\documentclass[aps,prd,onecolumn,tightenlines,nofootinbib,notitlepage,11pt,floatfix]{revtex4-2}
\usepackage{amsmath, amssymb, amsthm}
\usepackage{relsize}
\usepackage{enumitem}
\usepackage{booktabs}
\usepackage{longtable}
\usepackage{graphicx,lmodern,microtype,xcolor}\usepackage[T1]{fontenc}\definecolor{linkblue}{RGB}{22,64,128}
% revtex4-2 redefines \label, so nameref warns while loading and puts the
% kernel \label back, which is what hyperref needs. Mute nameref's warnings
% during the load only.
\makeatletter
\let\SGM@PackageWarning\PackageWarning
\def\SGM@nameref{nameref}
\def\PackageWarning#1#2{\def\SGM@pkg{#1}%
  \ifx\SGM@pkg\SGM@nameref\else\SGM@PackageWarning{#1}{#2}\fi}
\usepackage[colorlinks=true, urlcolor=linkblue, linkcolor=linkblue, citecolor=linkblue, bookmarks, psdextra]{hyperref}
\let\PackageWarning\SGM@PackageWarning
\makeatother

% Include the supplemental material after the bibliography. Change
% \supplementtrue to \supplementfalse to build the main text alone: the
% document then ends after the bibliography, the introduction drops its
% sentence about the supplement, and references cited only in the supplement
% drop out of the bibliography.
\newif\ifsupplement
\supplementtrue \newif\ifanonymous \anonymousfalse % \anonymoustrue restores double-anonymous review

% Lead each bibliography entry with its BibTeX key, which is what \citep takes
% in the source. Wrapping \@lbibitem at \begin{document} catches whatever
% definition natbib and hyperref have settled on by then.
\makeatletter
\AtBeginDocument{%
  \let\SGM@lbibitem\@lbibitem
  % BibTeX-key prefix disabled for clean preprint output.
}
\makeatother

% English patterns leave this German journal name unbreakable, which strands a
% loose line in the Madelung entry once the key precedes it.
\hyphenation{Zeit-schrift}\setlength{\emergencystretch}{2em}\tolerance=1000\def\bibpreamble{\raggedright}\linespread{1.08}

% Table spacing
\setlength{\tabcolsep}{7.5pt}
\renewcommand{\arraystretch}{1.18}\setlist{itemsep=2pt,topsep=4pt}
% revtex4-2 breaks the array package's >{} column prefix, so set every
% paragraph cell ragged right here instead. Its table-preamble builder also
% expands \@endpbox, which longtable v4.27 redefines, so a p column in a
% longtable fails with "Incomplete \iffalse"; keep \@endpbox unexpanded.
% At a longtable page break it reserves room for the foot but not for the
% depth of the last row, which can overfill the page; reserve \maxdepth too.
\makeatletter
\g@addto@macro\@arrayparboxrestore{\raggedright\let\\\tabularnewline}
% Stylish compact preprint abstract with thin accent rules and justified small text.
% Fits cleanly on Page 1 ahead of Section I.
\def\frontmatter@title@format{\LARGE\bfseries\centering\linespread{1.15}\selectfont\parskip\z@skip}\def\frontmatter@abstractheading{\begingroup\centering\small\bfseries Abstract\par\vspace{3pt}\endgroup}\def\frontmatter@abstractfont{\footnotesize\linespread{1.04}\selectfont\leftskip=0.04\textwidth\rightskip=\leftskip}\def\prep@absbox{\vspace{-4pt}\noindent\hspace{0.04\textwidth}\textcolor{black!25}{\rule{0.92\textwidth}{0.5pt}}\par\vspace{4pt}}\def\post@absbox{\vspace{2pt}\noindent\hspace{0.04\textwidth}\textcolor{black!25}{\rule{0.92\textwidth}{0.5pt}}\par\vspace{6pt}}
\appdef\@mkpream@relax{\let\@endpbox\relax}
\def\LT@adj{%
 \setbox\z@\vbox{\null}\dimen@-\ht\z@
 \setbox\z@\vbox{\unvbox\z@\LT@bot}\advance\dimen@\ht\z@
 \advance\dimen@\maxdepth
 \global\advance\vsize-\dimen@
}
\makeatother

\newtheoremstyle{indentedtheorem}%
  {\topsep}{\topsep}{\itshape}{\parindent}{\bfseries}{:}{.5em}{}
\theoremstyle{indentedtheorem}
\newtheorem{theorem}{Theorem}
\newtheorem{lemma}{Lemma}
\newtheorem{corollary}{Corollary}
\newtheorem{definition}{Definition}
\newtheorem{axiom}{Axiom}
\newtheorem{proposition}{Proposition}
\newtheorem{remark}{Remark}

\begin{document}

\title{Spectral Gradient Mechanics:\texorpdfstring{\\}{ }%
force as the gradient of wave function frequency,\texorpdfstring{\\}{ }%
and time dilation as the origin of gravity and inertia}
\newcommand{\companionauthor}{\ifanonymous Anonymous\else L.~Hutchison\fi}
\ifanonymous\else\author{Luke Hutchison}
\email{luke.hutch@alum.mit.edu}\noaffiliation\fi
\date{\today}

%===============================================================================
\begin{abstract}
%===============================================================================
In flat-spacetime quantum mechanics, a stationary state has one coordinate angular frequency $\omega=E/\hbar$ over its whole spatial extent, and its rest term $\omega_0=mc^2/\hbar$ is taken uniform everywhere, so $\nabla(\omega_0\alpha)\equiv\mathbf{0}$. In general relativity, clocks at different positions run at different rates $\alpha(\mathbf{x})=d\tau/dt$: a stationary metric still admits modes of one coordinate phase rate $E/\hbar$, whereas the local proper frequency $E/(\hbar\alpha(\mathbf{x}))$ and the rest-energy term $\omega_0\alpha(\mathbf{x})$ of the phase-space frequency function $\omega(\mathbf{x},\mathbf{p})$ vary with position ($\nabla(\omega_0\alpha)\neq\mathbf{0}$). The de~Broglie momentum-phase relation $\mathbf{p}=\hbar\nabla\theta$ treats the phase of the wave function as a field, and the time derivative of this relation along the particle's trajectory in a smooth field, with the gradient taken at fixed momentum, yields the force-frequency identity $\mathbf{F}=-\hbar\nabla\omega$ (which is also Hamilton's equation with the Hamiltonian $\hbar\omega$). Under the unnecessary assumption $\nabla(\omega_0\alpha)\equiv\mathbf{0}$ the rest term contributes zero to this identity. Instead, this paper demonstrates that the \emph{spectral gradient} $\nabla\omega$ of a particle is in general not zero, due at a minimum to the time dilation field; relaxing this single-frequency assumption resolves this longstanding tension between general relativity and quantum mechanics and explains the origin of all force laws other than those of the strong and weak interactions. The identity shows that gravity and inertia have a single origin, a gradient in the time dilation field (the static field for gravity, and the dynamic field, via Rindler time dilation, for inertia), which explains the equivalence principle. In a static field, the gravitational force on a particle at rest follows from the time dilation field alone. Planck's constant cancels from the gravitational force, and no combination of $c$ and Newton's constant $G$ has the dimension of action, so the constrained Newtonian and first post-Newtonian forces obtained in this way require neither a quantized gravitational field nor a graviton; configuration-space coupling yields an exact gravitationally induced entanglement purity law and a manifestly Hermitian $1$PN Einstein--Infeld--Hoffmann cross-momentum operator, whereas deterministic nonlinear mean-field gravity (both static Schr\"odinger--Newton and radiative) is excluded by finite-time state orthogonalization and Bell-pair gravitational-wave power signaling.
\end{abstract}

\maketitle

%===============================================================================
\section{Introduction}
\label{sec:intro}
%===============================================================================

Gravity is the one fundamental interaction without a complete quantum theory. The programs that seek one quantize the gravitational field, by canonical methods \citep{DeWitt1967Quantum}, in string theory \citep{Green1987Superstring,Polchinski1998String}, or in loop quantum gravity \citep{Ashtekar1986Variables,Rovelli2004Quantum}; a separate line treats gravity as emergent from thermodynamics or from entanglement entropy \citep{Jacobson1995Thermodynamics,Padmanabhan2010Thermodynamics,Verlinde2011Origin,Jacobson2016Entanglement,Verlinde2017Emergent}. Each of these programs locates the difficulty in the gravitational field, and whether that field must be quantized at all is unsettled: experiments that would detect entanglement generated by gravity alone have been proposed to decide it \citep{Bose2017Spin,Marletto2017Gravitationally}.

Quantum mechanics and general relativity also take inertia and its equality with gravitational mass as inputs. Inertia, the reaction $-m\mathbf{a}$ of a body of mass $m$ to its acceleration $\mathbf{a}$, enters quantum mechanics through the kinetic term of the Hamiltonian, as it enters Newtonian mechanics through the second law. Its origin has been sought in gravitational induction by the distant matter of the universe \citep{Sciama1953Origin} and in the dynamics of spatially extended trapped waves \citep{Jennison1977Approach,Jennison1983Wave}. The equality of gravitational and inertial mass, confirmed by the MICROSCOPE satellite at the level of $10^{-15}$ \citep{Touboul2017Microscope,Touboul2022Final}, is the founding hypothesis of general relativity \citep{Einstein1907Relativitatsprinzip}. In the Schr\"odinger equation gravity is a potential $mgz$, in a uniform field of strength $g$ at height $z$, whose coefficient is set equal to the inertial mass by postulate, and a transformation to accelerated coordinates adds the inertial potential in the same form \citep{Shukla2010Gauge}. A test with a single rubidium atom, split into a freely falling and a magnetically levitated wave packet, found their relative phase consistent with free fall in a uniform field \citep{Dobkowski2026Observation}; its interpretation as a test of the equivalence principle is disputed \citep{AsenbaumOverstreet2025Magnetic,Dobkowski2025Reply}. The measured phase does not contain the speed of light, so the test probes gravity only in its Newtonian limit.

The two theories treat time differently, which is the \emph{problem of time} \citep{Isham1993Canonical,Kuchavr1992Time,Anderson2012Problem}: quantum mechanics evolves a state in one external time parameter, while general relativity is invariant under diffeomorphisms of spacetime and has no such parameter. The Planck--Einstein relation $E=\hbar\omega$ \citep{Planck1901Radiation,Einstein1905Photoelectric} and the de~Broglie relation $\mathbf{p}=\hbar\nabla\theta$ \citep{deBroglie1925Thesis} give the wave function of every massive particle a frequency and a momentum, where $\hbar$ is the reduced Planck constant, $E$ and $\omega$ are the energy and the angular frequency, $\mathbf{p}$ is the momentum, and $\theta$ is the phase of the wave function. Quantum mechanics gives each stationary state one phase rate $E/\hbar$ at every point of the wave function's spatial extent, which in general is all of space; Hamilton--Jacobi theory and the Wentzel--Kramers--Brillouin approximation likewise separate the phase with one constant $E$ \citep{Wentzel1926WKB,Kramers1926WKB,Brillouin1926WKB,Goldstein2002Classical}. The frequency function $H(\mathbf{x},\mathbf{p})/\hbar$, where $H$ is the Hamiltonian, varies with position through the potential, but its rest term $mc^2/\hbar$, where $c$ is the speed of light, has one value at every point (the \emph{single-frequency assumption}). In flat spacetime the rest energy contributes only a phase factor that is the same at every point and changes no probability, and nonrelativistic quantum mechanics omits it for that reason.

General relativity does not allow this assumption. In a static field the rate of proper time $\tau$ relative to coordinate time $t$ at position $\mathbf{x}$ is the time-dilation factor $\alpha(\mathbf{x})=d\tau/dt$, whose variation has been resolved across one millimeter of height \citep{Bothwell2022Resolving}. In coordinate time the rest energy at $\mathbf{x}$ is therefore $mc^2\alpha(\mathbf{x})$, and the rest frequency there is $mc^2\alpha(\mathbf{x})/\hbar$. A stationary metric still admits modes of one coordinate frequency $E/\hbar$, but a clock held at $\mathbf{x}$ measures the frequency of such a mode as $E/(\hbar\alpha(\mathbf{x}))$.

In the frame that accelerates with a particle, the rate of proper time varies across the wave function by the time dilation of special relativity, so the rest frequency has a gradient across the state while the acceleration lasts. Where the rest frequency differs between two points of the wave function, the phases there drift apart and the phase gradient changes in time; by the de~Broglie relation the momentum changes with it. A stationary state keeps one coordinate frequency because other terms of its frequency function compensate this variation: the kinetic term along the momentum field of a falling particle, or a potential that holds the particle in place. The rest-energy term still has a nonzero gradient at fixed momentum, and minus this gradient is a force: it accelerates a falling particle, and the force of a holding potential balances it.

A change in a source also cannot reach the far parts of a state faster than $c$. An instantaneous update of the whole extent would let a controllable energy change send a signal to arbitrary distance, and no Lorentz-invariant dynamics does so; the correlations that violate Bell's inequalities \citep{Bell1964EPR,Aspect1982Experimental} send no signal, since the outcome statistics on each side do not depend on the setting chosen at the other. Nonrelativistic quantum mechanics nonetheless sets the potential of a source, such as the Coulomb potential, at every point of a state at the same instant. Without the single-frequency assumption the rest frequency, and with it the frequency function, is a field over the whole extent of the wave function and can follow the rate that general relativity gives the clock at each point. With the assumption, the two theories disagree about the frequency of one particle's wave function before the gravitational field is quantized.

General relativity treats gravity and inertia as one geometry. The weak-field derivation of Newton's law from time dilation acting on the rest frequency $mc^2/\hbar$, also called the Compton frequency, has been found independently several times \citep{Evans2001Matter,Wootters2003Why,Mueller2014Quantum,Xu2019DeBroglie,Czarnecki2021Gravitational}. Inertia has been obtained from the frequency change that accompanies motion \citep{Schlegel1954Wave}. For an extended charge, the propagation of light in accelerated and supported frames gives inertial and gravitational self-forces with the same mass \citep{Petkov1999Acceleration}. The emergent-gravity program derives Newton's law and inertia from a single thermodynamic principle \citep{Verlinde2011Origin}. The Compton-frequency derivations use the static gradient alone. The Dirac Hamiltonian of a uniformly accelerated frame contains a rest-energy term that varies with position \citep{HehlNi1990Inertial}, and the inertial and tidal potentials of a localized system have been built from its rest-mass term times the local redshift factor \citep{Perche2022Localized}. This paper derives both the static and the Rindler gradient from the same rest frequency through the identity $\mathbf{F}=-\hbar\nabla\omega$. Their common coefficient gives the equality of gravitational and inertial mass and free fall, and the second derivatives of the same frequency field give the tidal difference between gravity and inertia.

The Dirac spinor has been written as the square root of a density, with a phase factor, times a rotation \citep{Hestenes1967Real,Hestenes1990Zitterbewegung}, and gravity has been formulated as a gauge theory in which rotations act on that spinor from the left and phase rotations from the right \citep{LasenbyDoranGull1998Gravity}. The Aharonov--Bohm phase of a charge that encircles a magnetic flux through a region free of field \citep{Aharonov1959Significance} has been given local accounts that add degrees of freedom beyond the particle's own: the charges of the solenoid \citep{Vaidman2012Role} or the quantized electromagnetic field \citep{Marletto2020Aharonov}.

Section~\ref{sec:identity} states the axioms, derives the identity by four routes, and shows that it holds, in expectation values, for every interaction with a quantum Hamiltonian. Section~\ref{sec:grav_inertia} obtains Newton's law from Gauss's law for the rate of clocks through the static gradient, and the inertial reaction from the dynamic gradient; from their common coefficient it obtains the equality of gravitational and inertial mass and free fall, and from the second derivatives of the frequency the tidal difference between the two. In the same section the balance of an applied force against the inertial reaction gives Newton's second law, and the dispersion relation of a massless wave gives the deflection of light and the Shapiro delay. Section~\ref{sec:potentials} adds the gauge potentials and, from the companion paper \citep{HutchisonWavefn2026}, the exact frequency equation of the Dirac equation and its expansion into the common force laws, and shows that the identity holds at every spin. Section~\ref{sec:canonical} recovers the Lagrangian, Hamiltonian, and conservation laws, the angular gradient, and a local account of the Aharonov--Bohm phase. Sections~\ref{sec:experiment} and~\ref{sec:scope} treat experiment and prior work, and Section~\ref{sec:discussion} gives three corollaries: the reach of a particle's couplings, the cancellation of Planck's constant from the gravitational force, and static fields, gravity among them, as functions on the configuration space of all particles. \ifsupplement The supplemental material after the bibliography gives the notation and conventions and, in full, the derivations that the main text compresses.\fi

%===============================================================================
\section{The force-frequency identity}
\label{sec:identity}
%===============================================================================

At each point of the wave function the frequency is the sum of three terms: the rest term $mc^2\alpha/\hbar$, where $\alpha$ is the local time-dilation factor; a kinetic term from the motion; and terms from the potentials. The time-dilation factor and the potentials vary with position, so the frequency is a field over the spatial extent of the wave function. The phase of one particle's state depends on fields across macroscopic distances: in atom interferometry the two wave packets of one atom, separated by as much as 16~cm, accumulate a relative phase from the spacetime curvature across that separation \citep{Asenbaum2017Phase}.

\begin{axiom}[Frequency is a field, measured only at the position of the particle]\label{ax:local_frequency}
A wave function's frequency is a field $\omega(\mathbf{x})$ over the wave function's spatial extent; at each point where the amplitude is nonzero, $\omega(\mathbf{x})$ is the rate $-\partial_t\theta$, at fixed position, at which the phase $\theta$ decreases under the equation of motion, in the gauge of the stated potentials.
\end{axiom}

\noindent For a particle of charge $q$ in an electromagnetic field with scalar potential $\Phi$, a gauge transformation by a function $\lambda(\mathbf{x},t)$ multiplies the wave function by $e^{iq\lambda/\hbar}$, where $i$ is the imaginary unit. It adds $-q\,\partial_t\lambda/\hbar$ to $\omega(\mathbf{x})$ and $-q\,\partial_t\lambda$ to $q\Phi$, so $\hbar\omega(\mathbf{x})-q\Phi$ is gauge invariant.

The \emph{far tail} of the state is the region beyond a few uncertainty widths of the particle's position, in which the probability of finding the particle is negligible. To bound the effect of a change of frequency in the far tail, write the normalized state at one instant as $\psi=\psi_\text{core}+\psi_\text{tail}$, where $\psi_\text{tail}$ is its part in the far tail and $\psi_\text{core}$ is its part elsewhere. In Dirac notation, $\langle\phi_1|\phi_2\rangle=\int\phi_1^*\phi_2\,d^3x$ is the inner product of two wave functions, where $\phi_1^*$ is the complex conjugate of $\phi_1$. The norm is $\|\phi\|=\langle\phi|\phi\rangle^{1/2}$, and the probability of finding the particle in the far tail is $P_\text{tail}=\|\psi_\text{tail}\|^2$. A change of frequency there changes the phase of $\psi_\text{tail}$ and leaves its modulus unchanged, giving the state $\psi'=\psi_\text{core}+\psi'_\text{tail}$ with $\|\psi'_\text{tail}\|^2=P_\text{tail}$. The part $\psi_\text{core}$ vanishes in the far tail and both tails vanish outside it, so $\langle\psi|\psi'\rangle=1-P_\text{tail}+\langle\psi_\text{tail}|\psi'_\text{tail}\rangle$, and by the Cauchy--Schwarz inequality the last term has magnitude at most $P_\text{tail}$. For $P_\text{tail}\le 1/2$ the overlap therefore obeys $|\langle\psi|\psi'\rangle|\ge 1-2P_\text{tail}$. Hence $1-|\langle\psi|\psi'\rangle|^2\le 1-(1-2P_\text{tail})^2=4P_\text{tail}(1-P_\text{tail})\le 4P_\text{tail}$, and
\begin{equation}
    \sqrt{1-|\langle\psi|\psi'\rangle|^2} \le 2\sqrt{P_\text{tail}} .
    \label{eq:tail_distance}
\end{equation}
The left side is the \emph{trace distance} $d_\text{tr}$ between the two pure states. (By comparison, truncating $\psi$ to the core and renormalizing to $\psi_\text{box}\equiv\psi_\text{core}/\sqrt{1-P_\text{tail}}$ gives $|\langle\psi|\psi_\text{box}\rangle|=\sqrt{1-P_\text{tail}}$ and the exact trace distance $d_\text{tr}(\psi,\psi_\text{box})=\sqrt{P_\text{tail}}$, half the phase-perturbation bound; under causal finite-speed evolution over a duration $T$, padding the core by $cT$ preserves the bound.) The probability of a measurement outcome is $\langle\psi|\Pi|\psi\rangle$, where $\Pi$ is the projector onto the states of that outcome, so its change is the trace of $\Pi$ times the operator $|\psi\rangle\langle\psi|-|\psi'\rangle\langle\psi'|$. That operator acts in the plane of $\psi$ and $\psi'$, its trace is zero, and the trace of its square is $2-2|\langle\psi|\psi'\rangle|^2$, so its eigenvalues are $\pm d_\text{tr}$. The change is therefore $d_\text{tr}$ times the difference of the expectations of $\Pi$ in the two eigenvectors. Each expectation lies between $0$ and $1$, so the magnitude of the change is at most $d_\text{tr}$, which by Eq.~\eqref{eq:tail_distance} is at most $2\sqrt{P_\text{tail}}$. This bound holds at one instant for every outcome, including the outcomes of momentum and phase measurements. It is below a tolerance $\varepsilon$ when $P_\text{tail}\le\varepsilon^2/4$; for a Gaussian state and $\varepsilon=10^{-2}$, the far tail is then the region beyond about $4.2$ standard deviations on either side of the mean in one dimension, and beyond a radius of about $4.9$ standard deviations in three. The expectation of an unbounded observable, such as the energy, can change by more, since a small probability can multiply a large eigenvalue.

\begin{axiom}[The Planck--Einstein relation]\label{ax:planck}
The energy of a stationary quantum state is the reduced Planck constant $\hbar$ times the angular frequency of its wave function \citep{Planck1901Radiation,Einstein1905Photoelectric},
\begin{equation}
    E = \hbar \, \omega .
    \label{eq:planck}
\end{equation}
In a superposition of stationary states each component has its own frequency.
\end{axiom}

Equating $E=\hbar\omega$ with the rest energy $E=mc^2$ \citep{Einstein1905Inertia} gives a particle of mass $m$ the rest frequency
\begin{equation}
    \omega_0 = mc^2/\hbar ,
    \label{eq:rest_freq}
\end{equation}
the frequency of its wave function at rest and free of every interaction. Since the frequency is a field (Axiom~\ref{ax:local_frequency}), Eq.~\eqref{eq:planck} holds at each point, $E(\mathbf{x})=\hbar\,\omega(\mathbf{x})$, so an interaction that shifts the energy at a point shifts the frequency there away from $\omega_0$ by the energy shift divided by $\hbar$.

Throughout, an overdot is a total time derivative along the particle's trajectory, $\dot{\mathcal{X}}\equiv d\mathcal{X}/dt$ for any quantity $\mathcal{X}$, so the velocity is $\mathbf{v}=\dot{\mathbf{x}}$; the partial derivative $\partial_t$ is the rate at a fixed position. In both, $t$ is the coordinate time of the frame in which the fields are written. An unqualified \emph{frequency} or \emph{rate} is measured per unit of that $t$, and an unqualified \emph{energy} is $\hbar$ times such a frequency. A quantity read on a clock carried at the point itself is always called \emph{proper}.

\begin{axiom}[The de~Broglie relation]\label{ax:debroglie}
A particle's momentum is $\hbar$ times the spatial gradient of the phase $\theta(\mathbf{x},t)$ of its wave function \citep{deBroglie1925Thesis},
\begin{equation}
    \mathbf{p} = \hbar \, \nabla\theta ,
    \label{eq:debroglie}
\end{equation}
and its angular frequency is the negative time rate of that phase,
\begin{equation}
    \omega = -\partial_t\theta .
    \label{eq:phase_freq}
\end{equation}
The phase gradient is the \emph{wavevector} $\mathbf{k}\equiv\nabla\theta$, whose magnitude is the inverse of the reduced de~Broglie wavelength $\lambdabar_\mathrm{dB}=\hbar/|\mathbf{p}|$, and Eq.~\eqref{eq:debroglie} is equivalently
\begin{equation}
    \mathbf{p} = \hbar \, \mathbf{k} .
    \label{eq:debroglie_k}
\end{equation}
\end{axiom}

\noindent The spatial gradient of the phase is a momentum field $\mathbf{p}(\mathbf{x})=\hbar\nabla\theta$, a different object from the state's Fourier content, whose spread the uncertainty relation bounds. The negative time rate is the \emph{phase rate} $\omega(\mathbf{x})=-\partial_t\theta$, the frequency field of Axiom~\ref{ax:local_frequency}, a coordinate frequency because the derivative is taken in $t$. The de~Broglie relation holds pointwise, as the canonical momentum per unit probability density $\rho$. It holds for a wave function $\psi$ of one component, with $\rho=|\psi|^2$, and for a wave function of several components $\rho^{1/2}e^{i\theta}\xi$ with a constant unit spinor $\xi$. Where $\xi$ varies, the momentum per unit density gains $-i\hbar\,\xi^\dagger\nabla\xi$, where $\xi^\dagger$ is the conjugate transpose of $\xi$.

\begin{axiom}[Hamiltonian dynamics]\label{ax:hamilton}
A particle with Hamiltonian $H(\mathbf{x},\mathbf{p})$ moves by Hamilton's canonical equations \citep{Hamilton1835Second,Goldstein2002Classical}. In particular the momentum changes at a rate equal to the negative spatial gradient of the Hamiltonian taken at fixed momentum, $\dot{\mathbf{p}}=-\nabla_{\mathbf{x}}H|_{\mathbf{p}}$, and force is by definition that rate, $\mathbf{F}=\dot{\mathbf{p}}$.
\end{axiom}

\noindent The kinetic energy in the Hamiltonian makes the frequency depend on momentum as well as on position. The frequency is therefore also a function on phase space, the \emph{frequency function} $\omega(\mathbf{x},\mathbf{p})=H(\mathbf{x},\mathbf{p})/\hbar$ given by Eq.~\eqref{eq:planck}. The function is fixed by the Hamiltonian; the state determines only the points at which it is evaluated, such as the particle's mean position and mean momentum, which are jointly well defined.

The \emph{quantum potential} of the wave function $\psi$ is $Q=-(\hbar^2/2m)\,\nabla^2|\psi|/|\psi|$ \citep{Madelung1927Quantentheorie,Bohm1952Suggested}, where $\nabla^2$ is the Laplacian. It depends on the amplitude alone and is negligible wherever the amplitude varies slowly on the scale of the de~Broglie wavelength. For the Hamiltonian $H=mc^2+\mathbf{p}^2/2m+V(\mathbf{x})$, with the rest energy of Eq.~\eqref{eq:rest_freq} and the potential energy $V$, the Schr\"odinger equation \citep{Schrodinger1926Quantization4} written in terms of the amplitude $|\psi|$ and the phase $\theta$ gives the phase rate as the frequency function evaluated on the state's own momentum field $\mathbf{p}(\mathbf{x})=\hbar\nabla\theta$ (Eq.~\eqref{eq:debroglie}), plus $Q/\hbar$:
\begin{equation}
    -\partial_t\theta = \omega\big(\mathbf{x},\mathbf{p}(\mathbf{x})\big) + Q(\mathbf{x})/\hbar .
    \label{eq:phase_rate_split}
\end{equation}
The phase rate $\omega(\mathbf{x})$ and the frequency function $\omega(\mathbf{x},\mathbf{p})$ share a symbol and differ in their arguments. Likewise $\mathbf{p}$ denotes either the particle's momentum or the momentum field $\mathbf{p}(\mathbf{x})$, and $\mathbf{x}$ either the particle's position or a point of the wave function's extent.

\begin{definition}[The spectral gradient]\label{def:spectral_gradient}
Throughout, $\nabla$ denotes the gradient of a function on phase space, such as the frequency function $\omega(\mathbf{x},\mathbf{p})$, with respect to position at fixed momentum,
\begin{equation}
    \nabla \equiv \nabla_{\mathbf{x}}\big|_{\mathbf{p}} \, .
    \label{eq:gradient_def}
\end{equation}
Applied to a function of position alone, such as a potential or the phase $\theta(\mathbf{x})$, the gradient $\nabla$ is the ordinary spatial gradient. The vector $\nabla\omega$, evaluated at the position and momentum of the particle, is the \emph{spectral gradient} of the wave function.
\end{definition}

\noindent In a stationary state of energy $E$ the phase separates as $\hbar\theta=S_0(\mathbf{x})-E\,t$, where $S_0(\mathbf{x})$ is a function of position alone, equal to Hamilton's abbreviated action \citep{Goldstein2002Classical} wherever $Q$ is negligible. The phase rate $-\partial_t\theta=E/\hbar$ is then the same at every point, and its spatial gradient is zero. The spectral gradient need not be zero: for $H=mc^2+\mathbf{p}^2/2m+V(\mathbf{x})$, Eq.~\eqref{eq:planck} gives $\hbar\omega=H$ and so $\nabla\omega=\nabla V/\hbar$, which is nonzero wherever $V$ varies.

\begin{axiom}[Static case: the rate of time varies with position]\label{ax:nonzero_spectral_gradient}
The rate of proper time relative to coordinate time varies continuously with position, as set by the distribution of mass, and its gradient is in general nonzero \citep{Einstein1907Relativitatsprinzip,Einstein1911Influence,Pound1960Apparent,Bothwell2022Resolving}.
\end{axiom}

\begin{axiom}[Causality: frequency changes propagate no faster than light]\label{ax:causal_propagation}
Information propagates no faster than $c$, the speed of light in vacuum \citep{Einstein1905Electrodynamics}. A localized change of frequency reaches a point at displacement $\delta\mathbf{x}$ through the phase no earlier than $|\delta\mathbf{x}|/c$. In the Dirac equation no solution propagates faster than $c$ \citep{Thaller1992Dirac}, and the front of the new phase, where the solution changes discontinuously, moves at the characteristic speed $c$.
\end{axiom}

\begin{axiom}[Local position invariance]\label{ax:lpi}
The outcome of any local nongravitational experiment is independent of where and when it is performed \citep{Will2014Confrontation}. In particular, the phase of a particle at rest advances at the same rate $\omega_0$ of Eq.~\eqref{eq:rest_freq} per unit of proper time at every location.
\end{axiom}

%-------------------------------------------------------------------------------
\subsection{Four derivations of the identity}
\label{sec:deriv}
%-------------------------------------------------------------------------------

%-------------------------------------------------------------------------------
\subsubsection{Route I: from force and energy}
\label{sec:route_energy}
%-------------------------------------------------------------------------------

The force on a particle is the negative spatial gradient of its energy $E(\mathbf{x},\mathbf{p})$ at fixed momentum (Eq.~\eqref{eq:gradient_def}),
\begin{equation}
    \mathbf{F} = -\nabla E \, .
    \label{eq:force_energy}
\end{equation}
Here $E$ takes the place of the function whose gradient is minus the force in Lagrange's analytical mechanics \cite{Lagrange1788Mecanique}. Lagrange derived from it the conservation of vis viva: the kinetic energy plus that function is constant. For central forces Helmholtz \cite{Helmholtz1847Erhaltung} called the constancy of this sum the conservation of force. Equation~\eqref{eq:force_energy} is the content of Hamilton's canonical equation (Axiom~\ref{ax:hamilton}) when the Hamiltonian equals the energy, $H=E$. Any part of the energy that depends on momentum alone has zero gradient at fixed momentum and drops out, so the force comes from the explicit position dependence of the energy. The Planck--Einstein relation (Eq.~\eqref{eq:planck}) identifies that energy with $\hbar$ times the frequency of the wave function,
\begin{align}
    H &= E \notag\\
      &= \hbar\,\omega(\mathbf{x},\mathbf{p}) .
    \label{eq:H_eq_hw}
\end{align}
Substituting Eq.~\eqref{eq:H_eq_hw} into Eq.~\eqref{eq:force_energy} gives
\begin{equation}
    \begin{gathered}
    \boxed{\; \mathbf{F} = -\hbar\nabla\omega \;} \, , \\[-3pt]
    \text{\small\itshape (The SGM identity)}
    \end{gathered}
    \label{eq:sgm}
\end{equation}
where $\mathbf{F}$ is the force on the particle (Axiom~\ref{ax:hamilton}), $\omega(\mathbf{x},\mathbf{p})$ is the frequency function of Eq.~\eqref{eq:H_eq_hw}, and $\nabla$ is its gradient in position at fixed momentum (Eq.~\eqref{eq:gradient_def}), evaluated at the particle's position. With the force written as the rate of change of momentum (Axiom~\ref{ax:hamilton}), the identity reads
\begin{equation}
    \begin{gathered}
    \boxed{\; \dot{\mathbf{p}} = -\hbar\nabla\omega \;} \, . \\[-3pt]
    \text{\small\itshape (The SGM identity, momentum form)}
    \end{gathered}
    \label{eq:sgm_as_dot_p}
\end{equation}
Here $\mathbf{p}$ is the particle's canonical momentum, the momentum paired with the position $\mathbf{x}$ in Hamilton's equations (Axiom~\ref{ax:hamilton}).

\begin{definition}[Spectral gradient mechanics (SGM)]\label{def:sgm_def}
Spectral gradient mechanics (SGM) is the identification of the spectral gradient as the origin of force laws. Equation~\eqref{eq:sgm}, equivalently Eq.~\eqref{eq:sgm_as_dot_p}, is the \emph{SGM identity}, or force-frequency identity.
\end{definition}

%-------------------------------------------------------------------------------
\subsubsection{Route II: from momentum and phase}
\label{sec:route_phase}
%-------------------------------------------------------------------------------

The identity also follows from the de~Broglie relation $\mathbf{p}=\hbar\nabla\theta$ (Eq.~\eqref{eq:debroglie}) together with the dispersion relation, the frequency function $\omega(\mathbf{x},\mathbf{p})$ of Eq.~\eqref{eq:H_eq_hw}. Differentiating the de~Broglie relation in time at fixed position gives $\partial_t\mathbf{p}=\hbar\,\partial_t\nabla_{\mathbf{x}}\theta$, where $\nabla_{\mathbf{x}}$ is the spatial gradient at fixed time. On a twice continuously differentiable ($\theta\in C^2$) phase field that gradient and the time derivative commute,
\begin{equation}
    \partial_t\nabla_{\mathbf{x}}\theta = \nabla_{\mathbf{x}}\partial_t\theta \, ,
    \label{eq:mixed_orders}
\end{equation}
and the time derivative of the phase is the negative of the frequency field, $\omega(\mathbf{x})=-\partial_t\theta$ (Eq.~\eqref{eq:phase_freq}), giving
\begin{equation}
    \partial_t\mathbf{p} = -\hbar\,\nabla_{\mathbf{x}}\omega .
    \label{eq:dp_dt_local}
\end{equation}

Equation~\eqref{eq:dp_dt_local} holds at a fixed point of space, and the gradient in it is the full gradient of the frequency field $\omega(\mathbf{x})$. Where the quantum potential $Q$ is negligible, that field equals $\omega(\mathbf{x},\mathbf{p}(\mathbf{x}))$, the frequency function evaluated on the state's own momentum field (Eq.~\eqref{eq:phase_rate_split}). The chain rule splits that gradient into the spectral gradient and a convective term,
\begin{equation}
    \partial_{x_i}\,\omega(\mathbf{x},\mathbf{p}(\mathbf{x})) = \partial_{x_i}\omega\big|_{\mathbf{p}} + (\partial_{x_i}p_j)\,\partial_{p_j}\omega ,
    \label{eq:chain_convective}
\end{equation}
with the index $j$ summed over the three spatial axes. In the convective term, $\partial_{x_i}p_j=\hbar\,\partial_{x_i}\partial_{x_j}\theta$ is $\hbar$ times the Hessian of the phase (Eq.~\eqref{eq:debroglie}), and $\hbar\,\partial_{p_j}\omega=\partial H/\partial p_j$ is the group velocity \citep{deBroglie1925Thesis}, which by Hamilton's first canonical equation (Axiom~\ref{ax:hamilton}) is the particle velocity $\dot{x}_j$. The trajectory is the curve $\mathbf{x}(t)$ with this velocity, and along it the momentum is the value $\mathbf{p}(\mathbf{x}(t),t)$ of the momentum field. The force is the rate of the momentum along the trajectory, the rate at a fixed point plus the change due to the particle's motion through the momentum field,
\begin{equation}
    \dot{\mathbf{p}} = \partial_t\mathbf{p} + (\dot{\mathbf{x}}\cdot\nabla)\,\mathbf{p} .
    \label{eq:total_momentum_rate}
\end{equation}
Substituting Eq.~\eqref{eq:dp_dt_local} for $\partial_t\mathbf{p}$ in Eq.~\eqref{eq:total_momentum_rate}, splitting its gradient by Eq.~\eqref{eq:chain_convective}, and writing each derivative of the momentum field as $\hbar$ times the Hessian of the phase gives, in components,
\begin{equation}
    \dot{p}_i = -\hbar\,\partial_{x_i}\omega\big|_{\mathbf{p}} - \hbar\,\dot{x}_j\,\partial_{x_i}\partial_{x_j}\theta + \hbar\,\dot{x}_j\,\partial_{x_j}\partial_{x_i}\theta .
    \label{eq:route2_components}
\end{equation}
The second term is the convective term of Eq.~\eqref{eq:chain_convective}, with the sign of Eq.~\eqref{eq:dp_dt_local}; the third is the term $(\dot{\mathbf{x}}\cdot\nabla)\,\mathbf{p}$ of Eq.~\eqref{eq:total_momentum_rate}. Since $\partial_{x_i}\partial_{x_j}\theta=\partial_{x_j}\partial_{x_i}\theta$ for $\theta\in C^2$, the two cancel at every particle velocity, leaving $\dot{p}_i=-\hbar\,\partial_{x_i}\omega|_{\mathbf{p}}$. The derivative at fixed momentum is the gradient $\nabla$ of Eq.~\eqref{eq:gradient_def}, so the three components together are Eq.~\eqref{eq:sgm_as_dot_p}.

In the instantaneous rest frame, $\dot{\mathbf{x}}=\mathbf{0}$, the second and third terms of Eq.~\eqref{eq:route2_components} vanish separately. In a stationary state in which $Q$ is negligible the frequency field is uniform at $E/\hbar$, so the left side of Eq.~\eqref{eq:chain_convective} vanishes and the spectral gradient, which gives the force, equals minus the convective term.

The derivation equated the phase rate $-\partial_t\theta$ with the function $\omega(\mathbf{x},\mathbf{p}(\mathbf{x}))$, which holds where the quantum potential $Q$ is negligible. For an exact state the phase rate is the function plus $Q/\hbar$ (Eq.~\eqref{eq:phase_rate_split}), and the gradient of $Q$ adds to the force. In the ground state of a one-dimensional harmonic oscillator of angular frequency $\omega_\text{osc}$, for example, the wave function is real, so its phase is uniform in space and the momentum field is zero. By Eq.~\eqref{eq:phase_rate_split} the phase rate is then $(mc^2+V+Q)/\hbar$, which equals $\omega_0+\omega_\text{osc}/2$ at every point, so the force obtained from the phase vanishes while $-\nabla V$ does not. In that state the identity at fixed momentum still holds for the mean momentum and the mean force, which are both zero.

The steps from Eq.~\eqref{eq:debroglie} to Eq.~\eqref{eq:sgm_as_dot_p} are the standard derivation, in geometrical optics, of the eikonal ray equation \citep{Whitham1974Linear},
\begin{equation}
    \dot{\mathbf{k}} = -\nabla\omega ,
    \label{eq:eikonal_ray}
\end{equation}
with the gradient at fixed wavevector (Eq.~\eqref{eq:gradient_def}), as in ray tracing \citep{Dodin2009Langmuir}. Equation~\eqref{eq:eikonal_ray} holds for every wave that has a phase field and a dispersion relation $\omega(\mathbf{x},\mathbf{k})$, light included \citep{Evans1986Optics}, wherever the wavelength is short compared with the length over which the medium varies. Multiplying Eq.~\eqref{eq:eikonal_ray} by $\hbar$ gives Eq.~\eqref{eq:sgm_as_dot_p}, since $\hbar\dot{\mathbf{k}}=\dot{\mathbf{p}}$ by $\mathbf{p}=\hbar\mathbf{k}$ (Eq.~\eqref{eq:debroglie_k}), and a gradient at fixed wavevector is a gradient at fixed momentum. Planck's constant enters the identity only through these conversions of wave quantities into mechanical ones, $\mathbf{p}=\hbar\mathbf{k}$ and $E=\hbar\omega$ (Eq.~\eqref{eq:planck}).

%-------------------------------------------------------------------------------
\subsubsection{Route III: from the Euler--Lagrange equation}
\label{sec:route_lagrangian}
%-------------------------------------------------------------------------------

For a Lagrangian $L(\mathbf{x},\dot{\mathbf{x}})$ the Euler--Lagrange equation \citep{Euler1744Methodus,Lagrange1788Mecanique},
\begin{equation}
    \frac{d}{dt}\frac{\partial L}{\partial\dot{\mathbf{x}}} = \frac{\partial L}{\partial\mathbf{x}} ,
    \label{eq:euler_lagrange}
\end{equation}
equates the time rate of the canonical momentum $\mathbf{p}=\partial L/\partial\dot{\mathbf{x}}$ with a derivative at fixed velocity. The Legendre transform $H(\mathbf{x},\mathbf{p})=\mathbf{p}\cdot\dot{\mathbf{x}}-L(\mathbf{x},\dot{\mathbf{x}})$ relates this derivative of $L$ at fixed velocity to the derivative of the Hamiltonian $H$ of Eq.~\eqref{eq:H_eq_hw} at fixed momentum. Inverting $\mathbf{p}=\partial L/\partial\dot{\mathbf{x}}$ gives the velocity as a function $\dot{\mathbf{x}}(\mathbf{x},\mathbf{p})$, so at fixed momentum the velocity still changes with position. Differentiating the transform with respect to the position component $x_i$ at fixed momentum therefore gives three terms by the chain rule, one from each place the velocity appears and one from the explicit position dependence of $L$, with the index $j$ summed:
\begin{equation}
    \frac{\partial H}{\partial x_i}\bigg|_{\mathbf{p}} = p_j\,\frac{\partial\dot{x}_j}{\partial x_i}\bigg|_{\mathbf{p}} - \frac{\partial L}{\partial x_i}\bigg|_{\dot{\mathbf{x}}} - \frac{\partial L}{\partial\dot{x}_j}\,\frac{\partial\dot{x}_j}{\partial x_i}\bigg|_{\mathbf{p}} .
    \label{eq:legendre_chain}
\end{equation}
The first and third cancel by $p_j=\partial L/\partial\dot{x}_j$, leaving, for each component,
\begin{equation}
    \frac{\partial L}{\partial\mathbf{x}}\bigg|_{\dot{\mathbf{x}}} = -\frac{\partial H}{\partial\mathbf{x}}\bigg|_{\mathbf{p}} .
    \label{eq:legendre_bridge}
\end{equation}
The Euler--Lagrange equation and Eq.~\eqref{eq:legendre_bridge}, with $H=\hbar\omega$ (Eq.~\eqref{eq:H_eq_hw}), then give $\dot{\mathbf{p}}=\partial L/\partial\mathbf{x}|_{\dot{\mathbf{x}}}=-\partial H/\partial\mathbf{x}|_{\mathbf{p}}=-\hbar\nabla\omega$, the SGM identity (Eq.~\eqref{eq:sgm_as_dot_p}). The Euler--Lagrange equation holds on the path of stationary action, which is the path the particle follows \citep{Goldstein2002Classical}; Eq.~\eqref{eq:legendre_bridge} is an identity wherever the Legendre transform exists.

%-------------------------------------------------------------------------------
\subsubsection{Route IV: from the Schr\"odinger equation by Ehrenfest's theorem}
\label{sec:route_schrodinger}
%-------------------------------------------------------------------------------

Write the state as a wave function $\psi(\mathbf{x},t)$, normalized so that $\int|\psi|^2\,d^3x=1$, evolving under the Schr\"odinger equation $i\hbar\,\partial_t\psi=\hat{H}\psi$ \citep{Schrodinger1926Quantization4} with the Hamiltonian operator
\begin{equation}
    \hat{H}=\frac{\hat{\mathbf{p}}^2}{2m}+V(\mathbf{x}) ,
    \label{eq:schrodinger_hamiltonian}
\end{equation}
where $\hat{\mathbf{p}}=-i\hbar\nabla$ is the momentum operator. The expectation value of an operator $\hat{A}$ in the state is $\langle\hat{A}\rangle=\int\psi^*\hat{A}\psi\,d^3x$. The frequency operator is the Hamiltonian divided by $\hbar$, $\hat{\omega}=\hat{H}/\hbar$ (Eq.~\eqref{eq:planck}).

The \emph{commutator} is $[\hat{H},\hat{A}]=\hat{H}\hat{A}-\hat{A}\hat{H}$. For an operator without explicit time dependence the evolution law of expectation values \citep{Ehrenfest1927Bemerkung,Sakurai2020Modern} follows from the product rule (first line), the Schr\"odinger equation and its complex conjugate (second), the Hermiticity of $\hat{H}$ applied to the first integral (third), and the definition of the commutator (fourth):
\begin{align}
    \frac{d}{dt}\langle\hat{A}\rangle &= \int(\partial_t\psi^*)\,\hat{A}\psi\,d^3x + \int\psi^*\hat{A}\,(\partial_t\psi)\,d^3x \notag\\
    &= \frac{i}{\hbar}\int(\hat{H}\psi)^*\hat{A}\psi\,d^3x - \frac{i}{\hbar}\int\psi^*\hat{A}\hat{H}\psi\,d^3x \notag\\
    &= \frac{i}{\hbar}\int\psi^*\big(\hat{H}\hat{A}-\hat{A}\hat{H}\big)\psi\,d^3x \notag\\
    &= \frac{i}{\hbar} \, \big\langle[\hat{H},\hat{A}]\big\rangle .
    \label{eq:ehrenfest_commutator}
\end{align}
The derivation requires that the state be normalizable, that $\psi$ and $\hat{A}\psi$ lie in the domain of $\hat{H}$, and that $\hat{H}\psi$ lie in the domain of $\hat{A}$, so that the expectation values are finite and $\hat{H}$ can act on either factor of the integrand with no boundary term.

Take $\hat{A}=\hat{\mathbf{p}}$. The kinetic term of $\hat{H}$ is built from $\hat{\mathbf{p}}$ and commutes with it, so the whole commutator comes from the potential. Applied to the state, the commutator gives
\begin{align}
    [\hat{H},\hat{\mathbf{p}}] \, \psi &= V \, (-i\hbar\nabla\psi) + i\hbar\,\nabla(V\psi) \notag\\
    &= i\hbar\,(\nabla V) \, \psi ,
    \label{eq:commutator_V}
\end{align}
because $\nabla(V\psi)=V\,\nabla\psi+(\nabla V)\,\psi$ and the $V\,\nabla\psi$ terms cancel. Substituting Eq.~\eqref{eq:commutator_V} into Eq.~\eqref{eq:ehrenfest_commutator}, with $(i/\hbar)(i\hbar)=-1$, gives Ehrenfest's theorem \citep{Ehrenfest1927Bemerkung},
\begin{equation}
    \frac{d}{dt}\langle\hat{\mathbf{p}}\rangle = -\langle\nabla V\rangle .
    \label{eq:ehrenfest}
\end{equation}
For the frequency function $\omega(\mathbf{x},\mathbf{p})=H(\mathbf{x},\mathbf{p})/\hbar$ of this Hamiltonian, $H=\mathbf{p}^2/2m+V$, the kinetic term has no position dependence at fixed momentum, so $\hbar\nabla\omega=\nabla V$ (Eq.~\eqref{eq:gradient_def}) at every momentum. The operator of the spectral gradient is $\nabla\hat{\omega}\equiv(i/\hbar)[\hat{\mathbf{p}},\hat{\omega}]$, which by Eq.~\eqref{eq:commutator_V} multiplies the wave function by $\nabla V/\hbar$. Equation~\eqref{eq:ehrenfest} is therefore the SGM identity written for expectation values,
\begin{equation}
    \frac{d}{dt}\langle\hat{\mathbf{p}}\rangle = -\hbar\,\langle\nabla\hat{\omega}\rangle .
    \label{eq:sgm_expectation}
\end{equation}
With $\nabla\hat{\omega}$ defined by the commutator, Eq.~\eqref{eq:sgm_expectation} is also Eq.~\eqref{eq:ehrenfest_commutator} for $\hat{A}=\hat{\mathbf{p}}$ and $\hat{H}=\hbar\hat{\omega}$, so it holds for any Hamiltonian.

The particle's momentum and position are the expectation values $\mathbf{p}=\langle\hat{\mathbf{p}}\rangle$ and $\mathbf{x}=\langle\hat{\mathbf{x}}\rangle$, where $\hat{\mathbf{x}}$ is the position operator. The Taylor expansion of $\nabla V(\hat{\mathbf{x}})$ about $\mathbf{x}$, averaged over the state, has a first-order term proportional to $\langle\hat{\mathbf{x}}-\mathbf{x}\rangle$, which is zero, so $\langle\nabla V\rangle=\nabla V(\mathbf{x})$ up to terms of second order in the width of the state, and exactly when $V$ is at most quadratic in position. To that order, with $\hbar\,\nabla\omega=\nabla V$ at the mean position, Eq.~\eqref{eq:sgm_expectation} is the SGM identity, Eq.~\eqref{eq:sgm_as_dot_p}.

For this Hamiltonian Eq.~\eqref{eq:sgm} assigns to each point $\mathbf{x}$ the force $-\nabla V(\mathbf{x})$, and by Eq.~\eqref{eq:sgm_expectation} the force on a particle in a given state is the average of that force over the state, weighted by $|\psi|^2$. In a bound stationary state every expectation value is constant in time, so $\langle\hat{\mathbf{p}}\rangle$ is constant and Eq.~\eqref{eq:sgm_expectation} gives $\langle\nabla\hat{\omega}\rangle=0$ exactly, while $\nabla V$ stays nonzero across the region the state occupies. A classical bound orbit that neither decays nor precesses returns to the same momentum once per revolution, so the force averaged over the period is zero.

This route covers the static case only, a static profile of $\alpha$ or a static potential. The Schr\"odinger equation assumes $c\to\infty$, which sets the light-flight time across the state to zero, contrary to Axiom~\ref{ax:causal_propagation}. The route gives the rate of change of the momentum of a state accelerated by an applied potential. It does not give the frequency gradient that the acceleration itself produces across the state, which is the rest frequency times the gradient of the time-dilation factor of the accelerating frame.

%-------------------------------------------------------------------------------
\subsection{The single-frequency assumption}
\label{sec:single_frequency}
%-------------------------------------------------------------------------------

The identity follows from relations that were all in place by 1926. Route~I uses only the Planck--Einstein relation (Eq.~\eqref{eq:planck}) and Hamilton's canonical equations (Axiom~\ref{ax:hamilton}). Route~II is the time derivative of the de~Broglie relation (Eq.~\eqref{eq:debroglie}), with the dispersion relation of Eq.~\eqref{eq:H_eq_hw}, where the quantum potential $Q$ is negligible. On a smooth phase field the time and space derivatives commute (Eq.~\eqref{eq:mixed_orders}), and the symmetry of the Hessian of the phase cancels the two terms that contain the particle velocity (Eq.~\eqref{eq:route2_components}). In geometrical optics the same steps give the eikonal ray equation (Eq.~\eqref{eq:eikonal_ray}).

A stationary state has the phase rate $E/\hbar$ at every point (the paragraph after Definition~\ref{def:spectral_gradient}), so the gradient of its frequency field is zero, and by Eq.~\eqref{eq:dp_dt_local} its momentum field does not change at any fixed point. The time derivative of the de~Broglie relation at a fixed point therefore shows no force. In a stationary state in which $Q$ is negligible, the spectral gradient equals minus the convective term at every point (Section~\ref{sec:route_phase}), and the force appears only along the trajectory.

Quantum mechanics also takes the rest term of the frequency to be $\omega_0$ at every point of the wave function's extent, so that $\alpha=1$ everywhere as in an inertial frame of flat spacetime, and in the Schr\"odinger equation a uniform gravitational field is the potential $mgz$ (Section~\ref{sec:intro}). This is the single-frequency assumption of Section~\ref{sec:intro}, and it sets the spectral gradient of the rest term to zero at every point:
\begin{equation}
    \nabla(\omega_0\,\alpha) = \mathbf{0} .
    \label{eq:single_frequency}
\end{equation}
No postulate of quantum mechanics states Eq.~\eqref{eq:single_frequency}. It enters when the rest energy is omitted as a phase factor (Section~\ref{sec:intro}), since that factor is the same at every point only where Eq.~\eqref{eq:single_frequency} holds. \citet{Dirac1947Principles} states that the constant term $mc^2$ ``can still have no physical effects, since the Hamiltonian in the quantum theory \ldots\ is undefined to the extent of an arbitrary additive real constant'', and \citet{Feynman1965Lectures3}, obtaining Newton's second law from the phase that a slowly varying potential adds, let the internal energy remain ``the same in both regions''. Separating the phase of a stationary state with one constant $E$ does not depend on the assumption, since a stationary metric also admits modes of one coordinate frequency while $\alpha$ varies (Section~\ref{sec:intro}).

Under Eq.~\eqref{eq:single_frequency} the rest term contributes nothing to the spectral gradient. For $\hat{H}=\hat{\mathbf{p}}^2/2m+V$ the spectral gradient is $\nabla V/\hbar$, and the expectation-value form of the identity, Eq.~\eqref{eq:sgm_expectation}, is Ehrenfest's theorem (Eq.~\eqref{eq:ehrenfest}). Within quantum mechanics the identity therefore gives no force that the Hamiltonian does not already give. This is why a relation that follows in a few lines from century-old physics has been applied to the forces of the potentials, and to the rest term only in particular settings, but has not been used as a general law of force. A standard text states that, while the de~Broglie relations connect wavelength and momentum, ``no such simple relations exist between the wave properties of matter and force'' \citep{Bohm1951Quantum}. The derivations of Newton's law from the Compton frequency cited in Section~\ref{sec:intro} set the assumption aside for the static field of a mass, \citet{HehlNi1990Inertial} for an accelerated, rotating frame, and \citet{Perche2022Localized} for a localized system on a timelike trajectory in curved spacetime. Axiom~\ref{ax:nonzero_spectral_gradient} replaces Eq.~\eqref{eq:single_frequency} in a static field, as the time dilation of special relativity does in an accelerated frame (Section~\ref{sec:intro}). In both cases the rate of proper time per unit coordinate time varies with position, so the time-dilated rest frequency has in general a nonzero spectral gradient and contributes $-\hbar\nabla(\omega_0\,\alpha)$ to the force.

%-------------------------------------------------------------------------------
\subsection{A universal force law}
\label{sec:universal_law}
%-------------------------------------------------------------------------------

Each route starts from an equation of standard mechanics (Eqs.~\eqref{eq:force_energy}, \eqref{eq:debroglie}, \eqref{eq:euler_lagrange}, and~\eqref{eq:ehrenfest_commutator}) and holds for any interaction that provides what the route requires. Route~I requires an energy that depends on position; Route~II, a smooth phase with negligible quantum potential $Q$; Route~III, a Lagrangian with a Legendre transform; and Route~IV, a Hamiltonian and a normalizable state on which the operators and their products act with no boundary term, with $\hat{H}=\hat{\mathbf{p}}^2/2m+V$ for the reduction to $-\nabla V$. Every interaction that quantum mechanics describes enters through the Hamiltonian, and the Hamiltonian divided by $\hbar$ is the frequency function (Eq.~\eqref{eq:H_eq_hw}) or, as an operator, the frequency operator $\hat{\omega}$ (Section~\ref{sec:route_schrodinger}), the identification that all four routes use.

External classical fields, the gravitational field among them, enter the Hamiltonian as position-dependent coefficients. A state that follows one level adiabatically, in an eigenbasis that does not vary in space, experiences the force $-\hbar\nabla\omega$, where $\omega$ is the frequency of that level, and a state on several levels experiences the force of each level weighted by its probability. An eigenbasis that turns in space adds a Berry connection to the projected dynamics, and the force then includes a term from its curvature \citep{Berry1984Quantal,Sundaram1999Wave}.

When the interaction depends on position alone and the kinetic term does not depend explicitly on time, the rate of the canonical momentum is the observable force. The reaction of a constraint, such as the force that holds a particle on a circle, is included when the constraint is part of the Hamiltonian, as the limit of a stiff confining potential. For an ideal wall, on which $\psi=0$, it enters instead as a boundary term that the domain conditions of Route~IV exclude. Computing $d\langle\hat{\mathbf p}\rangle/dt$ from the Schr\"odinger equation and integrating by parts over the region inside the walls leaves a force directed away from the wall and equal to $\hbar^2/2m$ times the squared magnitude of the derivative of $\psi$ normal to the wall, integrated over the wall. In a rectangular box with side $L_\text{box}$ along $x$, the normalized eigenstate with $n$ half-wavelengths along $x$ has energy $\pi^2\hbar^2n^2/(2mL_\text{box}^2)$ plus terms that do not depend on $L_\text{box}$, and at the wall $x=L_\text{box}$ this force has magnitude $\pi^2\hbar^2n^2/(mL_\text{box}^3)$, which is minus the rate at which the energy changes as that wall moves outward.

\begin{theorem}[Universality of the force-frequency law]\label{thm:universal}
Any interaction described by a Hermitian quantum Hamiltonian $\hat{H}=\hbar\hat{\omega}$ in the unweighted inner product drives the canonical momentum at the rate $\dot{\mathbf{p}}=-\hbar\nabla\omega$ in expectation values (Eq.~\eqref{eq:sgm_expectation}) on every normalizable state $\psi$ such that $\psi,\hat{\mathbf{p}}\psi\in\operatorname{dom}(\hat{H})$ and $\hat{H}\psi\in\operatorname{dom}(\hat{\mathbf{p}})$ with vanishing boundary terms; when a spin coupling makes the Hamiltonian a matrix, $\nabla\hat{\omega}$ is the matrix gradient. For interactions that depend on position alone, when the kinetic term depends on momentum alone (with no explicit time dependence, position-dependent metric, or position-dependent effective mass), that rate is the observable force.
\end{theorem}

Each route takes the evolution to be generated by a Hamiltonian. The reduced evolution of an open system is not generated by one \citep{Lindblad1976Generators}, and a Lindblad term in it can change $\langle\hat{\mathbf{p}}\rangle$ without changing $\omega(\mathbf{x},\mathbf{p})$. The force of the environment is then the spectral gradient of the closed system of particle and environment, by Route~IV applied to that system, with the gradient taken in the particle's coordinates and the coordinates of the environment held fixed.

%-------------------------------------------------------------------------------
\subsection{The frequency field about the particle}
\label{sec:taylor_expansion}
%-------------------------------------------------------------------------------

At fixed momentum (Eq.~\eqref{eq:gradient_def}), the frequency field at the displacement $\delta\mathbf{x}$ from the particle has the Taylor series
\begin{equation}
    \omega(\mathbf{x}+\delta\mathbf{x},\mathbf{p}) = \omega + \nabla\omega\cdot\delta\mathbf{x} + \tfrac{1}{2}\,(\partial_i\partial_j\omega)\,\delta x_i\,\delta x_j + \cdots ,
    \label{eq:omega_expansion}
\end{equation}
where $\partial_i$ is the derivative with respect to the position component $x_i$ at fixed momentum, with every derivative evaluated at the particle, the indices $i,j$ running over the three spatial axes, and repeated indices summed. The first-order coefficient is the spectral gradient of Eq.~\eqref{eq:sgm}, and the second-order coefficient is half the Hessian $\partial_i\partial_j\omega$. Since $F_j=-\hbar\,\partial_j\omega$, where $F_j$ is the component of $\mathbf{F}$ along the $j$th axis, the force Jacobian is $\partial_i F_j=-\hbar\,\partial_i\partial_j\omega$. It gives the variation of the force across the spatial extent of the wave function.

%===============================================================================
\section{Gravity and inertia}
\label{sec:grav_inertia}
%===============================================================================

%-------------------------------------------------------------------------------
\subsection{Gravity from the static gradient}
\label{sec:gravity}
%-------------------------------------------------------------------------------

A particle of mass $m$ at rest oscillates, per unit of its proper time, at its rest frequency $\omega_0=mc^2/\hbar$ (Eq.~\eqref{eq:rest_freq}). This oscillation is de~Broglie's internal clock at the Compton angular frequency \citep{Compton1923Quantum,deBroglie1925Thesis}, the periodic phenomenon from which he derived the matter wave. Proper-time differences can be read interferometrically with internal clock states \citep{Zych2011Quantum}; interpreting an atom interferometer as a clock at the Compton frequency is disputed \citep{Lan2013Clock,Wolf2011,SinhaSamuel2011}.

In a stationary spacetime a clock held at fixed spatial coordinates $\mathbf{x}$ measures a proper time $\tau$ that advances relative to the coordinate time $t$ at the rate \citep{Einstein1915GR,MTW1973Gravitation}
\begin{align}
    \alpha(\mathbf{x}) &\equiv \frac{d\tau}{dt} \notag\\
                       &= \sqrt{-g_{00}(\mathbf{x})} ,
    \label{eq:alpha_def}
\end{align}
where $g_{00}$ is the time-time component of the metric, negative in the signature $(-,+,+,+)$, with time coordinate $x^0=ct$. The \emph{time-dilation factor} $\alpha$ is dimensionless, equal to $1$ far from all masses and less than $1$ near a mass. In a static metric $\alpha$ equals the lapse of the $3+1$ split of general relativity; where the metric has cross terms $g_{0i}$ the two differ, and $\alpha$ always denotes $\sqrt{-g_{00}}$. Here $\alpha(\mathbf{x})$ is a smooth classical field, as in general relativity.

In the particle's local rest frame the time variable is $\tau$ and the energy is the rest energy $mc^2$, so Eqs.~\eqref{eq:planck} and~\eqref{eq:phase_freq} evaluated in that frame give the proper rate $-\partial_\tau\theta=\omega_0$. Local position invariance (Axiom~\ref{ax:lpi}) makes that rest-frame rate $\omega_0$ at every location, so $\theta=-\omega_0\,\tau$ up to a constant. For a static particle $d\tau=\alpha\,dt$ (Eq.~\eqref{eq:alpha_def}), so the phase advances in coordinate time as $\theta=-\omega_0\,\alpha\,t$, and the frequency $\omega=-\partial_t\theta$ is the redshifted value
\begin{equation}
    \begin{gathered}
    \omega(\mathbf{x}) = \omega_0\,\alpha(\mathbf{x}) , \\[-3pt]
    \text{\small\itshape (The time-dilated rest frequency)}
    \end{gathered}
    \label{eq:omega_alpha}
\end{equation}
which is the zeroth-order coefficient of Eq.~\eqref{eq:omega_expansion} and the value $\omega(\mathbf{x},\mathbf{0})$ of the frequency function at zero momentum. In a stationary state the phase rate is the uniform $E/\hbar$ (Section~\ref{sec:identity}), while the frequency function at fixed momentum varies with position as $\omega_0\,\alpha(\mathbf{x})$ does.

By $E=\hbar\omega$ (Eq.~\eqref{eq:planck}) and $\hbar\,\omega_0=mc^2$ (Eq.~\eqref{eq:rest_freq}), the coordinate frequency of Eq.~\eqref{eq:omega_alpha} gives the energy at $\mathbf{x}$ of a particle at rest there, conserved because the metric is static,
\begin{align}
    E(\mathbf{x}) &= \hbar \, \omega_0 \, \alpha(\mathbf{x}) \notag\\
    &= mc^2 \, \alpha(\mathbf{x}) .
    \label{eq:energy_alpha_massive}
\end{align}

To first order in the gravitational potential, Eq.~\eqref{eq:omega_alpha} gives different coordinate rates to the rest-energy phases of two narrow, separated wave packets of one wave function at different heights, such as the two arms of an atom interferometer, so the packets accumulate a relative phase \citep{Asenbaum2017Phase}. The motion of the packets and the beam splitters add further terms to it.

Clocks at rest at $\mathbf{x}\pm\ell\hat{\mathbf{u}}$, where $\hat{\mathbf{u}}$ is a unit vector and $\ell>0$, tick at $\omega_0\,\alpha(\mathbf{x}\pm\ell\hat{\mathbf{u}})$ in the common coordinate time in which their phases are compared, so
\begin{align}
    \hat{\mathbf{u}}\cdot\nabla\omega(\mathbf{x},\mathbf{0}) &= \lim_{\ell\to0}\frac{\omega(\mathbf{x}+\ell\hat{\mathbf{u}},\mathbf{0})-\omega(\mathbf{x}-\ell\hat{\mathbf{u}},\mathbf{0})}{2\ell} \notag\\
    &= \omega_0\,\hat{\mathbf{u}}\cdot\nabla\alpha(\mathbf{x}) ,
    \label{eq:gradient_span}
\end{align}
nonzero wherever $\alpha$ varies, although the proper frequency $\omega_0$ is the same at both ends of every span. The gradient $\nabla\omega=\omega_0\,\nabla\alpha$ is the \emph{static spectral gradient}.

For any profile of $\alpha$, substituting Eqs.~\eqref{eq:rest_freq} and~\eqref{eq:omega_alpha} into Eq.~\eqref{eq:sgm} gives the gravitational force on a massive particle at zero momentum in a static metric exactly,
\begin{equation}
    \mathbf{F}_\text{grav} = -mc^2\,\nabla\alpha .
    \label{eq:geometric_force}
\end{equation}
For a particle of speed $v\ll c$ the same force holds to zeroth order in $v^2/c^2$. Since $mc^2$ is constant, Eq.~\eqref{eq:geometric_force} is $\mathbf{F}_\text{grav}=-\nabla E$ (Eq.~\eqref{eq:force_energy}) for the energy of Eq.~\eqref{eq:energy_alpha_massive}, with the gradient taken at fixed momentum (Eq.~\eqref{eq:gradient_def}).

Equation~\eqref{eq:geometric_force} needs a field law that gives $\alpha$ from the masses. The field law taken here is Gauss's law for gravity \citep{Gauss1813Theoria,MTW1973Gravitation}, written for the time-dilation factor. Let $\partial\mathcal{V}$ be any closed surface, the boundary of a region $\mathcal{V}$, with outward area element $d\boldsymbol{\mathcal{S}}$ of magnitude $d\mathcal{S}$, and let $M_\text{enc}$ be the mass enclosed. Then
\begin{equation}
    \oint_{\partial\mathcal{V}}\nabla\alpha\cdot d\boldsymbol{\mathcal{S}} = \frac{4\pi G \, M_\text{enc}}{c^2} ,
    \label{eq:gauss_alpha}
\end{equation}
where $G$ is the gravitational constant and the dot is the scalar product of the gradient with the area element. Equation~\eqref{eq:gauss_alpha} is the premise of this subsection. Its source-free part states that the flux of $\nabla\alpha$ through any surface that encloses no mass is zero, a conservation law that contains neither $G$ nor any mass. Its source part states that the flux out of a mass is proportional to that mass, with the coupling $4\pi G/c^2$, whose value is measured. Comparisons of clocks test the source-free part directly; a test of the source part also needs the enclosed mass. Neither part contains a force. In the weak field of slowly moving matter, with the gradient taken in the coordinates $\mathbf{x}$, Eq.~\eqref{eq:gauss_alpha} is the integral form of the general-relativistic equation that fixes $\alpha$ in a static metric from the energy density and stress of the matter \citep{Gourgoulhon2007Formalism}.

For a spherically symmetric mass $M$, $\alpha$ depends only on the distance $r$ from the center, so on a sphere of radius $r$ the integrand of Eq.~\eqref{eq:gauss_alpha} is the constant $d\alpha/dr$ times $d\mathcal{S}$, and the sphere encloses $M_\text{enc}=M$:
\begin{equation}
    \frac{d\alpha}{dr} \, 4\pi r^2 = \frac{4\pi G M}{c^2} ,
    \label{eq:gauss_sphere}
\end{equation}
and dividing by $4\pi r^2$ gives
\begin{equation}
    \frac{d\alpha}{dr} = \frac{GM}{c^2 \, r^2} .
    \label{eq:dalpha_dr}
\end{equation}
The slope falls as the inverse square of the distance because a conserved flux is spread over a growing spherical area.

Coordinate time is normalized to the proper time of a distant static observer, so $\alpha(\infty)=1$. Integrating Eq.~\eqref{eq:dalpha_dr} over the radius $r'$ from $r$ to infinity gives the change in $\alpha$ along the radial path,
\begin{align}
    \alpha(\infty) - \alpha(r) &= \int_r^\infty \frac{GM}{c^2 \, r'^2} \, dr' \notag\\
    &= \frac{GM}{c^2 \, r} .
    \label{eq:alpha_deficit}
\end{align}
This dimensionless deficit is the fraction by which a clock at radius $r$ runs slow against the distant one; it is positive because $\alpha$ rises outward. With $\alpha(\infty)=1$, Eq.~\eqref{eq:alpha_deficit} gives the weak-field time-dilation factor
\begin{equation}
    \alpha(r) \approx 1 - \frac{GM}{c^2 \, r} ,
    \label{eq:alpha_weak}
\end{equation}
which is the first-order expansion of the Schwarzschild solution of Einstein's field equations \citep{Einstein1915GR,Schwarzschild1916Metric,MTW1973Gravitation}, and clocks measure its gradient directly \citep{Pound1960Apparent,Bothwell2022Resolving}. The integration in Eq.~\eqref{eq:alpha_deficit} is exact, so the approximation in Eq.~\eqref{eq:alpha_weak} comes only from the weak-field form of the premise, Eq.~\eqref{eq:gauss_alpha}.

At Earth's surface, where $GM/r^2$ is the surface gravity $g$, Eq.~\eqref{eq:dalpha_dr} gives $d\alpha/dr=g/c^2\approx1.1\times10^{-16}$ per meter, so two clocks a millimeter apart differ in rate by about one part in $10^{19}$, the shift resolved across a single atomic sample \citep{Bothwell2022Resolving}. Across the reduced Compton wavelength $\hbar/(mc)$ of an electron, of mass $m$, the factor $\alpha$ changes by the fraction $g\hbar/(mc^3)\approx4\times10^{-29}$.

Equations~\eqref{eq:omega_alpha} and~\eqref{eq:alpha_weak} give the coordinate frequency of a test particle of mass $m$ as $\omega(r)\approx\omega_0\,(1-GM/c^2r)$, whose gradient, with $\omega_0=mc^2/\hbar$ (Eq.~\eqref{eq:rest_freq}), is
\begin{align}
    \nabla\omega &= \omega_0\,\frac{GM}{c^2 \, r^2}\,\hat{\mathbf{r}} \notag\\
    &= \frac{mGM}{\hbar r^2}\,\hat{\mathbf{r}} ,
    \label{eq:grav_gradient}
\end{align}
where $\hat{\mathbf{r}}$ is the outward unit radial vector. Substituting Eq.~\eqref{eq:grav_gradient} into the identity $\mathbf{F}=-\hbar\,\nabla\omega$ (Eq.~\eqref{eq:sgm}) gives the gravitational force on the particle in the weak field outside a spherical mass $M$ \citep{Newton1687Philosophiae},
\begin{equation}
    \begin{gathered}
    \boxed{\;\mathbf{F}_\text{grav} = -\frac{GMm}{r^2}\,\hat{\mathbf{r}}\;} \, . \\[-3pt]
    \text{\small\itshape (Newton's law of universal gravitation)}
    \end{gathered}
    \label{eq:newton}
\end{equation}
The frequency falls toward the mass, so its gradient points outward and the force points toward the mass.

Dividing Eq.~\eqref{eq:geometric_force} by $m$ gives the gravitational acceleration $\mathbf{g}=\mathbf{F}_\text{grav}/m$ of a slow particle, the same for every mass:
\begin{equation}
    \mathbf{g} = -c^2\,\nabla\alpha .
    \label{eq:g_grad_alpha}
\end{equation}
In the Newtonian limit the metric component is $g_{00}=-(1+2\Phi_g/c^2)$, where $\Phi_g$ is the Newtonian gravitational potential \citep{MTW1973Gravitation}. By Eq.~\eqref{eq:alpha_def}, $\alpha=\sqrt{1+2\Phi_g/c^2}$, which to first order in $\Phi_g/c^2$ is $\alpha=1+\Phi_g/c^2$. Equation~\eqref{eq:g_grad_alpha} is then $\mathbf{g}=-\nabla\Phi_g$, and Eq.~\eqref{eq:gauss_alpha} is Gauss's law for gravity in its usual form, $\oint_{\partial\mathcal{V}}\mathbf{g}\cdot d\boldsymbol{\mathcal{S}}=-4\pi G \, M_\text{enc}$.

Equation~\eqref{eq:alpha_weak} also fixes the energy of Eq.~\eqref{eq:energy_alpha_massive} at radius $r$ as $E(r)\approx mc^2-GMm/r$, whose second term is the Newtonian gravitational potential energy $m \, \Phi_g$. That energy is the amount by which time dilation lowers the particle's rest energy at $r$ as measured by a distant observer, while the local rest energy stays $mc^2$.

In Eqs.~\eqref{eq:gauss_alpha} to~\eqref{eq:newton} the field law gives $\alpha$, and the force is $-\hbar$ times the gradient of the frequency $\omega_0 \, \alpha$ (Eqs.~\eqref{eq:omega_alpha} and~\eqref{eq:sgm}); no law of gravitational force enters. Gauss's law was first obtained from Newton's law, and outside a spherical source Eq.~\eqref{eq:gauss_alpha} is equivalent to Eq.~\eqref{eq:dalpha_dr}. But Eq.~\eqref{eq:dalpha_dr} is a law for the gradient of the rate of clocks. Its inverse square comes from the conservation of flux in three dimensions, and its one constant is the flux $4\pi GM/c^2$ out of the mass, given by the source part of Eq.~\eqref{eq:gauss_alpha}. The identity then gives the force the same inverse-square law and the factor $m$ through $\omega_0=mc^2/\hbar$. The coupling $G$ is measured, as the vacuum permittivity is in Coulomb's law, which follows in the same way from Gauss's law for the electric field $\mathbf{E}$ and the force $q\mathbf{E}$.

Outside a spherically symmetric mass, with the cosmological term zero, the vacuum field equations of general relativity, which contain neither $G$ nor any source strength, give the exterior Schwarzschild metric \citep{Schwarzschild1916Metric,MTW1973Gravitation}. Its time-dilation factor is $\alpha=\sqrt{1-R_s/r}$, normalized to $\alpha(\infty)=1$. The one constant $R_s$ is the Schwarzschild radius, and $r$ is the areal radius: the sphere through the point has area $4\pi r^2$. A local ruler measures the proper radial distance $r_\text{p}$, which satisfies $dr_\text{p}=dr/\alpha$, and $\alpha^2=1-R_s/r$ gives $2\alpha\,d\alpha/dr=R_s/r^2$, so the gradient of $\alpha$ along that ruler is
\begin{align}
    \frac{d\alpha}{dr_\text{p}} &= \alpha \, \frac{d\alpha}{dr} \notag\\
    &= \frac{R_s}{2r^2} ,
    \label{eq:dalpha_dl}
\end{align}
and its flux through the sphere of area $4\pi r^2$ is
\begin{equation}
    \frac{d\alpha}{dr_\text{p}} \, 4\pi r^2 = 2\pi R_s
    \label{eq:gauss_exact}
\end{equation}
at every radius of the exterior geometry and at every field strength. Equation~\eqref{eq:gauss_exact} is the exact exterior Komar--Gauss flux law $\oint_{S^2(r)}(\nabla_\text{prop}\alpha)\cdot d\boldsymbol{\mathcal{S}} = 2\pi R_s = 4\pi GM/c^2$, that is, Eq.~\eqref{eq:gauss_alpha} with the gradient taken along proper distance and $2\pi R_s$ in place of $4\pi GM/c^2$. If $|\mathbf{x}|$ is the areal radius $r$, the only approximation in Eq.~\eqref{eq:gauss_alpha} is taking the gradient along $r$ instead of along proper distance.

A comparison of clocks measures $R_s$ \citep{Pound1960Apparent,Bothwell2022Resolving}. The ratio of the rates of two clocks at rest at distinct areal radii $r_1$ and $r_2$, both larger than $R_s$, is $\alpha(r_1)/\alpha(r_2)=[(1-R_s/r_1)/(1-R_s/r_2)]^{1/2}$, which is monotonic in $R_s$ and so gives one value of $R_s$. Far from the mass, where $\alpha\approx1$, Eq.~\eqref{eq:geometric_force} gives the force $-(mc^2R_s/2r^2)\,\hat{\mathbf{r}}$ on a particle at rest, which is Eq.~\eqref{eq:newton} with $GM=R_sc^2/2$, so Eq.~\eqref{eq:geometric_force} relates two quantities observed independently of each other. The force on a slow massive particle depends on the metric only through $g_{00}$, the component those clocks measure; the spatial components enter at order $v^2/c^2$, in the deflection of light rather than in Newton's law.

Close to the mass Eq.~\eqref{eq:geometric_force} needs no approximation and exhibits a three-force strong-field triad. With $\alpha=\sqrt{1-R_s/r}$ and $GM=R_sc^2/2$, the coordinate-gradient force $\mathbf{F}_\text{SGM}\equiv -mc^2\nabla_r\alpha = -(GMm/(r^2\alpha))\,\hat{\mathbf{r}}$ of Eq.~\eqref{eq:geometric_force}, the local proper weight $\mathbf{F}_\text{prop}\equiv -mc^2\nabla_\text{prop}\ln\alpha = -(GMm/(r^2\alpha))\,\hat{\mathbf{r}}$ read by an accelerometer held at rest~\citep{Wald1984General}, and the force $\mathbf{F}_\infty\equiv\alpha\,\mathbf{F}_\text{prop}=-mc^2\nabla_\text{prop}\alpha = -(GMm/r^2)\,\hat{\mathbf{r}}$ transmitted to infinity by a massless rope~\citep{Wald1984General} satisfy $\mathbf{F}_\text{SGM}=\mathbf{F}_\text{prop}=\alpha^{-1}\mathbf{F}_\infty$ at every radius $r>R_s$. The equality $\mathbf{F}_\text{SGM}=\mathbf{F}_\text{prop}$ holds in Schwarzschild areal coordinates specifically because $dr_\text{p}=dr/\alpha$ (so $g_{00}g_{rr}=-1$ and $d\ln\alpha/dr_\text{p}=d\alpha/dr$ by Eq.~\eqref{eq:dalpha_dl}); on a general static spatial slice with spatial metric $h_{jk}$ (companion paper~\citep{HutchisonWavefn2026}), the proper weight satisfies $m\,a_\text{prop}=\alpha^{-1}\sqrt{h^{jk}F_j F_k}$ with $F_j=-mc^2\,\partial_j\alpha$.

\begin{corollary}[From metric to force]\label{cor:causal_gravity}
In a static metric, time dilation alone gives the gravitational force on a particle at rest. With Gauss's law for the time-dilation factor (Eq.~\eqref{eq:gauss_alpha}) as the field law, Eq.~\eqref{eq:geometric_force} gives Newton's law; with the vacuum field equations and the clock constant $R_s$ it gives the same law far outside a spherical mass, where $\alpha\approx1$. The inverse square follows from the conservation of flux and the proportionality to $m$ from the identity; the equality of $R_s$ with the value $2GM/c^2$ that orbits give is a prediction. The Newtonian gravitational field is $-c^2$ times the gradient of the rate of proper time, which clocks measure (Eq.~\eqref{eq:g_grad_alpha}), and its potential is $\Phi_g=c^2(\alpha-1)$ to first order.
\end{corollary}

General relativity gives the same acceleration from the geodesic equation, whose limit for a slow particle in a static weak field is $d^2\mathbf{x}/dt^2=-c^2\nabla\alpha$ \citep{MTW1973Gravitation}, so Corollary~\ref{cor:causal_gravity} agrees with it there. The identity gives the acceleration without the geodesic equation, from the gradient of the rest frequency of the particle's own wave function (Eq.~\eqref{eq:omega_alpha}).

%-------------------------------------------------------------------------------
\subsection{Inertia from the dynamic gradient}
\label{sec:inertia}
%-------------------------------------------------------------------------------

The acceleration of a particle is a dynamic source of a spectral gradient, in addition to the static one of Section~\ref{sec:gravity}. A free state in uniform motion has one frequency at every point in its rest frame, so this source can act only while the motion changes.

The wave function has nonzero spatial extent, centered on the particle's mean position $\mathbf{x}$. Throughout this section $\mathbf{a}$ is the proper acceleration of the state's mean position, $a=|\mathbf{a}|$ its magnitude, $\hat{\mathbf{a}}=\mathbf{a}/a$ its direction, and $\delta x_\parallel=\hat{\mathbf{a}}\cdot\delta\mathbf{x}$ the component of a displacement $\delta\mathbf{x}$ along that direction. Under a uniform applied force every point of the state has the proper acceleration $\mathbf{a}$.

Exact rigidity is Born rigidity \citep{Born1909Theorie,Moller1952Theory,Rindler1966Kruskal}, in which a point at proper distance $X$ from a common horizon has proper acceleration $c^2/X$. The mean position is at $X=c^2/a$, so the proper acceleration at the displacement $\delta\mathbf{x}$ is $a/(1+a \, \delta x_\parallel/c^2)$. The uniform and the Born-rigid profiles differ at first order in $a \, \delta x_\parallel/c^2$, which is $1.1\times10^{-16}$ per meter of $\delta x_\parallel$ at $a=g$. They agree at an instant when the state is at rest. Both then have zero velocity at every point, so the rate of proper time at each point depends only on the position of that point in the co-accelerating coordinates. The velocities that build up after that instant change the rate only at second order in $v/c$.

The Rindler metric gives the offset in the rate of time across an accelerating frame. In the co-accelerating coordinates $(t,\delta\mathbf{x})$, where $t$ is the proper time of the mean position, its line element is $ds^2=-(1+\mathbf{a}\cdot\delta\mathbf{x}/c^2)^2c^2\,dt^2+|d(\delta\mathbf{x})|^2$ \citep{Rindler1966Kruskal,MTW1973Gravitation}, so its time-dilation factor $\sqrt{-g_{00}}$ (Eq.~\eqref{eq:alpha_def}) at displacement $\delta\mathbf{x}$ is
\begin{equation}
    \alpha = 1 + \delta\alpha ,
    \label{eq:alpha_rindler}
\end{equation}
where the \emph{fractional frequency shift}
\begin{equation}
    \delta\alpha = \frac{\mathbf{a}\cdot\delta\mathbf{x}}{c^2}
    \label{eq:rapidity_gain}
\end{equation}
is the fraction by which the rate of proper time at $\delta\mathbf{x}$ exceeds that at the mean position. The factor is exact for a Born-rigid state at every separation with $\mathbf{a}\cdot\delta\mathbf{x}=a \, \delta x_\parallel>-c^2$ \citep{Moller1952Theory,Rindler1966Kruskal,MTW1973Gravitation}; \citet{Einstein1907Relativitatsprinzip} obtained it to first order in $a\,\delta x_\parallel/c^2$, and \citet{Einstein1911Influence} obtained it again as the Doppler shift \citep{Doppler1842Farbige} of light exchanged between the ends of an accelerated box. By Axiom~\ref{ax:lpi} the intrinsic part of the phase, the part due to the rest energy, changes by $-\omega_0$ per unit of proper time at every point of the state. By Eq.~\eqref{eq:alpha_rindler} that proper time advances at $\alpha$ per unit of $t$, so in the co-accelerating coordinates the intrinsic part of the phase at $\mathbf{x}+\delta\mathbf{x}$ has the frequency
\begin{equation}
    \omega(\mathbf{x}+\delta\mathbf{x}) = \omega_0 \, (1+\delta\alpha) .
    \label{eq:rindler_rate}
\end{equation}
Its departure from the rest rate, $\omega_0 \, \delta\alpha$, depends on the force only through the acceleration. Equation~\eqref{eq:rindler_rate} is $\omega=\omega_0\,\alpha$ of Eq.~\eqref{eq:omega_alpha} with the Rindler factor in place of the static one.

Differencing this rate (Eq.~\eqref{eq:rindler_rate}) across the opposite displacements $\pm\delta\mathbf{x}$ gives, per unit separation along $\hat{\mathbf{d}}=\delta\mathbf{x}/|\delta\mathbf{x}|$,
\begin{equation}
    \frac{\omega(\mathbf{x}+\delta\mathbf{x})-\omega(\mathbf{x}-\delta\mathbf{x})}{2\,|\delta\mathbf{x}|} = \frac{\omega_0\,\mathbf{a}\cdot\hat{\mathbf{d}}}{c^2} \, .
    \label{eq:front_behind}
\end{equation}
In the expansion of Eq.~\eqref{eq:omega_expansion} the symmetric difference cancels every even-order term, including the rest frequency $\omega_0=\omega(\mathbf{x})$ and the second-order tidal term. After division by $2|\delta\mathbf{x}|$ the odd terms beyond the first are at most of order $|\delta\mathbf{x}|^2$, so the quotient converges to the directional derivative $\nabla\omega\cdot\hat{\mathbf{d}}$ of the rate as $|\delta\mathbf{x}|\to0$. That rate is a function of displacement alone, so its ordinary gradient is the fixed-momentum spectral gradient of Definition~\ref{def:spectral_gradient}, here in the particle's accelerating rest frame, where the Rindler factor is the time-dilation factor. The Rindler factor (Eq.~\eqref{eq:alpha_rindler}) is linear in $\delta\mathbf{x}$, so the quotient already has its limiting value at every separation out to the distance $c^2/a$ behind the particle at which the factor reaches zero, and the result is independent of the size of the wave function on that scale.

Equation~\eqref{eq:front_behind} holds for every direction $\hat{\mathbf{d}}$, so the component of $\nabla\omega$ along each is $(\omega_0\,\mathbf{a}/c^2)\cdot\hat{\mathbf{d}}$. With $\omega_0=mc^2/\hbar$ (Eq.~\eqref{eq:rest_freq}) the gradient is the \emph{dynamic spectral gradient}
\begin{align}
    \nabla\omega &= \frac{\omega_0\,\mathbf{a}}{c^2} \notag\\
    &= \frac{m\,\mathbf{a}}{\hbar} ,
    \label{eq:inertia_gradient}
\end{align}
and substituting it into $\mathbf{F}=-\hbar\,\nabla\omega$ (Eq.~\eqref{eq:sgm}) gives the inertial reaction $\mathbf{F}_\text{in}$,
\begin{equation}
    \begin{gathered}
    \boxed{\;\mathbf{F}_\text{in} = -m\,\mathbf{a}\;} \, . \\[-3pt]
    \text{\small\itshape (The inertial reaction)}
    \end{gathered}
    \label{eq:inertia}
\end{equation}
The proportionality to $\mathbf{a}$ comes from the Rindler factor, which is linear in $\mathbf{a}$. The particle enters only through its rest rate $\omega_0$, and no kinetic term appears, so the coefficient is the rest mass $m=\hbar \, \omega_0/c^2$, the inertia of the rest energy \citep{Einstein1905Inertia}. The derivation uses only the proper rest rate $\omega_0$ at every point of the state (Axiom~\ref{ax:lpi}) and the Rindler factor, which follows from special relativity for worldlines in flat spacetime. It takes the acceleration as input and gives the force, with no gravitational field and no law for the applied force.

The gradient of the Rindler factor (Eq.~\eqref{eq:alpha_rindler}) with respect to the displacement $\delta\mathbf{x}$ is $\nabla\alpha=\nabla(1+\mathbf{a}\cdot\delta\mathbf{x}/c^2)=\mathbf{a}/c^2$. Substituting $\mathbf{a}=c^2\,\nabla\alpha$ into Eq.~\eqref{eq:inertia} gives
\begin{equation}
    \mathbf{F}_\text{in} = -mc^2\,\nabla\alpha .
    \label{eq:inertia_geometric}
\end{equation}

Let a uniform applied force $\mathbf{F}$ derive from the potential $V=-\mathbf{F}\cdot\mathbf{x}$. By Eq.~\eqref{eq:planck} a potential shifts the coordinate rate by $V/\hbar$, so in the inertial frame in which the particle is momentarily stationary the phase at the displacement $\delta\mathbf{x}$ has the frequency, written with the subscript st for stationary,
\begin{equation}
    \omega_\text{st}(\mathbf{x}+\delta\mathbf{x}) = \omega_0 \left(1-\frac{\mathbf{F}\cdot\delta\mathbf{x}}{mc^2}\right) ,
    \label{eq:detuned_rate}
\end{equation}
exact because the potential is linear. Because Eq.~\eqref{eq:detuned_rate} is expressed in the momentarily comoving inertial time $t_\text{in}$, transforming the phase-rate covector $\omega=-\partial_t\theta$ to the Rindler coordinate time $t_\text{R}$ multiplies the entire inertial frequency (rest plus potential) by $\partial t_\text{in}/\partial t_\text{R}=1+\delta\alpha$ (which equals $1$ at $\delta\mathbf{x}=\mathbf{0}$ and differs from an unredshifted coordinate-frame potential only at $O(|\delta\mathbf{x}|^2)$), giving the product
\begin{equation}
    \omega_\text{acc}(\mathbf{x}+\delta\mathbf{x}) = \omega_0 \left(1-\frac{\mathbf{F}\cdot\delta\mathbf{x}}{mc^2}\right) (1+\delta\alpha) .
    \label{eq:flat_total}
\end{equation}
To first order in $\delta\mathbf{x}$ this rate is $\omega_0[1+(\mathbf{a}-\mathbf{F}/m)\cdot\delta\mathbf{x}/c^2]$, and by Eq.~\eqref{eq:sgm} its gradient exerts the force $\mathbf{F}-m\mathbf{a}$. In its accelerating rest frame the particle's momentum stays zero, so this force vanishes:
\begin{equation}
    \mathbf{F} = m\,\mathbf{a} .
    \label{eq:second_law}
\end{equation}
This is d'Alembert's principle, the vanishing of the sum of the applied and inertial forces \citep{DAlembert1743Traite,Goldstein2002Classical}, stated as a condition on a frequency field: the shift $-\omega_0\delta\alpha$ that the applied potential gives the rate cancels the Rindler shift $+\omega_0\delta\alpha$, and the frequency field in the accelerating rest frame is uniform to first order in $\delta\mathbf{x}$. Of the two frequency fields, only the inertial-frame field of Eq.~\eqref{eq:detuned_rate} is the frequency function of Definition~\ref{def:spectral_gradient}, and that field generates the motion: the canonical momentum of Eq.~\eqref{eq:sgm_as_dot_p} changes at the rate $\dot{\mathbf{p}}=\mathbf{F}$. The Rindler shift enters only through the time of the accelerating frame.

The dynamic spectral gradient (Eq.~\eqref{eq:inertia_gradient}) comes from the relativity of simultaneity. In the inertial frame in which the state was momentarily at rest, once the state has gained a small velocity $\mathbf{v}$, a clock moving with it at the displacement $\delta\mathbf{x}$ reads, to first order in $v/c$, the time $t-\mathbf{v}\cdot\delta\mathbf{x}/c^2$, so the rest phase there is $-\omega_0\,(t-\mathbf{v}\cdot\delta\mathbf{x}/c^2)$. Its wavevector $\omega_0\,\mathbf{v}/c^2$ grows at the rate $\omega_0\,\mathbf{a}/c^2$, which is the dynamic spectral gradient.

This gradient is present at an instant when the state is at rest and its energy is not changing, so it is not the gradient at a front by which a change of frequency spreads across the state (Axiom~\ref{ax:causal_propagation}). A frequency that changes at the rate $\partial_t\omega$ and spreads at the speed $c$ along a unit vector $\hat{\mathbf{n}}$, $\omega(t-\hat{\mathbf{n}}\cdot\delta\mathbf{x}/c)$, has the gradient $-(\hat{\mathbf{n}}/c)\,\partial_t\omega$; at rest the force does work at the rate $\mathbf{F}\cdot\mathbf{v}=0$, so the energy does not change and the front gradient is zero.

%-------------------------------------------------------------------------------
\subsection{Gravity and inertia are one mechanism}
\label{sec:duality}
%-------------------------------------------------------------------------------

The gravitational force (Eq.~\eqref{eq:geometric_force}) and the inertial force (Eq.~\eqref{eq:inertia_geometric}) have one form, $\mathbf{F} = -mc^2\,\nabla\alpha$, and differ in the origin of the time-dilation factor, and through it in the second derivatives of $\alpha$, the tidal terms. For gravity, $\alpha$ is a static field over position, sourced by a gravitating mass. For inertia, $\alpha$ is the Rindler factor over displacement, produced by the particle's own acceleration (Eq.~\eqref{eq:alpha_rindler}), and its gradient is nonzero only while the motion changes (Section~\ref{sec:inertia}).

Both forces act through the multiplicative form $\omega=\omega_0 \, \alpha$, with the same factor $\omega_0=mc^2/\hbar$ (Eq.~\eqref{eq:rest_freq}) in Eqs.~\eqref{eq:grav_gradient} and~\eqref{eq:inertia_gradient}, so both are proportional to the mass $m=\hbar \, \omega_0/c^2$, the rest energy divided by $c^2$. The gravitational mass $m_g$ is the coefficient of $-c^2\nabla\alpha$ in Eq.~\eqref{eq:geometric_force}, and the inertial mass $m_i$ is the coefficient of $-\mathbf{a}$ in Eq.~\eqref{eq:inertia}. Hence $m_g=m_i$, because the metric multiplies the rate of every clock by the same factor $\alpha$ (Eq.~\eqref{eq:alpha_def}).

For a composite body the rest energy is the total internal energy, binding and kinetic included, and the body's overall frequency is that total over $\hbar$ (Axiom~\ref{ax:lpi}). Neither derivation used any property of the state beyond the rest rate $\omega_0$, so the coupling applies unchanged to the center of mass of a composite body small compared with the distance over which $\nabla\alpha$ varies (for a relativistic bound state, the virial balance between internal kinetic energy coupled through $\alpha_n W^{-2}$ and internal binding energy cancels the spatial-curvature parameter $\gamma_\text{PPN}$ from the total weight, preserving universal coupling to $\alpha$); tidal, multipole and spin forces on an extended body are separate.

The MICROSCOPE satellite measured the E\"otv\"os parameter of titanium and platinum, the fractional difference of their free-fall accelerations, as $[-1.5\pm2.3\,(\text{stat})\pm1.5\,(\text{syst})]\times10^{-15}$, which with the two uncertainties combined in quadrature lies 0.5 standard deviations from the value $0$ that $m_g=m_i$ gives \citep{Touboul2017Microscope,Touboul2022Final}. Neutron interferometry in accelerated motion tests $m_g=m_i$ at the level of the matter-wave phase \citep{Bonse1983Measurement}.

Because $\omega_0$ is positive for antiparticles as for particles, antihydrogen falls downward like hydrogen. The ALPHA-g experiment measured its downward acceleration as $(0.75\pm0.13\pm0.16)\,g$, where the first uncertainty is statistical and systematic and the second is from simulation; with the two combined in quadrature, the value lies 1.2 standard deviations from $g$ and 8.5 from the value $-g$ of repulsion \citep{Anderson2023Observation}.

The nonrelativistic theory does not show this mechanism, because its reduction removes both factors of the inertial gradient. Equation~\eqref{eq:inertia_gradient} is the product of $\omega_0=mc^2/\hbar$, which grows as $c^2$, and the fractional frequency shift per unit separation, $a/c^2$. The reduction removes the rest phase $e^{-imc^2t/\hbar}$, and with it $\omega_0$, and takes $c$ to infinity, which takes $\delta\alpha$ of Eq.~\eqref{eq:rapidity_gain} to zero. Their product $mc^2\,\delta\alpha=m\,\mathbf{a}\cdot\delta\mathbf{x}$ remains only as the inertial potential, whose coefficient $m$ is a parameter of the equation. Gravity likewise remains only as the potential $mgz$, whose coefficient is set equal to the inertial mass by postulate. The transformation $\psi(\mathbf{x},t)=\psi_{\mathbf{a}}(\tilde{\mathbf{x}},t)\exp[\tfrac{i}{\hbar}(m\,\mathbf{a}\cdot\tilde{\mathbf{x}}\,t+\tfrac16 m\,a^2t^3)]$ to the coordinates $\tilde{\mathbf{x}}=\mathbf{x}-\tfrac12\mathbf{a}t^2$ adds $+m\,\mathbf{a}\cdot\tilde{\mathbf{x}}$ exactly and removes a uniform field in the falling frame \citep{Shukla2010Gauge}. In neither potential is the coefficient the rest frequency times a time-dilation factor, so the two forces cannot be identified there as one mechanism.

The free Schr\"odinger kernel from a source point $\mathbf{x}_\text{s}$, $(m/2\pi i\hbar t)^{3/2}\exp[im|\mathbf{x}-\mathbf{x}_\text{s}|^2/(2\hbar t)]$, has the same modulus at every separation, so the equation also contains no light-flight time. The static force is unaffected, but the inertial derivation needs a finite $c$.

%-------------------------------------------------------------------------------
\subsection{Free fall and the equivalence principle}
\label{sec:freefall}
%-------------------------------------------------------------------------------

A particle in free fall has both gradients at once: the static gradient $\nabla\omega_\text{grav}=\omega_0 \, \nabla\alpha$ (Eq.~\eqref{eq:omega_alpha}) and a dynamic gradient $\nabla\omega_\text{in}$ of the form of Eq.~\eqref{eq:inertia_gradient}. Since the dynamic gradient depends on the force only through the acceleration (Section~\ref{sec:inertia}), it can be evaluated for a freely falling particle before asking what force acts.

Equation~\eqref{eq:inertia_gradient} was derived for a momentarily resting state, so it is evaluated here in coordinates that move with the particle. The static metric is $ds^2=-\alpha^2c^2\,dt^2+\gamma_{ij}\,dx^idx^j$, where $\gamma_{ij}(\mathbf{x})$ is its spatial part and $\gamma^{ij}$ its inverse. Let $\mathbf{x}_0(t)$ be the particle's position in the static coordinates and $\mathbf{a}_c=\ddot{\mathbf{x}}_0$ its coordinate acceleration, which is nonzero although the proper acceleration vanishes on the geodesic. Here $\mathbf{x}$ is again a general point, not the mean position of Section~\ref{sec:inertia}. Define coordinates $(\tilde{t},\delta\mathbf{x})$ that move with the particle by $\mathbf{x}=\mathbf{x}_0(\tilde{t})+\delta\mathbf{x}$ and
\begin{equation}
    t = \tilde{t} + \frac{\gamma_{ij}(\mathbf{x}_0)\,\dot{x}_0^{\,j}(\tilde{t})\,\delta x^i}{\alpha(\mathbf{x}_0)^2\,c^2} ,
    \label{eq:moving_sync}
\end{equation}
which makes each slice of constant $\tilde{t}$ orthogonal to the particle's worldline $\delta\mathbf{x}=\mathbf{0}$, on which $\tilde{t}=t$. At an instant when the coordinate velocity $\dot{\mathbf{x}}_0$ vanishes, the differentials are $d\mathbf{x}=d(\delta\mathbf{x})$ and $dt=(1+\gamma_{ij}a_c^{\,j}\delta x^i/(\alpha_0^2c^2))\,d\tilde{t}$, where $\alpha_0=\alpha(\mathbf{x}_0)$ and $\gamma_{ij}$ is evaluated at $\mathbf{x}_0$. Substituting them leaves the spatial part of the metric unchanged, adds no cross term, and multiplies its time-time component by the square of the factor in $dt$, so the time-dilation factor $\sqrt{-g_{00}}$ in the moving coordinates is
\begin{equation}
    \tilde{\alpha} = \alpha(\mathbf{x}_0+\delta\mathbf{x}) \, \Big(1+\frac{\gamma_{ij}\,a_c^{\,j}\,\delta x^i}{\alpha_0^2\,c^2}\Big) ,
    \label{eq:moving_frame_alpha}
\end{equation}
the static factor times a Rindler factor of the form of Eq.~\eqref{eq:alpha_rindler}. By the product rule its gradient with respect to $\delta x^i$ at the particle is $\partial_i\alpha+\gamma_{ij}\,a_c^{\,j}/(\alpha_0\,c^2)$. Multiplied by $\omega_0$, it is the sum of the static gradient $\nabla\omega_\text{grav}=\omega_0\,\nabla\alpha$ and a dynamic gradient $\nabla\omega_\text{in}$ with components $\omega_0\,\gamma_{ij}a_c^{\,j}/(\alpha_0c^2)$, which in flat spacetime ($\alpha_0=1$, $\gamma_{ij}=\delta_{ij}$) is Eq.~\eqref{eq:inertia_gradient} with $\mathbf{a}=\mathbf{a}_c$.

As for the applied force of Eq.~\eqref{eq:second_law}, the particle's momentum in its accelerating rest frame stays zero, so the sum of the two gradients vanishes,
\begin{equation}
    \nabla\omega_\text{grav}+\nabla\omega_\text{in} = \mathbf{0} ,
    \label{eq:geodesic_balance}
\end{equation}
which is d'Alembert's principle: on the realized path the gravitational force and the inertial reaction sum to zero \citep{DAlembert1743Traite,Lagrange1788Mecanique,Goldstein2002Classical}. Substituting the two gradients into Eq.~\eqref{eq:geodesic_balance}, multiplying by $\alpha_0c^2\gamma^{ki}/\omega_0$, and summing over $i$ gives the acceleration
\begin{equation}
    a_c^{\,k} = -c^2 \, \alpha_0 \, \gamma^{ki} \, \partial_i\alpha ,
    \label{eq:freefall_accel}
\end{equation}
exact at an instant of rest, and the coordinate acceleration of the geodesic there \citep{MTW1973Gravitation}. To first order in the potential, $\alpha_0\approx1$ and $\gamma^{ki}\approx\delta^{ki}$, so $\mathbf{a}_c=-c^2\nabla\alpha$, which the weak-field factor of Eq.~\eqref{eq:alpha_weak} evaluates to $-GM \, \hat{\mathbf{r}}/r^2$, the acceleration of Newtonian free fall. A particle held at rest against gravity has zero coordinate acceleration, so its dynamic gradient in these coordinates vanishes, and it feels the uncanceled static gradient as its weight (Eq.~\eqref{eq:newton}).

The cancellation by itself is not a prediction, since substituting $\mathbf{a}=\mathbf{F}/m$ into a reaction term $-m\mathbf{a}$ makes the sum vanish for every force law. The prediction is that the two coefficients, derived separately from the static metric and from the Rindler factor of the particle's own acceleration, are both equal to $\omega_0$. That shared value cancels out of Eq.~\eqref{eq:geodesic_balance}, so Eq.~\eqref{eq:freefall_accel} does not contain the mass: every body falls with the same acceleration. This is the weak equivalence principle, the consequence of $m_g=m_i$ that free-fall experiments test. With two unequal coefficients their ratio would remain in Eq.~\eqref{eq:freefall_accel}, and the fall would depend on the body. Metric theories of gravity postulate the universality of free fall: test bodies follow geodesics of one metric, to which all nongravitational fields couple in the same way \citep{Will2014Confrontation}. Here it follows because the common coefficient of the two frequency gradients is the proper rest frequency $\omega_0$, the same at every location by Axiom~\ref{ax:lpi}.

Equation~\eqref{eq:freefall_accel} sets the gradient of $\tilde{\alpha}$ (Eq.~\eqref{eq:moving_frame_alpha}) to zero, so the frequency field of a freely falling particle, in its own falling frame, has zero gradient at its position, which is the equivalence principle stated for the frequency field. The static gradient $\nabla\omega_\text{grav}$ is still present in Eq.~\eqref{eq:geodesic_balance}; only its sum with the inertial gradient $\nabla\omega_\text{in}$ vanishes. The Rindler factor belongs to the particle's accelerating rest frame (Section~\ref{sec:inertia}), so in the static coordinates the fixed-momentum gradient contains only the static term, and the canonical momentum evolves at the rate $-\hbar\nabla\omega$ of Eq.~\eqref{eq:sgm_as_dot_p}, which is Newton's law (Eq.~\eqref{eq:newton}). By Ehrenfest's theorem (Eq.~\eqref{eq:sgm_expectation}) the balance also holds for the force averaged over the state, so the mean position of a wave packet across which the force varies negligibly falls along the same worldline; across a larger packet the variation of the force, the tidal field, adds corrections of finite size.

General relativity obtains that worldline as force-free geodesic flow: the particle extremizes proper time in the given metric and feels no invariant force. Equation~\eqref{eq:freefall_accel} gives the same worldline from the cancellation of two spectral gradients. Both accounts use the same metric and give the same path, in the way that Newtonian gravity gives the weak-field trajectories with an inverse-square force.

\begin{remark}[Tidal curvature distinguishes gravity from inertia]\label{rem:rank}
The shared spectral gradient $\nabla\omega=\omega_0 \, \nabla\alpha$ is the first-order coefficient of the frequency field (Eq.~\eqref{eq:omega_expansion}), and at that order gravity and inertia coincide in both the force and the mass. They separate at second order, in the Hessian $\omega_0\,\partial_i\partial_j\alpha$, the curvature of the frequency field. To first order in the potential this Hessian is $\omega_0\,R_{0i0j}$, where $R_{0i0j}$ are the tidal components of the Riemann curvature tensor, the index $0$ labels the time coordinate $ct$, and $i,j$ label the space coordinates \citep{MTW1973Gravitation}. For the metric component $g_{00}=-\alpha^2$ with no cross term and a flat spatial part, $R_{0i0j}=\alpha\,\partial_i\partial_j\alpha$. The Rindler factor (Eq.~\eqref{eq:alpha_rindler}) is linear in position, so its Hessian and its Riemann tensor vanish, and a change to an inertial frame removes the factor: uniform acceleration produces a spectral gradient but no tidal force. The weak-field factor (Eq.~\eqref{eq:alpha_weak}) is nonlinear, so its Hessian is nonzero; the tidal components $R_{0i0j}$ do not vanish in the Newtonian limit and produce tidal forces that persist in every frame, since a tensor that is nonzero in one frame is nonzero in every frame.

The local force depends only on the shared first-order coefficient (Eq.~\eqref{eq:geometric_force}), so a uniform gravitational field and a uniformly accelerating frame give the particle the same frequency field and the same force, which is the equivalence principle of \citet{Einstein1907Relativitatsprinzip} as it applies to the particle. In the experiment of \citet{Dobkowski2026Observation} (Section~\ref{sec:intro}), the wave packet held static has the static spectral gradient and the freely falling packet, in its own frame, has none. The Hessian does not separate the two inverse-square forces: the Newtonian and Coulomb potentials are both proportional to $1/r$, whose Hessian has the eigenvalues $2/r^3$ along the radial direction and $-1/r^3$ along each of the two directions perpendicular to it, so the tidal tensor of a point mass is proportional to that of a point charge.
\end{remark}

%-------------------------------------------------------------------------------
\subsection{Light}
\label{sec:light}
%-------------------------------------------------------------------------------

The wavevector form of the identity, the eikonal ray equation (Eq.~\eqref{eq:eikonal_ray}), applies to light as well as matter where the wavelength is short compared with the length over which $\alpha$ varies. For a massive particle the time-dilation factor multiplies the rest frequency, $\omega=\omega_0\,\alpha$ (Eq.~\eqref{eq:omega_alpha}). Light has no rest frequency for $\alpha$ to multiply, and $\alpha$ acts on it through its dispersion relation instead. In isotropic coordinates the weak-field metric of a static mass is, to first order in the potential,
\begin{equation}
    ds^2 = -\alpha^2 \, c^2\,dt^2 + \alpha^{-2}\,d\mathbf{x}^2
    \label{eq:isotropic}
\end{equation}
\citep{MTW1973Gravitation}. With $x^0=ct$, the covariant wavevector is $k_\mu=(-\omega/c,\mathbf{k})$, and $g^{\mu\nu}$ is the inverse of the metric $g_{\mu\nu}$. For a wave of mass $m$, the mass-shell condition $g^{\mu\nu}k_\mu k_\nu=-(mc/\hbar)^2$, with the Greek indices summed over $0$ to $3$, gives for this metric the dispersion relation
\begin{equation}
    \hbar\,\omega = \alpha\,\sqrt{m^2c^4+\alpha^2c^2\hbar^2|\mathbf{k}|^2} .
    \label{eq:isotropic_dispersion}
\end{equation}
As $|\mathbf{k}|\to0$ it gives $\omega=\omega_0\,\alpha$, and with it the gravitational force of Eq.~\eqref{eq:geometric_force}. At $m=0$ it gives
\begin{equation}
    \omega = \frac{c\,|\mathbf{k}|}{n_\text{r}} ,
    \label{eq:light_dispersion}
\end{equation}
where $n_\text{r}=1/\alpha^2$. This is the dispersion relation of light in a medium of refractive index $n_\text{r}$ \citep{Eddington1920,Gordon1923Refraction,deFelice1971Gravitational}. For a particle at rest $\alpha$ enters the dispersion relation once, through the rest frequency; for light it enters twice, once through the time part of the metric and once through the spatial part. For a body of speed $v$ the spatial part enters at order $v^2/c^2$, which is why a slowly moving body falls almost entirely because of the time dilation \citep{Gould2016Why}.

The gradient of Eq.~\eqref{eq:light_dispersion} at fixed wavevector gives, by Eq.~\eqref{eq:eikonal_ray}, the spectral gradient of the ray,
\begin{align}
    \hbar\,\dot{\mathbf{k}} &= -\hbar\,\nabla\omega \notag\\
    &= \frac{\hbar \, \omega}{n_\text{r}}\,\nabla n_\text{r} ,
    \label{eq:light_gradient}
\end{align}
with $\nabla n_\text{r}=-2\nabla\alpha/\alpha^3$. Light therefore deflects toward the region where clocks run slower, because $n_\text{r}$ rises where $\alpha$ falls.

Along the ray $\omega$ is constant, since the index is static: with $\dot{\mathbf{x}}=\partial\omega/\partial\mathbf{k}$ and Eq.~\eqref{eq:eikonal_ray}, $\dot{\omega}=\nabla\omega\cdot\dot{\mathbf{x}}+\dot{\mathbf{k}}\cdot\partial\omega/\partial\mathbf{k}=0$. Let $s_\text{ray}$ be the path length along the ray. The group velocity $(c/n_\text{r})\,\mathbf{k}/|\mathbf{k}|$ is parallel to $\mathbf{k}$, and $|\mathbf{k}|=n_\text{r}\omega/c$ by Eq.~\eqref{eq:light_dispersion}. The speed along the ray and the wavevector are therefore
\begin{align*}
    \frac{ds_\text{ray}}{dt} &= \frac{c}{n_\text{r}} , \\
    \mathbf{k} &= \frac{n_\text{r}\,\omega}{c}\,\frac{d\mathbf{x}}{ds_\text{ray}} .
\end{align*}
Substituting both into Eq.~\eqref{eq:light_gradient}, with $\omega$ constant along the ray, gives
\[
    \frac{\hbar\,\omega}{n_\text{r}}\,\frac{d}{ds_\text{ray}}\!\left(n_\text{r} \, \frac{d\mathbf{x}}{ds_\text{ray}}\right) = \frac{\hbar \, \omega}{n_\text{r}}\,\nabla n_\text{r} ,
\]
and dividing by $\hbar\omega/n_\text{r}$ gives the ray equation of geometrical optics,
\begin{equation}
    \frac{d}{ds_\text{ray}}\!\left(n_\text{r} \, \frac{d\mathbf{x}}{ds_\text{ray}}\right) = \nabla n_\text{r} ,
    \label{eq:optical_path}
\end{equation}
\citep{Evans1986Optics}, the differential form of Fermat's principle \citep{Fermat1662Refractions}. Each factor of $\alpha$ in $n_\text{r}$ contributes one half of the deflection $4GM/(c^2b_\text{imp})$ of a ray at impact parameter $b_\text{imp}$ \citep{MTW1973Gravitation}, and the same index gives the Shapiro delay \citep{Shapiro1964Fourth,Bertotti2003Test}.

Maxwell's equations in a curved spacetime are the flat-space equations in a medium whose relative permittivity and permeability matrices $\varepsilon_r^{ij}$ and $\mu_r^{ij}$ are equal, with
\begin{equation}
    \varepsilon_r^{ij} = -\frac{\sqrt{-\det(g_{\mu\nu})}\,g^{ij}}{g_{00}} \, ,
    \label{eq:constitutive_map}
\end{equation}
where $\det(g_{\mu\nu})$ is the determinant of the metric. Where $g_{0i}\neq0$ the medium also couples the electric and magnetic fields, through $g_{0i}/g_{00}$ \citep{Gordon1923Refraction,Plebanski1960Electromagnetic,Leonhardt2006General}. For the metric of Eq.~\eqref{eq:isotropic}, $\sqrt{-\det(g_{\mu\nu})}=\alpha^{-2}$, $g^{ij}=\alpha^2\delta^{ij}$ and $g_{00}=-\alpha^2$, so $\varepsilon_r^{ij}=\mu_r^{ij}=\varepsilon_r\,\delta^{ij}$ with $\varepsilon_r=\mu_r=1/\alpha^2$, and each response contributes one factor $1/\alpha$ to the index $\sqrt{\varepsilon_r\mu_r}=n_\text{r}$.

A medium with these responses has the wave impedance $\mathcal{Z}_0\sqrt{\mu_r/\varepsilon_r}$, where $\mathcal{Z}_0=\sqrt{\mu_0/\varepsilon_0}$ is the impedance of free space and $\varepsilon_0$ and $\mu_0$ are its permittivity and permeability. Since $\mu_r=\varepsilon_r$, the impedance is $\mathcal{Z}_0$ at every point, and the medium reflects no wave at normal incidence, however fast its index changes. For every metric with $g_{0i}=0$ the two responses are equal matrices, so along each direction the two polarizations have the same index, even where the medium is anisotropic and that index depends on the direction, as in Schwarzschild coordinates.

Einstein's 1911 calculation \cite{Einstein1911Influence} slowed light to $c(1+\Phi_g/c^2)$, an index of $1/\alpha$ to first order, and Huygens' construction \citep{Huygens1690Traite} then gave $2GM/(c^2b_\text{imp})$, half the measured deflection. For the time part of Eq.~\eqref{eq:isotropic} alone, $g_{00}=-\alpha^2$ with $g_{ij}=\delta_{ij}$, Eq.~\eqref{eq:constitutive_map} gives $\varepsilon_r=\mu_r=1/\alpha$, the index of 1911. The spatial part multiplies each response by a further $1/\alpha$, so the missing half enters the electric and the magnetic response equally.

For a static, spherically symmetric vacuum the field equations fix the spatial part as a function of $\alpha$: $g_{rr}=-1/g_{00}$ exactly in Schwarzschild coordinates. The full index is then a function of $\alpha$ alone, $1/\alpha^2$ to first order and $4/[\alpha(1+\alpha)^2]$ exactly in isotropic coordinates \citep{MTW1973Gravitation}, and in flat coordinate space the whole deflection follows from the gradient of $\alpha$ through Eq.~\eqref{eq:light_gradient}. Lensing measures only the product $\varepsilon_r\mu_r=n_\text{r}^2$; the equal split, and with it the impedance $\mathcal{Z}_0$, comes from Eq.~\eqref{eq:constitutive_map}.

%===============================================================================
\section{Force laws from the Dirac equation}
\label{sec:potentials}
%===============================================================================

The companion paper \citep{HutchisonWavefn2026} obtains the frequency of a spin-$\tfrac{1}{2}$ particle in electromagnetic potentials and a stationary metric from the minimally coupled Dirac equation.

%-------------------------------------------------------------------------------
\subsection{Additive potentials}
\label{sec:em}
%-------------------------------------------------------------------------------

Gravity and inertia multiply the rest frequency by $\alpha$ (Eq.~\eqref{eq:omega_alpha}). A gauge potential instead adds a term to the energy, and so to the frequency. For a particle of charge $q$ in the scalar potential $\Phi$ and the vector potential $\mathbf{A}$, the \emph{kinematic momentum} is $\boldsymbol{\pi}\equiv\mathbf{p}-q\mathbf{A}$ and the minimally coupled frequency is
\begin{equation}
    \omega = \omega_0 + \frac{|\mathbf{p}-q\mathbf{A}|^2/2m + q\Phi}{\hbar} .
    \label{eq:omega_em}
\end{equation}
Hamilton's velocity is $\mathbf{v}=\hbar\nabla_{\mathbf{p}}\omega=\boldsymbol{\pi}/m$, where $\nabla_{\mathbf{p}}$ is the gradient in momentum at fixed position, and the gradient at fixed $\mathbf{p}$ is $\hbar\,\partial_i\omega=q\,\partial_i\Phi-q\,v_j\,\partial_iA_j$, summed over $j$. By Theorem~\ref{thm:universal} (Eq.~\eqref{eq:sgm_as_dot_p}) the canonical momentum therefore changes at the rate $\dot{p}_i=-q\,\partial_i\Phi+q\,v_j\,\partial_iA_j$, which depends on the gauge. The observable force is the rate of $\boldsymbol{\pi}$,
\begin{align}
    \dot{\boldsymbol{\pi}} &= -\hbar\nabla\omega-q\,\dot{\mathbf{A}} \notag\\
    &= q \, (\mathbf{E}+\mathbf{v}\times\mathbf{B}) ,
    \label{eq:lorentz} \\[-3pt]
    \multispan2{\hfil\text{\small\itshape (The Lorentz force)}\hfil} \notag
\end{align}
where $\mathbf{E}=-\nabla\Phi-\partial_t\mathbf{A}$ is the electric field, $\mathbf{B}=\nabla\times\mathbf{A}$ the magnetic field, and $\dot{\mathbf{A}}=\partial_t\mathbf{A}+(\mathbf{v}\cdot\nabla)\mathbf{A}$ the rate of change of the potential along the path. The second line follows because $v_j\,\partial_iA_j-v_j\,\partial_jA_i=(\mathbf{v}\times\mathbf{B})_i$, so the result depends on the potentials only through $\mathbf{E}$ and $\mathbf{B}$.

An additive term has no factor $\omega_0$ in its gradient, so its force is proportional to the charge and independent of the mass. The gradient of the geometric term $\omega_0\alpha$ is $\omega_0\nabla\alpha$, with $\omega_0=mc^2/\hbar$, so its force is proportional to the mass and every body has the same free-fall acceleration (Eq.~\eqref{eq:freefall_accel}).

A uniform additive potential shifts the frequency of every level of a system of net charge $q_\text{tot}$ by the same $q_\text{tot} \, \Phi/\hbar$. The shift cancels from every transition frequency, so it leaves the rate of every clock unchanged; it appears only in the phase, and an interferometer detects it only as a difference between arms held at different potentials \citep{Aharonov1959Significance}. A clock inside a charged conducting shell, where $\Phi$ is uniform, keeps its transition frequencies, while a gravitational potential multiplies them by $\alpha$. Table~\ref{tab:classes} summarizes the classes.

\begin{table}[htbp]
\centering
\fontsize{10}{11.5}\selectfont
\caption{The two coupling classes and the Lorentz-scalar coupling, distinguished by how each acts on the rest frequency $\omega_0=mc^2/\hbar$ (Eq.~\eqref{eq:rest_freq}). The multiplicative form is Eq.~\eqref{eq:omega_alpha} and the additive form Eq.~\eqref{eq:omega_em}. A Lorentz-scalar potential $V_\text{sc}$ changes the rest mass $m_0$ to $m=m_0+V_\text{sc}/c^2$; it belongs to the geometric class only if it changes every mass by the same fraction.}
\label{tab:classes}
\begin{tabular}{lll}
\toprule
\textbf{Coupling} & \textbf{Frequency at rest} & \textbf{Mass-independent acceleration} \\
\midrule
Geometric (gravity, inertia, rotation) & $\omega=\omega_0\,\alpha$ & yes, $\mathbf{F}\propto m$ \\
Charge (electromagnetic, strong) & $\omega=\omega_0+V/\hbar$ & no \\
Lorentz scalar & $\omega=\omega_0+V_\text{sc}/\hbar$ & only if $V_\text{sc}\propto m_0$ \\
\bottomrule
\end{tabular}
\end{table}

\begin{remark}[Gauge invariance fixes the additive form]\label{rem:additivity}
In the coordinates $x^\mu=(ct,\mathbf{x})$, with $\partial_\mu\equiv\partial/\partial x^\mu$, the four-gradient of the phase is
\begin{equation}
    \partial_\mu\theta = (-\omega/c,\ \mathbf{k})
    \label{eq:phase_fourvector}
\end{equation}
by Eqs.~\eqref{eq:phase_freq} and~\eqref{eq:debroglie_k}. A gauge transformation $\psi\to e^{iq\lambda/\hbar}\psi$ changes the four-potential $A_\mu=(-\Phi/c,\mathbf{A})$ to $A_\mu+\partial_\mu\lambda$ and leaves $\hbar\,\partial_\mu\theta-qA_\mu$ unchanged.

Suppose the potential entered as $\omega=\omega_0\,f(\Phi)$, with $f$ differentiable and $f(0)=1$. The gauge function $\lambda=-t\,\delta\Phi$, for a constant $\delta\Phi$, raises $\Phi$ by $\delta\Phi$ and $\omega$ by $q\,\delta\Phi/\hbar$, so $\omega_0\,[f(\Phi+\delta\Phi)-f(\Phi)]=q\,\delta\Phi/\hbar$ for all $\Phi$ and $\delta\Phi$. Differentiating at $\delta\Phi=0$ gives $f'=q/mc^2$, so $\omega=\omega_0+q\Phi/\hbar$: the potential adds. Terms built from the field strength, such as an anomalous magnetic moment, are gauge invariant by themselves and contain no bare $A_\mu$, so they do not change the additive form. The factor $\alpha$ adds no phase of its own: it converts coordinate time to proper time, $d\tau=\alpha\,dt$ (Eq.~\eqref{eq:alpha_def}), and so scales the rate at which the existing phase advances.
\end{remark}

%-------------------------------------------------------------------------------
\subsection{The spinor as a quaternion}
\label{sec:spinor}
%-------------------------------------------------------------------------------

The quaternion units $e_1$, $e_2$, $e_3$ multiply as
\begin{equation}
    e_j\,e_k = -\delta_{jk} + \epsilon_{jkl}\,e_l ,
    \label{eq:quaternion_units}
\end{equation}
summed over $l$, where $\delta_{jk}$ is the Kronecker delta and $\epsilon_{jkl}$ the Levi-Civita symbol. A quaternion is $Y=Y_0+\mathbf{Y}\cdot\mathbf{e}$, with a real number $Y_0$, a real vector $\mathbf{Y}$, and $\mathbf{e}\equiv(e_1,e_2,e_3)$. Its conjugate is $\bar{Y}=Y_0-\mathbf{Y}\cdot\mathbf{e}$ and its squared norm $|Y|^2=Y\bar{Y}$. The two complex components of a spin-$\tfrac{1}{2}$ wave function $(\psi_1,\psi_2)$ are four real numbers, and they form one quaternion,
\begin{equation}
    \psi = \operatorname{Re}\psi_1 - (\operatorname{Im}\psi_1)\,e_3 + e_2\,\big[\operatorname{Re}\psi_2 - (\operatorname{Im}\psi_2)\,e_3\big] .
    \label{eq:spinor_quaternion}
\end{equation}
In this form multiplication by $i$ is right multiplication by $-e_3$, and the Pauli matrices $\sigma_k$ act from the left,
\begin{equation}
    \sigma_k\,\psi = -e_k\,\psi\,e_3 ,
    \label{eq:pauli_quaternion}
\end{equation}
so $-i\sigma_k$ is left multiplication by $e_k$ \citep{HutchisonWavefn2026}. The density is $\rho=|\psi|^2=|\psi_1|^2+|\psi_2|^2$, and the spin is $\mathbf{S}=\tfrac{1}{2}\hbar\boldsymbol{\sigma}$.

Every nonzero quaternion is its norm times a unit quaternion, so
\begin{equation}
    \psi = \rho^{1/2}\,\mathcal{R}_\psi ,
    \label{eq:polar_form}
\end{equation}
with $\mathcal{R}_\psi\bar{\mathcal{R}}_\psi=1$. A unit quaternion $\mathcal{R}$ rotates a vector $\mathbf{w}$ by $\mathbf{w}\cdot\mathbf{e}\mapsto\mathcal{R}\,(\mathbf{w}\cdot\mathbf{e})\,\bar{\mathcal{R}}$, and $\mathcal{R}=\exp(\tfrac{1}{2}\vartheta\,\hat{\mathbf{n}}\cdot\mathbf{e})$ rotates it by the angle $\vartheta$ about the unit vector $\hat{\mathbf{n}}$. The spin points along the unit vector $\hat{\mathbf{s}}$ given by
\begin{equation}
    \hat{\mathbf{s}}\cdot\mathbf{e} = \mathcal{R}_\psi\,e_3\,\bar{\mathcal{R}}_\psi ,
    \label{eq:spin_rotor}
\end{equation}
so left multiplication of $\psi$ by $\mathcal{R}$ rotates the spin. A phase factor $e^{-i\phi}$ is right multiplication by $\exp(e_3\phi)$, so a stationary state of frequency $\omega$ is
\begin{equation}
    \psi(\mathbf{x},t) = \psi_0(\mathbf{x})\,\exp(e_3\,\omega t) .
    \label{eq:stationary_time}
\end{equation}
The wave function is therefore a density and a frame at each point: the frame $\mathcal{R}_\psi$ rotates $e_3$ into the spin axis $\mathcal{R}_\psi e_3\bar{\mathcal{R}}_\psi$, and as the phase advances its two transverse axes $\mathcal{R}_\psi e_1\bar{\mathcal{R}}_\psi$ and $\mathcal{R}_\psi e_2\bar{\mathcal{R}}_\psi$ turn about the spin axis at the rate $2\omega$.

Because a change of phase acts on the right, the potentials of the charge class enter the covariant derivatives as right multiplications. The gauge-covariant derivatives are
\begin{align}
    D_t\psi &= \partial_t\psi - \frac{V}{\hbar}\,\psi\,e_3 ,
    \label{eq:phase_connection_time} \\
    D_k\psi &= \partial_k\psi + \frac{q}{\hbar}\,A_k\,\psi\,e_3 ,
    \label{eq:right_connection}
\end{align}
where $V$ is the additive potential energy ($q\Phi$ for a charge), $\partial_k\equiv\partial/\partial x_k$, and $k=1,2,3$. With $\hat{\boldsymbol{\pi}}=-i\hbar\nabla-q\mathbf{A}$, Eq.~\eqref{eq:right_connection} reads $D_k\psi=-\hbar^{-1}(\hat{\pi}_k\psi)\,e_3$. The fields are commutators of these derivatives:
\begin{align}
    [D_1,D_2]\,\psi &= \frac{q}{\hbar}\,B_3\,\psi\,e_3 ,
    \label{eq:phase_curvature} \\
    [D_t,D_k]\,\psi &= \frac{1}{\hbar}\,\big(\partial_kV+q\,\partial_tA_k\big)\,\psi\,e_3 ,
    \label{eq:time_space_curvature}
\end{align}
with the cyclic permutations of the first, and for $V=q\Phi$ the right side of the second is $-(q/\hbar)E_k\,\psi\,e_3$. With $\mathbf{D}\equiv(D_1,D_2,D_3)$ and $\mathbf{D}^2\equiv D_1^2+D_2^2+D_3^2$, Eqs.~\eqref{eq:quaternion_units} and~\eqref{eq:phase_curvature} give
\begin{equation}
    (\mathbf{e}\cdot\mathbf{D})^2\,\psi = -\mathbf{D}^2\psi + \frac{q}{\hbar}\,(\mathbf{e}\cdot\mathbf{B})\,\psi\,e_3 ,
    \label{eq:fueter_square}
\end{equation}
which in Pauli form is $(\boldsymbol{\sigma}\cdot\hat{\boldsymbol{\pi}})^2=\hat{\boldsymbol{\pi}}^2-q\hbar\,\boldsymbol{\sigma}\cdot\mathbf{B}$, the magnetic spin coupling with g-factor $g_s=2$ \citep{Dirac1928Quantum}.

The geometric class acts on the frame from the left and converts coordinate time to proper time. A stationary metric with conformally flat spatial slices is
\begin{equation}
    ds^2 = -\alpha_n^2c^2\,dt^2 + W^2\,|d\mathbf{x}+\mathbf{A}_g\,dt|^2 ,
    \label{eq:lapse_shift_metric}
\end{equation}
where the \emph{lapse} $\alpha_n>0$ converts coordinate time to the proper time of the observer whose four-velocity is normal to the slice, the \emph{shift} $\mathbf{A}_g$ is a velocity, the gravitomagnetic potential, and $W>0$ is the spatial scale factor \citep{Arnowitt2008Dynamics,MTW1973Gravitation}. Its time-time component gives
\begin{equation}
    \alpha = \sqrt{\alpha_n^2 - W^2|\mathbf{A}_g|^2/c^2} .
    \label{eq:alpha_lapse_shift}
\end{equation}
Flat spacetime has $\alpha_n=W=1$ and $\mathbf{A}_g=\mathbf{0}$. A frame rotating at the angular velocity $\boldsymbol{\Omega}$ has $\alpha_n=W=1$ and $\mathbf{A}_g=\boldsymbol{\Omega}\times\mathbf{x}$. The post-Newtonian metric has $\alpha_n=1+\Phi_g/c^2+(\beta_\text{PPN}-\tfrac{1}{2})\Phi_g^2/c^4$ and $W^{-1}=1+\gamma_\text{PPN}\Phi_g/c^2$, with the parameters $\beta_\text{PPN}=\gamma_\text{PPN}=1$ in general relativity, and a body of angular momentum $\mathbf{J}_\text{src}$ adds $\mathbf{A}_g=-(2G/c^2r^3)\,\mathbf{J}_\text{src}\times\mathbf{x}$ \citep{Will2014Confrontation}. An accelerating frame adds $\mathbf{a}\cdot\mathbf{x}/c^2$ to $\alpha_n$ (Eq.~\eqref{eq:alpha_rindler}).

The spinor is referred to the orthonormal axes of the metric and rescaled by $\alpha_n^{1/2}W^{3/2}$, so that its norm is $\int|\psi|^2\alpha_n^{-1}d^3x$. It then has the \emph{covariant derivative along the normal} to the slice
\begin{equation}
    D_n\psi = \frac{1}{c\,\alpha_n}\big[D_t\psi - \mathbf{A}_g\cdot\mathbf{D}\psi + \tfrac{1}{4}(\mathbf{e}\cdot\mathbf{B}_g)\,\psi - \tfrac{1}{2}(\nabla\cdot\mathbf{A}_g)\,\psi + \tfrac{1}{2}(\mathbf{A}_g\cdot\nabla\ln\alpha_n)\,\psi\big] ,
    \label{eq:normal_derivative}
\end{equation}
where $\mathbf{B}_g\equiv\nabla\times\mathbf{A}_g$. Its rotation term $\tfrac{1}{4}(\mathbf{e}\cdot\mathbf{B}_g)\,\psi$ turns the frame from the left. The curl $\mathbf{B}_g$ couples equally to the orbital angular momentum $\mathbf{L}=\mathbf{x}\times\hat{\boldsymbol{\pi}}$ and the spin $\mathbf{S}$, so a Dirac particle has no anomalous gravitomagnetic moment \citep{KobzarevOkun1963Gravitational,HehlNi1990Inertial}.

%-------------------------------------------------------------------------------
\subsection{The quaternionic wave equation}
\label{sec:wave_equation}
%-------------------------------------------------------------------------------

Define the operators
\begin{align}
    \mathsf{N}\psi &\equiv (D_n\psi)\,e_3 , \notag\\
    \mathsf{X}\psi &\equiv \mathbf{e}\cdot\mathbf{D}\psi , \notag\\
    \mathsf{X}_W\psi &\equiv W^{-1/2}\,\mathsf{X}\big(W^{-1/2}\psi\big) ,
    \label{eq:NX_operators}
\end{align}
and the inverse reduced Compton wavelength $\kappa\equiv mc/\hbar$. In flat spacetime $\mathsf{X}_W=\mathsf{X}=\boldsymbol{\sigma}\cdot\hat{\boldsymbol{\pi}}/\hbar$ and $\mathsf{N}=-(i\hbar\partial_t-V)/\hbar c$. For a free particle there, each of the four real components of $\psi$ obeys the Klein--Gordon equation \citep{Klein1926Quantum,Gordon1926Compton},
\begin{equation}
    \frac{1}{c^2}\,\partial_t^2\psi + (\mathbf{e}\cdot\nabla)^2\psi + \kappa^2\psi = 0 ,
    \label{eq:klein_gordon}
\end{equation}
since $(\mathbf{e}\cdot\nabla)^2=-\nabla^2$ by Eq.~\eqref{eq:quaternion_units}.

In the Weyl representation \citep{Peskin1995Introduction}, the Dirac equation \citep{Dirac1928Quantum} for the left- and right-handed spinors $\psi_L$ and $\psi_R$ is
\begin{equation}
    \begin{gathered}
    i\hbar\,\partial_t \begin{pmatrix}\psi_L\\ \psi_R\end{pmatrix} = \begin{pmatrix} q\Phi - c\,\boldsymbol{\sigma}\cdot\hat{\boldsymbol{\pi}} & mc^2 \\ mc^2 & q\Phi + c\,\boldsymbol{\sigma}\cdot\hat{\boldsymbol{\pi}} \end{pmatrix} \begin{pmatrix}\psi_L\\ \psi_R\end{pmatrix} . \\[-3pt]
    \text{\small\itshape (The Dirac equation)}
    \end{gathered}
    \label{eq:dirac_equation}
\end{equation}
With the chiral components $\varphi_R\equiv\psi_R$ and $\varphi_L\equiv-\psi_L$ written as quaternions, and with minimal coupling, $c^{-1}\partial_t\to D_n$ and $\nabla\to W^{-1/2}\mathbf{D}\,W^{-1/2}$, the Dirac equation in the metric of Eq.~\eqref{eq:lapse_shift_metric} becomes the pair \citep{HutchisonWavefn2026}
\begin{align}
    (\mathsf{N}+\mathsf{X}_W)\,\varphi_R &= \kappa\,\varphi_L ,
    \label{eq:weyl_pair_R} \\
    (\mathsf{N}-\mathsf{X}_W)\,\varphi_L &= \kappa\,\varphi_R .
    \label{eq:weyl_pair_L}
\end{align}
Minimal coupling is an assumption: it gives the particle no anomalous magnetic moment. For $m\neq0$ constant on the region, Eq.~\eqref{eq:weyl_pair_R} gives $\varphi_L$ from $\varphi_R$, and substituting it into Eq.~\eqref{eq:weyl_pair_L} gives, with $\psi\equiv\varphi_R$,
\begin{equation}
    \begin{gathered}
    \boxed{\; (\mathsf{N}-\mathsf{X}_W)(\mathsf{N}+\mathsf{X}_W) \, \psi = \kappa^2 \, \psi \;} \, . \\[-3pt]
    \text{\small\itshape (The quaternionic wave equation)}
    \end{gathered}
    \label{eq:pure_form}
\end{equation}
Its solutions correspond one to one with those of the Dirac equation. Expanding the operator product gives $\mathsf{N}^2-\mathsf{X}_W^2+[\mathsf{N},\mathsf{X}_W]$: the part $\mathsf{N}^2-\mathsf{X}_W^2$ is the Klein--Gordon operator of Eq.~\eqref{eq:klein_gordon} in covariant derivatives, with the magnetic spin coupling of Eq.~\eqref{eq:fueter_square} inside $\mathsf{X}^2$, and the commutator $[\mathsf{N},\mathsf{X}_W]$ couples the spin to the electric field and to the gradients of $\alpha_n$ and $W$. The conserved density of the Dirac equation is $(|\varphi_R|^2+|\varphi_L|^2)/\alpha_n$.

The coefficients $V$, $\mathbf{A}$, $\alpha_n$, $\mathbf{A}_g$, and $W$ are external fields. In a system of many particles in the mean-field approximation, each particle's wave equation contains the fields that Maxwell's equations and the field equations of gravity produce from the charge and mass densities of the others.

%-------------------------------------------------------------------------------
\subsection{The master frequency equation}
\label{sec:master}
%-------------------------------------------------------------------------------

For a stationary state, Eq.~\eqref{eq:stationary_time}, the time derivative becomes $\partial_t\psi=\psi\,e_3\,\omega$, and Eq.~\eqref{eq:pure_form} becomes an eigenvalue problem for $\omega$. In stationary fields, whenever the positive-frequency states form the graph of a bounded operator, a norm-preserving change of spinor separates them exactly from the negative-frequency states. The companion proves this graph condition when the terms of the Dirac Hamiltonian off its diagonal and the potentials are bounded operators whose norms sum to less than $mc^2$. Operators on a finite set of basis states are always bounded. The companion also proves the condition for a free particle, and assumes it for a particle in weak fields in a state built from momenta small compared with $mc$ \citep{HutchisonWavefn2026,Foldy1950Dirac,Eriksen1958Foldy}. Under the same graph hypothesis, the separation holds in every stationary metric, with any spatial metric on the slices \citep{HutchisonWavefn2026}. On the conformally flat slices of Eq.~\eqref{eq:lapse_shift_metric} the positive-frequency spinor $\chi$ then obeys the \emph{master frequency equation}
\begin{equation}
    \begin{gathered}
    \boxed{\;
    \hbar\,\omega \, \chi = \Big[V + \alpha_n\sqrt{\big(mc^2\big)^2 + c^2 \, W^{-1}\big(\hat{\boldsymbol{\pi}}^2 - q\hbar\,\boldsymbol{\sigma}\cdot\mathbf{B}\big)W^{-1}} - \hat{\Pi}_g - \tfrac{1}{4}\,\hbar\,\boldsymbol{\sigma}\cdot\mathbf{B}_g + \mathsf{C}\Big] \, \chi\;} \, . \\[-3pt]
    \text{\small\itshape (The master frequency equation)}
    \end{gathered}
    \label{eq:master}
\end{equation}
Here $V$ is the additive potential energy in coordinate time, and $m=m_0+V_\text{sc}/c^2$ is the mass with a Lorentz-scalar potential $V_\text{sc}$. The shift coupling is $\hat{\Pi}_g\equiv\alpha_n^{1/2}\,\tfrac{1}{2}\{\mathbf{A}_g,\hat{\boldsymbol{\pi}}\}\,\alpha_n^{-1/2}$, where $\{\cdot,\cdot\}$ is the anticommutator. With $\mathsf{G}\equiv m_0c^2(\alpha_n-1)$ and $\hat{V}\equiv V-\hat{\Pi}_g-\tfrac{1}{4}\hbar\,\boldsymbol{\sigma}\cdot\mathbf{B}_g$, the term $\mathsf{C}$ is a sum of commutators of $\boldsymbol{\sigma}\cdot\hat{\boldsymbol{\pi}}$, $W$, $\mathsf{G}$, and $\hat{V}$, and it vanishes when these operators commute \citep{HutchisonWavefn2026}.

The classical limit replaces each operator by its \emph{Weyl symbol}, the function of $\mathbf{x}$ and $\mathbf{p}$ that corresponds to it under symmetric ordering. At order $\hbar^0$ the spin terms and the symbol of $\mathsf{C}$ vanish \citep{HutchisonWavefn2026,PanatiSpohnTeufel2003Space}. For a particle at rest ($\dot{\mathbf{x}}=\mathbf{0}$) in any stationary field in which Eq.~\eqref{eq:master} holds, that equation at this order gives, with no expansion in $c^{-2}$ and with $\alpha$ from Eq.~\eqref{eq:alpha_lapse_shift},
\begin{equation}
    \hbar\omega = V + mc^2\alpha .
    \label{eq:static_master}
\end{equation}
At rest $\nabla_{\mathbf{p}}\omega=\mathbf{0}$, so the gradient of Eq.~\eqref{eq:static_master} along the rest momentum equals the gradient at fixed momentum, and the identity of Eq.~\eqref{eq:sgm} gives $\mathbf{F}=-\nabla V-c^2\nabla(m\alpha)$. For constant $m$ this is $\mathbf{F}=-\nabla V-mc^2\nabla\alpha$, the sum of the additive force of Section~\ref{sec:em} and the geometric force of Eq.~\eqref{eq:geometric_force}.

%-------------------------------------------------------------------------------
\subsection{The Taylor expansion and the force laws}
\label{sec:expansion_orders}
%-------------------------------------------------------------------------------

Expanding Eq.~\eqref{eq:master} in powers of $c^{-2}$ for a slow particle in weak fields gives its operator $\hat{h}$, the bracket of Eq.~\eqref{eq:master}, as a sum of eight terms \citep{HutchisonWavefn2026}:
\begin{equation}
    \begin{gathered}
    \boxed{\begin{aligned} \hat{h} = {} & \underbrace{mc^2 \, \alpha_n}_{\substack{= \hbar\,\omega_\text{geom} \\ \text{(geometric)}}} + \underbrace{V}_{\substack{= \hbar\,\omega_\text{ext} \\ \text{(extrinsic)}}} + \underbrace{\frac{\alpha_n}{4}\,\big\{m^{-1},W^{-1}\hat{\boldsymbol{\pi}}^2W^{-1}\big\}}_{\substack{= \hbar\,\omega_\text{kin} \\ \text{(kinetic)}}} - \underbrace{\hat{\Pi}_g}_{\substack{= \hbar\,\omega_\text{gm} \\ \text{(gravitomagnetic)}}} - \underbrace{\frac{\alpha_nq\hbar\,\boldsymbol{\sigma}\cdot\mathbf{B}}{2mW^2}}_{\substack{= \hbar\,\omega_\text{Pauli} \\ \text{(Pauli)}}} \\ & - \underbrace{\frac{\hbar\,\boldsymbol{\sigma}\cdot\mathbf{B}_g}{4}}_{\substack{= \hbar\,\omega_\text{rot} \\ \text{(spin-rotation)}}} - \underbrace{\frac{\hbar^4\mathsf{X}^4}{8m^3c^2}}_{\substack{= \hbar\,\omega_\text{rel} \\ \text{(relativistic)}}} + \underbrace{\frac{\hbar^2}{2m}\big(\tilde{\mathsf{X}}_W\tilde{\mathsf{X}}_W^\dagger-W^{-1}\mathsf{X}^2W^{-1}\big)}_{\substack{= \hbar\,\omega_\text{so} \\ \text{(Darwin \& spin-orbit)}}} + O(c^{-4}) \end{aligned}} \\[-3pt]
    \text{\small\itshape (The Taylor expansion of the frequency equation)}
    \end{gathered}
    \label{eq:nr_expansion}
\end{equation}
The last term is the leading part of $\mathsf{C}$. In it the dressed operators are $\tilde{\mathsf{X}}\equiv e^{-f/4mc^2}\,\mathsf{X}\,e^{f/4mc^2}$ and $\tilde{\mathsf{X}}_W\equiv W^{-1/2}\,\tilde{\mathsf{X}}\,W^{-1/2}$, with $f\equiv\hat{V}+\mathsf{G}-V_\text{sc}$, and $\dagger$ denotes the adjoint. The similarity transformation $\alpha_n^{-1/2}\hat{h}\,\alpha_n^{1/2}$ makes $\hat{h}$ Hermitian in the unweighted inner product. The leading Weyl symbol of each transformed term is $\hbar$ times one additive part of the frequency $\omega(\mathbf{x},\mathbf{p})$, and Hamilton's equations use these symbols; Eq.~\eqref{eq:ehrenfest_commutator} gives the exact quantum statement, for operator expectation values. Since Eq.~\eqref{eq:sgm} is linear in $\omega$, the forces of the parts add.

Table~\ref{tab:forces} lists the leading forces and spin precessions of the terms of Eq.~\eqref{eq:nr_expansion}. Velocity-dependent terms give their force through the rate of the kinematic momentum, as in Eq.~\eqref{eq:lorentz}, with the velocity $\dot{\mathbf{x}}=\boldsymbol{\pi}/m_0-\mathbf{A}_g$. A spin term gives the precession $d\mathbf{S}/dt=(i/\hbar)[\hat{h},\mathbf{S}]$ (Eq.~\eqref{eq:ehrenfest_commutator}).

The expansion assumes stationary fields, so the induction and Euler forces also require the time derivatives of the potentials in the momentum-velocity relation. The ponderomotive force is an average $\langle\cdot\rangle$ over one cycle of an oscillating field. The table uses the post-Newtonian lapse with the frame acceleration $\mathbf{a}$ of the accelerating frame added, $\alpha_n=1+\Phi_g/c^2+(\beta_\text{PPN}-\tfrac{1}{2})\Phi_g^2/c^4+\mathbf{a}\cdot\mathbf{x}/c^2$, and $V=q\Phi+V_\text{ext}$, where $V_\text{ext}$ is the potential energy of any other additive force, such as a trap.

\begin{table}[tbp]
\centering
\fontsize{9.5}{11}\selectfont
\caption{Selected leading force laws produced by the parts of the frequency in Eq.~\eqref{eq:nr_expansion} through order $c^{-2}$; $m_0$ is the rest mass and $\langle\cdot\rangle$ an average over one cycle. The force is $-\hbar\nabla\omega$ (Eq.~\eqref{eq:sgm}) for a position-dependent term and the rate of the kinematic momentum (Eq.~\eqref{eq:lorentz}) for a velocity-dependent one; the 1PN velocity row gives only the canonical-force contribution $-\nabla$ of its term at fixed $\mathbf{p}$. A spin term $\mathbf{S}\cdot\boldsymbol{\Omega}_S$ gives the precession $d\mathbf{S}/dt=\boldsymbol{\Omega}_S\times\mathbf{S}$. Each entry is derived, and the full catalog with the omitted cross terms is given, in \citet{HutchisonWavefn2026}.}
\label{tab:forces}
\renewcommand{\arraystretch}{1.0}
\setlength{\tabcolsep}{3pt}
\begin{tabular}{p{0.2\linewidth}lll}
\toprule
\textbf{Force law} & \textbf{Part} & \textbf{Term of }$\hbar\omega$ & \textbf{Force, canonical force, or precession} \\
\midrule
Newtonian gravity & geom & $m_0\Phi_g$ & $-m_0\nabla\Phi_g$ \\
Nonlinear gravity (1PN) & geom & $(\beta_\text{PPN}-\tfrac{1}{2})\,m_0\Phi_g^2/c^2$ & $-(2\beta_\text{PPN}-1)\,(m_0\Phi_g/c^2)\,\nabla\Phi_g$ \\
Inertial reaction & geom & $m_0\,\mathbf{a}\cdot\mathbf{x}$ & $-m_0\mathbf{a}$ \\
Lorentz-scalar force & geom & $\alpha_nV_\text{sc}$ & $-\nabla(\alpha_nV_\text{sc})$ \\
Coulomb force & ext & $q\Phi$ & $-q\nabla\Phi$ \\
Trap or other additive force & ext & $V_\text{ext}$ & $-\nabla V_\text{ext}$ \\
Magnetic Lorentz force, induction & kin & $|\hat{\boldsymbol{\pi}}|^2/2m_0$ & $q\,\dot{\mathbf{x}}\times\mathbf{B}-q\,\partial_t\mathbf{A}$ \\
Diamagnetic, ponderomotive & kin & $q^2|\mathbf{A}|^2/2m_0$ & $-(q^2/2m_0)\,\nabla\langle|\mathbf{A}|^2\rangle$ \\
Velocity-dependent gravity (1PN) & kin & $(\tfrac{1}{2}+\gamma_\text{PPN})\,(\Phi_g/c^2)\,|\hat{\boldsymbol{\pi}}|^2/m_0$ & $-(\tfrac{1}{2}+\gamma_\text{PPN})\,m_0\,(|\dot{\mathbf{x}}|^2/c^2)\,\nabla\Phi_g$ \\
Frame dragging (Lense--Thirring) & gm & $-\tfrac{1}{2}\{\mathbf{A}_g,\hat{\boldsymbol{\pi}}\}$ & $m_0\,\dot{\mathbf{x}}\times\mathbf{B}_g-m_0\,\partial_t\mathbf{A}_g$ \\
Coriolis, Euler, centrifugal & gm & $-\boldsymbol{\Omega}\cdot\mathbf{L}$ & $\begin{array}[t]{@{}l@{}} -2m_0\boldsymbol{\Omega}\times\dot{\mathbf{x}}-m_0\dot{\boldsymbol{\Omega}}\times\mathbf{x} \\ {}-m_0\boldsymbol{\Omega}\times(\boldsymbol{\Omega}\times\mathbf{x}) \end{array}$ \\
Stern--Gerlach force, Larmor precession & Pauli & $-(q/m_0)\,\mathbf{S}\cdot\mathbf{B}$ & $\nabla[(q/m_0)\,\mathbf{S}\cdot\mathbf{B}],\ \ \boldsymbol{\Omega}_S=-(q/m_0)\,\mathbf{B}$ \\
Spin-rotation (Mashhoon), Barnett & rot & $-\tfrac{1}{2}\,\mathbf{S}\cdot\mathbf{B}_g$ & $\nabla(\tfrac{1}{2}\,\mathbf{S}\cdot\mathbf{B}_g),\ \ \boldsymbol{\Omega}_S=-\tfrac{1}{2}\,\mathbf{B}_g$ \\
Relativistic mass increase & rel & $-|\hat{\boldsymbol{\pi}}|^4/8m_0^3c^2$ & $\dot{\mathbf{x}}=(\boldsymbol{\pi}/m_0)\,(1-|\boldsymbol{\pi}|^2/2m_0^2c^2)$ \\
Darwin term & so & $(\hbar^2/8m_0^2c^2)\,\nabla^2V$ & $-(\hbar^2/8m_0^2c^2)\,\nabla(\nabla^2V)$ \\
Spin-orbit, Thomas factor $\tfrac{1}{2}$ & so & $\mathbf{S}\cdot(\nabla V\times\hat{\boldsymbol{\pi}})/2m_0^2c^2$ & $\boldsymbol{\Omega}_S=(\nabla V\times\boldsymbol{\pi})/2m_0^2c^2$ \\
Geodetic precession & so & $(\tfrac{1}{2}+\gamma_\text{PPN})\,\mathbf{S}\cdot(\nabla\Phi_g\times\hat{\boldsymbol{\pi}})/m_0c^2$ & $\boldsymbol{\Omega}_S=(\tfrac{1}{2}+\gamma_\text{PPN})\,(\nabla\Phi_g\times\dot{\mathbf{x}})/c^2$ \\
Inertial spin-orbit & so & $\mathbf{S}\cdot(\mathbf{a}\times\hat{\boldsymbol{\pi}})/2m_0c^2$ & $\boldsymbol{\Omega}_S=(\mathbf{a}\times\dot{\mathbf{x}})/2c^2$ \\
Scalar spin-orbit & so & $-\mathbf{S}\cdot(\nabla V_\text{sc}\times\hat{\boldsymbol{\pi}})/2m_0^2c^2$ & $\boldsymbol{\Omega}_S=-(\nabla V_\text{sc}\times\boldsymbol{\pi})/2m_0^2c^2$ \\
\bottomrule
\end{tabular}
\end{table}

The Newtonian and frame-dragging rows of Table~\ref{tab:forces} have the form of electromagnetic rows with the charge $q$ replaced by the rest mass $m_0$, $\Phi$ by $\Phi_g$, and $\mathbf{A}$ by $\mathbf{A}_g$: Newtonian gravity pairs with the Coulomb force, and frame dragging with the magnetic Lorentz force and induction. The nonlinear-gravity and velocity-dependent-gravity rows have no electromagnetic row of the same form. The orbital couplings follow the same replacement. In uniform fields, with $\mathbf{A}=\tfrac{1}{2}\,\mathbf{B}\times\mathbf{x}$ and $\mathbf{A}_g=\tfrac{1}{2}\,\mathbf{B}_g\times\mathbf{x}$, the terms linear in the potentials are $-(q/2m_0)\,\mathbf{B}\cdot\mathbf{L}$ and $-\tfrac{1}{2}\,\mathbf{B}_g\cdot\mathbf{L}$. The spin couplings do not follow the replacement. The Pauli term is $-(q/2m_0)\,\mathbf{B}\cdot2\mathbf{S}$, with $g_s=2$ (Eq.~\eqref{eq:fueter_square}), while the spin-rotation term is $-\tfrac{1}{2}\,\mathbf{B}_g\cdot\mathbf{S}$, with the gravitomagnetic g-factor $1$ (Section~\ref{sec:spinor}). In a magnetic field the spin therefore precesses at $-(q/m_0)\,\mathbf{B}$, twice the Larmor rate $-(q/2m_0)\,\mathbf{B}$ of the orbit, while in a gravitomagnetic field spin and orbit both precess at $-\tfrac{1}{2}\,\mathbf{B}_g$. The replacement takes the spin-orbit term to $\mathbf{S}\cdot(\nabla\Phi_g\times\hat{\boldsymbol{\pi}})/2m_0c^2$, with the coefficient $\tfrac{1}{2}$ where the geodetic row has $\tfrac{1}{2}+\gamma_\text{PPN}$.

Each inertial row is a gravitational row with the potentials of an accelerating or a rotating frame. With $\Phi_g=\mathbf{a}\cdot\mathbf{x}$ the Newtonian row is the inertial reaction. With $\mathbf{A}_g=\boldsymbol{\Omega}\times\mathbf{x}$, whose curl is $\mathbf{B}_g=2\boldsymbol{\Omega}$, the frame-dragging force $m_0\,\dot{\mathbf{x}}\times\mathbf{B}_g-m_0\,\partial_t\mathbf{A}_g$ is the Coriolis force $-2m_0\boldsymbol{\Omega}\times\dot{\mathbf{x}}$ plus the Euler force $-m_0\dot{\boldsymbol{\Omega}}\times\mathbf{x}$, and the spin-rotation term is $-\boldsymbol{\Omega}\cdot\mathbf{S}$. The centrifugal force comes from the shift part of $\alpha$ (Eq.~\eqref{eq:alpha_lapse_shift}), which Eq.~\eqref{eq:nr_expansion} contains in its kinetic and gravitomagnetic terms rather than in $mc^2\alpha_n$. With $\boldsymbol{\pi}=m_0(\dot{\mathbf{x}}+\mathbf{A}_g)$ these two terms sum to $\tfrac{1}{2}m_0|\dot{\mathbf{x}}|^2-\tfrac{1}{2}m_0|\mathbf{A}_g|^2$. The second term is the part of order $c^0$ that the shift contributes to $m_0c^2\alpha$, and its negative gradient is $-m_0\boldsymbol{\Omega}\times(\boldsymbol{\Omega}\times\mathbf{x})$. With $\Phi_g=\mathbf{a}\cdot\mathbf{x}$ and $\gamma_\text{PPN}=0$ the geodetic row is the inertial spin-orbit row. The accelerating frame has flat space, $W=1$, and $\gamma_\text{PPN}$ enters only through $W$, so the frame has the part $\tfrac{1}{2}$ of the geodetic coefficient $\tfrac{1}{2}+\gamma_\text{PPN}$ and lacks the part $\gamma_\text{PPN}$, which comes from the curvature of space.
% Checked in verification/sgm_check.py, the blocks for eq:nr_expansion and tab:forces:
% B_g = 2 Omega, the Coriolis, Euler and centrifugal forces, and the geodetic row at gamma = 0.

%-------------------------------------------------------------------------------
\subsection{Maxwell's equations}
\label{sec:maxwell}
%-------------------------------------------------------------------------------

The Coulomb, induction, and magnetic Lorentz forces of Table~\ref{tab:forces} together are the Lorentz force $q(\mathbf{E}+\dot{\mathbf{x}}\times\mathbf{B})$ of Eq.~\eqref{eq:lorentz}, and its fields are the commutators of Eqs.~\eqref{eq:phase_curvature} and~\eqref{eq:time_space_curvature}. Let the index $a$ be $t$ or $1,2,3$. Since $D_a$ commutes with right multiplication by $e_3$, for a real function $f$ the commutator of $D_a$ with the operator $\psi\mapsto f\,\psi\,e_3$ is $\psi\mapsto(\partial_af)\,\psi\,e_3$. Each of the commutators in Eqs.~\eqref{eq:phase_curvature} and~\eqref{eq:time_space_curvature} is an operator of this form. Any three operators obey the Jacobi identity
\begin{equation}
    [D_a,[D_b,D_c]] + [D_b,[D_c,D_a]] + [D_c,[D_a,D_b]] = 0 .
    \label{eq:jacobi}
\end{equation}
For $a,b,c=1,2,3$ its left side applied to $\psi$ is $(q/\hbar)(\nabla\cdot\mathbf{B})\,\psi\,e_3$. For $a,b,c=t,1,2$ and $V=q\Phi$ it is $(q/\hbar)\big[\partial_tB_3+(\nabla\times\mathbf{E})_3\big]\,\psi\,e_3$, and $t,2,3$ and $t,3,1$ give the other two components. Since the left side of Eq.~\eqref{eq:jacobi} is zero, for $q\neq0$
\begin{align}
    \nabla\cdot\mathbf{B} &= 0 ,
    \label{eq:no_monopole} \\
    \nabla\times\mathbf{E}+\partial_t\mathbf{B} &= \mathbf{0} ,
    \label{eq:faraday}
\end{align}
the magnetic Gauss law and Faraday's law. They hold for any potentials $\Phi$ and $\mathbf{A}$: they are the Bianchi identity of the connection in Eqs.~\eqref{eq:phase_connection_time} and~\eqref{eq:right_connection}.

The other two of Maxwell's equations, here in flat spacetime and inertial coordinates, contain the charge density $\rho_q$ and the current density $\mathbf{J}_q$, the subscript $q$ for charge. The vacuum permittivity and permeability are $\varepsilon_0$ and $\mu_0=1/(\varepsilon_0c^2)$. A second imaginary unit $\mathsf{i}$, with $\mathsf{i}^2=-1$, commutes with $e_1$, $e_2$, and $e_3$. Form the biquaternion
\begin{equation}
    \mathcal{F} = \mathbf{E}\cdot\mathbf{e} + \mathsf{i}\,c\,\mathbf{B}\cdot\mathbf{e} .
    \label{eq:rs_biquaternion}
\end{equation}
It is $\sqrt{2/\varepsilon_0}$ times the vacuum \emph{Riemann--Silberstein vector} $\sqrt{\varepsilon_0/2}\,\mathbf{E}+\mathsf{i}\,\mathbf{B}/\sqrt{2\mu_0}$ written as a quaternion, since $\sqrt{2/\varepsilon_0}/\sqrt{2\mu_0}=c$ \citep{BialynickiBirula1996Photon}. For a real vector field $\mathbf{w}$, Eq.~\eqref{eq:quaternion_units} gives $(\mathbf{e}\cdot\nabla)(\mathbf{w}\cdot\mathbf{e})=-\nabla\cdot\mathbf{w}+(\nabla\times\mathbf{w})\cdot\mathbf{e}$: the scalar part is minus the divergence and the vector part is the curl. Applied to $\mathcal{F}$, with $\mathsf{i}^2=-1$, the operator $\mathbf{e}\cdot\nabla-(\mathsf{i}/c)\,\partial_t$ therefore gives the scalar part $-\nabla\cdot\mathbf{E}-\mathsf{i}\,c\,\nabla\cdot\mathbf{B}$ and the vector part $(\nabla\times\mathbf{E}+\partial_t\mathbf{B})\cdot\mathbf{e}+\mathsf{i}\,c\,(\nabla\times\mathbf{B}-c^{-2}\,\partial_t\mathbf{E})\cdot\mathbf{e}$. By Eqs.~\eqref{eq:no_monopole} and~\eqref{eq:faraday}, Gauss's law $\nabla\cdot\mathbf{E}=\rho_q/\varepsilon_0$, and the Amp\`ere--Maxwell law $\nabla\times\mathbf{B}-c^{-2}\,\partial_t\mathbf{E}=\mu_0\mathbf{J}_q$, the two parts equal $-\rho_q/\varepsilon_0$ and $\mathsf{i}\,\mu_0c\,\mathbf{J}_q\cdot\mathbf{e}$:
\begin{equation}
    \Big(\mathbf{e}\cdot\nabla-\frac{\mathsf{i}}{c}\,\partial_t\Big)\,\mathcal{F} = -\frac{\rho_q}{\varepsilon_0} + \mathsf{i}\,\mu_0c\,\mathbf{J}_q\cdot\mathbf{e} .
    \label{eq:maxwell_biquaternion}
\end{equation}
Conversely, since $\mathbf{E}$, $\mathbf{B}$, $\rho_q$, and $\mathbf{J}_q$ are real, the terms of the scalar part without $\mathsf{i}$ give Gauss's law and the coefficients of $\mathsf{i}$ give Eq.~\eqref{eq:no_monopole}, and in the vector part the terms without $\mathsf{i}$ give Eq.~\eqref{eq:faraday} and the coefficients of $\mathsf{i}$ give the Amp\`ere--Maxwell law. All four of Maxwell's equations are thus the four parts of one biquaternion equation. Silberstein wrote the two curl equations of the vacuum as one equation for the complex vector $\mathbf{E}+\mathsf{i}\mathbf{B}$ in units with $c=1$, with the divergence of that vector zero \citep{Silberstein1907Elektromagnetische}, and the four equations with sources as one quaternion equation \citep{Silberstein1912Quaternionic}.

For a free particle of zero mass in flat spacetime, $\mathsf{N}\psi=c^{-1}(\partial_t\psi)\,e_3$ and $\mathsf{X}=\mathbf{e}\cdot\nabla$ (Eq.~\eqref{eq:NX_operators}). Multiplication by $i$ is right multiplication by $-e_3$, so the operator of Eq.~\eqref{eq:weyl_pair_R} is $\mathsf{N}+\mathsf{X}=\mathbf{e}\cdot\nabla-(i/c)\,\partial_t$. This is the operator of Eq.~\eqref{eq:maxwell_biquaternion} with $i$ in place of $\mathsf{i}$. Both units commute with left multiplication by $e_k$, so the Weyl equation $(\mathsf{N}+\mathsf{X})\,\varphi_R=0$ and Maxwell's equations in vacuum differ in the object the operator acts on: a spinor, whose phase is a right multiplication, or the field $\mathcal{F}$. In matrix form the same comparison replaces the Pauli matrices of the Weyl equation with the spin-$1$ matrices of Maxwell's equations \citep{BialynickiBirula1996Photon}.

The conjugate operator $\mathbf{e}\cdot\nabla+(\mathsf{i}/c)\,\partial_t$ commutes with that of Eq.~\eqref{eq:maxwell_biquaternion}, and by Eq.~\eqref{eq:quaternion_units} the product of the two is $-\nabla^2+c^{-2}\partial_t^2$. Applying the conjugate operator to Eq.~\eqref{eq:maxwell_biquaternion} gives $(-\nabla^2+c^{-2}\partial_t^2)\,\mathcal{F}$ on the left, which has no scalar part, and the scalar part $-\mathsf{i}\,\mu_0c\,(\partial_t\rho_q+\nabla\cdot\mathbf{J}_q)$ on the right. The equation is therefore consistent only if
\begin{equation}
    \partial_t\rho_q+\nabla\cdot\mathbf{J}_q = 0 .
    \label{eq:charge_conservation}
\end{equation}
In vacuum each component of $\mathcal{F}$ obeys the wave equation with speed $c$, Eq.~\eqref{eq:klein_gordon} with $\kappa=0$. In the Lorenz gauge, $\nabla\cdot\mathbf{A}+c^{-2}\partial_t\Phi=0$, the conjugate operator applied to the potential biquaternion $-\Phi+\mathsf{i}\,c\,\mathbf{A}\cdot\mathbf{e}$ gives $\mathcal{F}$, so Eq.~\eqref{eq:maxwell_biquaternion} becomes $(-\nabla^2+c^{-2}\partial_t^2)(-\Phi+\mathsf{i}\,c\,\mathbf{A}\cdot\mathbf{e})=-\rho_q/\varepsilon_0+\mathsf{i}\,\mu_0c\,\mathbf{J}_q\cdot\mathbf{e}$. As in Eq.~\eqref{eq:pure_form}, the wave operator is a product of two first-order operators, here applied to the potentials that enter Eqs.~\eqref{eq:phase_connection_time} and~\eqref{eq:right_connection}.

The Dirac equation does not produce Gauss's law or the Amp\`ere--Maxwell law, which are the field's own equations. It supplies the source they require, since its density is conserved (Section~\ref{sec:wave_equation}), and $q$ times that density and its current satisfy Eq.~\eqref{eq:charge_conservation}.
% Checked in verification/sgm_check.py, the blocks for eq:jacobi to eq:charge_conservation:
% the Jacobi identity, the four parts of eq:maxwell_biquaternion, and the
% Riemann--Silberstein factors.

%-------------------------------------------------------------------------------
\subsection{Route IV for the Dirac equation, tensor rank, and spin}
\label{sec:dirac}
%-------------------------------------------------------------------------------

Route~IV (Section~\ref{sec:route_schrodinger}) holds line by line with the Dirac Hamiltonian \citep{Dirac1928Quantum} in place of the Schr\"odinger Hamiltonian of Eq.~\eqref{eq:schrodinger_hamiltonian}. Its kinetic term is linear in $\hat{\mathbf{p}}$ with constant matrix coefficients in an inertial frame, so it still drops out of the commutator, and its position-dependent terms multiply the spinor, so they obey the product rule of Eq.~\eqref{eq:commutator_V}, even though they are matrix valued.

The nonrelativistic reduction removes both factors of the inertial gradient, $\omega_0$ and $a/c^2$ (Section~\ref{sec:duality}), and the Dirac Hamiltonian keeps both. Its mass term contains the rest frequency $\omega_0$. The Hamiltonian contains $c$ at its finite value, so $a/c^2$, the fractional frequency shift per unit separation, stays nonzero, and its evolution kernel vanishes outside the light cone, as Axiom~\ref{ax:causal_propagation} requires \citep{Thaller1992Dirac}.

Hehl and Ni's Hamiltonian for a Dirac particle in a uniformly accelerated frame collects the rest-energy terms as $\beta \, mc^2 \, (1+\mathbf{a}\cdot\delta\mathbf{x}/c^2)$, the rest energy times the Rindler factor of Eq.~\eqref{eq:alpha_rindler}, where $\beta$ is the Dirac matrix that multiplies $mc^2$ in the Dirac Hamiltonian and separates the positive- and negative-energy branches \citep{HehlNi1990Inertial}. This term is the even part of the Hamiltonian before any Foldy--Wouthuysen expansion, and it is exact. On the branch where $\beta$ has the eigenvalue $+1$, this term divided by $\hbar$ is the frequency field $\omega_0 \, \alpha$ of Eq.~\eqref{eq:rindler_rate}, and Eq.~\eqref{eq:sgm} applied to its gradient gives the inertial reaction (Eq.~\eqref{eq:inertia}). In a static metric Eq.~\eqref{eq:master} gives the rest term the factor $\alpha=\sqrt{-g_{00}}$ (Eq.~\eqref{eq:static_master}), so the multiplicative form $\omega=\omega_0 \, \alpha$ that Sections~\ref{sec:gravity} and~\ref{sec:inertia} apply to the frequency field of Axiom~\ref{ax:local_frequency} follows from the Dirac equation.

For localized quantum systems on arbitrary timelike trajectories in curved spacetime, the same multiplicative form is the starting point: written in Fermi normal coordinates, the system's Hamiltonian is multiplied by the local redshift factor as the prescription for its time evolution, and the leading corrections, the inertial and tidal potentials, then arise as the product of the rest-mass term with that factor \citep{Perche2022Localized}.

The interactions differ in the tensor rank of the field that mediates them: rank $1$ for the electromagnetic four-potential, whose time component is $-\Phi/c$ (Remark~\ref{rem:additivity}) and whose space part is $\mathbf{A}$, and rank $2$ for the metric, which enters through the lapse $\alpha_n$ and through the $g_{0i}$ that supply $\mathbf{A}_g$. The frequency has no free index. A field of any rank enters Eq.~\eqref{eq:master} only as a scalar, the potential $V$, the lapse $\alpha_n$ or the scale factor $W$, through a scalar contraction with the particle's state, $\boldsymbol{\pi}^2$, $\mathbf{A}_g\cdot\boldsymbol{\pi}$, $\boldsymbol{\sigma}\cdot\mathbf{B}$, or $\boldsymbol{\sigma}\cdot\mathbf{B}_g$, or through the commutators of these that make up $\mathsf{C}$ (Section~\ref{sec:master}), and combines with $\omega_0$ additively or multiplicatively. The rank of the field equals the rank of its source, the four-current for a rank-$1$ field and the stress-energy tensor for a rank-$2$ field. The rank does not change the rate $-\hbar\nabla\omega$ of Theorem~\ref{thm:universal}.

The identity uses only $H=\hbar\omega$ and Hamilton's equation, so it holds at every spin. At leading WKB order the part of Eq.~\eqref{eq:master} without $\boldsymbol{\sigma}$ is $\hbar\omega=V+\alpha_n\sqrt{(mc^2)^2+c^2W^{-2}\boldsymbol{\pi}^2}-\mathbf{A}_g\cdot\boldsymbol{\pi}$, the relativistic dispersion relation with minimal coupling in the metric, the same at every spin. A particle of spin $0$ obeys Eq.~\eqref{eq:klein_gordon} with each derivative replaced by its covariant derivative. In a nonuniform background its operators do not commute, so an exact positive-frequency reduction of that equation needs a separate treatment. Light obeys Eq.~\eqref{eq:eikonal_ray} with its own dispersion relation.

The two terms in $\boldsymbol{\sigma}$, the magnetic term with $g_s=2$ and the rotation term $-\tfrac{1}{4}\,\hbar\,\boldsymbol{\sigma}\cdot\mathbf{B}_g$, are those of spin $\tfrac{1}{2}$, since Eq.~\eqref{eq:pure_form} is an equation for one quaternion, whose four real numbers are the two complex components of a spin-$\tfrac{1}{2}$ state (Section~\ref{sec:spinor}). A rotating frame couples to the total angular momentum (Section~\ref{sec:spinor}), so the rotation term keeps the form $-\tfrac{1}{2}\,\mathbf{B}_g\cdot\mathbf{S}$ at every spin, where $\mathbf{S}$ is the spin of the particle. In a frame rotating at $\boldsymbol{\Omega}$ the shift is $\boldsymbol{\Omega}\times\mathbf{x}$, whose curl is $\mathbf{B}_g=2\boldsymbol{\Omega}$, so the rotation term is $-\boldsymbol{\Omega}\cdot\mathbf{S}$ and the shift term $-\mathbf{A}_g\cdot\boldsymbol{\pi}=-\tfrac{1}{2}\,\mathbf{B}_g\cdot\mathbf{L}$ is $-\boldsymbol{\Omega}\cdot\mathbf{L}$ (Section~\ref{sec:expansion_orders}). An observer rotating at the angular velocity $\boldsymbol{\Omega}$ therefore finds each energy of a Dirac particle lowered by $\hbar M_\Omega\,\Omega$ before multiplication by the observer's time-dilation factor; \citet{Mashhoon1998Helicity} finds the same for scalar and vector fields. Here $\Omega=|\boldsymbol{\Omega}|$, and $\hbar M_\Omega$ is the component of the total angular momentum $\mathbf{L}+\mathbf{S}$ along $\boldsymbol{\Omega}$. Receivers of the Global Positioning System record the case of circularly polarized light, whose component is $\pm\hbar$, as phase wrap-up \citep{Ashby2003Relativity}.
% spin scope: identity at every spin; sigma terms spin-1/2. Checked in verification/sgm_check.py,
% the block for this paragraph: a rotating observer lowers each energy by hbar M Omega.

%===============================================================================
\section{Lagrangian, Hamiltonian, and Noether structure}
\label{sec:canonical}
%===============================================================================

%-------------------------------------------------------------------------------
\subsection{Lagrangian form}
\label{sec:lagrangian}
%-------------------------------------------------------------------------------

Route~III (Section~\ref{sec:route_lagrangian}) holds for any Lagrangian related to $\hbar\omega$ by a Legendre transform. At Newtonian order, with constant rest mass and the spin terms left out, the Lagrangian of the particle treated here, the Legendre transform of $\hbar$ times the frequency of Eq.~\eqref{eq:nr_expansion}, is
\begin{equation}
    \begin{gathered}
    L = \tfrac{1}{2}m|\dot{\mathbf{x}}|^2 - U(\mathbf{x},\dot{\mathbf{x}}) , \\[-3pt]
    \text{\small\itshape (The slow-particle Lagrangian)}
    \end{gathered}
    \label{eq:sgm_lagrangian}
\end{equation}
where the \emph{generalized potential} \citep{Goldstein2002Classical} is $U(\mathbf{x},\dot{\mathbf{x}})=mc^2(\alpha-1)+q\Phi-q\,\dot{\mathbf{x}}\cdot\mathbf{A}-m\,\dot{\mathbf{x}}\cdot\mathbf{A}_g$. Its first term, $\hbar\omega_0(\alpha-1)$, is $\hbar$ times the excess of the time-dilated rest frequency $\omega_0\,\alpha$ of Eq.~\eqref{eq:omega_alpha} over the rest frequency; to first order in $\Phi_g/c^2$ it is the gravitational potential energy $m \, \Phi_g$ (Section~\ref{sec:gravity}). The next two are the scalar and velocity-dependent potential energies of Eq.~\eqref{eq:omega_em} with $\mathbf{v}=\dot{\mathbf{x}}$. The last comes from $-\mathbf{A}_g\cdot\boldsymbol{\pi}$, the Weyl symbol of the gravitomagnetic term of Eq.~\eqref{eq:nr_expansion}.

The shift term follows from the Legendre transform. Hamilton's equation $\dot{x}_j=\partial H/\partial p_j$ (Axiom~\ref{ax:hamilton}) applied to the kinetic and gravitomagnetic terms of Eq.~\eqref{eq:nr_expansion} gives $\dot{\mathbf{x}}=\alpha_nW^{-2}\boldsymbol{\pi}/m-\mathbf{A}_g$. The inverted transform $L=\mathbf{p}\cdot\dot{\mathbf{x}}-H$ of Route~III then has the kinetic term $(mW^2/2\alpha_n)\,|\dot{\mathbf{x}}+\mathbf{A}_g|^2$ in place of $\tfrac{1}{2}m|\dot{\mathbf{x}}|^2$. Its term $(mW^2/2\alpha_n)\,|\mathbf{A}_g|^2$ combines with $-mc^2\alpha_n$ into $-mc^2\alpha$ by Eq.~\eqref{eq:alpha_lapse_shift}, up to terms of order $|\mathbf{A}_g|^4/c^4$. Up to relative corrections of order $|W^2/\alpha_n-1|$, its cross term is $m\,\dot{\mathbf{x}}\cdot\mathbf{A}_g$ and its remaining term is $\tfrac{1}{2}m|\dot{\mathbf{x}}|^2$.

Equation~\eqref{eq:legendre_bridge} holds for a velocity-independent and a velocity-dependent $U$ alike, so the two differ only in the canonical momentum. Without the shift and the vector potential, and at this order, the kinetic energy $|\mathbf{p}|^2/2m$ has no explicit position dependence, so $U(\mathbf{x})=\hbar\omega-\hbar\omega_0-|\mathbf{p}|^2/2m$ and $\nabla U=\hbar\nabla\omega$. For a velocity-independent $U$ the canonical momentum is $m\dot{\mathbf{x}}$, and Eq.~\eqref{eq:euler_lagrange} gives $m\ddot{\mathbf{x}}=-\nabla U$, which by $\nabla U=\hbar\nabla\omega$ is
\begin{equation}
    m\ddot{\mathbf{x}} = -\hbar\nabla\omega ,
    \label{eq:el_sgm}
\end{equation}
the identity of Eq.~\eqref{eq:sgm} with $\mathbf{F}=m\ddot{\mathbf{x}}$. When $U$ contains the velocity-dependent term $-q\,\dot{\mathbf{x}}\cdot\mathbf{A}$, the canonical momentum $\partial L/\partial\dot{\mathbf{x}}$ is $m\dot{\mathbf{x}}+q\mathbf{A}$, and the left side of Eq.~\eqref{eq:euler_lagrange} contains the extra term $q\dot{\mathbf{A}}=q[\partial_t\mathbf{A}+(\dot{\mathbf{x}}\cdot\nabla)\mathbf{A}]$, the rate of change of $q\mathbf{A}$ along the path. Moving that term to the right side gives, without the shift, $m\ddot{\mathbf{x}}=-mc^2\nabla\alpha+q(\mathbf{E}+\dot{\mathbf{x}}\times\mathbf{B})$, the gravitational force and the Lorentz force (Section~\ref{sec:em}). The shift term has the same form with $q\mathbf{A}$ replaced by $m\mathbf{A}_g$, so the same steps give the force $m\,\dot{\mathbf{x}}\times\mathbf{B}_g-m\,\partial_t\mathbf{A}_g$.

In coordinates $\mathbf{x}$ that rotate at a constant angular velocity $\boldsymbol{\Omega}$, the inertial position of a point is $\mathbf{x}$ turned through the angle $\Omega t$ about $\boldsymbol{\Omega}$. Its inertial displacement is $d\mathbf{x}+\boldsymbol{\Omega}\times\mathbf{x}\,dt$ turned through the same angle, which has the same length. The flat metric therefore takes the form of Eq.~\eqref{eq:lapse_shift_metric} with $\alpha_n=1$ and $\mathbf{A}_g=\boldsymbol{\Omega}\times\mathbf{x}$, whose curl is $2\boldsymbol{\Omega}$ for constant $\boldsymbol{\Omega}$. The shift term then gives the Coriolis force $-2m\,\boldsymbol{\Omega}\times\dot{\mathbf{x}}$, and the factor $\alpha=\sqrt{1-|\boldsymbol{\Omega}\times\mathbf{x}|^2/c^2}$ in the first term of $U$ gives the centrifugal force $-m\,\boldsymbol{\Omega}\times(\boldsymbol{\Omega}\times\mathbf{x})$. These are the Coriolis and centrifugal forces of Table~\ref{tab:forces}.

Substituting $\mathbf{p}=\hbar\mathbf{k}=\hbar\nabla\theta$ and $H=\hbar\omega=-\hbar\,\partial_t\theta$ (Axioms~\ref{ax:planck} and~\ref{ax:debroglie}) into the inverse Legendre transform $L=\mathbf{p}\cdot\dot{\mathbf{x}}-H$ gives $L=\hbar(\mathbf{k}\cdot\dot{\mathbf{x}}-\omega)$. The particle's own phase advances along its path at the rate $d\theta/dt=\partial_t\theta+\dot{\mathbf{x}}\cdot\nabla\theta=-\omega+\mathbf{k}\cdot\dot{\mathbf{x}}$, the same quantity, so the Lagrangian is $\hbar$ times that rate,
\begin{equation}
    L = \hbar\,\frac{d\theta}{dt} ,
    \label{eq:L_phase_rate}
\end{equation}
an energy, and the action $S=\int L\,dt$ along the path is $\hbar$ times the phase accumulated along it,
\begin{equation}
    S = \hbar\,\Delta\theta .
    \label{eq:action_phase}
\end{equation}
Equation~\eqref{eq:L_phase_rate} holds along the classical path, the one on which Eq.~\eqref{eq:euler_lagrange} holds, where $\hbar\theta$ is Hamilton's principal function; on a neighboring path the action is the integral of $L$ along that path. Equation~\eqref{eq:sgm_lagrangian} measures $U$ from the rest energy and so omits the constant $-\hbar\omega_0$ that Eq.~\eqref{eq:L_phase_rate} keeps. A constant added to $L$ does not change Eq.~\eqref{eq:euler_lagrange}, but the identification of action with phase needs that constant.

For the free relativistic particle the identification holds exactly. Its phase advances at the rest rate in proper time, $d\theta/d\tau=-\omega_0$ (Eq.~\eqref{eq:rest_freq} and Axiom~\ref{ax:lpi}), so $\hbar\,d\theta/dt=-\hbar\omega_0\,d\tau/dt=-mc^2\sqrt{1-|\dot{\mathbf{x}}|^2/c^2}$, the free relativistic Lagrangian, and $S=-mc^2\!\int d\tau$: the accumulated phase is $\Delta\theta=-\omega_0\,\Delta\tau$.

By Eq.~\eqref{eq:action_phase}, Hamilton's principle of stationary action is a condition on the phase of the wave. Neighboring paths whose accumulated phases differ by much less than a radian add in phase, and paths whose phases differ by many radians cancel, so the paths that contribute lie near the one on which the phase is stationary. Hamilton's principle is therefore Huygens's construction applied to the matter wave \citep{Huygens1690Traite,Schrodinger1926Quantization}.

%-------------------------------------------------------------------------------
\subsection{Hamiltonian form: dual frequency gradients}
\label{sec:hamiltonian}
%-------------------------------------------------------------------------------

With the Hamiltonian $H=\hbar \, \omega(\mathbf{x},\mathbf{p})$ of Eq.~\eqref{eq:H_eq_hw}, Hamilton's canonical equations read
\begin{align}
    \dot{\mathbf{p}} &= -\hbar\,\nabla\omega , \label{eq:ham_p}\\
    \dot{\mathbf{x}} &= +\hbar\,\nabla_{\mathbf{p}}\,\omega , \label{eq:ham_x} \\[-3pt]
    \multispan2{\hfil\text{\small\itshape (Hamilton's equations)}\hfil} \notag
\end{align}
The first is the SGM identity (Eq.~\eqref{eq:sgm}); the second gives the group velocity, $\dot{\mathbf{x}}=c^2\mathbf{p}/E$ for a free relativistic particle with $E=c\sqrt{\mathbf{p}^2+m^2c^2}$. The opposite signs follow from the symplectic form $\sum_j dp_j\wedge dx_j$. The flow $X=(\dot{\mathbf{x}},\dot{\mathbf{p}})$ of $H=\hbar\omega$ is defined by $\big(\sum_j dp_j\wedge dx_j\big)(X,\cdot)=-dH$. The antisymmetry of the form gives $\sum_j(\dot{p}_j\,dx_j-\dot{x}_j\,dp_j)=-dH$, so matching coefficients gives $\dot{p}_j=-\partial H/\partial x_j$ and $\dot{x}_j=+\partial H/\partial p_j$.

%-------------------------------------------------------------------------------
\subsection{Noether: the conservation laws}
\label{sec:conservation}
%-------------------------------------------------------------------------------

Along the flow generated by the Hamiltonian $\hbar\omega$, the energy $E=\hbar\omega$ (Eq.~\eqref{eq:planck}) changes at the total rate $dE/dt=\dot{\mathbf{x}}\cdot\nabla(\hbar\omega)+\dot{\mathbf{p}}\cdot\nabla_{\mathbf{p}}(\hbar\omega)+\hbar\,\partial_t\omega$, where the frequency function may now depend explicitly on time. By Eqs.~\eqref{eq:ham_p} and~\eqref{eq:ham_x} the first two terms are $(\hbar\nabla_{\mathbf{p}}\omega)\cdot(\hbar\nabla\omega)-(\hbar\nabla\omega)\cdot(\hbar\nabla_{\mathbf{p}}\omega)$, which cancel because the dot product is symmetric and the two terms have opposite signs, leaving
\begin{equation}
    \frac{dE}{dt} = +\hbar\,\frac{\partial\omega}{\partial t} ,
    \label{eq:sgm_temporal}
\end{equation}
which relates the rate of the energy to the time derivative of the frequency, as Eq.~\eqref{eq:sgm_as_dot_p} relates the rate of the momentum to the gradient of the frequency. Differentiating the Legendre transform of Section~\ref{sec:route_lagrangian} with respect to time instead of position gives, as in Eq.~\eqref{eq:legendre_chain},
\begin{equation}
    \frac{\partial H}{\partial t}\bigg|_{\mathbf{x},\mathbf{p}} = p_j\,\frac{\partial\dot{x}_j}{\partial t} - \frac{\partial L}{\partial t}\bigg|_{\mathbf{x},\dot{\mathbf{x}}} - \frac{\partial L}{\partial\dot{x}_j}\,\frac{\partial\dot{x}_j}{\partial t} ,
    \label{eq:legendre_time}
\end{equation}
with $j$ summed over the three axes and $\partial\dot{x}_j/\partial t$ taken at fixed $\mathbf{x}$ and $\mathbf{p}$. Its first and third terms cancel by $p_j=\partial L/\partial\dot{x}_j$, leaving $\partial L/\partial t|_{\mathbf{x},\dot{\mathbf{x}}}=-\partial H/\partial t|_{\mathbf{x},\mathbf{p}}$, so $dE/dt=-\partial L/\partial t$, and a Lagrangian with no explicit time dependence conserves the energy.

The phase four-gradient of Eq.~\eqref{eq:phase_fourvector}, multiplied by $\hbar$, is the four-momentum $p_\mu=\hbar\,\partial_\mu\theta=(-E/c,\ \mathbf{p})$ (Eqs.~\eqref{eq:debroglie} and~\eqref{eq:phase_freq}), and with $\partial_0=c^{-1}\partial_t$,
\begin{equation}
    \begin{gathered}
    \frac{dp_\mu}{dt} = -\hbar\,\partial_\mu\omega\big|_{\mathbf{p}} , \\[-3pt]
    \text{\small\itshape (The SGM identity, four-component form)}
    \end{gathered}
    \label{eq:sgm_fourvector}
\end{equation}
where $|_{\mathbf{p}}$ marks that the frequency function $\omega(\mathbf{x},\mathbf{p},t)$ is differentiated in $t$ and $\mathbf{x}$ at fixed momentum, as in Eq.~\eqref{eq:gradient_def}, while in $p_\mu=\hbar\,\partial_\mu\theta$ the derivative acts on the phase field $\theta(\mathbf{x},t)$. The three spatial components are Eq.~\eqref{eq:sgm_as_dot_p} and the time component is Eq.~\eqref{eq:sgm_temporal}, each written, like every rate in this paper, in the coordinates of the chosen frame. The force and the spectral gradient depend on that frame, while the predicted invariants, such as orbits, deflection angles, and redshifts, do not. Neither side is a four-vector: $\omega=-c\,\partial_0\theta$ is $-c$ times the time component of the phase gradient, not a scalar, and $d/dt$ is the rate in coordinate time, so the equation holds component by component in the chosen frame. Equation~\eqref{eq:sgm_fourvector} extends the mixed-derivative statement of Route~II (Section~\ref{sec:route_phase}) to all four components of the phase gradient.

For a frequency that varies with position alone, $\partial_t\omega=0$ and Eq.~\eqref{eq:sgm_temporal} conserves the energy. Integrating the identity of Eq.~\eqref{eq:sgm} along a path from $\mathbf{x}_A$ to $\mathbf{x}_B$ then gives the work,
\begin{align}
    \mathcal{W} &= \int_{\mathbf{x}_A}^{\mathbf{x}_B}\mathbf{F}\cdot d\mathbf{x} \notag\\
    &= -\Delta U ,
    \label{eq:work_energy}
\end{align}
where $U(\mathbf{x})$ is the generalized potential of Eq.~\eqref{eq:sgm_lagrangian} when it does not depend on velocity, so that $\hbar\nabla\omega=\nabla U$ (Section~\ref{sec:lagrangian}) and the integral is $-\Delta U$ by the gradient theorem. Along the motion, Eq.~\eqref{eq:el_sgm} gives $\tfrac{d}{dt}(K+U)=\dot{\mathbf{x}}\cdot(m\ddot{\mathbf{x}}+\hbar\nabla\omega)=0$ with kinetic energy $K=\tfrac{1}{2}m|\dot{\mathbf{x}}|^2$, so $K+U$ is conserved and $\mathcal{W}=\Delta K$, the work-energy theorem. The work is therefore independent of the path: it depends only on the values of the single-valued function $U$ at $\mathbf{x}_A$ and $\mathbf{x}_B$. Integrating Eq.~\eqref{eq:sgm_as_dot_p} from a time $t_1$ to a time $t_2$ gives the change $\hbar\,\Delta\mathbf{k}$ of the canonical momentum $\mathbf{p}=\hbar\mathbf{k}$ (Eq.~\eqref{eq:debroglie_k}), the impulse $\mathbf{J}$,
\begin{align}
    \mathbf{J} &= \int_{t_1}^{t_2}\mathbf{F}\,dt \notag\\
    &= \Delta\mathbf{p} ,
    \label{eq:impulse_momentum}
\end{align}
the impulse-momentum theorem.

Momentum is conserved because the total frequency $\omega_\text{sys}=H_\text{sys}/\hbar$ of an isolated system, Eq.~\eqref{eq:H_eq_hw} applied to the system's Hamiltonian, depends only on relative positions, so its position gradients cancel, $\sum_i\nabla_i\,\omega_\text{sys}=\mathbf{0}$, where $\nabla_i$ is the gradient with respect to the position of particle $i$. Momentum and energy conservation are the Noether consequences of space- and time-translation invariance \citep{Noether1918Invariante,Goldstein2002Classical}.

%-------------------------------------------------------------------------------
\subsection{Rotation: the angular spectral gradient}
\label{sec:angular}
%-------------------------------------------------------------------------------

The spectral gradient of Definition~\ref{def:spectral_gradient} is the rate at which the frequency changes when the state is displaced, with the fields and the momentum held fixed. A rotation of the state by a small angle $\delta\vartheta$ about a unit vector $\hat{\mathbf{n}}$, the rotation of Section~\ref{sec:spinor}, moves the position to $\mathbf{x}+\delta\vartheta\,\hat{\mathbf{n}}\times\mathbf{x}$, turns the momentum to $\mathbf{p}+\delta\vartheta\,\hat{\mathbf{n}}\times\mathbf{p}$, and turns the spin expectation value to $\langle\boldsymbol{\sigma}\rangle+\delta\vartheta\,\hat{\mathbf{n}}\times\langle\boldsymbol{\sigma}\rangle$, the first-order change of a vector under the left multiplication by $\mathcal{R}$ of Section~\ref{sec:spinor}. The fields $V$, $\mathbf{A}$, $\alpha_n$, $\mathbf{A}_g$, $\mathbf{B}$, and $\mathbf{B}_g$ stay where they are. The rate at which the frequency changes per unit angle, $\partial\omega/\partial\vartheta$, with the fields held fixed, is the \emph{angular spectral gradient} about $\hat{\mathbf{n}}$. Unlike a displacement, a rotation turns the momentum with the position, keeping its magnitude $|\mathbf{p}|$ and its angle to $\mathbf{x}$, so the angular gradient is taken at a fixed axis of rotation and a fixed magnitude of the momentum, while the linear gradient is taken at a fixed momentum vector.

For the orbital part, the chain rule gives $\partial\omega/\partial\vartheta=(\hat{\mathbf{n}}\times\mathbf{x})\cdot\nabla\omega+(\hat{\mathbf{n}}\times\mathbf{p})\cdot\nabla_{\mathbf{p}}\omega$. By Eq.~\eqref{eq:ham_p} the first term is $-(\hat{\mathbf{n}}\times\mathbf{x})\cdot\dot{\mathbf{p}}/\hbar$, which by the cyclic property of the scalar triple product is $-\hat{\mathbf{n}}\cdot(\mathbf{x}\times\dot{\mathbf{p}})/\hbar$. By Eq.~\eqref{eq:ham_x} the second term is $(\hat{\mathbf{n}}\times\mathbf{p})\cdot\dot{\mathbf{x}}/\hbar$, which is $\hat{\mathbf{n}}\cdot(\mathbf{p}\times\dot{\mathbf{x}})/\hbar=-\hat{\mathbf{n}}\cdot(\dot{\mathbf{x}}\times\mathbf{p})/\hbar$. The two terms sum to $-\hat{\mathbf{n}}\cdot(\mathbf{x}\times\dot{\mathbf{p}}+\dot{\mathbf{x}}\times\mathbf{p})/\hbar$, and the bracket is the time derivative of $\mathbf{x}\times\mathbf{p}$ by the product rule. Where the vector potential vanishes, $\mathbf{p}=\boldsymbol{\pi}$ and $\mathbf{x}\times\mathbf{p}$ is the orbital angular momentum $\mathbf{L}=\mathbf{x}\times\boldsymbol{\pi}$ of Section~\ref{sec:spinor}, so $\hbar\,\partial\omega/\partial\vartheta=-\hat{\mathbf{n}}\cdot d\mathbf{L}/dt$. When in addition $\mathbf{p}=m\dot{\mathbf{x}}$, the product $\dot{\mathbf{x}}\times\mathbf{p}$ vanishes and $-\hbar$ times the first term alone is $\hat{\mathbf{n}}\cdot(\mathbf{x}\times(-\hbar\nabla\omega))$: the orbital angular gradient is the component along $\hat{\mathbf{n}}$ of the moment about the origin of the linear spectral gradient, and $-\hbar$ times it is the component along $\hat{\mathbf{n}}$ of the torque of the force $-\hbar\nabla\omega$ of Eq.~\eqref{eq:sgm}.

For the spin part, the Pauli and spin-rotation terms of Eq.~\eqref{eq:nr_expansion} are $\mathbf{b}\cdot\boldsymbol{\sigma}$. Here the vector $\mathbf{b}=-(\alpha_nq\hbar/2mW^2)\,\mathbf{B}-\tfrac{1}{4}\hbar\,\mathbf{B}_g$ has the dimension of energy and is taken to vary little across the state. By the product rule of the Pauli matrices (Section~\ref{sec:spinor}), $[\mathbf{b}\cdot\boldsymbol{\sigma},\sigma_k]=b_j(\sigma_j\sigma_k-\sigma_k\sigma_j)=2i\,\epsilon_{jkl}\,b_j\,\sigma_l$, summed over $j$ and $l$, and the terms of the Hamiltonian without $\boldsymbol{\sigma}$ commute with $\sigma_k$. Equation~\eqref{eq:ehrenfest_commutator} with $\hat{A}=\sigma_k$ then gives $d\langle\sigma_k\rangle/dt=(i/\hbar)\,2i\,\epsilon_{jkl}\,b_j\langle\sigma_l\rangle=-(2/\hbar)\,\epsilon_{jkl}\,b_j\langle\sigma_l\rangle$, and since $\epsilon_{jkl}=-\epsilon_{kjl}$ this is $(2/\hbar)\,(\mathbf{b}\times\langle\boldsymbol{\sigma}\rangle)_k$. With the expectation value $\mathbf{S}=\tfrac{1}{2}\hbar\langle\boldsymbol{\sigma}\rangle$ of the spin (Section~\ref{sec:spinor}), the Pauli and spin-rotation terms contribute $\mathbf{b}\times\langle\boldsymbol{\sigma}\rangle$ to $d\mathbf{S}/dt$. The rotation of the state changes the expectation value of $\mathbf{b}\cdot\boldsymbol{\sigma}$ by $\delta\vartheta\,\mathbf{b}\cdot(\hat{\mathbf{n}}\times\langle\boldsymbol{\sigma}\rangle)$, which by the cyclic property is $\delta\vartheta\,\hat{\mathbf{n}}\cdot(\langle\boldsymbol{\sigma}\rangle\times\mathbf{b})=-\delta\vartheta\,\hat{\mathbf{n}}\cdot d\mathbf{S}/dt$, so the spin part of $\hbar\,\partial\omega/\partial\vartheta$ is $-\hat{\mathbf{n}}\cdot d\mathbf{S}/dt$, which adds to the orbital part.

For the whole frequency operator, let $\hat{\mathbf{L}}\equiv\mathbf{x}\times\hat{\mathbf{p}}$, which equals $\mathbf{L}$ where the vector potential vanishes, and let $\mathsf{T}_\vartheta\equiv\exp[-i\vartheta\,\hat{\mathbf{n}}\cdot(\hat{\mathbf{L}}+\tfrac{1}{2}\hbar\boldsymbol{\sigma})/\hbar]$, which is unitary in the unweighted product and rotates $\mathbf{x}$, $\hat{\mathbf{p}}$, and $\boldsymbol{\sigma}$ by the angle $\vartheta$ about $\hat{\mathbf{n}}$. Let $\partial\hat{\omega}/\partial\vartheta$ be the derivative at $\vartheta=0$ of $\mathsf{T}_\vartheta^{-1}\,\hat{\omega}\,\mathsf{T}_\vartheta$, which is $(i/\hbar)\,[\hat{\mathbf{n}}\cdot(\hat{\mathbf{L}}+\tfrac{1}{2}\hbar\boldsymbol{\sigma}),\hat{\omega}]$. Where the vector potential vanishes, Eq.~\eqref{eq:ehrenfest_commutator} with $\hat{A}=\hat{\mathbf{n}}\cdot(\hat{\mathbf{L}}+\tfrac{1}{2}\hbar\boldsymbol{\sigma})$ and $\hat{H}=\hbar\hat{\omega}$ then gives, in expectation values (Eq.~\eqref{eq:sgm_expectation}),
\begin{equation}
    \begin{gathered}
    \boxed{\;\hat{\mathbf{n}}\cdot\frac{d}{dt}(\mathbf{L}+\mathbf{S}) = -\hbar\,\frac{\partial\omega}{\partial\vartheta}\;} . \\[-3pt]
    \text{\small\itshape (The SGM identity, angular form)}
    \end{gathered}
    \label{eq:sgm_angular}
\end{equation}
The displacement $\delta\mathbf{x}$ of Eq.~\eqref{eq:sgm_as_dot_p} is replaced by the angle $\delta\vartheta\,\hat{\mathbf{n}}$, and the momentum $\mathbf{p}$ by the total angular momentum $\mathbf{L}+\mathbf{S}$. The total frequency $\omega_\text{sys}$ of an isolated system depends only on the relative positions and the relative orientations of its parts, so a rotation of every part together about $\hat{\mathbf{n}}$ leaves it unchanged. Equation~\eqref{eq:sgm_angular} summed over the parts then conserves the component of the total $\mathbf{L}+\mathbf{S}$ along $\hat{\mathbf{n}}$, the Noether consequence of rotational invariance \citep{Noether1918Invariante,Goldstein2002Classical}, as the conservation of momentum is that of translational invariance (Section~\ref{sec:conservation}).

Section~\ref{sec:inertia} obtained the inertial reaction as $-\hbar$ times the gradient of the rate $\omega_0\,(1+\delta\alpha)$ that the accelerating frame assigns to the phase across the state (Eq.~\eqref{eq:rindler_rate}), with $\delta\alpha=\mathbf{a}\cdot\delta\mathbf{x}/c^2$. A moving or turning frame shifts the frequency by the Doppler effect. For a wave with the phase $\theta=\mathbf{k}\cdot\mathbf{x}-\omega t$, coordinates $\mathbf{x}'=\mathbf{x}-\mathbf{u}t$ that move at the constant velocity $\mathbf{u}$ give $\theta=\mathbf{k}\cdot\mathbf{x}'-(\omega-\mathbf{k}\cdot\mathbf{u})\,t$, so the frequency at fixed $\mathbf{x}'$ is $\omega-\mathbf{k}\cdot\mathbf{u}$, the first-order Doppler shift.

In coordinates that turn at the constant angular velocity $\boldsymbol{\Omega}$ about the axis $\hat{\mathbf{n}}=\boldsymbol{\Omega}/\Omega$ (Section~\ref{sec:lagrangian}), the azimuth about $\hat{\mathbf{n}}$ is $\varphi'=\varphi-\Omega t$, where $\varphi$ is the azimuth in the original coordinates. An orbital state whose component of $\mathbf{L}$ along $\hat{\mathbf{n}}$ is $\hbar m_l$, where $m_l$ is an integer, has the phase $m_l\varphi-\omega t=m_l\varphi'-(\omega-m_l\Omega)\,t$, so its frequency at fixed $\varphi'$ is $\omega-m_l\Omega$. The integer $m_l=\partial\theta/\partial\varphi$ is the phase per radian of azimuth, as the wavevector $\mathbf{k}=\nabla\theta$ is the phase per meter, and $m_l\Omega$ enters the frequency as $\mathbf{k}\cdot\mathbf{u}$ does.

For the spin, the turning coordinates refer the spinor to axes that turn at $\boldsymbol{\Omega}$, so in those coordinates the spinor is turned by the angle $-\Omega t$ about $\hat{\mathbf{n}}$, the left factor $\exp(-\tfrac{1}{2}\Omega t\,\hat{\mathbf{n}}\cdot\mathbf{e})$ of Section~\ref{sec:spinor}. Take the stationary state of Eq.~\eqref{eq:stationary_time} with its spin along $\hat{\mathbf{n}}$, and write its time-independent part in the polar form of Eq.~\eqref{eq:polar_form} as $\psi_0=\rho^{1/2}\mathcal{R}_{\psi_0}$, so that $\psi=\rho^{1/2}\mathcal{R}_{\psi_0}\exp(e_3\,\omega t)$ and $\mathcal{R}_\psi=\mathcal{R}_{\psi_0}\exp(e_3\,\omega t)$. Since $\exp(e_3\,\omega t)$ commutes with $e_3$, Eq.~\eqref{eq:spin_rotor} gives $\hat{\mathbf{n}}\cdot\mathbf{e}=\mathcal{R}_{\psi_0}\,e_3\,\bar{\mathcal{R}}_{\psi_0}$. In each power of $\mathcal{R}_{\psi_0}\,e_3\,\bar{\mathcal{R}}_{\psi_0}$, the factors $\bar{\mathcal{R}}_{\psi_0}\mathcal{R}_{\psi_0}$ between adjacent copies equal $1$. The exponential series therefore gives $\exp(-\tfrac{1}{2}\Omega t\,\hat{\mathbf{n}}\cdot\mathbf{e})\,\mathcal{R}_{\psi_0}=\mathcal{R}_{\psi_0}\exp(-\tfrac{1}{2}\Omega t\,e_3)$. In the turning coordinates the state is
\begin{equation}
    \exp\big(-\tfrac{1}{2}\Omega t\,\hat{\mathbf{n}}\cdot\mathbf{e}\big)\,\psi = \rho^{1/2}\,\mathcal{R}_{\psi_0}\exp\big(e_3\,(\omega-\tfrac{1}{2}\Omega)\,t\big) ,
    \label{eq:spin_doppler}
\end{equation}
which has the frequency $\omega-\tfrac{1}{2}\Omega$. If the spin is against $\hat{\mathbf{n}}$, then $\hat{\mathbf{n}}\cdot\mathbf{e}=-\mathcal{R}_{\psi_0}\,e_3\,\bar{\mathcal{R}}_{\psi_0}$ and the frequency is $\omega+\tfrac{1}{2}\Omega$. The transverse axes of the frame $\mathcal{R}_\psi$ turn about the spin axis at the rate $2\omega$ (Section~\ref{sec:spinor}); in the turning coordinates they turn at $2\omega-\Omega$ when the spin is along $\hat{\mathbf{n}}$, twice the frequency $\omega-\tfrac{1}{2}\Omega$. The spin has the angular wavenumber $\pm\tfrac{1}{2}$, its component along $\hat{\mathbf{n}}$ divided by $\hbar$. Orbital and spin shifts together lower the frequency by $\boldsymbol{\Omega}\cdot(\mathbf{L}+\mathbf{S})/\hbar$. Equation~\eqref{eq:nr_expansion} gives the same term with $\mathbf{A}_g=\boldsymbol{\Omega}\times\mathbf{x}$ and $\mathbf{B}_g=2\boldsymbol{\Omega}$ (Section~\ref{sec:dirac}), so a free particle in the turning coordinates has
\begin{equation}
    \hbar\,\omega = mc^2 + \frac{\boldsymbol{\pi}^2}{2m} - \boldsymbol{\Omega}\cdot(\mathbf{L}+\mathbf{S}) .
    \label{eq:omega_rotating}
\end{equation}
Here $\boldsymbol{\pi}=\mathbf{p}$, since $\mathbf{A}=\mathbf{0}$, and Eq.~\eqref{eq:ham_x} applied to the terms $\boldsymbol{\pi}^2/2m-\mathbf{A}_g\cdot\boldsymbol{\pi}$ gives $\dot{\mathbf{x}}=\boldsymbol{\pi}/m-\mathbf{A}_g$, so $\boldsymbol{\pi}=m(\dot{\mathbf{x}}+\boldsymbol{\Omega}\times\mathbf{x})$, which is $m$ times the velocity in the original coordinates resolved on the turning axes (Section~\ref{sec:lagrangian}), and $\mathbf{L}+\mathbf{S}$ is the angular momentum of the original coordinates resolved on the turning axes.

The kinetic term of Eq.~\eqref{eq:omega_rotating} depends only on $|\boldsymbol{\pi}|$, which a rotation does not change. A rotation by $\delta\vartheta$ about any unit vector $\hat{\mathbf{n}}'$ turns $\mathbf{L}+\mathbf{S}$ by $\delta\vartheta\,\hat{\mathbf{n}}'\times(\mathbf{L}+\mathbf{S})$, so $\hbar\,\partial\omega/\partial\vartheta=-\boldsymbol{\Omega}\cdot(\hat{\mathbf{n}}'\times(\mathbf{L}+\mathbf{S}))$. By the cyclic property this is $-\hat{\mathbf{n}}'\cdot((\mathbf{L}+\mathbf{S})\times\boldsymbol{\Omega})$. Equation~\eqref{eq:sgm_angular} then gives $\hat{\mathbf{n}}'\cdot d(\mathbf{L}+\mathbf{S})/dt=\hat{\mathbf{n}}'\cdot((\mathbf{L}+\mathbf{S})\times\boldsymbol{\Omega})$ for every $\hat{\mathbf{n}}'$, so in the turning coordinates
\begin{equation}
    \frac{d}{dt}(\mathbf{L}+\mathbf{S}) = (\mathbf{L}+\mathbf{S})\times\boldsymbol{\Omega} .
    \label{eq:gyroscopic}
\end{equation}
The right side is the \emph{gyroscopic torque}, $-\hbar$ times the angular spectral gradient of the rotational Doppler shift, as the inertial reaction is $-\hbar$ times the gradient of the rate in the accelerating frame (Section~\ref{sec:inertia}). For the spin it comes from the spin-rotation term alone. A vector that is constant in the original coordinates is seen from axes turning at $\boldsymbol{\Omega}$ to change at the rate $-\boldsymbol{\Omega}\times\mathbf{w}$, for any vector $\mathbf{w}$, so a rate in the original coordinates is the rate in the turning coordinates plus $\boldsymbol{\Omega}\times\mathbf{w}$ \citep{Goldstein2002Classical}. Applied to $\mathbf{L}+\mathbf{S}$ with Eq.~\eqref{eq:gyroscopic}, the rate in the original coordinates is $(\mathbf{L}+\mathbf{S})\times\boldsymbol{\Omega}+\boldsymbol{\Omega}\times(\mathbf{L}+\mathbf{S})=\mathbf{0}$.

A spin held at rest in coordinates that turn at $\boldsymbol{\Omega}$, as by the gimbal of a gyroscope, has $d\mathbf{S}/dt=\mathbf{0}$ in those coordinates. With an applied torque $\mathbf{N}_\text{app}$ \citep{Goldstein2002Classical} added to the gyroscopic torque of Eq.~\eqref{eq:gyroscopic},
\begin{equation}
    \mathbf{N}_\text{app} + \mathbf{N}_\text{in} = \mathbf{0} ,
    \label{eq:torque_balance}
\end{equation}
where $\mathbf{N}_\text{in}=\mathbf{S}\times\boldsymbol{\Omega}$ is the gyroscopic torque. In the original coordinates the spin turns at $\boldsymbol{\Omega}\times\mathbf{S}$, which by Eq.~\eqref{eq:torque_balance} equals the applied torque, the law of the precessing gyroscope. In coordinates that turn with the held spin, the applied and gyroscopic torques cancel (Eq.~\eqref{eq:torque_balance}), as the inertial reaction and the gravitational force cancel in the falling frame (Eq.~\eqref{eq:geodesic_balance}).

The orbital part of the gyroscopic torque is the moment of the linear inertial reactions. An element of mass $m$ at rest at $\mathbf{x}$ in coordinates whose angular velocity $\boldsymbol{\Omega}(t)$ changes at the rate $\dot{\boldsymbol{\Omega}}$ has, in the original coordinates, the acceleration $\mathbf{a}=\dot{\boldsymbol{\Omega}}\times\mathbf{x}+\boldsymbol{\Omega}\times(\boldsymbol{\Omega}\times\mathbf{x})$, and its dynamic spectral gradient is $m\mathbf{a}/\hbar$ (Eq.~\eqref{eq:inertia_gradient}). Its angular gradient about $\hat{\mathbf{n}}$ is the moment of that gradient, $(\hat{\mathbf{n}}\times\mathbf{x})\cdot m\mathbf{a}/\hbar=\hat{\mathbf{n}}\cdot(\mathbf{x}\times m\mathbf{a})/\hbar$ by the cyclic property of the scalar triple product, and summed over the elements of a body, $\hbar$ times it is $\hat{\mathbf{n}}\cdot\sum m\,\mathbf{x}\times\mathbf{a}$. The inertia tensor $\mathcal{I}$ is the linear map $\mathcal{I}\mathbf{w}=\sum m\,\mathbf{x}\times(\mathbf{w}\times\mathbf{x})$ \citep{Goldstein2002Classical}. The first term of $\mathbf{a}$ contributes $\mathcal{I}\dot{\boldsymbol{\Omega}}$, and the second contributes $\sum m\,\mathbf{x}\times(\boldsymbol{\Omega}\times(\boldsymbol{\Omega}\times\mathbf{x}))=\boldsymbol{\Omega}\times\mathcal{I}\boldsymbol{\Omega}$, since $\mathbf{x}\times(\boldsymbol{\Omega}\times\mathbf{w})=\boldsymbol{\Omega}\times(\mathbf{x}\times\mathbf{w})-\mathbf{w}\times(\mathbf{x}\times\boldsymbol{\Omega})$ by the Jacobi identity of the cross product and, with $\mathbf{w}=\boldsymbol{\Omega}\times\mathbf{x}$, the last term vanishes because $\mathbf{x}\times\boldsymbol{\Omega}=-\mathbf{w}$. The inertial torque on the body is therefore $\mathbf{N}_\text{in}=-\mathcal{I}\dot{\boldsymbol{\Omega}}-\boldsymbol{\Omega}\times\mathcal{I}\boldsymbol{\Omega}$, and the balance of Eq.~\eqref{eq:torque_balance} gives
\begin{equation}
    \mathbf{N}_\text{app} = \mathcal{I}\dot{\boldsymbol{\Omega}} + \boldsymbol{\Omega}\times\mathcal{I}\boldsymbol{\Omega} ,
    \label{eq:euler_rigid}
\end{equation}
Euler's equations of a rigid body \citep{Goldstein2002Classical}, for a body that turns with the coordinates.

The body's orbital angular momentum is $\mathbf{L}=\sum m\,\mathbf{x}\times(\boldsymbol{\Omega}\times\mathbf{x})=\mathcal{I}\boldsymbol{\Omega}$, so the second term of $\mathbf{N}_\text{in}$ is $\mathbf{L}\times\boldsymbol{\Omega}$, the orbital part of the gyroscopic torque of Eq.~\eqref{eq:gyroscopic}: the same torque follows from the angular gradient of $-\boldsymbol{\Omega}\cdot\mathbf{L}$ and from the moments of the centrifugal reactions. Of the two reactions $-m\mathbf{a}$, the second, $-m\,\boldsymbol{\Omega}\times(\boldsymbol{\Omega}\times\mathbf{x})$, is the centrifugal force that the factor $\alpha$ of Section~\ref{sec:lagrangian} gives, the gradient of the time-dilation factor. The term $\dot{\boldsymbol{\Omega}}\times\mathbf{x}$ has the curl $2\dot{\boldsymbol{\Omega}}$, as $\boldsymbol{\Omega}\times\mathbf{x}$ has the curl $2\boldsymbol{\Omega}$, so it is not the gradient of any scalar; it is the force $-m\,\partial_t\mathbf{A}_g$ of Section~\ref{sec:lagrangian}, with $\partial_t\mathbf{A}_g=\dot{\boldsymbol{\Omega}}\times\mathbf{x}$.

The spin has no torque proportional to $\dot{\boldsymbol{\Omega}}$. A body that turns with the coordinates has $\mathbf{L}=\mathcal{I}\boldsymbol{\Omega}$, so a change of its rate of rotation requires the torque $\mathcal{I}\dot{\boldsymbol{\Omega}}$. The spin $\mathbf{S}$ is not $\mathcal{I}\boldsymbol{\Omega}$ for any moment of inertia. Under the spin-rotation term it only turns, at the rate $\mathbf{S}\times\boldsymbol{\Omega}$ of Eq.~\eqref{eq:gyroscopic}, which keeps its length. A change of $\boldsymbol{\Omega}$ changes the term $-\boldsymbol{\Omega}\cdot\mathbf{S}$ at the rate $\hbar\,\partial\omega/\partial t=-\dot{\boldsymbol{\Omega}}\cdot\mathbf{S}$, which by Eq.~\eqref{eq:sgm_temporal} changes the energy, while the torque from its angular gradient remains $\mathbf{S}\times\boldsymbol{\Omega}$.

The linear shift $\mathbf{k}\cdot\mathbf{u}$ vanishes for a state at rest, $\mathbf{k}=\mathbf{0}$, and the inertial gradient of Section~\ref{sec:inertia} arises instead from the time dilation of the accelerating frame (Eq.~\eqref{eq:alpha_rindler}). The rotational shift of the spin is nonzero at rest, at first order in $\Omega$, since its angular wavenumber $\pm\tfrac{1}{2}$ is a property of the state, independent of its motion, so the gyroscopic torque on a spin requires no time dilation. A rotation turns $\mathbf{L}+\mathbf{S}$, so $-\boldsymbol{\Omega}\cdot(\mathbf{L}+\mathbf{S})$ has an angular gradient. The linear shift $\mathbf{k}\cdot\mathbf{u}$ has no position gradient, because a translation leaves $\mathbf{k}$ unchanged, and a uniform velocity produces no force.

The angular gradient acts on the phase as the linear gradient does. Along the trajectory $\mathbf{x}(t)$ of Section~\ref{sec:route_phase}, the wavevector $\mathbf{k}=\mathbf{p}/\hbar$ (Eq.~\eqref{eq:debroglie}) changes at the rate $\dot{\mathbf{k}}=-\nabla\omega$ by Eq.~\eqref{eq:sgm_as_dot_p}, with the gradient at fixed momentum: a linear spectral gradient changes the wavevector of the particle. The derivative along the azimuth about $\hat{\mathbf{n}}$ is $\partial/\partial\varphi=(\hat{\mathbf{n}}\times\mathbf{x})\cdot\nabla$, and by the cyclic property $\partial\theta/\partial\varphi=\hat{\mathbf{n}}\cdot(\mathbf{x}\times\nabla\theta)$, which where $\mathbf{A}$ vanishes is $\hat{\mathbf{n}}\cdot(\mathbf{x}\times\mathbf{k})$, the component of the orbital angular momentum along $\hat{\mathbf{n}}$ at the point $\mathbf{x}$, divided by $\hbar$. For a state whose component of $\mathbf{L}$ along $\hat{\mathbf{n}}$ is $\hbar m_l$ it is the integer $m_l$ of the rotational Doppler shift, and otherwise it varies from point to point and need not be an integer. Along the trajectory its rate is $\hat{\mathbf{n}}\cdot(\dot{\mathbf{x}}\times\mathbf{k})+\hat{\mathbf{n}}\cdot(\mathbf{x}\times\dot{\mathbf{k}})$, and by the cyclic property the second term is $-(\hat{\mathbf{n}}\times\mathbf{x})\cdot\nabla\omega$. The first term vanishes when $\dot{\mathbf{x}}$ is parallel to $\mathbf{k}$. This holds where $\mathbf{A}$ and $\mathbf{A}_g$ vanish and the spin terms are left out, since the frequency of Eq.~\eqref{eq:nr_expansion} then depends on the momentum only through $|\mathbf{p}|^2$ and Eq.~\eqref{eq:ham_x} gives $\dot{\mathbf{x}}\propto\mathbf{p}$. A spin-orbit term adds to $\dot{\mathbf{x}}$ the part $\mathbf{S}\times\nabla V/2m^2c^2$, which is not parallel to $\mathbf{p}$. For a frequency of the form $\omega(\mathbf{x},|\mathbf{p}|^2)$,
\begin{equation}
    \frac{d}{dt}\,\frac{\partial\theta}{\partial\varphi} = -\frac{\partial\omega}{\partial\varphi} .
    \label{eq:winding_rate}
\end{equation}
The right side is the moment of the linear spectral gradient, which is the first term of the orbital angular gradient. Along the trajectory through each point of the state where $Q$ is negligible, an angular spectral gradient therefore changes the winding $\partial\theta/\partial\varphi$ of the phase about the axis.

Consider a body whose rate of rotation $\Omega$ about $\hat{\mathbf{n}}$ changes at the rate $\dot\Omega$, and a particle of mass $m$ in it at the distance $r_\perp=|\hat{\mathbf{n}}\times\mathbf{x}|$ from the axis. Its acceleration $\mathbf{a}=\dot\Omega\,\hat{\mathbf{n}}\times\mathbf{x}+\boldsymbol{\Omega}\times(\boldsymbol{\Omega}\times\mathbf{x})$ has the azimuthal component $\dot\Omega\,r_\perp$, since the centripetal term is perpendicular to $\hat{\mathbf{n}}\times\mathbf{x}$. In the original coordinates the net force $m\mathbf{a}$ that the rest of the body and the applied torque exert on the particle is $-\hbar\nabla\omega$ by Eq.~\eqref{eq:sgm}, so $\hbar\,\partial\omega/\partial\varphi=-(\hat{\mathbf{n}}\times\mathbf{x})\cdot m\mathbf{a}=-m\dot\Omega\,r_\perp^2$, and by Eq.~\eqref{eq:winding_rate} the phase of the particle winds at the rate $d(\partial\theta/\partial\varphi)/dt=m\dot\Omega\,r_\perp^2/\hbar$. A change $\delta\Omega$ of the rate of rotation adds $m r_\perp^2\,\delta\Omega/\hbar$ to $\partial\theta/\partial\varphi$, which is the change of the particle's component of $\mathbf{L}$ along $\hat{\mathbf{n}}$, divided by $\hbar$. The definition of $\mathcal{I}$ gives $\hat{\mathbf{n}}\cdot\mathcal{I}\hat{\mathbf{n}}=\sum m\,\hat{\mathbf{n}}\cdot(\mathbf{x}\times(\hat{\mathbf{n}}\times\mathbf{x}))$, and by the cyclic property each term is $m\,|\hat{\mathbf{n}}\times\mathbf{x}|^2=m r_\perp^2$, so $\hbar$ times the winding rate summed over the body is $\dot\Omega\,\hat{\mathbf{n}}\cdot\mathcal{I}\hat{\mathbf{n}}=\hat{\mathbf{n}}\cdot\mathcal{I}\dot{\boldsymbol{\Omega}}$, the component along the axis of the first term of Eq.~\eqref{eq:euler_rigid}. A rotational impulse on a body therefore winds the phase of every particle in it about the axis by the amount that the particle's share of the change of $\mathbf{L}$ requires, as a linear impulse changes the wavevector of every particle by its share of the change of $\mathbf{p}$.

The spins of the particles take no part in this winding. Under the spin-rotation coupling alone, which gives no torque proportional to $\dot{\boldsymbol{\Omega}}$, a spin stays fixed in the original coordinates and turns at $\mathbf{S}\times\boldsymbol{\Omega}$ relative to the body (Eq.~\eqref{eq:gyroscopic}). Seen from the body, the two levels of a spin are shifted by $\mp\tfrac{1}{2}\Omega$ (Eq.~\eqref{eq:spin_doppler}). Its precession frequency, the difference of the two level frequencies, is therefore shifted by $\Omega$, whatever its gyromagnetic ratio. A magnetic field along $\hat{\mathbf{n}}$, of strength $\Omega$ divided by that ratio, would produce the same shift. \citet{Wood2017Magnetic} observed this shift for the $^{13}$C nuclear spins in a rotating diamond. The coupling of the spin of a free particle to rotation has been measured with neutrons, in polarimetry \citep{Demirel2015Measurement} and in interferometry \citep{Danner2020Spin}, and the orbital coupling $-\boldsymbol{\Omega}\cdot\mathbf{L}$ as the phase shift that the rotation of the Earth gives a neutron in an interferometer \citep{Werner1979Effect}.

Within a body the spin and the orbital angular momentum exchange angular momentum only through the terms that couple them, the spin-orbit coupling of Section~\ref{sec:expansion_orders} and the interactions of the spins with the lattice of a solid, and Eq.~\eqref{eq:sgm_angular} conserves their sum. Through these couplings a rotating body becomes magnetized \citep{Barnett1915}, and a change of magnetization makes a body rotate \citep{EinsteinDeHaas1915Experimental}. Table~\ref{tab:linear_angular} compares the linear and angular mechanisms.

\begin{table}[tbp]
\centering
\caption{The linear and angular spectral gradients. The linear column is from Sections~\ref{sec:inertia}, \ref{sec:freefall}, and~\ref{sec:conservation}; the angular column is from Section~\ref{sec:angular}.}
\label{tab:linear_angular}
\begin{tabular}{lll}
\toprule
\textbf{Quantity} & \textbf{Linear} & \textbf{Angular} \\
\midrule
Change of the state & displacement $\delta\mathbf{x}$ & rotation $\delta\vartheta\,\hat{\mathbf{n}}$ \\
Conjugate quantity & momentum $\mathbf{p}$ & angular momentum $\mathbf{L}+\mathbf{S}$ \\
Identity & $\dot{\mathbf{p}}=-\hbar\nabla\omega$ & $\hat{\mathbf{n}}\cdot d(\mathbf{L}+\mathbf{S})/dt=-\hbar\,\partial\omega/\partial\vartheta$ \\
Phase per unit change & $\mathbf{k}=\nabla\theta$ & $m_l=\partial\theta/\partial\varphi$; spin $\pm\tfrac{1}{2}$ \\
Rate of the phase per unit change & $\dot{\mathbf{k}}=-\nabla\omega$ & $d(\partial\theta/\partial\varphi)/dt=-\partial\omega/\partial\varphi$ \\
Doppler shift of $\omega$ & $-\mathbf{k}\cdot\mathbf{u}$ & $-\boldsymbol{\Omega}\cdot(\mathbf{L}+\mathbf{S})/\hbar$ \\
Inertial reaction & $\mathbf{F}_\text{in}=-m\mathbf{a}$ & $\mathbf{N}_\text{in}=(\mathbf{L}+\mathbf{S})\times\boldsymbol{\Omega}$, $-\mathcal{I}\dot{\boldsymbol{\Omega}}$ \\
Balance & $\mathbf{F}+\mathbf{F}_\text{in}=\mathbf{0}$ & $\mathbf{N}_\text{app}+\mathbf{N}_\text{in}=\mathbf{0}$ \\
Conserved under & translation & rotation \\
\bottomrule
\end{tabular}
\end{table}
% angular form checked in verification/sgm_check.py, the blocks for eq:sgm_angular to
% tab:linear_angular: the orbital angular gradient, the spin frequency shift, the gyroscopic torque,
% Euler's equations, eq:winding_rate, and the azimuthal moment m Omegadot r_perp^2.

%-------------------------------------------------------------------------------
\subsection{The Aharonov--Bohm effect}
\label{sec:ab}
%-------------------------------------------------------------------------------

The additive coupling of Eq.~\eqref{eq:omega_em} puts the potential itself into the frequency and, through the Legendre relation of Eq.~\eqref{eq:action_phase} and the four-gradient $p_\mu=\hbar\,\partial_\mu\theta$ of Eq.~\eqref{eq:sgm_fourvector}, into the phase accumulated along a path. Where the static potentials give $\Phi=0$ and $\mathbf{B}=\mathbf{0}$ but $\mathbf{A}\neq\mathbf{0}$, the force at the particle's position vanishes, since the electric field $\mathbf{E}=-\nabla\Phi-\partial_t\mathbf{A}$ and the magnetic force $q\,\mathbf{v}\times\mathbf{B}$ both vanish (Eq.~\eqref{eq:lorentz}). The canonical momentum is $\hbar\nabla\theta=\boldsymbol{\pi}+q\mathbf{A}$, where $\boldsymbol{\pi}$ is the kinematic momentum of Section~\ref{sec:em} (Eq.~\eqref{eq:debroglie}), so each path acquires, in addition to the kinetic phase of $\boldsymbol{\pi}$, the potential phase $(q/\hbar)\int\mathbf{A}\cdot d\boldsymbol{\ell}$, where $d\boldsymbol{\ell}$ is the path element. For two paths enclosing magnetic flux $\Phi_B$ the potential phases differ by $q/\hbar$ times the loop integral of $\mathbf{A}$, which by Stokes's theorem is the enclosed flux,
\begin{align}
    \Delta\theta_\text{AB} &= \frac{q}{\hbar}\oint \mathbf{A}\cdot d\boldsymbol{\ell} \notag\\
    &= \frac{q}{\hbar}\,\Phi_B ,
    \label{eq:ab} \\[-3pt]
    \multispan2{\hfil\text{\small\itshape (The Aharonov--Bohm phase)}\hfil} \notag
\end{align}
the Aharonov--Bohm phase \citep{Aharonov1959Significance}, a geometric phase \citep{Berry1984Quantal}. Experiment confirms it \citep{Chambers1960Shift}, including with the magnetic field confined inside a toroidal magnet by a superconductor and so excluded from both electron paths \citep{Tonomura1986Evidence}, and confirms the absence of a force: electrons passing a macroscopic solenoid arrive with no time delay, which rules out a semiclassical force of the size needed to produce the shift \citep{Caprez2007Macroscopic}. A gravitational Aharonov--Bohm effect, an interference shift induced by the potential of a source mass beyond any classical-force contribution, has been proposed as a force-free redshift measurement \citep{Hohensee2012Force} and observed in atom interferometry \citep{Overstreet2022Observation}.

With $\Phi=0$ and $\mathbf{B}=\mathbf{0}$ the potential does no work. At fixed velocity $\mathbf{v}=\boldsymbol{\pi}/m$, Eq.~\eqref{eq:omega_em} gives $\hbar\omega=mc^2+\tfrac{1}{2}m|\mathbf{v}|^2$, which does not contain $\mathbf{A}$, so $\omega$ equals its value at $\mathbf{A}=\mathbf{0}$; the measured absence of a time delay agrees with this \citep{Caprez2007Macroscopic}. The vector potential changes only the wavevector $\nabla\theta=(\boldsymbol{\pi}+q\mathbf{A})/\hbar$ (Eq.~\eqref{eq:debroglie}). The phase along each arm is the integral of that wavevector along the arm, so it depends only on the potential along that arm. Aharonov and Bohm's own calculation is local in the same way: the potential multiplies the field-free solution by a phase factor whose gradient is $(q/\hbar)\mathbf{A}$ \citep{Aharonov1959Significance}. The pointwise rates depend on the gauge, but the two-path difference that the fringe records is the gauge-invariant $\Delta\theta_\text{AB}$ of Eq.~\eqref{eq:ab}. A quaternion $\psi$ transported with zero component of $\mathbf{D}\psi$ along the path (Eq.~\eqref{eq:right_connection}) returns around the loop multiplied on the right by $\exp(-e_3\,\Delta\theta_\text{AB})$, the \emph{holonomy} of the gauge-covariant derivative.

Equation~\eqref{eq:pure_form} contains the potentials only through that derivative. A field strength is the commutator of two of its components (Eqs.~\eqref{eq:phase_curvature} and~\eqref{eq:time_space_curvature}). Parallel transport along the path solves a first-order equation, which forms no commutator, so it involves no field strength at the points of the path. Where $\mathbf{B}=\mathbf{0}$ the derivative is \emph{flat}, so the phase vanishes around every loop that can be contracted inside the field-free region. The region outside a solenoid is multiply connected, and around a loop that encircles the flux the common value $q\Phi_B/\hbar$ can be nonzero. It is a multiple of $2\pi$ when the flux is a multiple of $h/q$, where $h=2\pi\hbar$ \citep{Aharonov1959Significance}.

In coordinates that rotate at $\boldsymbol{\Omega}$, the shift $\mathbf{A}_g=\boldsymbol{\Omega}\times\mathbf{x}$ enters the phase as $\mathbf{A}$ does. There $\hbar\nabla\theta=\boldsymbol{\pi}=m(\dot{\mathbf{x}}+\mathbf{A}_g)$ (Section~\ref{sec:angular}), so each arm acquires the phase $(m/\hbar)\int\mathbf{A}_g\cdot d\boldsymbol{\ell}$ in addition to the phase of $m\dot{\mathbf{x}}$. Let a loop enclose the vector area $\boldsymbol{\mathcal{A}}$, oriented with the loop. By Stokes's theorem and $\nabla\times\mathbf{A}_g=2\boldsymbol{\Omega}$, the phase of the arm that runs in the sense of the loop exceeds that of the other arm by $2m\boldsymbol{\Omega}\cdot\boldsymbol{\mathcal{A}}/\hbar$. This is the Sagnac phase, observed first for light \citep{Sagnac1913Effect}; for neutrons it is the phase shift from the rotation of the Earth that \citet{Werner1979Effect} measured (Section~\ref{sec:angular}). Unlike the Aharonov--Bohm phase it comes with a force, since $\mathbf{B}_g=2\boldsymbol{\Omega}$ is nonzero on the paths and gives the Coriolis force of Table~\ref{tab:forces}.
% Checked in verification/sgm_check.py, the blocks for eq:ab.

%===============================================================================
\section{Consistency with experiment}
\label{sec:experiment}
%===============================================================================

A light-pulse atom interferometer splits an atomic wave function along two paths of vertical separation $\Delta z(t)$, the height of the upper path above the lower, in a gravitational field of strength $g$, enclosing the spacetime area $\mathcal{A}_{zt}=\int\Delta z\,dt$. The standard Feynman path-integral calculation, to first order in $g$ along the unperturbed paths, gives the phase difference $\Delta\theta_g=-mg \, \mathcal{A}_{zt}/\hbar$, the phase of the upper arm less that of the lower \citep{Anandan1977Gravitational,Kasevich1992Measurement,Storey1994Feynman,Peters1999Measurement}. In Earth's field the gravitational frequency gradient of Eq.~\eqref{eq:grav_gradient} is $\partial_z\omega=\omega_0 \, g/c^2=mg/\hbar$, and, since each arm's phase advances at the rate $-\omega$ (Eq.~\eqref{eq:phase_freq}), the frequency difference between the arms, $\partial_z\omega\,\Delta z(t)$, accumulates in time to
\begin{align}
    \Delta\theta_g &= -\,\partial_z\omega\,\mathcal{A}_{zt} \notag\\
    &= -\frac{mg}{\hbar}\,\mathcal{A}_{zt} ,
    \label{eq:interferometer} \\[-3pt]
    \multispan2{\hfil\text{\small\itshape (The phase of a light-pulse atom interferometer)}\hfil} \notag
\end{align}
identical to the path-integral phase. Stodolsky obtained the equality in equivalent form for matter and light waves in gravitational fields \citep{Stodolsky1979Matter}. Equation~\eqref{eq:interferometer} is evaluated, like the path-integral phase, on the unperturbed paths, along which the phases imprinted by the beam splitters cancel between the arms and the whole difference is the propagation phase. Along the actual paths, which fall, the total phase difference divides differently between the propagation phase and the pulse phases. In the closed geometry of three pulses separated by equal intervals $T_\text{p}$, the first pulse gives one arm the vertical recoil velocity $v_\text{r}$ and the second takes it from that arm and gives it to the other, and the propagation phase difference is zero. Integrating the free fall of each arm shows that the kinetic part $\int\tfrac{1}{2}m\dot{z}^2\,dt$ and the potential part $\int mgz\,dt$ of the action each differ between the arms by $mgv_\text{r}T_\text{p}^2$, for any $g$, so their difference cancels, and the phases imprinted by the pulses make up the whole of the same total.

Colella, Overhauser, and Werner first resolved this phase with neutrons, in an interferometer cut from a single crystal of silicon \citep{Colella1975Observation}. For their geometry Eq.~\eqref{eq:interferometer} gives $59.6$~radians, the value they predicted; they measured $54.3$~radians and ascribed the difference to bending of the interferometer under its own weight, and the sign of the measured phase shift confirmed that gravity enters the Schr\"odinger equation as the potential $+mgz$.

Peters, Chung, and Chu measured $g$ by the same phase with a fountain of cesium atoms \citep{Peters1999Measurement}. There the photon recoil that splits each atom sets the arm separation in inverse proportion to the atom's mass, so the mass cancels from the phase and the instrument measures $g$ alone; they determined $g$ at the site to three parts in $10^{9}$. The phase from the spacetime curvature across the arms, which vanishes in a uniform field and depends on where the atom is, has since been resolved as well \citep{Asenbaum2017Phase}.

Both experiments predicted the fringe phase from the weight: the potential $m \, g \, z$, whose sign the neutron measurement confirmed, gives the phase of Eq.~\eqref{eq:interferometer}. Equation~\eqref{eq:sgm} uses the same data in the opposite direction, to obtain the force from the phase. Dividing the measured phase by the enclosed spacetime area gives minus the vertical frequency gradient, and multiplying by $\hbar$ gives the force on the interfering particle. For the neutron interferometer the enclosed spacetime area is $3.83\times10^{-7}$~m\,s, the $10.5$~cm$^2$ area enclosed by the two beam paths divided by the neutron speed $2.74\times10^{3}$~m/s, and $\hbar$ times the predicted $59.6$~radians divided by it is $1.64\times10^{-26}$~N, the weight of one neutron. For the cesium interferometer the enclosed area is inversely proportional to the atom's mass while the phase is independent of it, and the same quotient is the weight of one atom, $m_\text{Cs}\,g=2.16\times10^{-24}$~N, where $m_\text{Cs}$ is the mass of a cesium-133 atom.

The time-dilation factor $\alpha$ is measured directly, with flown atomic clocks \citep{Hafele1972Clocks,Hafele1972Observed}, with a rocket-borne hydrogen maser \citep{Vessot1980Test}, and with optical clocks across height differences from tens of centimeters \citep{Chou2010Optical} down to a single millimeter \citep{Bothwell2022Resolving}; gravitational bound states of ultracold neutrons confirm the same potential $mgz$ \citep{Nesvizhevsky2002Quantum}. Whether an atom-interferometer phase tests a gravitational redshift at the Compton frequency, rather than the universality of free fall, is disputed \citep{Mueller2010Precision,Wolf2011,SinhaSamuel2011,Hohensee2012Comment,Wolf2012Reply,Schleich2013Representation,Roura2020Gravitational}. Spectral Gradient Mechanics reproduces the measured phase either way and does not by itself settle that question.

The shift $\mathbf{A}_g=-(2G/c^2r^3)\,\mathbf{J}_\text{src}\times\mathbf{x}$ of the rotating Earth (Eq.~\eqref{eq:lapse_shift_metric}) is measured through its orbital coupling in Eq.~\eqref{eq:nr_expansion}. The orbital coupling gives the frame-dragging force of Table~\ref{tab:forces}, which drags the orbits of satellites: a measurement on two Earth satellites found $99\pm5$ percent of the value of general relativity, with a total uncertainty allowance of $10$ percent \citep{Ciufolini2004Confirmation}. For a spinning body small compared with the distance over which $\mathbf{B}_g$ varies, the same coupling summed over its parts contains the term $-\tfrac{1}{2}\mathbf{B}_g\cdot\mathbf{L}$, where $\mathbf{L}$ is the angular momentum of the body about its center, so the axis of a gyroscope turns at the rate $-\tfrac{1}{2}\mathbf{B}_g$: the gyroscopes of Gravity Probe B drifted by $37.2\pm7.2$~mas/yr, against the $39.2$~mas/yr of general relativity \citep{Everitt2011Gravity}. The spin coupling $-\tfrac{1}{4}\hbar\,\boldsymbol{\sigma}\cdot\mathbf{B}_g=-\tfrac{1}{2}\mathbf{B}_g\cdot\mathbf{S}$ has the same coefficient $-\tfrac{1}{2}$, and it has been measured in rotating coordinates, where $\mathbf{B}_g=2\boldsymbol{\Omega}$ (Section~\ref{sec:angular}).
% Checked in verification/sgm_check.py, the blocks for eq:interferometer: the body sum gives -B_g.L/2.

The additive frequency contribution $q\Phi/\hbar$ is measured directly in the alternating-current Josephson effect. A superconducting junction held at a potential difference $\Delta\Phi$ oscillates at the frequency $f_J=(2e/h)\,\Delta\Phi$, where $e$ is the elementary charge and $2e$ the Cooper-pair charge, which is the frequency difference that Eq.~\eqref{eq:omega_em} assigns to the pair wave functions on the two sides \citep{Josephson1962Possible,Shapiro1963Josephson}. The relation is exact in the 2019 SI, where $e$ and $h$ are fixed, and it realizes the volt standard \citep{BIPM2019Ampere,CODATA2022Constants}: a Josephson voltmeter measures $\Delta\Phi$ by counting $f_J$.

Optical clocks at different heights measure the multiplicative coupling in the same way: the factor $\alpha$ multiplies the frequency of their transition as it multiplies $\omega_0$ in Eq.~\eqref{eq:omega_alpha} \citep{Chou2010Optical,Bothwell2022Resolving}. In each class the angular beat frequency is the coupling constant times the potential difference, divided by $\hbar$: the charge $2e$ times $\Delta\Phi$ across the junction, and the mass $m$ times the gravitational potential difference $g\,\Delta z$ between interferometer arms (Eq.~\eqref{eq:interferometer}).
%===============================================================================
\section{Scope and relation to prior work}
\label{sec:scope}
%===============================================================================

The optical-mechanical analogy, in which a trajectory follows the rays of a wave through a refractive medium, is due to Hamilton \citep{Synge1937Geometrical}, and its general-relativistic form has been worked out for the paths of light and of the planets \citep{Evans1996Optical}. The identity's own wavevector form, the eikonal ray equation (Eq.~\eqref{eq:eikonal_ray}), is the standard equation of ray tracing in wave physics \citep{Whitham1974Linear,Dodin2009Langmuir}. Light deflection as a vacuum refractive index $1/\alpha^2$ traces to Eddington \cite{Eddington1920} and to the optical, or Gordon, metric \citep{Gordon1923Refraction}, and the gravitational matter-wave phase was given by Stodolsky \cite{Stodolsky1979Matter} and has been observed from the Colella--Overhauser--Werner neutron-interferometry experiment onward \citep{Colella1975Observation,Greenberger1979Coherence}.

\citet{Gould2016Why} has shown classically, from the geodesic alone, that a released body falls mainly because of the time dilation rather than the spatial curvature, and the five derivations cited in Section~\ref{sec:intro} obtain the same result from the Compton frequency \citep{Wootters2003Why,Mueller2014Quantum}. The accounts of inertia cited there differ from Section~\ref{sec:inertia} in what produces the reaction. Schlegel derives Newton's second law from the de~Broglie and Planck--Einstein relations and the Lorentz transformation, with inertia as the change of frequency, and so of energy, that the time transformation gives a moving matter wave as a whole \citep{Schlegel1954Wave}. Petkov obtains the inertial and the gravitational self-force of an extended charge, with one mass, from the anisotropic propagation of light in accelerated and supported frames; that mass is the energy of the charge's own electromagnetic field divided by $c^2$ \citep{Petkov1999Acceleration}.

In Section~\ref{sec:inertia} the reaction comes from the particle's own wave function: while the particle accelerates, the rate of proper time in its rest frame varies across its wave function by the Rindler factor, so the frequency of its rest phase varies from one side of the wave function to the other. That gradient does not depend on the size of the wave function or on any boundary, and its coefficient is the particle's whole rest frequency $\omega_0=mc^2/\hbar$ (Eq.~\eqref{eq:rest_freq}).

The gauge theory of gravity of \citet{LasenbyDoranGull1998Gravity} (Section~\ref{sec:intro}) writes the Dirac spinor as an even element of a Clifford algebra, and its commutator of two internal covariant derivatives likewise places the field on the left of the spinor and the fixed unit on the right; Section~\ref{sec:spinor} makes the same left-right split for one chiral component. In Hestenes's factorization of the Dirac wave function (Section~\ref{sec:intro}) the phase is the orientation of the two transverse axes about the spin axis \citep{Hestenes1990Zitterbewegung}.

Feynman's derivation of the Lorentz force from the commutation relations of position and velocity has been extended to curved spacetime, where for a particle whose force depends on its position and velocity alone it admits only a force independent of the mass, from a scalar or a gauge field, and one proportional to the mass, from gravity \citep{Tanimura1992Relativistic}. The gauge and gravitational terms are the two classes of Table~\ref{tab:classes}. A scalar term independent of the mass acts as a Lorentz-scalar potential (Table~\ref{tab:classes}) that changes the masses of different particles by different fractions, and so lies outside both classes.

In the additive coupling of Table~\ref{tab:classes}, the Aharonov--Bohm phase is the potential's addition to the phase gradient $(\omega,\mathbf{k})$ at each point of the wave function's extent, accumulated locally along each path in every gauge (Section~\ref{sec:ab}). This paper shows that the same addition to the frequency field gives the force law of Eq.~\eqref{eq:sgm}. The other local accounts cited in Section~\ref{sec:intro} add degrees of freedom instead of using the particle's frequency. With the solenoid's charges treated quantum mechanically, the electron's own field acts locally on the source and the phase develops in the source's state \citep{Vaidman2012Role}, an account found incomplete for an electron separated from the solenoid by superconducting shielding \citep{Aharonov2016Nonlocality}. With the electromagnetic field quantized, the phase is generated by entanglement between the charge and the field \citep{Marletto2020Aharonov}.

The rest-frequency phase that gives gravity and inertia (Sections~\ref{sec:gravity} and~\ref{sec:inertia}) has no observable effect in Galilean quantum mechanics, where Bargmann's superselection rule decouples the overall mass phase \citep{Bargmann1954Unitary}. It becomes observable through relativity: it is de~Broglie's internal clock, which runs at the particle's proper time.

Pilot-wave theory \citep{deBroglie1927Wave,Bohm1952Suggested}, whose realist ontology the present framework shares, guides the particle by the quantum potential $Q$ of Section~\ref{sec:identity} \citep{Madelung1927Quantentheorie}, built from the wave function's amplitude $|\psi|$; $Q$ supplies the quantum corrections to the motion and is negligible where the amplitude varies slowly on the scale of the de~Broglie wavelength. The spectral gradient is built instead from the frequency function of Definition~\ref{def:spectral_gradient}, a function on phase space that contains no amplitude. For a particle at rest in a gravitational field the metric fixes the gradient as $\nabla\omega=\omega_0 \, \nabla\alpha$, and the force $-\hbar\nabla\omega$ does not vanish in the classical limit, since $\hbar\nabla\omega=mc^2 \, \nabla\alpha$ contains no $\hbar$ (Eq.~\eqref{eq:rest_freq}).

The phase rate $-\partial_t\theta$ of an exact state is the frequency function plus $Q/\hbar$ (Section~\ref{sec:route_phase}). The relation $\mathbf{F}=-\hbar\nabla\omega$ is Hamilton's canonical equation evaluated through $E=\hbar\omega$, and with the quantum potential added to the frequency, $\hbar\omega=H+Q$, it is the force law of de~Broglie--Bohm \citep{Bohm1952SuggestedB}.

Emergent-gravity programs likewise leave the gravitational field unquantized, but they derive gravity from entropy: from the thermodynamics of horizons and holographic screens in the earlier work \citep{Jacobson1995Thermodynamics,Padmanabhan2010Thermodynamics,Verlinde2011Origin}, and from the entanglement entropy of the vacuum in the later \citep{Jacobson2016Entanglement,Verlinde2017Emergent}. Here the force is $-\hbar$ times the gradient of one particle's own oscillation rate, with no average over the microstates of a horizon or of the vacuum.

The framework treats particles in stationary external fields and metrics, the case $\partial_t\omega=0$ in which Eq.~\eqref{eq:sgm_temporal} conserves the energy, together with constrained static and first post-Newtonian ($1$PN) moving-source many-body interactions on the matter Hilbert space. Its inertial cases are uniform acceleration and rigid rotation. It does not quantize dynamical gravitational radiation, and fully time-dependent strong-field metrics lie outside it, except for the response of a gravitational-wave detector and the exclusion of deterministic radiative mean-field gravity.

A wave whose transverse-traceless metric perturbation at the detector is $h^\text{TT}_{ij}(t)$ changes the time-dilation factor of the detector's freely falling frame to $\alpha=1+\tfrac{1}{2}R_{0i0j}\,x_ix_j$, summed over $i$ and $j$, to second order in the position $\mathbf{x}$ relative to the frame's origin \citep{Rakhmanov2014Fermi}, where $R_{0i0j}=-\ddot{h}^\text{TT}_{ij}/(2c^2)$ and each dot is a derivative with respect to $t$. For masses at rest at separations much smaller than the wavelength, $-mc^2\nabla\alpha$ is then the force with components $\tfrac{1}{2}m\,\ddot{h}^\text{TT}_{ij}\,x_j$, summed over $j$. It shifts the coordinate $x_i$ of a mass at $\mathbf{x}$ by $\tfrac{1}{2}h^\text{TT}_{ij}\,x_j$, the displacement that interferometric detectors measure \citep{Flanagan2005Basics}.

%===============================================================================
\section{Corollaries of the derived laws}
\label{sec:discussion}
%===============================================================================

%-------------------------------------------------------------------------------
\subsection{Particles are as large as the universe}
\label{sec:reach}
%-------------------------------------------------------------------------------

The bound eigenfunctions of the Schr\"odinger equation reach every distance. The hydrogen ground state $\psi_{100}\propto e^{-r/a_0}$, where $a_0$ is the Bohr radius, is nonzero at every finite $r$. Every hydrogenic eigenfunction is a polynomial in $r$ times a decaying exponential in $r$ times a spherical harmonic \citep{Sakurai2020Modern}; its nodes are surfaces of zero volume, so its amplitude is nonzero \emph{generically}, off those nodes, at every distance. The free Schr\"odinger kernel of Section~\ref{sec:duality} has the modulus $(m/2\pi\hbar t)^{3/2}$ at every separation, so a state confined at one instant generally has unbounded support an instant later, although interference nodes can remain.

The Dirac kernel, whose nonrelativistic limit is the Schr\"odinger kernel, vanishes outside the light cone \citep{Thaller1992Dirac}, and so bounds the spread. Amplitude that starts inside a bounded region therefore stays within the reach that causal contact allows, a distance $c$ times the time the state has evolved since it was prepared, and a change in a state spreads at most $c$ times the time since the change was made. A state confined to a bounded region mixes both signs of energy: a state built from positive-energy solutions of the free Dirac equation alone, as the state of a free electron is, has infinite support at every instant \citep{Hegerfeldt1998Instantaneous}, so its reach is unbounded from the start.

For matter whose wave functions have evolved since the early universe, that reach is the present particle horizon, about $46$~Glyr: the farthest position from which light emitted at $t=0$ has reached an observer here now. It is about three times the Hubble radius $R_H=c/H_0$, where $H_0$ is the Hubble constant \citep{DavisLineweaver2004Expanding}. It limits how far a change in a state can have spread, while the support of a positive-energy state, and with it the frequency field of Axiom~\ref{ax:local_frequency}, can already be unbounded. For any Hamiltonian bounded below, the probability of finding a particle in a given region is either nonzero at almost every instant or zero at every instant \citep[Theorem~2]{Hegerfeldt1998Instantaneous}, so a state with positive energy that is absent from a region for an interval of time is absent from that region at every time.

For a bound eigenstate with energy below the continuum threshold, whose tail falls off exponentially, every moment of $|\psi|^2$ is finite. In particular the spread $\Delta x$, the standard deviation of the coordinate $x$, is finite, and the probability of finding the particle beyond a few multiples of $\Delta x$ falls exponentially. The Kennard inequality bounds the product of the spreads from below, $\Delta x\,\Delta p\geq\hbar/2$, where $\Delta p$ is the standard deviation of the momentum component $p_x$ \citep{Sakurai2020Modern}. The hydrogen ground state, with $\langle x^2\rangle=a_0^2$ and $\langle p_x^2\rangle=\hbar^2/(3a_0^2)$, lies near that bound, at $\Delta x\,\Delta p=\hbar/\sqrt{3}$. A normalized state can therefore have a tail that is nonzero at every distance causal contact allows, while its probability stays concentrated near its mean. Below the threshold the spread is finite; at the threshold the second moment need not be finite. In units with $\hbar=m=1$ and $x$ a pure number, the state $\psi=[\pi(1+x^2)]^{-1/2}$ on a line is normalized and solves the Schr\"odinger equation with energy $0$ in the potential $V=(2x^2-1)/[2(1+x^2)^2]$, which tends to $0$ at large $|x|$. This eigenstate at the threshold has an infinite second moment, and so an infinite spread.

By the derivation from the Dirac equation in Section~\ref{sec:potentials}, each term of the Taylor expansion (Eq.~\eqref{eq:nr_expansion}) is the part of the frequency at a point $\mathbf{x}$ that one coupling contributes, evaluated with the fields that the other particles produce at $\mathbf{x}$, and it is defined wherever the amplitude is nonzero. Every coupling of the particle therefore acts over the whole region where its Dirac wave function is nonzero, which is unbounded for a state built from positive-energy solutions of the free Dirac equation \citep{Hegerfeldt1998Instantaneous} and, generically, for a bound state, whose tail reaches every distance.

\begin{proposition}[Particles are as large as the universe]\label{prop:reach}
In a potential without impenetrable hard walls (where unique continuation forbids a bound eigenfunction from vanishing identically on any open set), a bound state generically has nonzero amplitude at arbitrarily large distance, and a state built from positive-energy solutions of the free Dirac equation, as the state of a free electron is, has nonzero amplitude at arbitrarily large distance at every instant. For matter whose wave functions have evolved since the early universe, the reach that causal contact allows the state is bounded only by the particle horizon; a change made in the state at a later time reaches only the points inside the future light cone of the change. For a state of finite second moment, such as a bound eigenstate below the continuum threshold, the probability of finding the particle stays concentrated near its mean position, with a finite spread $\Delta x$ whose product with the momentum spread $\Delta p$ obeys the Kennard inequality $\Delta x\,\Delta p\geq\hbar/2$.
\end{proposition}

\begin{remark}[Each particle can act on any other]\label{rem:reach}
Proposition~\ref{prop:reach} concerns the amplitude of each state. In the mean-field approximation of Section~\ref{sec:wave_equation}, one particle acts on another only through the coefficients that the field laws produce from its charge and mass densities, so whether it acts at a given distance is a property of those laws. A coefficient at a point is the integral of the whole density against the Green function of the law. For a spherically symmetric density, the part that lies farther from its center than the point contributes zero to the gradient of $\alpha$ there, by Eq.~\eqref{eq:gauss_sphere} for the sphere through the point, where $M$ is the mass inside that sphere; Gauss's law for the electric field gives the same for $\Phi$. For a static mass density $\rho_m(\mathbf{y})$ in the weak field, the solution of Gauss's law for gravity (Eq.~\eqref{eq:gauss_alpha}) that tends to $\alpha=1$ far from all masses (Eq.~\eqref{eq:alpha_def}) is the integral of the density against the Green function $-1/(4\pi|\mathbf{x}-\mathbf{y}|)$ of the Laplacian,
\begin{equation}
    \alpha(\mathbf{x}) - 1 = -\frac{G}{c^2}\int\frac{\rho_m(\mathbf{y})}{|\mathbf{x}-\mathbf{y}|} \, d^3y ,
    \label{eq:static_density_clock_field}
\end{equation}
where $\mathbf{y}$ is the integration point and $d^3y$ the volume element at it. The integrand is positive wherever $\rho_m$ is, so the static mass density of a particle lowers $\alpha$ below 1 at every distance from it, and the charge density of a charged particle, which has one sign, makes $\Phi$ nonzero at every distance.

A change of either density reaches another particle only within the causal future of the change (Axiom~\ref{ax:causal_propagation}), and there only where the change alters the field: a rearrangement of a spherically symmetric mass at fixed total mass leaves the field outside it unchanged (Eq.~\eqref{eq:gauss_sphere}). No known law cuts off the interaction between wave functions at a maximum distance, so each particle can in principle act on any other within causal contact, and each pair interacts across the extents of both wave functions. Since every coupling of a particle acts wherever its wave function is nonzero, every particle is as large as the universe, in the sense that its couplings act over the whole of the universe within its causal contact. In a universe with a cosmological event horizon, such as de Sitter space \citep{Gibbons1977Cosmological}, the causal future of a change is not the whole universe.

This classical Gauss--Birkhoff exterior reach should be distinguished from holographic boundary completeness in perturbative quantum gravity \citep{Bousso2002,Marolf2009,Laddha2021,Raju2022}. In classical general relativity, a spherically symmetric rearrangement at fixed ADM mass leaves every exterior multipole of $\alpha$ unchanged, whereas in perturbative quantum gravity the asymptotic ADM Hamiltonian and boundary algebra at spatial or null infinity encode fine-grained information about interior quantum excitations \citep{Marolf2009,Laddha2021,Raju2022}. Remark~\ref{rem:reach} relies strictly on the classical exterior field laws and causal propagation, without invoking quantum-gravitational boundary completeness.
\end{remark}

%-------------------------------------------------------------------------------
\subsection{Planck's constant cancels from the gravitational force}
\label{sec:hbar_cancels}
%-------------------------------------------------------------------------------

In the classical limit $\hbar$ enters the motion only through the product $\hbar\nabla\omega$. The rest frequency $\omega_0=mc^2/\hbar$ (Eq.~\eqref{eq:rest_freq}) contains $\hbar$, and so does the frequency $\omega=\omega_0 \, \alpha$ (Eq.~\eqref{eq:omega_alpha}) at which the state oscillates. Substituting both into Eq.~\eqref{eq:sgm} gives
\begin{align}
    \mathbf{F} &= -\hbar \, \nabla\!\left(\frac{m c^2}{\hbar} \, \alpha\right) \notag\\
    &= -m c^2 \, \nabla\alpha \, .
    \label{eq:hbar_cancels_identity}
\end{align}
The scalar $mc^2/\hbar$ is constant, so it comes out of the gradient, and the two factors of $\hbar$ cancel exactly, for every $\alpha$ and every $m$: the force is $-\hbar$ times the gradient of the frequency of a quantum phase, and the product contains no $\hbar$. The corresponding classical trajectory contains no explicit $\hbar$, while the state that follows it remains quantum mechanical \citep{Ehrenfest1927Bemerkung}. No quantization of the gravitational field is needed for any effect derived here.

The $\hbar$ in Eq.~\eqref{eq:sgm} belongs to the matter: it is the action per radian of the particle's own phase, the constant that converts the mass to the rest frequency $\omega_0=mc^2/\hbar$ (Eq.~\eqref{eq:rest_freq}). The field $\alpha$ is a ratio of clock rates and contains no constant of action.

The action of a field configuration enters a quantum phase only after division by a constant of the same dimension. Take the gravitational sector to consist of the metric $g_{\mu\nu}$, whose time-time component $g_{00}$ gives the time-dilation factor $\alpha=\sqrt{-g_{00}}$ (Eq.~\eqref{eq:alpha_def}), together with the constants $c$ and $G$. In the coordinates $x^\mu=(ct,\mathbf{x})$ every coordinate has the dimension of length, as the line element $ds$ does, so every component of the metric is dimensionless, and so is the factor $\alpha=\sqrt{-g_{00}}$, a ratio of proper time to coordinate time. No power or gradient of a dimensionless field acquires the dimension of action.

Suppose that for some rational $n_c$ and $n_G$ the product $c^{n_c} \, G^{n_G}$ had the dimension $\mathsf{M} \, \mathsf{L}^2 \, \mathsf{T}^{-1}$ of action, where $\mathsf{M}$, $\mathsf{L}$, and $\mathsf{T}$ are the dimensions of mass, length, and time. Since $c$ has dimension $\mathsf{L} \, \mathsf{T}^{-1}$ and $G$ has dimension $\mathsf{M}^{-1} \, \mathsf{L}^3 \, \mathsf{T}^{-2}$, the product has dimension $\mathsf{M}^{-n_G} \, \mathsf{L}^{n_c+3n_G} \, \mathsf{T}^{-n_c-2n_G}$. Matching its exponents of $\mathsf{M}$, $\mathsf{L}$, and $\mathsf{T}$ to those of action requires
\begin{align}
    -n_G &= 1 \, , \notag\\
    n_c+3n_G &= 2 \, , \notag\\
    -n_c-2n_G &= -1 \, . \notag
\end{align}
The first condition gives $n_G=-1$, the second then gives $n_c=5$, and at those values the left side of the third is $-3$. The three conditions are inconsistent, so no such pair exists.

\begin{proposition}[The gravitational sector contains no constant of action]\label{prop:no_grav_action}
Let the gravitational sector consist of the metric $g_{\mu\nu}$, whose time-time component $g_{00}$ gives the time-dilation factor $\alpha=\sqrt{-g_{00}}$ (Eq.~\eqref{eq:alpha_def}), together with the constants $c$ and $G$. Then every component of the metric in the coordinates $x^\mu=(ct,\mathbf{x})$ is dimensionless, no product of powers of $c$ and $G$ has the dimension of action, and Planck's constant cancels identically from the force that Eq.~\eqref{eq:sgm} assigns to a gradient of $\alpha$ (Eq.~\eqref{eq:hbar_cancels_identity}). The sector therefore supplies no constant of dimension action by which a value of that dimension could be divided to give a phase.
\end{proposition}

By Proposition~\ref{prop:no_grav_action}, gravity is unquantized in this framework: the gravitational sector supplies no constant of action, so it defines no quantum of the gravitational field. The argument rests on two premises: that a quantum amplitude of the gravitational field requires a constant of dimension action that belongs to the gravitational sector itself, rather than one taken from the matter on which the field acts, and that the gravitational sector consists of the metric, $c$, and $G$. The first premise is not derived here, and a dimensional argument does not supply it: the free electromagnetic field, whose field equations contain only $c$ and the vacuum permittivity $\varepsilon_0$, contains no constant of dimension action and is quantized with the universal $\hbar$. The Einstein--Hilbert action $S_\text{EH}=(c^3/16\pi G)\int R\,\sqrt{-\det g_{\mu\nu}}\,d^4x$, where $R$ is the scalar curvature of the metric and $d^4x$ the volume element in the coordinates $x^\mu$, has the dimension of action, but it is a functional of the metric and not a constant: a sum over metrics weighted by $\exp(iS_\text{EH}/\hbar)$ divides it by $\hbar$, a constant from outside the sector. Under the first premise, the $\hbar$ of Eq.~\eqref{eq:rest_freq} belongs to matter and enters the coupling to a field only through the identity.

This also clarifies why ``everything else in nature is quantum'' does not force a particle-scale graviton. General relativity treats spacetime itself—and with it the time-dilation factor $\alpha(\mathbf{x})=\sqrt{-g_{00}}$—as a smooth, continuous manifold; its field equations and local proper-time structure rely on that continuous geometry, and in the present framework gravity is a \emph{consequence} of the gradient of time dilation (Corollary~\ref{cor:causal_gravity}, Eq.~\eqref{eq:geometric_force}) rather than the reverse. Moreover, standard quantum mechanics and the constant $\hbar$ of Eq.~\eqref{eq:rest_freq} operate at the \emph{particle (Compton) scale} $\hbar/(mc)$—roughly $2.10\times 10^{-16}~\text{m}$ for a proton ($m_p$) and $3.86\times 10^{-13}~\text{m}$ for an electron ($m_e$)—not at the \emph{Planck scale} $\ell_P\equiv\sqrt{\hbar G/c^3}\approx 1.62\times 10^{-35}~\text{m}$. The ratio of the particle scale to the Planck length is
\begin{equation}
    \frac{\hbar/(m_p c)}{\ell_P} = \frac{m_P}{m_p} \approx 1.30\times 10^{19} , \qquad \frac{\hbar/(m_e c)}{\ell_P} = \frac{m_P}{m_e} \approx 2.39\times 10^{22} ,
    \label{eq:compton_planck_ratio}
\end{equation}
where $m_P\equiv\sqrt{\hbar c/G}\approx 2.18\times 10^{-8}~\text{kg}$ is the Planck mass—roughly $20$ orders of magnitude separating particle-scale quantization from the Planck scale. Even if spacetime geometry possesses a discrete or quantized microstructure at the Planck scale $\ell_P$, any resulting Planck-scale quantization of gravity would enter only because time dilation $\alpha$ itself is quantized at the Planck scale, yet $\alpha(\mathbf{x})$ would remain indistinguishable from a smooth continuous field across all particle-scale experiments.

For the static field the first premise concerns $\alpha$: a quantum amplitude for $\alpha$ would need a constant of action from the gravitational sector. Optical clocks resolve its change across a height difference of $0.33$~m \citep{Chou2010Optical} and measure a linear frequency gradient across a single millimeter \citep{Bothwell2022Resolving}. A smooth measured profile of $\alpha$ is compatible both with a classical field and with the expectation value of a quantized one, so clock comparisons establish the redshift profile but do not by themselves test field quantization.

Some factors of the metric can be quantized: \citet{DeWitt1967Quantum} quantized the radius of a closed Friedmann universe, a scale factor of its spatial metric, and wrote its wave function as a function of that radius. The lapse $\alpha_n$ of Eq.~\eqref{eq:lapse_shift_metric}, which equals $\alpha$ where the shift vanishes (Eq.~\eqref{eq:alpha_lapse_shift}), differs from a scale factor: in the decomposition of the metric into the lapse, the shift, and the spatial metric, the Lagrangian of general relativity contains no time derivative of the lapse or the shift; they are Lagrange multipliers, and the spatial metric and its conjugate momentum are the complete set of initial data \citep{Arnowitt2008Dynamics,Gourgoulhon2007Formalism}. The momenta conjugate to the lapse and the shift therefore vanish identically \citep{DeWitt1967Quantum,Gourgoulhon2007Formalism}, and in the quantum theory this is a condition on the state: the wave function depends on the spatial metric alone, and that of the quantized Friedmann universe ``cannot depend on'' the lapse \citep{DeWitt1967Quantum}.

On slices whose extrinsic curvature has zero trace, the field equations give the lapse by an elliptic equation with no time derivative of the lapse. Its source is the energy density and stress of the matter, and in the Newtonian limit it reduces to the Poisson equation for the gravitational potential \citep{Gourgoulhon2007Formalism}. The slices of a static metric have zero extrinsic curvature, so for a static metric that equation, with the spatial metric and the boundary condition $\alpha_n\to1$ at infinity, fixes $\alpha_n$ from the matter. This application to the static constrained sector is made here, and its weak-field form is Gauss's law (Eq.~\eqref{eq:gauss_alpha}). The static time-dilation factor is therefore fixed by its source and has no conjugate momentum to quantize (analogous to the Coulomb-gauge scalar potential $\Phi$). For a source in a superposition of configurations, this paper takes the $\alpha$ of each configuration to be the static solution for that configuration.

The second premise restricts a constant of action of the gravitational sector to one built from the metric, $c$, and $G$. In four dimensions the only divergence-free tensor of two indices built from the metric and its first two derivatives is a constant multiple of the Einstein tensor plus a constant multiple of the metric \citep{Lovelock1972Four}, so the field equations contain $G$, $c$, and one further constant, the coefficient of the metric, which is the cosmological term. Gauss's law (Eq.~\eqref{eq:gauss_alpha}) and the vacuum field equations with the cosmological term zero (Corollary~\ref{cor:causal_gravity}) contain no constant beyond $G$ and $c$. The same field equations have gravitational waves as solutions, which are the transverse-traceless part of the metric and leave $\alpha$ equal to $1$ in the coordinates adapted to them \citep{Flanagan2005Basics}, so a sector of $\alpha$ alone would omit them. Like the free electromagnetic field, gravitational waves have degrees of freedom of their own, so whether they are quantized depends on the first premise and not on dimensions.

Any further constant in the field law from which, with $c$ and $G$, a constant of dimension action follows would defeat the argument. The cosmological term is such a constant: it has the dimension of an inverse area, and $c^3$ divided by $G$ and by it has the dimension of action. Moved to the source side of the field equations, the cosmological term is the stress of a vacuum whose energy density is uniform and whose pressure is minus that density, and it is equivalent to a vacuum mass density equal to $c^2/(8\pi G)$ times the cosmological term \citep{Carroll2001Cosmological}. The argument therefore treats the measured cosmological term as that vacuum mass density among the sources, as it treats a mass $m$, and not as a constant of the sector.

%-------------------------------------------------------------------------------
\subsection{A static field is a function on the configuration space of all particles}
\label{sec:fields_gravity}
%-------------------------------------------------------------------------------

The derivation from the Dirac equation (Section~\ref{sec:potentials}) shows that a particle is acted on only through the frequency of its own wave function at each point $\mathbf{x}$ of its extent, each coupling entering as one additive part of the $c^{-2}$ expansion, Eq.~\eqref{eq:nr_expansion}. Equation~\eqref{eq:static_density_clock_field} gives the coefficient at $\mathbf{x}$ as an integral over the points $\mathbf{y}$ of the mass densities of the other particles, each point contributing through the kernel $-G/(c^2|\mathbf{x}-\mathbf{y}|)$, a function only of the point $\mathbf{x}$ of one wave function and the point $\mathbf{y}$ of another.

Consider two particles of masses $m_1$ and $m_2$ with normalized densities $\rho_1$ and $\rho_2$, both static and with no shift, so that the lapse $\alpha_n$ equals $\alpha$. Each density is the conserved density $(|\varphi_R|^2+|\varphi_L|^2)/\alpha_n$ of Section~\ref{sec:wave_equation}, not the part $\rho=|\psi|^2$ contributed by $\varphi_R$. The weight $1/\alpha_n$ differs from 1 at first order in the potential, so it changes the pair energy only at second order in $G$. The geometric term $mc^2\alpha_n$ of Eq.~\eqref{eq:nr_expansion} for particle 1, weighted by $\rho_1(\mathbf{x})$ and integrated over $\mathbf{x}$, contains the part $-G\,m_1m_2\iint\rho_1(\mathbf{x})\,\rho_2(\mathbf{y})\,|\mathbf{x}-\mathbf{y}|^{-1}\,d^3x\,d^3y$ that particle 2 produces. This energy is unchanged when the particles are exchanged, and it couples every point of one density to every point of the other, across the spatial extents of both. It is the expectation value, in a state that is a product of one wave function for each particle, of the pair term $-G\,m_1m_2/|\mathbf{x}_1-\mathbf{x}_2|$, where $\mathbf{x}_1$ and $\mathbf{x}_2$ are the positions of the two particles.

A point $\mathbf{y}$ with mass $dm$ contributes $(G\,dm/c^2)\,(\mathbf{x}-\mathbf{y})/|\mathbf{x}-\mathbf{y}|^3$ to $\nabla\alpha$ at $\mathbf{x}$, and the flux of that contribution through a closed surface is $G\,dm/c^2$ times the solid angle the surface subtends at $\mathbf{y}$: $4\pi$ if $\mathbf{y}$ lies inside, and zero if it lies outside, since every ray from an outside point crosses the surface inward as often as outward. The gradient $\nabla\alpha$ on the surface is the sum of these contributions over all points $\mathbf{y}$, inside and outside the volume, so its flux is the sum of their fluxes. Summed over every point of every wave function, the flux of $\nabla\alpha$ through the surface is therefore $4\pi G/c^2$ times the enclosed mass, however the mass inside and outside the volume is distributed, which is Gauss's law (Eq.~\eqref{eq:gauss_alpha}).

The derivation needs the kernel to fall as $1/|\mathbf{x}-\mathbf{y}|$: the flux of its gradient through a sphere about $\mathbf{y}$ is then the same at every radius, while for a kernel of finite range, such as $e^{-r/\lambda_Y}/r$ where $r=|\mathbf{x}-\mathbf{y}|$, it is $4\pi(1+r/\lambda_Y)e^{-r/\lambda_Y}$ in magnitude and falls with the radius, so the flux is not a function of the enclosed mass alone. The derivation takes two inputs from outside the pairwise sum: the form $1/|\mathbf{x}-\mathbf{y}|$ of the kernel, which the vacuum field equations give with no reference to a source (Corollary~\ref{cor:causal_gravity}), and its coefficient $G/c^2$, the measured strength of Section~\ref{sec:gravity}.

The Coulomb kernel $1/(4\pi\varepsilon_0|\mathbf{x}-\mathbf{y}|)$ gives the electric potential from the charge densities in the same way. In the Lorenz gauge of Section~\ref{sec:maxwell}, a change of a localized source reaches $\mathbf{x}$ through the retarded Green function, one propagation time $|\mathbf{x}-\mathbf{y}|/c$ after it occurs at $\mathbf{y}$ (Axiom~\ref{ax:causal_propagation}). Where $|\mathbf{x}-\mathbf{y}|/c$ is short compared with the time over which $\rho_q$ changes, expanding $\rho_q(\mathbf{y},t-|\mathbf{x}-\mathbf{y}|/c)$ in powers of $1/c$ shows that the delay changes the electric field, which is the same in every gauge, only at order $c^{-2}$ \citep{Jackson1999Classical,LandauLifshitz1975Classical}. The term of first order adds $-(4\pi\varepsilon_0c)^{-1}\,d\big(\int\rho_q\,d^3y\big)/dt$ to $\Phi$, which is the same at every $\mathbf{x}$ and so has no gradient; it is zero, since the total charge is conserved. The next term of the expansion, and $\partial_t\mathbf{A}$, are of order $c^{-2}$, since the magnetic constant is $\mu_0=1/(\varepsilon_0c^2)$.

Applied to the joint wave function, the Coulomb pair term entangles the particles it couples. Two ions held about $27~\mu\text{m}$ apart in separate potential wells, and coupled through the Coulomb interaction between them, have been prepared in an entangled spin state with fidelity $0.82(1)$ \citep{Wilson2014Tunable}. A potential computed from the expectation values of the charge densities could not produce this result: it acts on each ion through a function of that ion's position alone, so a product state stays a product state. The Coulomb term is fixed by the charges through Gauss's law, and it produces entanglement whether or not the radiation field is quantized \citep{Anastopoulos2018Comment}.
% Checked in verification/sgm_check.py, the block for eq:alpha_configuration: the 1/c expansion of the retarded integrand.

Like the equation for the electric potential in the Coulomb gauge, $\nabla\cdot\mathbf{A}=0$, the field equation for the lapse contains no time derivative, and the momentum conjugate to the lapse vanishes (Section~\ref{sec:hbar_cancels}), so the static constraint fixes $\alpha$ from the masses once a source configuration and boundary conditions are specified. The many-body prescription used here applies that solution configuration by configuration. The joint wave function has as its arguments the positions of all the particles; at each configuration, $\alpha$ is the static solution for the masses at the positions that configuration assigns, and it acts on the joint wave function as multiplication by that function of the positions. Let the other particles have masses $m_j$ at the positions $\mathbf{x}_j$ for $j=2,3,\ldots$, so that their mass density is concentrated at those points. Equation~\eqref{eq:static_density_clock_field} then gives, in that configuration,
\begin{equation}
    \alpha(\mathbf{x}) - 1 = -\frac{G}{c^2}\sum_{j\geq2}\frac{m_j}{|\mathbf{x}-\mathbf{x}_j|} ,
    \label{eq:alpha_configuration}
\end{equation}
and at $\mathbf{x}=\mathbf{x}_1$ the single-particle geometric term $m_1c^2(\alpha_1(\mathbf{x}_1)-1)$ contains the sum $-\sum_{j\geq2}G\,m_1m_j/|\mathbf{x}_1-\mathbf{x}_j|$. Summing $m_i c^2(\alpha_i(\mathbf{x}_i)-1)$ over all $N$ particles counts each unordered pair $\{i,j\}$ twice, exactly as summing $q_i\Phi_i(\mathbf{x}_i)$ does in electrostatics, so the total configuration-space gravitational interaction operator on the $N$-particle Hilbert space $\mathcal{H}_\text{mat}$ carries the standard factor $\tfrac{1}{2}$:
\begin{equation}
    H_\text{grav,pair} = \frac{1}{2}\sum_{i=1}^{N} m_i c^2\,\big(\alpha_i(\mathbf{x}_i)-1\big) = -\sum_{1\le i<j\le N}\frac{G\,m_i m_j}{|\mathbf{x}_i-\mathbf{x}_j|} ,
    \label{eq:h_grav_pair}
\end{equation}
each pair contributing $-G\,m_i m_j/|\mathbf{x}_i-\mathbf{x}_j|$ once to the joint energy. In a product state of one wave function per particle, the expectation value of the $(1,2)$ term in Eq.~\eqref{eq:h_grav_pair} equals the two-particle pair energy above, while differentiating $H_\text{grav,pair}$ with respect to $\mathbf{x}_1$ recovers $-m_1c^2\nabla_{\mathbf{x}_1}\alpha_1$ without any double-counting factor.

\begin{corollary}[Fields arise from the wave functions of all particles]\label{cor:fields}
In a static weak field, the time-dilation factor $\alpha$ and the electric potential $\Phi$ to which a particle's wave equation couples at a point are fixed by the positions of all the other particles and have no degrees of freedom of their own. On the joint wave function of all particles each acts as a function of those positions, the sum over the other particles of the Green function of its field law (Eq.~\eqref{eq:alpha_configuration} for $\alpha$), and the gravitational part of the energy is, to first order in $G$, a sum of pair terms, each a function of the positions of two particles, the form that the Coulomb term of the many-particle Hamiltonian has. Every pair of particles interacts through this function, across the extents of both wave functions. A field is the part of this function to which one particle's wave equation couples, additively or multiplicatively (Table~\ref{tab:classes}), and in the mean-field approximation of Section~\ref{sec:wave_equation} it is the integral of the densities of the other particles against the Green function (Eq.~\eqref{eq:static_density_clock_field}), not a superposition of the tails of their wave functions. For moving sources the field equations of general relativity and Maxwell's equations, taken as input, give $\alpha$, $\Phi$, $\mathbf{A}$, and $\mathbf{A}_g$ from the densities and currents of all particles, except for free radiation, which has degrees of freedom of its own.
\end{corollary}

By Corollary~\ref{cor:fields} the pair term of two masses entangles them as the Coulomb pair term entangles the ions, so two masses, each prepared in a spatial superposition and coupled only through gravity, become entangled, which is the gravitationally induced entanglement proposed as a test by \citet{Bose2017Spin} and \citet{Marletto2017Gravitationally}.

The gravitational force is linear in the gradient of the time-dilation factor $\alpha(\mathbf{x})$ (Eq.~\eqref{eq:geometric_force}), and by Corollary~\ref{cor:fields} that gradient at a point of a particle's wave function is the sum of the gradients of the contributions of the other particles (Eq.~\eqref{eq:alpha_configuration}). By Remark~\ref{rem:reach} every other particle within causal contact contributes to it, and a change in the density of one of them changes it no earlier than the time a signal at the speed $c$ takes to cross the separation (Axiom~\ref{ax:causal_propagation}). The identity (Eq.~\eqref{eq:sgm}) therefore gives the force on each particle as a sum over its interactions, through the wave functions, with all particles within its causal contact:

\begin{corollary}[Gravity is due only to the interaction between wave functions]\label{cor:gravity}
Every interaction of a particle with other particles, gravity included, acts on it only through the coupling between its wave function and the fields over its extent, which the field laws give from the densities and currents of the other particles within its causal contact. A change in one of those densities changes the force no earlier than the time a signal at the speed $c$ takes to cross the separation. For gravity the field is the time-dilation factor $\alpha$, the rate of proper time, its static law is Eq.~\eqref{eq:gauss_alpha}, and the force is $-mc^2\nabla\alpha$ at each point of the wave function.
\end{corollary}

A central force delayed by the propagation time would be aberrated: it would no longer point toward the center of mass, and the orbits would be unstable. In general relativity velocity-dependent terms in the interaction cancel that aberration almost exactly, a cancellation that follows from conservation laws and the quadrupole nature of gravitational radiation, while for a scalar interaction that is not Lorentz invariant conservation laws place no restriction on the radiation, and aberration is expected at all orders \citep{Carlip2000Aberration}. The field equation that gives $\alpha$ from a moving source is not derived here, so it is not shown that the delay in $\alpha$ leaves the orbits unaberrated; Eq.~\eqref{eq:gauss_alpha} is the static law. The force on a particle at rest in a static metric is $-mc^2\nabla\alpha$ (Eq.~\eqref{eq:geometric_force}), and the velocity-dependent effects, among them the full geodetic precession, need the rest of the metric.

\begin{remark}[Test of the ontology of $\alpha$]\label{rem:lapse_test}
Corollary~\ref{cor:fields} predicts that the time-dilation factor $\alpha$ follows its source into quantum superposition. The alternative is a mean-field $\alpha$, computed from the expectation values of the mass densities. Interferometry with a superposed test particle in the field of a classical source (Section~\ref{sec:experiment}) does not distinguish the two: for a classical source they make the same prediction, and they differ only when the source itself is superposed. Gravitationally induced entanglement distinguishes them. By the argument that Section~\ref{sec:fields_gravity} gives for a potential computed from the expectation values of the charge densities, a field computed from the expectation values of the mass densities cannot entangle the masses, so an observation of the entanglement would exclude the mean-field $\alpha$.

\citet{PageGeilker1981Indirect} excluded one form of the mean field when they let a quantum event choose the positions of the source masses of a torsion balance, and the field they measured followed the realized positions alone, with no contribution from the alternatives. That result excludes an $\alpha$ computed from the average over macroscopic alternatives that never collapse, which is the semiclassical Einstein equation without collapse of the wave function; it does not exclude an $\alpha$ computed from the mass configuration after collapse, and it is not a test of gravitational entanglement.

A pair of masses $m_1$ and $m_2$ held at separation $r_{ab}$ in branch $|a,b\rangle$ ($a,b\in\{0,1\}$) has joint gravitational interaction energy $-G m_1 m_2 / r_{ab}$ by Eq.~\eqref{eq:h_grav_pair}, so over a holding time $\tau$ that branch accumulates the gravitational phase $\theta_{ab} = +G m_1 m_2 \tau / (\hbar r_{ab})$ (about $0.8$~radian for $10^{-14}$~kg masses at $200~\mu\text{m}$ over $\tau=2.5$~s \citep{Bose2017Spin}). The entanglement is governed by the phase combination $\Delta\theta_\text{grav}\equiv\theta_{00}+\theta_{11}-\theta_{01}-\theta_{10}$. Consider a protocol that keeps the branches coherent over the holding time, in a geometry for which the weak-field pair terms predict that $\Delta\theta_\text{grav}$ differs from zero modulo $2\pi$ by more than the resolution of the measurement. A null result in such a protocol would exclude Corollary~\ref{cor:fields}, and with it the application of the constraint that fixes the lapse separately in each mass configuration.

In each branch of the superposition, the static $\alpha$ is fixed by that branch's mass configuration through Eq.~\eqref{eq:gauss_alpha} and has no degrees of freedom of its own. The phase that entangles the masses is an energy multiplied by a time and divided by the $\hbar$ of the matter, the constant that the first premise of the argument assigns to matter, so the predicted entanglement is consistent with Proposition~\ref{prop:no_grav_action}. \citet{Anastopoulos2018Comment} argue likewise from the constraint structure of general relativity: in the weak-field, nonrelativistic limit the Newtonian interaction is fixed by the scalar constraint, as the Coulomb term is fixed by Gauss's law, so the entanglement it produces does not by itself show that the gravitational radiation field is quantized. \citet{Bose2017Spin} and \citet{Marletto2017Gravitationally} take the entanglement instead as evidence that gravity is quantum, \citet{Christodoulou2019Possibility} describe the effect as a superposition of proper times, and \citet{DanielsonSatishchandranWald2022} note that enforcing relativistic causality when the superpositions are recombined rapidly links the Newtonian phase to radiative modes; within the static constrained sector treated here, Eq.~\eqref{eq:h_grav_pair} generates entanglement directly on $\mathcal{H}_\text{mat}$.
\end{remark}

The contrast between configuration-space coupling (Corollary~\ref{cor:fields}) and deterministic mean-field gravity (Remark~\ref{rem:lapse_test}) directly answers the standard question for a quantum source in spatial superposition: when an atom is ``in two places at once,'' which place does its gravity pull from, and how do probabilities remain conserved at $100\%$? Under Eq.~\eqref{eq:h_grav_pair}, the atom pulls from neither a single collapsed site nor an expectation-value midpoint; instead, the static constraint assigns to each configuration $|a,b\rangle$ of the joint matter wave function $|\Psi\rangle\in\mathcal{H}_\text{mat}$ the lapse $\alpha$ sourced by the masses at those branch positions. Because $H_\text{grav,pair}$ is a real multiplication operator on configuration space, it is strictly Hermitian on $\mathcal{H}_\text{mat}$, so $\exp(-i H_\text{grav,pair}\tau/\hbar)$ is strictly unitary: total probability stays identically $1$ ($\langle\Psi(\tau)|\Psi(\tau)\rangle=1$) and overlaps between distinct states are preserved, while configuration-dependent gravitational phases entangle the branches without invoking a graviton. By contrast, deterministic mean-field gravity sourced by the averaged mass density cannot entangle and violates linear unitarity by driving non-orthogonal states to orthogonality in finite time. This contrast is exact and quantitative.

\begin{proposition}[Quantitative BMV purity and two-site Schr\"odinger--Newton orthogonalization]\label{prop:config_space_gravity}
For two masses coupled through static weak-field gravity:
\begin{enumerate}[label=(\roman*)]
\item \emph{Configuration-space BMV entanglement purity:} Under the unitary operator $\exp(-i H_\text{grav,pair}\tau/\hbar)$ of Eq.~\eqref{eq:h_grav_pair}, the initial product state $|\Psi(0)\rangle = \tfrac{1}{2}(|0\rangle+|1\rangle)\otimes(|0\rangle+|1\rangle)$ evolves to $|\Psi(\tau)\rangle = \tfrac{1}{2}\sum_{a,b\in\{0,1\}} e^{i\theta_{ab}}|a,b\rangle$ with $\theta_{ab} = G m_1 m_2 \tau / (\hbar r_{ab})$. Tracing out the second mass gives the exact reduced-state purity of the reduced density operator $\hat{\rho}_1(\tau)\equiv\operatorname{tr}_2|\Psi(\tau)\rangle\langle\Psi(\tau)|$,
\begin{equation}
    \operatorname{tr}\big(\hat{\rho}_1(\tau)^2\big) = 1 - \frac{1}{2}\sin^2\!\left(\frac{\Delta\theta_\text{grav}}{2}\right), \qquad \Delta\theta_\text{grav} \equiv \theta_{00}+\theta_{11}-\theta_{01}-\theta_{10} ,
    \label{eq:gie_phase}
\end{equation}
which reaches maximal entanglement ($\operatorname{tr}(\hat{\rho}_1^2)=\tfrac{1}{2}$) at $\Delta\theta_\text{grav}=\pi\pmod{2\pi}$ while conserving global unitarity on $\mathcal{H}_\text{mat}$.
\item \emph{Two-site Schr\"odinger--Newton mean-field orthogonalization:} By contrast, under the deterministic mean-field Schr\"odinger--Newton evolution on two localized sites $|0\rangle, |1\rangle$ with self-coupling energy $U_\text{self}$ and cross-site coupling energy $U_\text{cross}$ ($U_\text{self}>U_\text{cross}>0$)~\citep{Bahrami2014},
\begin{equation}
    i\hbar\,\frac{dc_a}{dt} = -U_\text{self} |c_a|^2 c_a - U_\text{cross} |c_{1-a}|^2 c_a \qquad (a\in\{0,1\}) ,
    \label{eq:sn_mean_field}
\end{equation}
the site probabilities $p_a=|c_a|^2$ remain constant while the relative phase shifts at a state-dependent rate. Two initially non-orthogonal pure states $|\psi^{(A)}(0)\rangle = \sqrt{1-p_A}|0\rangle + \sqrt{p_A}|1\rangle$ and $|\psi^{(B)}(0)\rangle = \sqrt{1-p_B}|0\rangle + \sqrt{p_B}|1\rangle$ evolve with squared overlap
\begin{equation}
    \big|\langle\psi^{(A)}(t)\mid\psi^{(B)}(t)\rangle\big|^2 = (1-p_A)(1-p_B) + p_A p_B + 2\sqrt{p_A(1-p_A)p_B(1-p_B)}\,\cos(\Delta\omega_\text{SN}\,t) ,
    \label{eq:sn_overlap_drift}
\end{equation}
where $\Delta\omega_\text{SN} = \frac{2(U_\text{self}-U_\text{cross})}{\hbar}(p_B-p_A)$. In particular, for $(p_A,p_B)=(\tfrac{1}{4},\tfrac{3}{4})$ the overlap $\big|\langle\psi^{(A)}(t)\mid\psi^{(B)}(t)\rangle\big| = \tfrac{\sqrt{3}}{2}\big|\cos\big(\frac{(U_\text{self}-U_\text{cross})t}{2\hbar}\big)\big|$ drops from $\sqrt{3}/2$ to exact orthogonality $0$ at the finite time $t_*^\text{SN} = \pi\hbar/(U_\text{self}-U_\text{cross})$, violating unitarity and enabling superluminal Gisin collapse signaling~\citep{Gisin1990,Bahrami2014}.
\end{enumerate}
\end{proposition}
\begin{proof}
For (i), $\hat{\rho}_1(\tau) = \operatorname{tr}_2|\Psi(\tau)\rangle\langle\Psi(\tau)|$ has diagonal entries $\tfrac{1}{2}$ and off-diagonal entry $(\hat{\rho}_1)_{01} = \tfrac{1}{4}\big(e^{i(\theta_{00}-\theta_{10})} + e^{i(\theta_{01}-\theta_{11})}\big)$, so $|(\hat{\rho}_1)_{01}|^2 = \tfrac{1}{4}\cos^2(\Delta\theta_\text{grav}/2)$ and $\operatorname{tr}(\hat{\rho}_1^2) = \tfrac{1}{2} + 2|(\hat{\rho}_1)_{01}|^2 = 1 - \tfrac{1}{2}\sin^2(\Delta\theta_\text{grav}/2)$. For (ii), Eq.~\eqref{eq:sn_mean_field} gives $\frac{d}{dt}|c_a|^2 = 0$ and phase rates $\dot{\phi}_1-\dot{\phi}_0 = (U_\text{self}-U_\text{cross})(2p-1)/\hbar$, so the relative phase difference between states $A$ and $B$ grows at $\Delta\omega_\text{SN} = 2(U_\text{self}-U_\text{cross})(p_B-p_A)/\hbar$; expanding $|\langle\psi^{(A)}(t)|\psi^{(B)}(t)\rangle|^2$ gives Eq.~\eqref{eq:sn_overlap_drift}, which vanishes at $t_*^\text{SN}=\pi\hbar/(U_\text{self}-U_\text{cross})$ for $(p_A,p_B)=(\tfrac{1}{4},\tfrac{3}{4})$.
\end{proof}

%-------------------------------------------------------------------------------
\subsection{Moving sources, radiative mean-field exclusion, and horizon saturation}
\label{sec:dynamical_gravity}
%-------------------------------------------------------------------------------

Three exact results extend Section~\ref{sec:fields_gravity} beyond static weak fields while keeping the gravitational constraint sector unquantized on $\mathcal{H}_\text{mat}$.

First, for slowly moving masses at first post-Newtonian ($1$PN) order, eliminating the constrained gravitomagnetic shift yields the Einstein--Infeld--Hoffmann (EIH) cross-momentum interaction directly on the $N$-particle matter Hilbert space $\mathcal{H}_\text{mat}$ without introducing graviton degrees of freedom.

\begin{proposition}[Manifestly Hermitian $1$PN Einstein--Infeld--Hoffmann cross-momentum operator]\label{prop:eih_1pn}
Define the real symmetric $3\times 3$ spatial tensor kernel
\begin{equation}
    \mathcal{W}_{ij}^{ab}(\mathbf{r}) \equiv \frac{G}{2c^2|\mathbf{r}|}\!\left(7\,\delta^{ab} + \frac{r^a r^b}{|\mathbf{r}|^2}\right) \qquad (a,b\in\{1,2,3\},\; i\neq j) .
    \label{eq:eih_kernel}
\end{equation}
Then the $1$PN EIH cross-momentum Hamiltonian operator on $\mathcal{H}_\text{mat}$,
\begin{equation}
    H_\text{EIH}^\text{cross} \equiv \frac{1}{2}\sum_{1\le i<j\le N}\sum_{a,b=1}^{3} \Big(\hat{p}_{i,a}\,\mathcal{W}_{ij}^{ab}(\mathbf{x}_i-\mathbf{x}_j)\,\hat{p}_{j,b} + \hat{p}_{j,b}\,\mathcal{W}_{ij}^{ab}(\mathbf{x}_i-\mathbf{x}_j)\,\hat{p}_{i,a}\Big) ,
    \label{eq:eih_1pn}
\end{equation}
is manifestly Hermitian on $\mathcal{H}_\text{mat}$ (with no graviton Fock space), has the exact classical $1$PN EIH symbol $+\sum_{i<j}\frac{G}{2c^2 r_{ij}}\big(7\mathbf{p}_i\cdot\mathbf{p}_j + (\mathbf{p}_i\cdot\hat{\mathbf{r}}_{ij})(\mathbf{p}_j\cdot\hat{\mathbf{r}}_{ij})\big)$, and yields positive Lense--Thirring frame dragging $\dot{\mathbf{x}}_1|_{\mathbf{p}_1=\mathbf{0}} = +\frac{G}{2c^2 r}\big(7\mathbf{p}_2 + (\mathbf{p}_2\cdot\hat{\mathbf{r}})\hat{\mathbf{r}}\big)$ (reducing to $+\frac{7G}{2c^2 r}\mathbf{p}_2$ for $\mathbf{p}_2\perp\mathbf{r}$).
\end{proposition}
\begin{proof}
Because $\hat{p}_{i,a}$ and $\hat{p}_{j,b}$ are Hermitian and act on distinct particles $i\neq j$, taking the adjoint of $\hat{p}_{i,a}\mathcal{W}_{ij}^{ab}\hat{p}_{j,b}$ reverses factor order to $\hat{p}_{j,b}\mathcal{W}_{ij}^{ab}\hat{p}_{i,a}$ (since $\mathcal{W}_{ij}^{ab}$ is real and symmetric in $(a,b)$), making each summand of Eq.~\eqref{eq:eih_1pn} identically Hermitian. Under the Legendre transform with $\mathbf{v}_i=\mathbf{p}_i/m_i+O(c^{-2})$, the $1$PN EIH Lagrangian cross-velocity term $-\sum_{i<j}\frac{G m_i m_j}{2c^2 r_{ij}}\big(7\mathbf{v}_i\cdot\mathbf{v}_j+(\mathbf{v}_i\cdot\hat{\mathbf{r}}_{ij})(\mathbf{v}_j\cdot\hat{\mathbf{r}}_{ij})\big)$ flips sign to $+H_\text{EIH}^\text{cross}$, and differentiating its symbol with respect to $\mathbf{p}_1$ at $\mathbf{p}_1=\mathbf{0}$ yields $\dot{\mathbf{x}}_1|_{\mathbf{p}_1=\mathbf{0}} = +\frac{G}{2c^2 r}(7\mathbf{p}_2+(\mathbf{p}_2\cdot\hat{\mathbf{r}})\hat{\mathbf{r}})$.
\end{proof}

Second, although free transverse-traceless gravitational radiation $h_{ij}^\text{TT}$ carries independent dynamical degrees of freedom (Section~\ref{sec:hbar_cancels}), coupling matter to a deterministic nonlinear radiative mean field $h_{ij}^\text{TT}[\langle\hat{T}_{\mu\nu}\rangle]$ fails even when the static mass density $T_{00}$ is identical across branches.

\begin{proposition}[Exclusion of deterministic nonlinear radiative mean-field gravity]\label{prop:rad_meanfield_signaling}
Consider two orthogonal matter states $|0\rangle, |1\rangle$ with identical static mass-energy density $T_{00}^{(0)}(\mathbf{x})\equiv T_{00}^{(1)}(\mathbf{x})$ (so the static Schr\"odinger--Newton shift $\Delta\omega_\text{SN}$ of Eq.~\eqref{eq:sn_overlap_drift} vanishes identically) but orthogonal $+/\times$ traceless quadrupole tensors $E^{(0)}=\operatorname{diag}(1,-1,0)$ and $E^{(1)}=\bigl(\begin{smallmatrix}0&1&0\\1&0&0\\0&0&0\end{smallmatrix}\bigr)$ ($\operatorname{tr}(E^{(a)}E^{(b)})=2\delta_{ab}$). Under deterministic nonlinear radiative mean-field coupling $h_{ij}^\text{TT}[\langle\hat{T}_{\mu\nu}\rangle]$:
\begin{enumerate}[label=(\roman*)]
\item \emph{Radiative pure-state orthogonalization:} Near-zone time-symmetric TT self-coupling produces diagonal energies $E_0^\text{rad}(p)=-\kappa_\text{rad}(1-p)$ and $E_1^\text{rad}(p)=-\kappa_\text{rad}p$ ($\kappa_\text{rad}>0$), driving $|\psi^{(A)}(0)\rangle = \sqrt{1-p_A}|0\rangle + \sqrt{p_A}|1\rangle$ and $|\psi^{(B)}(0)\rangle = \sqrt{1-p_B}|0\rangle + \sqrt{p_B}|1\rangle$ to obey
\begin{equation}
    \big|\langle\psi^{(A)}(t)\mid\psi^{(B)}(t)\rangle\big|^2 = (1-p_A)(1-p_B) + p_A p_B + 2\sqrt{p_A(1-p_A)p_B(1-p_B)}\,\cos(\Delta\omega_\text{rad}\,t) ,
    \label{eq:rad_overlap_drift}
\end{equation}
with $\Delta\omega_\text{rad} = \frac{2\kappa_\text{rad}}{\hbar}(p_B-p_A)$, orthogonalizing $(p_A,p_B)=(\tfrac{1}{4},\tfrac{3}{4})$ at finite time $t_*^\text{rad}=\pi\hbar/\kappa_\text{rad}$~\citep{Polchinski1991}.
\item \emph{Superluminal Bell gravitational-wave power signaling:} If Alice and Bob share $|\Phi^+\rangle=\tfrac{1}{\sqrt{2}}(|00\rangle+|11\rangle)$ and Bob drives oscillating $+/\times$ quadrupoles $\hat{I}_{ij}^\text{TF}(t) = \mathcal{Q}_0\cos(\Omega_\text{quad}t)\big(E_{ij}^{(0)}|0\rangle\langle 0| + E_{ij}^{(1)}|1\rangle\langle 1|\big)$, then under collapsed-branch sourcing Bob's ensemble-average quadrupole power $P_\text{GW} = \frac{G}{5c^5}\langle\dddot{I}_{ij}^\text{TF}\dddot{I}_{ij}^\text{TF}\rangle_t$ equals $P_\text{GW}^{(Z)} = P_0 \equiv \frac{G\mathcal{Q}_0^2\Omega_\text{quad}^6}{5c^5}$ when Alice measures $Z$ versus $P_\text{GW}^{(X)}=\tfrac{1}{2}P_0$ when Alice measures $X$, producing the superluminal power bias
\begin{equation}
    \Delta P_\text{GW} \equiv P_\text{GW}^{(Z)} - P_\text{GW}^{(X)} = \frac{1}{2}\,P_0 = \frac{G\,\mathcal{Q}_0^2\,\Omega_\text{quad}^6}{10\,c^5} > 0 .
    \label{eq:gw_power_bias}
\end{equation}
Replacing collapsed-branch sourcing by uncollapsed reduced-density-matrix sourcing removes $\Delta P_\text{GW}$ while retaining pure-state orthogonalization~\eqref{eq:rad_overlap_drift}~\citep{Polchinski1991} and remaining excluded macroscopically by \citet{PageGeilker1981Indirect}.
\end{enumerate}
\end{proposition}
\begin{proof}
Direct trace evaluation gives $\operatorname{tr}(E^{(a)})=0$, $\operatorname{tr}(E^{(a)}E^{(b)})=2\delta_{ab}$, and $\operatorname{tr}\big((\tfrac{1}{2}(E^{(0)}+E^{(1)}))^2\big)=1=\tfrac{1}{2}\operatorname{tr}((E^{(0)})^2)$. Averaging $\sin^2(\Omega_\text{quad}t)$ over a cycle gives $\tfrac{1}{2}$, so each $Z$-branch radiates $\frac{G\mathcal{Q}_0^2\Omega_\text{quad}^6}{10c^5}\cdot 2 = P_0$, whereas each $X$-branch $|\pm\rangle$ has expectation quadrupole $\tfrac{1}{2}(E^{(0)}+E^{(1)})$ radiating $\frac{G\mathcal{Q}_0^2\Omega_\text{quad}^6}{10c^5}\cdot 1 = \tfrac{1}{2}P_0$, proving Eq.~\eqref{eq:gw_power_bias}. Equation~\eqref{eq:rad_overlap_drift} follows identically to Proposition~\ref{prop:config_space_gravity}(ii) with $U_\text{self}-U_\text{cross}$ replaced by $\kappa_\text{rad}$.
\end{proof}
\noindent Proposition~\ref{prop:rad_meanfield_signaling} excludes deterministic nonlinear mean-field coupling for gravitational radiation; whether dynamical $h_{ij}^\text{TT}$ modes are themselves quantized or coupled through a linear stochastic channel remains open outside the constrained static and $1$PN sectors derived here.

Third, at the strong-field horizon boundary of Section~\ref{sec:gravity}, the lapse profile $\alpha(r)=\sqrt{1-R_s/r}$ connects the matter-side Bekenstein bound directly to the black-hole area law. With the Planck length $\ell_P = \sqrt{\hbar G/c^3}$ of Eq.~\eqref{eq:compton_planck_ratio}, a static spherical mass of energy $E=Mc^2$, areal radius $r>R_s=2GM/c^2$, and surface area $A(r)=4\pi r^2$ satisfies the exact identity
\begin{equation}
    \frac{S_\text{Bek}}{k_B} \equiv \frac{2\pi r E}{\hbar c} = \frac{A(r)}{4\ell_P^2}\big(1-\alpha(r)^2\big) \xrightarrow{\;r\to R_s^+\;} \frac{A(R_s)}{4\ell_P^2} ,
    \label{eq:bek_bh_sat}
\end{equation}
since $1-\alpha(r)^2 = R_s/r = 2GE/(rc^4)$ and $A(r)/(4\ell_P^2) = \pi r^2 c^3/(\hbar G)$. Equation~\eqref{eq:bek_bh_sat} shows that $\hbar$ enters the horizon saturation value $A(R_s)/(4\ell_P^2)$ directly from the matter wavelength scale $2\pi r E/(\hbar c)$ at $\alpha\to 0$, without quantizing the static lapse itself, while any microscopic interpretation of $A/(4\ell_P^2)$ as holographic boundary entropy belongs to quantum gravity beyond the classical exterior field laws~\citep{Bousso2002,Marolf2009}.

%===============================================================================
\section{Conclusion}
\label{sec:conclusion}
%===============================================================================

Quantum mechanics in flat spacetime gives the rest frequency one value over the whole extent of a state, and general relativity makes the rate of proper time vary with position. Dropping this single-frequency assumption makes the frequency a field over the wave function, and four derivations, from the Planck--Einstein relation, the de~Broglie relation, Hamilton's principle, and Ehrenfest's theorem, give the identity $\mathbf{F}=-\hbar\nabla\omega$ (Eq.~\eqref{eq:sgm}), with the gradient taken at fixed momentum. Every interaction with a quantum Hamiltonian drives the canonical momentum at this rate in expectation values (Theorem~\ref{thm:universal}).

In a static metric, the gravitational force on a particle at rest follows from the time-dilation factor alone, without the geodesic equation (Corollary~\ref{cor:causal_gravity}). With Gauss's law for $\alpha$ (Eq.~\eqref{eq:gauss_alpha}) as the field law, the static gradient of the time-dilated rest frequency is Newtonian gravity in the weak field (Eq.~\eqref{eq:newton}), with $\mathbf{g}=-c^2\nabla\alpha$ (Corollary~\ref{cor:causal_gravity}), and the dynamic gradient from the Rindler factor is the inertial reaction $-m\mathbf{a}$ (Eq.~\eqref{eq:inertia}). Both have the coefficient $\omega_0$, so $m_g=m_i$ and every body falls with the same acceleration (Section~\ref{sec:freefall}); the two separate in the Hessian of $\alpha$, which for gravity is the tidal curvature and for uniform acceleration is zero (Remark~\ref{rem:rank}). In the accelerating rest frame the applied force cancels the inertial reaction, which is Newton's second law (Eq.~\eqref{eq:second_law}). The dispersion relation of a wave in the field of a static mass (Eq.~\eqref{eq:isotropic_dispersion}) gives the rest frequency $\omega_0\alpha$ at zero wavevector and, at zero mass, gives light the refractive index $1/\alpha^2$ in the weak field, and with it the full deflection of light and the Shapiro delay (Section~\ref{sec:light}). The nonrelativistic reduction removes both $\omega_0$ and the time-dilation shift and keeps each force only as a given potential (Section~\ref{sec:duality}), while the Dirac equation keeps both (Section~\ref{sec:dirac}).

Written as one quaternion \citep{HutchisonWavefn2026}, the spinor of a spin-$\tfrac12$ particle couples to the additive and the multiplicative class through two covariant derivatives (Table~\ref{tab:classes}), and in a stationary metric and fields in which the positive- and negative-frequency states separate, the Dirac equation reduces exactly to a frequency equation. On conformally flat spatial slices this is the master frequency equation (Eq.~\eqref{eq:master}), whose expansion in $c^{-2}$ gives the terms that produce the common force laws and spin precessions (Table~\ref{tab:forces}). The Lagrangian is $\hbar$ times the rate at which a path advances the phase (Eq.~\eqref{eq:action_phase}). Momentum and the total angular momentum, orbital and spin, are conserved because the total frequency of an isolated system does not change under a common displacement or rotation (Eq.~\eqref{eq:sgm_angular}); in turning coordinates the angular gradient gives the gyroscopic torque and Euler's equations (Eqs.~\eqref{eq:gyroscopic} and~\eqref{eq:euler_rigid}). The Aharonov--Bohm phase is accumulated locally along each path, with no degree of freedom beyond the particle's (Section~\ref{sec:ab}).

Every coupling of a particle acts wherever its wave function is nonzero, so its reach is limited only by causal contact (Proposition~\ref{prop:reach}). Planck's constant cancels from the gravitational force, and no product of powers of $c$ and $G$ has the dimension of action (Proposition~\ref{prop:no_grav_action}), so on the two premises stated in Section~\ref{sec:hbar_cancels} the constrained gravitational sector is unquantized and no graviton is exchanged. The static time-dilation factor, fixed continuously by its source with no degrees of freedom of its own at the particle (Compton) scale—roughly $20$ orders of magnitude above the Planck scale $\ell_P$ (Eq.~\eqref{eq:compton_planck_ratio})—acts on the joint wave function as the Hermitian pair operator $H_\text{grav,pair}$ (Eq.~\eqref{eq:h_grav_pair}, Corollary~\ref{cor:fields}), conserving total probability unitarily and giving the exact BMV entanglement purity of Eq.~\eqref{eq:gie_phase} while deterministic Schr\"odinger--Newton and radiative mean fields are excluded by finite-time state orthogonalization and Bell gravitational-wave power signaling (Propositions~\ref{prop:config_space_gravity} and~\ref{prop:rad_meanfield_signaling}). At $1$PN order, moving sources couple via the manifestly Hermitian tensor-sandwich operator $H_\text{EIH}^\text{cross}$ on $\mathcal{H}_\text{mat}$ alone (Proposition~\ref{prop:eih_1pn}), and at the strong-field horizon the matter-side Bekenstein bound saturates $A/(4\ell_P^2)$ identically as $\alpha(r)\to 0$ (Eq.~\eqref{eq:bek_bh_sat}).

The four derivations are century-old mathematics, and the local equivalence of gravity and acceleration is Einstein's hypothesis. The derivation of Newton's law from the rest frequency has been found independently several times, and inertia has been obtained before from the frequency change that accompanies motion (Section~\ref{sec:intro}). This paper derives the static and the inertial gradient from the one rest frequency $\omega_0$ through one relation, and obtains from their common coefficient the equality $m_g=m_i$ and free fall, and from the second derivatives of the frequency the tidal difference between gravity and inertia, where metric theories of gravity postulate free fall along geodesics \citep{Will2014Confrontation}. The identity acts on the effective frequency at each point: the particle's time-dilated rest frequency plus the frequencies of its motion, its spin and the potentials of the other particles.

For two states that agree outside the far tail of the wave function (Axiom~\ref{ax:local_frequency}) and differ in it only in phase, the probabilities of any outcome at any one instant differ by at most $2\sqrt{P_\text{tail}}$ (Eq.~\eqref{eq:tail_distance}). The matter-wave phase of Eq.~\eqref{eq:interferometer} agrees with the cesium-fountain measurement; the first neutron measurement fell $9\%$ short, which its authors ascribed to bending of the interferometer under its own weight.

The inertial $\alpha$ is derived, from special relativity with the acceleration given (Section~\ref{sec:inertia}); the gravitational $\alpha$ is not. Its radial dependence follows from the dilution of a conserved flux over a growing spherical area (Eq.~\eqref{eq:gauss_alpha}), but the strength of that flux is the measured constant $G$ (Corollary~\ref{cor:causal_gravity}), and the identity would convert any profile of $\alpha$ to a force in the same way. While Section~\ref{sec:dynamical_gravity} extends the constrained matter-only coupling to $1$PN cross-momentum interactions and excludes deterministic nonlinear radiative mean fields, the full relativistic field equations for moving sources (including $2.5$PN radiation reaction and all-orders aberration cancellation), quantized dynamical radiation, and fully time-dependent strong-field metrics remain outside the framework (Section~\ref{sec:scope}).

%===============================================================================
% Bibliography
%===============================================================================
\bibliographystyle{apsrev4-2}
\bibliography{references-sgm}

% Without the supplement the document ends here. \expandafter closes the
% conditional before \SGMendmain runs, so no \if is left open; the argument
% sits inside the macro, so nothing is left to typeset when the switch is on.
\newcommand{\SGMendmain}{\end{document}}
\ifsupplement\else\expandafter\SGMendmain\fi

\clearpage

%===============================================================================
% Supplemental material.
%
% Numbering restarts here with an S prefix. Each \theH... form is hyperref's
% anchor name for the same counter: the restarted counters would otherwise
% generate anchors that collide with the main text's, and hyperref discards a
% duplicate anchor, which silently breaks the links that point to it. They are
% set with \def because hyperref leaves \theHsubsection and \theHremark
% undefined, so \renewcommand would fail on them.
%===============================================================================
\setcounter{section}{0}
\renewcommand{\thesection}{S\arabic{section}}
\setcounter{equation}{0}
\renewcommand{\theequation}{S\arabic{equation}}
\setcounter{remark}{0}
\renewcommand{\theremark}{S\arabic{remark}}
\def\theHsection{S\arabic{section}}
\def\theHsubsection{S\arabic{section}.\arabic{subsection}}
\def\theHequation{S\arabic{equation}}
\def\theHremark{S\arabic{remark}}
\setcounter{table}{0}
\renewcommand{\thetable}{S\arabic{table}}
\def\theHtable{S\arabic{table}}

\begin{center}
{\Large\bfseries Supplemental Material}
\end{center}

\vspace{1ex}

\noindent Sections, equations, and tables numbered with a leading S belong to this supplemental material; every other number refers to the main text.

%===============================================================================
\section{Notation and conventions}
\label{sup:notation}
%===============================================================================

The metric signature is $(-,+,+,+)$ and the coordinates are $x^\mu=(ct,\mathbf{x})$. The component $g_{00}$ is dimensionless, with $g_{00}=-1$ in inertial Cartesian coordinates in flat spacetime. As in Section~\ref{sec:identity}, frequencies, rates, and energies are measured per unit of the coordinate time $t$ unless they are called proper. The rest frequency $\omega_0=mc^2/\hbar$ (Eq.~\eqref{eq:rest_freq}) is the proper rate of the Compton oscillation and is the same at every location. Its coordinate rate for a particle at rest, $\omega=\omega_0 \, \alpha$ (Eq.~\eqref{eq:omega_alpha}), varies with position, while the phase rate $-\partial_t\theta$ of a stationary state is uniform. The gradient $\nabla$ acts on position at fixed momentum (Eq.~\eqref{eq:gradient_def}).

Bold letters are triples: three-vectors such as $\mathbf{A}$ and $\mathbf{B}$, the quaternion units $\mathbf{e}=(e_1,e_2,e_3)$, and the Pauli matrices $\boldsymbol{\sigma}$. Unless otherwise indicated, the italic letter with the subscript $k$ is the $k$th Cartesian component of a bold letter, as $a_k$ of $\mathbf{a}$. A hat on a bold letter marks a unit vector, as in $\hat{\mathbf{a}}=\mathbf{a}/a$, except on $\mathbf{x}$, $\mathbf{p}$, $\boldsymbol{\pi}$, and $\mathbf{L}$, where it marks an operator. A hat on a letter that is not bold marks an operator, as in $\hat{H}$ and $\hat{\omega}$. The kinetic orbital angular momentum $\mathbf{L}=\mathbf{x}\times\hat{\boldsymbol{\pi}}$ carries no hat, and the canonical one, $\hat{\mathbf{L}}=\mathbf{x}\times\hat{\mathbf{p}}$, does (Section~\ref{sec:angular}). Other operators with letters of their own, such as $\mathbf{S}$, $\mathsf{N}$, $\mathsf{X}$, and $D_k$, carry no hat.

Sans-serif letters denote the operators $\mathsf{N}$ and $\mathsf{X}$ of the quaternionic wave equation, the lapse energy $\mathsf{G}$ and the commutator term $\mathsf{C}$ of the master frequency equation, the second imaginary unit $\mathsf{i}$ of the biquaternions, the rotation operator $\mathsf{T}_\vartheta$, the dimensions $\mathsf{M}$, $\mathsf{L}$, and $\mathsf{T}$, and the matrix $\mathsf{D}$ of their exponents.

All equations are in SI units. Table~\ref{tab:notation} lists each symbol of the main text with its meaning, the equation or section that defines it, and its SI units. The symbols are grouped by the section that defines them formally, in the order of the text. For an operator the units column gives the units by which the operator multiplies the function it acts on, and ``as $\psi$'' stands for the units of the wave function, $\mathrm{m^{-3/2}}$ for a normalized state. The units column is empty for a generic symbol, such as an arbitrary operator or quaternion. A letter that has two meanings in the main text is listed once for each meaning. Primed variables of integration, such as $r'$, are not listed.

Four kinds of product occur besides the ordinary product of numbers.
\begin{enumerate}[label=(\roman*)]
\item Two quaternions written side by side are multiplied by the product of Eq.~\eqref{eq:quaternion_units}, in the order in which they are written. The product is associative and not commutative: $e_je_k=-e_ke_j$ for $j\neq k$. A quaternion written to the right of $\psi$, as in $\psi\,e_3$, multiplies it from the right, and the imaginary unit $i$ acting on $\psi$ is right multiplication by $-e_3$ (Section~\ref{sec:spinor}). The shorthand $\mathbf{Y}\cdot\mathbf{e}=\mathbf{e}\cdot\mathbf{Y}=Y_1e_1+Y_2e_2+Y_3e_3$ is a quaternion, a sum over the three units, and not the dot product of two vectors. The unit $\mathsf{i}$ of the biquaternions commutes with every quaternion and is distinct from the $i$ that acts on $\psi$ (Section~\ref{sec:maxwell}).
\item The shorthand $\boldsymbol{\sigma}\cdot\mathbf{Y}=\sigma_1Y_1+\sigma_2Y_2+\sigma_3Y_3$ is a $2\times2$ matrix, applied to the two components $(\psi_1,\psi_2)$ of the spinor written as a column (Eq.~\eqref{eq:pauli_quaternion}). When the $Y_k$ are operators, as in $\boldsymbol{\sigma}\cdot\hat{\boldsymbol{\pi}}$, it is a $2\times2$ matrix of operators.
\item A product of two operators is their composition, the right factor applied first (Section~\ref{sec:wave_equation}). The commutator is $[\hat{A},\hat{B}]=\hat{A}\hat{B}-\hat{B}\hat{A}$ and the anticommutator $\{\hat{A},\hat{B}\}=\hat{A}\hat{B}+\hat{B}\hat{A}$.
\item Between two spatial vectors, $\mathbf{a}\cdot\mathbf{b}$ and $\mathbf{a}\times\mathbf{b}$ are the dot and cross products. A repeated index on the components of spatial vectors and tensors is summed from $1$ to $3$. Indices that label particles are summed only where a $\sum$ is written. Repeated spacetime indices $\mu$ and $\nu$, as in $g^{\mu\nu}k_\mu k_\nu$, are summed from $0$ to $3$.
\end{enumerate}

{\small
\begin{longtable}{@{}p{2.5cm}p{6.6cm}p{3.0cm}p{2.2cm}@{}}
\caption{Symbols of the main text.}\label{tab:notation}\\
\toprule
\textbf{Symbol} & \textbf{Meaning} & \textbf{Defined in} & \textbf{Units} \\
\midrule
\endfirsthead
\toprule
\textbf{Symbol} & \textbf{Meaning} & \textbf{Defined in} & \textbf{Units} \\
\midrule
\endhead
\bottomrule
\endfoot
\multicolumn{4}{@{}l}{\itshape Section~\ref{sec:intro}} \\*
$g$ & strength of a uniform gravitational field; Earth's surface gravity & Section~\ref{sec:intro} & $\mathrm{m\,s^{-2}}$ \\
$z$ & height & Section~\ref{sec:intro} & $\mathrm{m}$ \\
\addlinespace
\multicolumn{4}{@{}l}{\itshape Section~\ref{sec:identity}} \\*
$\mathbf{x}$, $t$ & position; coordinate time of the frame in which the fields are written & Section~\ref{sec:identity} & $\mathrm{m}$; $\mathrm{s}$ \\
$\omega(\mathbf{x})$ & frequency field of the wave function, its phase rate $-\partial_t\theta$ & Axiom~\ref{ax:local_frequency}; Eq.~\eqref{eq:phase_freq} & $\mathrm{rad\,s^{-1}}$ \\
$\theta$ & phase of the wave function & Axiom~\ref{ax:local_frequency} & $\mathrm{rad}$ \\
$q$, $\Phi$, $\lambda$ & charge; electric scalar potential; gauge function, $\psi\to e^{iq\lambda/\hbar}\psi$ & Section~\ref{sec:identity} & $\mathrm{C}$; $\mathrm{V}$; $\mathrm{V\,s}$ \\
$i$ & imaginary unit; from Section~\ref{sec:spinor} on, acting on the quaternion $\psi$, right multiplication by $-e_3$ & Section~\ref{sec:identity} & $1$ \\
$\psi$ & wave function; from Section~\ref{sec:spinor} on, a two-component spinor written as one quaternion; from Eq.~\eqref{eq:pure_form} on, the chiral component $\varphi_R$ & Section~\ref{sec:identity}; Eq.~\eqref{eq:spinor_quaternion}; Eq.~\eqref{eq:pure_form} & $\mathrm{m^{-3/2}}$ \\$\psi_\text{core}$, $\psi_\text{tail}$, $\psi_\text{box}$ & parts of a normalized state $\psi=\psi_\text{core}+\psi_\text{tail}$ outside and inside the far tail; renormalized core $\psi_\text{core}/\sqrt{1-P_\text{tail}}$ & Section~\ref{sec:identity} & as $\psi$ \\
$T$ & evolution duration for causal core padding $cT$ & Section~\ref{sec:identity} & $\mathrm{s}$ \\
$\langle\phi_1|\phi_2\rangle$, $\|\phi\|$ & inner product $\int\phi_1^*\phi_2\,d^3x$ of two wave functions $\phi_1$ and $\phi_2$; norm $\langle\phi|\phi\rangle^{1/2}$ & Section~\ref{sec:identity} & \\
$\psi^*$ & complex conjugate of the scalar wave function & Section~\ref{sec:identity} & as $\psi$ \\
$P_\text{tail}$ & probability of finding the particle in the far tail of the state & Eq.~\eqref{eq:tail_distance} & $1$ \\
$\psi'$, $\psi'_\text{tail}$ & the state $\psi_\text{core}+\psi'_\text{tail}$ after a change of frequency in the far tail; its part there & Section~\ref{sec:identity} & as $\psi$ \\
$d_\text{tr}$ & trace distance $\sqrt{1-|\langle\psi|\psi'\rangle|^2}$ between two pure states & Eq.~\eqref{eq:tail_distance} & $1$ \\
$\Pi$ & projector onto the states of a measurement outcome & Section~\ref{sec:identity} & $1$ \\
$\varepsilon$ & tolerance on $d_\text{tr}$ & Section~\ref{sec:identity} & $1$ \\
$\hbar$ & reduced quantum of action & Axiom~\ref{ax:planck} & $\mathrm{J\,s}$ \\
$E$, $E(\mathbf{x})$ & energy, $E=\hbar\omega$; energy at $\mathbf{x}$, $\hbar\,\omega(\mathbf{x})$ & Eq.~\eqref{eq:planck}; Section~\ref{sec:identity} & $\mathrm{J}$ \\
$m$, $\omega_0$ & rest mass; rest frequency $mc^2/\hbar$ & Eq.~\eqref{eq:rest_freq} & $\mathrm{kg}$; $\mathrm{rad\,s^{-1}}$ \\
$\mathcal{X}$, $\dot{\mathcal{X}}$, $\partial_t$ & any quantity; its total time derivative along the trajectory; rate at a fixed position & Section~\ref{sec:identity} & any; units of $\mathcal{X}$ times $\mathrm{s^{-1}}$; $\mathrm{s^{-1}}$ \\
$\mathbf{v}$, $v$ & velocity $\dot{\mathbf{x}}$; speed & Section~\ref{sec:identity} & $\mathrm{m\,s^{-1}}$ \\
$\mathbf{p}$ & momentum $\hbar\nabla\theta$ & Eq.~\eqref{eq:debroglie} & $\mathrm{kg\,m\,s^{-1}}$ \\
$\mathbf{k}$ & wavevector $\nabla\theta$, with components $k_j$ & Eq.~\eqref{eq:debroglie_k} & $\mathrm{m^{-1}}$ \\
$\lambdabar_\mathrm{dB}$ & reduced de~Broglie wavelength $\hbar/|\mathbf{p}|$ & Eq.~\eqref{eq:debroglie_k} & $\mathrm{m}$ \\
$\rho$, $\xi$, $\dagger$ & probability density, $|\psi|^2$ for one component; unit spinor of a wave function of several components; conjugate transpose, or adjoint & Section~\ref{sec:identity} & $\mathrm{m^{-3}}$; $1$; \\
$H$ & Hamiltonian $H(\mathbf{x},\mathbf{p})$ & Axiom~\ref{ax:hamilton} & $\mathrm{J}$ \\
$\nabla_{\mathbf{x}}$, $|_{\mathbf{p}}$ & gradient with respect to position; evaluated at fixed momentum & Axiom~\ref{ax:hamilton} & $\mathrm{m^{-1}}$; $1$ \\
$\mathbf{F}$ & force $\dot{\mathbf{p}}$, equal to $-\hbar\nabla\omega$ by Eq.~\eqref{eq:sgm} & Axiom~\ref{ax:hamilton} & $\mathrm{N}$ \\
$\omega(\mathbf{x},\mathbf{p})$ & frequency function $H/\hbar$ on phase space; equal to the phase rate on the momentum field where $Q$ is negligible & Section~\ref{sec:identity}; Eq.~\eqref{eq:phase_rate_split} & $\mathrm{rad\,s^{-1}}$ \\
$Q$ & quantum potential, $-(\hbar^2/2m)\nabla^2|\psi|/|\psi|$ & Section~\ref{sec:identity} & $\mathrm{J}$ \\
$\nabla^2$ & Laplacian & Section~\ref{sec:identity} & $\mathrm{m^{-2}}$ \\
$V$ & potential energy & Section~\ref{sec:identity} & $\mathrm{J}$ \\
$\nabla$ & gradient in position at fixed momentum & Eq.~\eqref{eq:gradient_def} & $\mathrm{m^{-1}}$ \\
$\nabla\omega$ & spectral gradient & Definition~\ref{def:spectral_gradient} & $\mathrm{rad\,s^{-1}\,m^{-1}}$ \\
$S_0$ & spatial part of $\hbar\theta$ in a stationary state; equal to Hamilton's abbreviated action where the quantum potential $Q$ is negligible & Section~\ref{sec:identity} & $\mathrm{J\,s}$ \\$c$ & speed of causality & Axiom~\ref{ax:causal_propagation} & $\mathrm{m\,s^{-1}}$ \\
$\delta\mathbf{x}$ & a displacement across the wave function's extent & Axiom~\ref{ax:causal_propagation} & $\mathrm{m}$ \\\addlinespace
\multicolumn{4}{@{}l}{\itshape Sections~\ref{sec:deriv}, \ref{sec:universal_law}, and~\ref{sec:taylor_expansion}} \\*
$\partial_{x_i}$, $\partial_{p_j}$ & derivative with respect to the component $x_i$ of position; with respect to the component $p_j$ of momentum & Eq.~\eqref{eq:chain_convective} & $\mathrm{m^{-1}}$; $\mathrm{kg^{-1}\,m^{-1}\,s}$ \\
$\dot{x}_j$ & components of the group velocity & Eq.~\eqref{eq:chain_convective} & $\mathrm{m\,s^{-1}}$ \\
$\omega_\text{osc}$ & angular frequency of a harmonic oscillator & Section~\ref{sec:route_phase} & $\mathrm{rad\,s^{-1}}$ \\
$L$ & Lagrangian $L(\mathbf{x},\dot{\mathbf{x}})$ & Eq.~\eqref{eq:euler_lagrange} & $\mathrm{J}$ \\
$\hat{H}$, $\hat{\mathbf{p}}$, $\hat{\mathbf{x}}$ & Hamiltonian operator $\hat{\mathbf{p}}^2/2m+V$; momentum operator $-i\hbar\nabla$; position operator & Section~\ref{sec:route_schrodinger} & $\mathrm{J}$; $\mathrm{kg\,m\,s^{-1}}$; $\mathrm{m}$ \\
$\hat{A}$, $\langle\hat{A}\rangle$ & an operator; its expectation value $\int\psi^*\hat{A}\psi\,d^3x$ & Section~\ref{sec:route_schrodinger} & \\
$\hat{\omega}$ & frequency operator $\hat{H}/\hbar$ & Section~\ref{sec:route_schrodinger} & $\mathrm{rad\,s^{-1}}$ \\
$[\hat{A},\hat{B}]$ & commutator $\hat{A}\hat{B}-\hat{B}\hat{A}$ of two operators & Eq.~\eqref{eq:ehrenfest_commutator} & \\
$\nabla\hat{\omega}$ & operator of the spectral gradient, $(i/\hbar)[\hat{\mathbf{p}},\hat{\omega}]$ & Section~\ref{sec:route_schrodinger} & $\mathrm{rad\,s^{-1}\,m^{-1}}$ \\$F_j$, $\partial_i$ & components of $\mathbf{F}$; derivative in $x_i$ at fixed momentum & Eq.~\eqref{eq:omega_expansion} & $\mathrm{N}$; $\mathrm{m^{-1}}$ \\
\addlinespace
\multicolumn{4}{@{}l}{\itshape Section~\ref{sec:gravity}} \\*
$\tau$, $\alpha$ & proper time; time-dilation factor $d\tau/dt=\sqrt{-g_{00}}$ & Eq.~\eqref{eq:alpha_def} & $\mathrm{s}$; $1$ \\
$g_{00}$ & time component of the metric, signature $(-,+,+,+)$ & Eq.~\eqref{eq:alpha_def} & $1$ \\
$\partial_\tau$ & rate in the proper time of the particle's rest frame & Section~\ref{sec:gravity} & $\mathrm{s^{-1}}$ \\
$\ell$, $\hat{\mathbf{u}}$ & distance of each of two clocks from $\mathbf{x}$; unit vector along their separation & Section~\ref{sec:gravity} & $\mathrm{m}$; $1$ \\$\mathbf{F}_\text{grav}$ & gravitational force $-mc^2\nabla\alpha$ & Eq.~\eqref{eq:geometric_force} & $\mathrm{N}$ \\
$\mathcal{V}$, $\partial\mathcal{V}$, $d\boldsymbol{\mathcal{S}}$, $d\mathcal{S}$ & a volume; its boundary; the outward area element; its magnitude & Eq.~\eqref{eq:gauss_alpha} & $\mathrm{m^3}$ for $\mathcal{V}$, $\mathrm{m^2}$ for $d\boldsymbol{\mathcal{S}}$ and $d\mathcal{S}$ \\
$G$, $M_\text{enc}$ & gravitational constant; enclosed mass & Eq.~\eqref{eq:gauss_alpha} & $\mathrm{m^3\,kg^{-1}\,s^{-2}}$; $\mathrm{kg}$ \\
$r$, $M$ & radial distance; source mass & Eq.~\eqref{eq:gauss_sphere} & $\mathrm{m}$; $\mathrm{kg}$ \\
$\hbar/(mc)$ & reduced Compton wavelength & Section~\ref{sec:gravity} & $\mathrm{m}$ \\
$\hat{\mathbf{r}}$ & outward radial unit vector & Eq.~\eqref{eq:grav_gradient} & $1$ \\
$\Phi_g$, $\mathbf{g}$ & Newtonian potential, with $\alpha=1+\Phi_g/c^2$ to first order in $\Phi_g/c^2$; gravitational field $-c^2\nabla\alpha$, which is $-\nabla\Phi_g$ to the same order & Eq.~\eqref{eq:g_grad_alpha} & $\mathrm{m^2\,s^{-2}}$; $\mathrm{m\,s^{-2}}$ \\
$R_s$ & Schwarzschild radius, $GM=R_sc^2/2$ & Section~\ref{sec:gravity} & $\mathrm{m}$ \\
$r_\text{p}$ & proper radial distance, $dr_\text{p}=dr/\alpha$ & Eq.~\eqref{eq:dalpha_dl} & $\mathrm{m}$ \\
$r_1$, $r_2$ & areal radii of two clocks at rest & Section~\ref{sec:gravity} & $\mathrm{m}$ \\
$\mathbf{F}_\text{SGM}$, $\mathbf{F}_\text{prop}$, $\mathbf{F}_\infty$ & coordinate-gradient force $-mc^2\nabla_r\alpha$; local proper weight $-mc^2\nabla_\text{prop}\ln\alpha$; force at infinity $\alpha\mathbf{F}_\text{prop}$ & Section~\ref{sec:gravity} & $\mathrm{N}$ \\
$h_{jk}$, $h^{jk}$, $a_\text{prop}$ & spatial metric of a static slice; its inverse; proper acceleration of a static observer & Section~\ref{sec:gravity} & $1$; $1$; $\mathrm{m\,s^{-2}}$ \\
\addlinespace
\multicolumn{4}{@{}l}{\itshape Section~\ref{sec:inertia}} \\*
$\mathbf{a}$, $a$, $\hat{\mathbf{a}}$, $\delta x_\parallel$ & proper acceleration; its magnitude; its direction $\mathbf{a}/a$; component of $\delta\mathbf{x}$ along $\mathbf{a}$ & Section~\ref{sec:inertia} & $\mathrm{m\,s^{-2}}$; $\mathrm{m\,s^{-2}}$; $1$; $\mathrm{m}$ \\
$X$ & proper distance of a point of a Born-rigid state from the common horizon & Section~\ref{sec:inertia} & $\mathrm{m}$ \\
$ds$ & spacetime line element & Section~\ref{sec:inertia} & $\mathrm{m}$ \\$\delta\alpha$ & fractional frequency shift $\mathbf{a}\cdot\delta\mathbf{x}/c^2$, the fraction by which the rate of proper time at $\delta\mathbf{x}$ in the accelerating frame exceeds that at the mean position & Eq.~\eqref{eq:rapidity_gain} & $1$ \\
$\hat{\mathbf{d}}$ & unit vector $\delta\mathbf{x}/|\delta\mathbf{x}|$ & Eq.~\eqref{eq:front_behind} & $1$ \\
$\mathbf{F}_\text{in}$ & inertial reaction $-m\mathbf{a}$ & Eq.~\eqref{eq:inertia} & $\mathrm{N}$ \\
$\omega_\text{st}$, $\omega_\text{acc}$ & coordinate rate of the phase in the momentary inertial rest frame; in the accelerating rest frame & Eqs.~\eqref{eq:detuned_rate} and~\eqref{eq:flat_total} & $\mathrm{rad\,s^{-1}}$ \\
$t_\text{in}$, $t_\text{R}$ & momentarily comoving inertial time; Rindler coordinate time & Section~\ref{sec:inertia} & $\mathrm{s}$ \\
$\hat{\mathbf{n}}$ & direction in which a change of frequency spreads; from Section~\ref{sec:spinor} on, the axis of a rotation & Section~\ref{sec:inertia} & $1$ \\
\addlinespace
\multicolumn{4}{@{}l}{\itshape Section~\ref{sec:duality}} \\*
$m_g$, $m_i$ & gravitational mass, the coefficient of $-c^2\nabla\alpha$ in $\mathbf{F}_\text{grav}$; inertial mass, the coefficient of $-\mathbf{a}$ in $\mathbf{F}_\text{in}$ & Section~\ref{sec:duality} & $\mathrm{kg}$ \\
$\tilde{\mathbf{x}}$, $\psi_{\mathbf{a}}$ & position in coordinates accelerating at $\mathbf{a}$; wave function in those coordinates & Section~\ref{sec:duality} & $\mathrm{m}$; as $\psi$ \\
$\mathbf{x}_\text{s}$ & source point of the free Schr\"odinger evolution kernel & Section~\ref{sec:duality} & $\mathrm{m}$ \\
\addlinespace
\multicolumn{4}{@{}l}{\itshape Section~\ref{sec:freefall}} \\*
$\nabla\omega_\text{grav}$, $\nabla\omega_\text{in}$ & static gradient $\omega_0\nabla\alpha$; dynamic gradient with components $\omega_0\gamma_{ij}a_c^j/(\alpha_0c^2)$, summed over $j$; $\omega_0\mathbf{a}_c/c^2$ in flat spacetime & Section~\ref{sec:freefall} & $\mathrm{rad\,s^{-1}\,m^{-1}}$ \\
$\gamma_{ij}$, $\gamma^{ij}$ & spatial part of the static metric; its inverse & Section~\ref{sec:freefall} & $1$ \\
$\mathbf{x}_0$, $\mathbf{a}_c$ & position and coordinate acceleration of the particle in the static coordinates & Section~\ref{sec:freefall} & $\mathrm{m}$; $\mathrm{m\,s^{-2}}$ \\$\tilde{t}$, $\tilde{\alpha}$ & time and time-dilation factor in the coordinates that move with the particle & Eqs.~\eqref{eq:moving_sync} and~\eqref{eq:moving_frame_alpha} & $\mathrm{s}$; $1$ \\
$\alpha_0$ & time-dilation factor at the particle, $\alpha(\mathbf{x}_0)$ & Section~\ref{sec:freefall} & $1$ \\
$R_{0i0j}$ & tidal components of the Riemann tensor & Remark~\ref{rem:rank} & $\mathrm{m^{-2}}$ \\
\addlinespace
\multicolumn{4}{@{}l}{\itshape Section~\ref{sec:light}} \\*
$g_{\mu\nu}$, $g^{\mu\nu}$, $k_\mu$ & metric components; inverse metric; covariant wavevector $(-\omega/c,\mathbf{k})$ & Section~\ref{sec:light} & $1$; $1$; $\mathrm{m^{-1}}$ \\
$n_\text{r}$ & refractive index $1/\alpha^2$ & Eq.~\eqref{eq:light_dispersion} & $1$ \\
$s_\text{ray}$ & path length along a ray & Section~\ref{sec:light} & $\mathrm{m}$ \\
$b_\text{imp}$ & impact parameter of a ray & Section~\ref{sec:light} & $\mathrm{m}$ \\
$\varepsilon_r^{ij}$, $\mu_r^{ij}$, $\varepsilon_r$, $\mu_r$ & relative permittivity and permeability matrices; their common diagonal value in an isotropic medium & Eq.~\eqref{eq:constitutive_map} & $1$ \\
$\det(g_{\mu\nu})$ & determinant of the metric & Eq.~\eqref{eq:constitutive_map} & $1$ \\
$\mathcal{Z}_0$, $\varepsilon_0$, $\mu_0$ & impedance $\sqrt{\mu_0/\varepsilon_0}$ of free space; vacuum permittivity; vacuum permeability $1/(\varepsilon_0c^2)$ & Section~\ref{sec:light} & $\mathrm{V\,A^{-1}}$; $\mathrm{F\,m^{-1}}$; $\mathrm{H\,m^{-1}}$ \\
\addlinespace
\multicolumn{4}{@{}l}{\itshape Section~\ref{sec:em}} \\*
$\mathbf{A}$ & vector potential & Eq.~\eqref{eq:omega_em} & $\mathrm{T\,m}$ \\
$\boldsymbol{\pi}$ & kinematic momentum $\mathbf{p}-q\mathbf{A}$ & Eq.~\eqref{eq:omega_em} & $\mathrm{kg\,m\,s^{-1}}$ \\
$\nabla_{\mathbf{p}}$ & gradient with respect to momentum & Section~\ref{sec:em} & $\mathrm{kg^{-1}\,m^{-1}\,s}$ \\
$\mathbf{E}$, $\mathbf{B}$ & electric field $-\nabla\Phi-\partial_t\mathbf{A}$; magnetic field $\nabla\times\mathbf{A}$ & Eq.~\eqref{eq:lorentz} & $\mathrm{V\,m^{-1}}$; $\mathrm{T}$ \\
$q_\text{tot}$ & net charge of a system & Section~\ref{sec:em} & $\mathrm{C}$ \\
$V_\text{sc}$, $m_0$ & Lorentz-scalar potential, which changes the mass to $m_0+V_\text{sc}/c^2$; rest mass without it & Table~\ref{tab:classes} & $\mathrm{J}$; $\mathrm{kg}$ \\
$x^\mu$, $\partial_\mu$ & coordinates $(ct,\mathbf{x})$; $\partial/\partial x^\mu$ & Remark~\ref{rem:additivity} & $\mathrm{m}$; $\mathrm{m^{-1}}$ \\
$A_\mu$ & four-potential $(-\Phi/c,\mathbf{A})$ & Remark~\ref{rem:additivity} & $\mathrm{T\,m}$ \\
$f$, $\delta\Phi$ & coupling function in $\omega=\omega_0\,f(\Phi)$; constant shift of $\Phi$ by the gauge function $\lambda=-t\,\delta\Phi$ & Remark~\ref{rem:additivity} & $1$; $\mathrm{V}$ \\
\addlinespace
\multicolumn{4}{@{}l}{\itshape Section~\ref{sec:spinor}} \\*
$e_1$, $e_2$, $e_3$, $\mathbf{e}$ & quaternion units, $e_je_k=-\delta_{jk}+\epsilon_{jkl}e_l$; the triple $(e_1,e_2,e_3)$ & Eq.~\eqref{eq:quaternion_units} & $1$ \\
$\delta_{jk}$, $\epsilon_{jkl}$ & Kronecker delta; Levi-Civita symbol & Section~\ref{sec:spinor} & $1$ \\$Y$, $Y_0$, $\mathbf{Y}$, $\bar{Y}$ & a quaternion $Y_0+\mathbf{Y}\cdot\mathbf{e}$; its scalar part; its vector part; its conjugate $Y_0-\mathbf{Y}\cdot\mathbf{e}$ & Section~\ref{sec:spinor} & \\
$\psi_1$, $\psi_2$ & the two complex components of the spinor & Section~\ref{sec:spinor} & as $\psi$ \\
$\operatorname{Re}$, $\operatorname{Im}$ & real and imaginary parts of a complex number, as of $\psi_1$ and $\psi_2$ & Eq.~\eqref{eq:spinor_quaternion} & $1$ \\
$\boldsymbol{\sigma}$, $\sigma_k$ & the Pauli matrices, $2\times2$ & Eq.~\eqref{eq:pauli_quaternion} & $1$ \\
$\mathbf{S}$ & spin $\tfrac{1}{2}\hbar\boldsymbol{\sigma}$ & Section~\ref{sec:spinor} & $\mathrm{J\,s}$ \\
$\mathcal{R}_\psi$ & the unit quaternion $\psi/\rho^{1/2}$, the frame of the spinor & Eq.~\eqref{eq:polar_form} & $1$ \\
$\mathcal{R}$, $\vartheta$ & unit quaternion $\exp(\tfrac{1}{2}\vartheta\,\hat{\mathbf{n}}\cdot\mathbf{e})$, which rotates a vector by the angle $\vartheta$ about the unit vector $\hat{\mathbf{n}}$ & Section~\ref{sec:spinor} & $1$; $\mathrm{rad}$ \\
$\mathbf{w}$ & any vector, rotated as the quaternion $\mathbf{w}\cdot\mathbf{e}$ & Section~\ref{sec:spinor} & any \\
$\hat{\mathbf{s}}$ & direction of the spin & Eq.~\eqref{eq:spin_rotor} & $1$ \\
$\phi$, $\psi_0$ & angle of the phase factor $\exp(e_3\phi)$; spatial factor of a stationary state & Eq.~\eqref{eq:stationary_time} & $\mathrm{rad}$; as $\psi$ \\$D_t$ & gauge-covariant derivative, time component, $\partial_t\psi-(V/\hbar)\,\psi\,e_3$ & Eq.~\eqref{eq:phase_connection_time} & $\mathrm{s^{-1}}$ \\
$D_k$ & gauge-covariant derivative, space components, $\partial_k\psi+(q/\hbar)A_k\,\psi\,e_3$ & Eq.~\eqref{eq:right_connection} & $\mathrm{m^{-1}}$ \\
$\partial_k$ & $\partial/\partial x_k$, $k=1,2,3$ & Section~\ref{sec:spinor} & $\mathrm{m^{-1}}$ \\
$\hat{\boldsymbol{\pi}}$ & kinematic momentum operator $-i\hbar\nabla-q\mathbf{A}$ & Section~\ref{sec:spinor} & $\mathrm{kg\,m\,s^{-1}}$ \\
$\mathbf{D}$, $\mathbf{D}^2$, $\mathbf{e}\cdot\mathbf{D}$, $\mathbf{e}\cdot\nabla$ & $(D_1,D_2,D_3)$; $D_1^2+D_2^2+D_3^2$; quaternion operator $\sum_ke_k\,D_k$; the same without a potential, $\sum_ke_k\,\partial_k$ & Eqs.~\eqref{eq:fueter_square} and~\eqref{eq:klein_gordon} & $\mathrm{m^{-1}}$; $\mathrm{m^{-2}}$; $\mathrm{m^{-1}}$; $\mathrm{m^{-1}}$ \\
$\mathbf{e}\cdot\mathbf{B}$ & quaternion $B_1e_1+B_2e_2+B_3e_3$, multiplied from the left & Eq.~\eqref{eq:fueter_square} & $\mathrm{T}$ \\
$g_s$ & spin g-factor & Eq.~\eqref{eq:fueter_square} & $1$ \\
$\alpha_n$, $\mathbf{A}_g$, $W$ & lapse; shift, the gravitomagnetic potential; spatial scale factor & Eq.~\eqref{eq:lapse_shift_metric} & $1$; $\mathrm{m\,s^{-1}}$; $1$ \\
$\boldsymbol{\Omega}$ & angular velocity of a rotating frame & Section~\ref{sec:spinor} & $\mathrm{rad\,s^{-1}}$ \\
$\beta_\text{PPN}$, $\gamma_\text{PPN}$ & post-Newtonian nonlinearity and spatial-curvature parameters & Section~\ref{sec:spinor} & $1$ \\
$\mathbf{J}_\text{src}$ & angular momentum of a rotating source & Section~\ref{sec:spinor} & $\mathrm{J\,s}$ \\
$D_n$ & covariant derivative along the normal to the slice & Eq.~\eqref{eq:normal_derivative} & $\mathrm{m^{-1}}$ \\
$\mathbf{B}_g$, $\mathbf{L}$ & gravitomagnetic field $\nabla\times\mathbf{A}_g$; orbital angular momentum $\mathbf{x}\times\hat{\boldsymbol{\pi}}$ & Section~\ref{sec:spinor} & $\mathrm{s^{-1}}$; $\mathrm{J\,s}$ \\\addlinespace
\multicolumn{4}{@{}l}{\itshape Section~\ref{sec:wave_equation}} \\*
$\mathsf{N}$, $\mathsf{X}$, $\mathsf{X}_W$ & operators $\mathsf{N}\psi=(D_n\psi)\,e_3$, $\mathsf{X}\psi=\mathbf{e}\cdot\mathbf{D}\,\psi$, and $W^{-1/2}\mathsf{X}(W^{-1/2}\psi)$ & Eq.~\eqref{eq:NX_operators} & $\mathrm{m^{-1}}$ \\
$\kappa$ & inverse reduced Compton wavelength $mc/\hbar$ & Section~\ref{sec:wave_equation} & $\mathrm{m^{-1}}$ \\
$\psi_L$, $\psi_R$ & left- and right-handed two-component spinors of the Dirac equation & Eq.~\eqref{eq:dirac_equation} & as $\psi$ \\
$\varphi_R$, $\varphi_L$ & chiral components $\psi_R$ and $-\psi_L$, written as quaternions & Section~\ref{sec:wave_equation} & as $\psi$ \\
\addlinespace
\multicolumn{4}{@{}l}{\itshape Section~\ref{sec:master}} \\*
$\chi$ & positive-frequency spinor & Eq.~\eqref{eq:master} & as $\psi$ \\
$\hat{\Pi}_g$ & shift coupling $\alpha_n^{1/2}\,\tfrac{1}{2}\{\mathbf{A}_g,\hat{\boldsymbol{\pi}}\}\,\alpha_n^{-1/2}$ & Eq.~\eqref{eq:master} & $\mathrm{J}$ \\
$\mathsf{G}$, $\hat{V}$ & lapse energy $m_0c^2(\alpha_n-1)$; potential-energy operator $V-\hat{\Pi}_g-\tfrac{1}{4}\hbar\,\boldsymbol{\sigma}\cdot\mathbf{B}_g$ & Eq.~\eqref{eq:master} & $\mathrm{J}$ \\
$\mathsf{C}$ & commutator term of the master frequency equation & Eq.~\eqref{eq:master} & $\mathrm{J}$ \\
\addlinespace
\multicolumn{4}{@{}l}{\itshape Section~\ref{sec:expansion_orders}} \\*
$\hat{h}$ & the bracket of Eq.~\eqref{eq:master}, an operator with eigenvalue $\hbar\omega$ & Eq.~\eqref{eq:nr_expansion} & $\mathrm{J}$ \\
$\omega_\text{geom}$, $\omega_\text{ext}$, $\omega_\text{kin}$, $\omega_\text{gm}$, $\omega_\text{Pauli}$, $\omega_\text{rot}$, $\omega_\text{rel}$, $\omega_\text{so}$ & geometric, extrinsic, kinetic, gravitomagnetic, Pauli, spin-rotation, relativistic, and Darwin and spin-orbit parts of the frequency & Eq.~\eqref{eq:nr_expansion} & $\mathrm{rad\,s^{-1}}$ \\
$\tilde{\mathsf{X}}$, $\tilde{\mathsf{X}}_W$, $f$ & dressed spatial operators $e^{-f/4mc^2}\,\mathsf{X}\,e^{f/4mc^2}$ and $W^{-1/2}\tilde{\mathsf{X}}W^{-1/2}$; dressing potential $\hat{V}+\mathsf{G}-V_\text{sc}$ & Section~\ref{sec:expansion_orders} & $\mathrm{m^{-1}}$; $\mathrm{m^{-1}}$; $\mathrm{J}$ \\
$V_\text{ext}$ & potential energy of any additive force other than the Coulomb force & Section~\ref{sec:expansion_orders} & $\mathrm{J}$ \\
$\boldsymbol{\Omega}_S$ & angular velocity of spin precession & Table~\ref{tab:forces} & $\mathrm{rad\,s^{-1}}$ \\
$\dot{\boldsymbol{\Omega}}$ & rate of change of $\boldsymbol{\Omega}$ & Table~\ref{tab:forces} & $\mathrm{rad\,s^{-2}}$ \\
\addlinespace
\multicolumn{4}{@{}l}{\itshape Section~\ref{sec:maxwell}} \\*
$\rho_q$, $\mathbf{J}_q$ & charge density; current density & Section~\ref{sec:maxwell} & $\mathrm{C\,m^{-3}}$; $\mathrm{A\,m^{-2}}$ \\
$\mathsf{i}$ & imaginary unit of the biquaternions, commuting with $e_1$, $e_2$, $e_3$ & Section~\ref{sec:maxwell} & $1$ \\
$\mathcal{F}$ & field biquaternion $\mathbf{E}\cdot\mathbf{e}+\mathsf{i}\,c\,\mathbf{B}\cdot\mathbf{e}$ & Eq.~\eqref{eq:rs_biquaternion} & $\mathrm{V\,m^{-1}}$ \\
\addlinespace
\multicolumn{4}{@{}l}{\itshape Section~\ref{sec:dirac}} \\*
$\beta$ & Dirac matrix that multiplies $mc^2$ in the Dirac Hamiltonian & Section~\ref{sec:dirac} & $1$ \\
$\Omega$, $M_\Omega$ & magnitude $|\boldsymbol{\Omega}|$; component of the total angular momentum $\mathbf{L}+\mathbf{S}$ along $\boldsymbol{\Omega}$ in units of $\hbar$ & Section~\ref{sec:dirac} & $\mathrm{rad\,s^{-1}}$; $1$ \\
\addlinespace
\multicolumn{4}{@{}l}{\itshape Section~\ref{sec:canonical}} \\*
$U(\mathbf{x},\dot{\mathbf{x}})$ & generalized potential, $mc^2(\alpha-1)+q\Phi-q\,\dot{\mathbf{x}}\cdot\mathbf{A}-m\,\dot{\mathbf{x}}\cdot\mathbf{A}_g$ & Eq.~\eqref{eq:sgm_lagrangian} & $\mathrm{J}$ \\
$S$ & action $\int L\,dt$ & Eq.~\eqref{eq:action_phase} & $\mathrm{J\,s}$ \\
$\Delta\theta$ & phase accumulated along a path, $S/\hbar$ & Eq.~\eqref{eq:action_phase} & $\mathrm{rad}$ \\
$X$, $\sum_j dp_j\wedge dx_j$ & phase-space flow $(\dot{\mathbf{x}},\dot{\mathbf{p}})$ generated by $H$; symplectic form $\sum_j dp_j\wedge dx_j$ & Section~\ref{sec:hamiltonian} & $\mathrm{m\,s^{-1}}$, $\mathrm{N}$; $\mathrm{J\,s}$ \\
$p_\mu$ & four-momentum $(-E/c,\mathbf{p})$ & Eq.~\eqref{eq:sgm_fourvector} & $\mathrm{kg\,m\,s^{-1}}$ \\
$\mathbf{x}_A$, $\mathbf{x}_B$ & endpoints of a path & Eq.~\eqref{eq:work_energy} & $\mathrm{m}$ \\
$\mathcal{W}$, $K$ & work; kinetic energy $\tfrac{1}{2}m|\dot{\mathbf{x}}|^2$ & Eq.~\eqref{eq:work_energy} & $\mathrm{J}$ \\$t_1$, $t_2$, $\mathbf{J}$ & the ends of a time interval; impulse & Eq.~\eqref{eq:impulse_momentum} & $\mathrm{s}$; $\mathrm{N\,s}$ \\
$\omega_\text{sys}$, $H_\text{sys}$, $\nabla_i$ & frequency and Hamiltonian of an isolated system; gradient in the position of particle $i$ & Section~\ref{sec:conservation} & $\mathrm{rad\,s^{-1}}$; $\mathrm{J}$; $\mathrm{m^{-1}}$ \\
\addlinespace
\multicolumn{4}{@{}l}{\itshape Section~\ref{sec:angular}} \\*
$\partial\omega/\partial\vartheta$ & angular spectral gradient about $\hat{\mathbf{n}}$ & Eq.~\eqref{eq:sgm_angular} & $\mathrm{s^{-1}}$ \\
$\mathbf{b}$ & coefficient of $\boldsymbol{\sigma}$ in the Pauli and spin-rotation terms of Eq.~\eqref{eq:nr_expansion} & Section~\ref{sec:angular} & $\mathrm{J}$ \\
$\hat{\mathbf{L}}$, $\mathsf{T}_\vartheta$ & orbital angular-momentum operator $\mathbf{x}\times\hat{\mathbf{p}}$; unitary rotation by the angle $\vartheta$ about $\hat{\mathbf{n}}$ & Section~\ref{sec:angular} & $\mathrm{J\,s}$; $1$ \\
$\varphi$, $\varphi'$, $m_l$ & azimuth about $\hat{\mathbf{n}}$ in the original and the turning coordinates; component of $\mathbf{L}$ along $\hat{\mathbf{n}}$ divided by $\hbar$ & Section~\ref{sec:angular} & $\mathrm{rad}$; $\mathrm{rad}$; $1$ \\
$\mathcal{R}_{\psi_0}$ & the unit quaternion $\psi_0/\rho^{1/2}$, the frame of the time-independent part $\psi_0$ of a stationary state & Section~\ref{sec:angular} & $1$ \\
$\mathbf{N}_\text{app}$, $\mathbf{N}_\text{in}$ & applied torque; inertial (gyroscopic) torque & Eq.~\eqref{eq:torque_balance} & $\mathrm{N\,m}$ \\
$\mathcal{I}$ & inertia tensor & Eq.~\eqref{eq:euler_rigid} & $\mathrm{kg\,m^2}$ \\
$r_\perp$ & distance from the axis of rotation & Section~\ref{sec:angular} & $\mathrm{m}$ \\
\addlinespace
\multicolumn{4}{@{}l}{\itshape Section~\ref{sec:ab}} \\*
$d\boldsymbol{\ell}$, $\Phi_B$, $\Delta\theta_\text{AB}$ & line element; magnetic flux; Aharonov--Bohm phase & Eq.~\eqref{eq:ab} & $\mathrm{m}$; $\mathrm{Wb}$; $\mathrm{rad}$ \\
$h$ & Planck constant $2\pi\hbar$ & Section~\ref{sec:ab} & $\mathrm{J\,s}$ \\
$\boldsymbol{\mathcal{A}}$ & vector area enclosed by a loop & Section~\ref{sec:ab} & $\mathrm{m^2}$ \\
\addlinespace
\multicolumn{4}{@{}l}{\itshape Section~\ref{sec:experiment}} \\*
$\Delta z$, $\mathcal{A}_{zt}$ & vertical separation of two paths; spacetime area $\int\Delta z\,dt$ & Eq.~\eqref{eq:interferometer} & $\mathrm{m}$; $\mathrm{m\,s}$ \\
$\Delta\theta_g$ & gravitational phase of the upper arm of an interferometer less that of the lower & Eq.~\eqref{eq:interferometer} & $\mathrm{rad}$ \\
$T_\text{p}$, $v_\text{r}$ & interval between the pulses of a light-pulse atom interferometer; vertical recoil velocity & Section~\ref{sec:experiment} & $\mathrm{s}$; $\mathrm{m\,s^{-1}}$ \\
$m_\text{Cs}$ & mass of a cesium-133 atom & Section~\ref{sec:experiment} & $\mathrm{kg}$ \\
$\Delta\Phi$, $f_J$ & potential difference across a junction; Josephson frequency $(2e/h)\,\Delta\Phi$ & Section~\ref{sec:experiment} & $\mathrm{V}$; $\mathrm{Hz}$ \\
$e$ & elementary charge & Section~\ref{sec:experiment} & $\mathrm{C}$ \\
\addlinespace
\multicolumn{4}{@{}l}{\itshape Sections~\ref{sec:scope} and~\ref{sec:discussion}} \\*
$h^\text{TT}_{ij}$ & transverse-traceless metric perturbation of a gravitational wave & Section~\ref{sec:scope} & $1$ \\
$a_0$, $\psi_{100}$ & Bohr radius; hydrogen ground state & Section~\ref{sec:reach} & $\mathrm{m}$; as $\psi$ \\
$R_H$, $H_0$ & Hubble radius $c/H_0$; Hubble constant & Section~\ref{sec:reach} & $\mathrm{m}$; $\mathrm{s^{-1}}$ \\
$\Delta x$, $\Delta p$ & standard deviations of the coordinate $x$ and of the momentum component $p_x$ & Section~\ref{sec:reach} & $\mathrm{m}$; $\mathrm{kg\,m\,s^{-1}}$ \\
$\rho_m$ & static mass density & Eq.~\eqref{eq:static_density_clock_field} & $\mathrm{kg\,m^{-3}}$ \\
$\mathbf{y}$, $d^3y$ & integration point; volume element at it & Eq.~\eqref{eq:static_density_clock_field} & $\mathrm{m}$; $\mathrm{m^3}$ \\
$n_c$, $n_G$ & exponents of $c$ and $G$ & Section~\ref{sec:hbar_cancels} & $1$ \\
$\mathsf{M}$, $\mathsf{L}$, $\mathsf{T}$ & dimensions of mass, length and time & Section~\ref{sec:hbar_cancels} & \\
$S_\text{EH}$, $R$ & Einstein--Hilbert action; scalar curvature of the metric & Section~\ref{sec:hbar_cancels} & $\mathrm{J\,s}$; $\mathrm{m^{-2}}$ \\
$d^4x$ & coordinate volume element & Section~\ref{sec:hbar_cancels} & $\mathrm{m^4}$ \\
$\ell_P$, $m_P$, $m_p$, $m_e$ & Planck length $\sqrt{\hbar G/c^3}$; Planck mass $\sqrt{\hbar c/G}$; proton mass; electron mass & Eq.~\eqref{eq:compton_planck_ratio} & $\mathrm{m}$; $\mathrm{kg}$; $\mathrm{kg}$; $\mathrm{kg}$ \\
$m_1$, $m_2$, $m_j$, $\rho_1$, $\rho_2$, $\mathbf{x}_1$, $\mathbf{x}_2$, $\mathbf{x}_j$ & masses, normalized densities, and positions of particles & Section~\ref{sec:fields_gravity} & $\mathrm{kg}$; $\mathrm{m^{-3}}$; $\mathrm{m}$ \\
$dm$ & mass at the integration point $\mathbf{y}$ & Section~\ref{sec:fields_gravity} & $\mathrm{kg}$ \\
$\lambda_Y$ & range of a Yukawa potential $e^{-r/\lambda_Y}/r$ & Section~\ref{sec:fields_gravity} & $\mathrm{m}$ \\
$\theta_{ab}$, $a$, $b$, $r_{ab}$ & gravitational phase of the branch with the first mass at position $a$ and the second at $b$; position labels in $\{0,1\}$; branch separation & Remark~\ref{rem:lapse_test} & $\mathrm{rad}$; $1$; $\mathrm{m}$ \\
$H_\text{grav,pair}$, $\mathcal{H}_\text{mat}$, $N$ & $N$-particle configuration-space gravitational pair operator; $N$-particle matter Hilbert space; particle count & Eq.~\eqref{eq:h_grav_pair} & $\mathrm{J}$; ; $1$ \\
$|\Psi(\tau)\rangle$, $\hat{\rho}_1(\tau)$, $\Delta\theta_\text{grav}$ & two-mass joint state; reduced density operator of mass $1$; entangling gravitational phase $\theta_{00}+\theta_{11}-\theta_{01}-\theta_{10}$ & Eq.~\eqref{eq:gie_phase} & $1$; $1$; $\mathrm{rad}$ \\
$U_\text{self}$, $U_\text{cross}$, $c_a$, $p_a$ & two-site Schr\"odinger--Newton self- and cross-site coupling energies; site amplitude and probability $|c_a|^2$ & Eq.~\eqref{eq:sn_mean_field} & $\mathrm{J}$; $\mathrm{J}$; $1$; $1$ \\
$|\psi^{(A)}\rangle$, $|\psi^{(B)}\rangle$, $p_A$, $p_B$ & pair of two-site states and their $|1\rangle$-site probabilities & Eq.~\eqref{eq:sn_overlap_drift} & $1$; $1$ \\
$\Delta\omega_\text{SN}$, $t_*^\text{SN}$ & Schr\"odinger--Newton relative phase-drift rate; finite orthogonalization time $\pi\hbar/(U_\text{self}-U_\text{cross})$ & Proposition~\ref{prop:config_space_gravity} & $\mathrm{rad\,s^{-1}}$; $\mathrm{s}$ \\
$\mathcal{W}_{ij}^{ab}$, $H_\text{EIH}^\text{cross}$, $r_{ij}$, $\hat{\mathbf{r}}_{ij}$ & $1$PN EIH symmetric spatial kernel; cross-momentum Hamiltonian operator; pair separation $|\mathbf{x}_i-\mathbf{x}_j|$ and unit vector & Eqs.~\eqref{eq:eih_kernel} and~\eqref{eq:eih_1pn} & $\mathrm{kg^{-1}}$; $\mathrm{J}$; $\mathrm{m}$; $1$ \\
$\hat{T}_{\mu\nu}$, $T_{00}^{(a)}$, $E^{(0)}$, $E^{(1)}$ & stress-energy operator; static mass-energy density of state $|a\rangle$; orthogonal traceless $+/\times$ quadrupole tensors & Proposition~\ref{prop:rad_meanfield_signaling} & $\mathrm{J\,m^{-3}}$; $\mathrm{J\,m^{-3}}$; $1$ \\
$\kappa_\text{rad}$, $E_a^\text{rad}$, $\Delta\omega_\text{rad}$, $t_*^\text{rad}$ & near-zone TT self-coupling energy; diagonal radiative shifts; radiative drift rate; orthogonalization time $\pi\hbar/\kappa_\text{rad}$ & Eq.~\eqref{eq:rad_overlap_drift} & $\mathrm{J}$; $\mathrm{J}$; $\mathrm{rad\,s^{-1}}$; $\mathrm{s}$ \\
$|\Phi^+\rangle$, $\hat{I}_{ij}^\text{TF}$, $\mathcal{Q}_0$, $\Omega_\text{quad}$ & Bell state; trace-free quadrupole operator; quadrupole amplitude; oscillation frequency & Proposition~\ref{prop:rad_meanfield_signaling} & $1$; $\mathrm{kg\,m^2}$; $\mathrm{kg\,m^2}$; $\mathrm{rad\,s^{-1}}$ \\
$P_\text{GW}$, $P_0$, $\Delta P_\text{GW}$ & quadrupole gravitational-wave power; $Z$-branch power $G\mathcal{Q}_0^2\Omega_\text{quad}^6/(5c^5)$; Bell power bias $P_\text{GW}^{(Z)}-P_\text{GW}^{(X)}$ & Eq.~\eqref{eq:gw_power_bias} & $\mathrm{W}$ \\
$S_\text{Bek}$, $k_B$, $A(r)$ & Bekenstein entropy bound; Boltzmann constant; sphere area $4\pi r^2$ & Eq.~\eqref{eq:bek_bh_sat} & $\mathrm{J\,K^{-1}}$; $\mathrm{J\,K^{-1}}$; $\mathrm{m^2}$ \\
\end{longtable}
}

%===============================================================================
\section{Route IV in detail}
\label{sup:routes}
%===============================================================================

%-------------------------------------------------------------------------------
\subsection{The evolution law for expectation values}
\label{sup:route_iv}
%-------------------------------------------------------------------------------

The third line of Eq.~\eqref{eq:ehrenfest_commutator}, in Route~IV (Section~\ref{sec:route_schrodinger}), moves $\hat{H}$ across the integrand and requires the domain conditions stated after that equation. Normalizability with $\psi$ in the domain of $\hat{H}$ does not supply them. In a box of length $L_\text{box}$ with hard walls and zero potential inside, the commutator gives zero force on an equal superposition of the two lowest modes, although the walls change its mean momentum through a boundary term.

With the walls at $x=0$ and $x=L_\text{box}$, the Schr\"odinger equation inside is $i\hbar\,\partial_t\psi=-(\hbar^2/2m)\,\partial_x^2\psi$, so the rate of change of $\langle\hat{p}\rangle=-i\hbar\int\psi^*\partial_x\psi\,dx$ is $(\hbar^2/2m)\int_0^{L_\text{box}}(\psi^*\,\partial_x^3\psi-\partial_x^2\psi^*\,\partial_x\psi)\,dx$. One integration by parts, with $\psi=0$ at the walls, gives $\int_0^{L_\text{box}}\psi^*\,\partial_x^3\psi\,dx=-\int_0^{L_\text{box}}\partial_x\psi^*\,\partial_x^2\psi\,dx$, and the integrand becomes $-\partial_x|\partial_x\psi|^2$, so \[ \frac{d\langle\hat{p}\rangle}{dt}=\frac{\hbar^2}{2m}\,\big[|\partial_x\psi(0)|^2-|\partial_x\psi(L_\text{box})|^2\big] , \] the push of the left wall less that of the right.

For $\psi=(u_1+u_2)/\sqrt{2}$ at $t=0$, where $u_n=\sqrt{2/L_\text{box}}\,\sin(n\pi x/L_\text{box})$ is the $n$th mode, the slopes at the two walls are $3\pi/L_\text{box}^{3/2}$ and $\pi/L_\text{box}^{3/2}$, and the mean momentum changes at the rate $(\hbar^2/2m)(9\pi^2-\pi^2)/L_\text{box}^3=4\pi^2\hbar^2/(mL_\text{box}^3)$, the force of the walls. The commutator misses this term because $\hat{p}\psi$ does not vanish at the walls, so $\hat{p}\psi$ lies outside the domain of $\hat{H}$.

%-------------------------------------------------------------------------------
\subsection{The centroid form and the pointwise form}
\label{sec:route_state}
%-------------------------------------------------------------------------------

Route~IV uses $\psi$ directly, not its polar form, so it holds at nodes and vortex lines, where the phase is undefined. Where the polar form $\psi=|\psi|e^{i\theta}$ exists, the real part of the Schr\"odinger equation reads $-\hbar\,\partial_t\theta=\mathbf{p}^2/2m+V+Q$ with the quantum potential $Q$ of Section~\ref{sec:identity} \citep{Madelung1927Quantentheorie,Bohm1952Suggested}. The pointwise frequency therefore contains $Q$, and its gradient contains $\nabla Q$, while the right side of Eq.~\eqref{eq:sgm_expectation} is $-\nabla V$ averaged over the state. Route~II includes $Q$ in its pointwise frequency, while the frequency $H(\mathbf{x},\mathbf{p})/\hbar$ of Route~I (Eq.~\eqref{eq:H_eq_hw}) does not; Route~III assumes a path and contains no $Q$.

For a real wave function the phase gradient vanishes, so at that instant the pointwise momentum and $\langle\hat{\mathbf{p}}\rangle$ are zero. The pointwise rate $\partial_t\mathbf{p}=-\nabla(V+Q)$ need not be zero. Its average over the state is $-\langle\nabla V\rangle$, the right side of Eq.~\eqref{eq:sgm_expectation}, because $|\psi|^2\,\partial_iQ=-(\hbar^2/4m)\,\partial_j(|\psi|^2\,\partial_i\partial_j\ln|\psi|^2)$, with $j$ summed, is a divergence, whose integral vanishes for a state that falls off at infinity.

For a potential at most quadratic in position, the expectation of the external force equals that force evaluated at the centroid, $\langle\nabla V\rangle=\nabla V(\langle\hat{\mathbf{x}}\rangle)$; the pointwise rate still contains $\nabla Q$. For a cubic potential $V\propto x^3$ and a state of spread $\Delta x$ centered at $\langle\hat{x}\rangle\ne0$, the expectation $\langle\nabla V\rangle$ and the value $\nabla V(\langle\hat{\mathbf{x}}\rangle)$ differ by the fraction $(\Delta x)^2/\langle\hat{x}\rangle^2$ for a state of any shape. The fraction contains no $\hbar$, so the difference comes from the spread of the state.

Route~IV applies also to the momentum conjugate to a parameter of the Hamiltonian. Let $\hat{H}(\varpi)$ depend on the position $\varpi$ of a heavy body, whose conjugate momentum $\hat{P}_\varpi$ satisfies $[\hat{H},\hat{P}_\varpi]=i\hbar\,\partial\hat{H}/\partial\varpi$, and let the state be a normalized eigenstate of $\hat{H}(\varpi)$ with eigenvalue $E(\varpi)$. Equation~\eqref{eq:ehrenfest_commutator} with $\hat{A}=\hat{P}_\varpi$ gives $d\langle\hat{P}_\varpi\rangle/dt=-\langle\partial\hat{H}/\partial\varpi\rangle$, and the Hellmann--Feynman theorem gives $\langle\partial\hat{H}/\partial\varpi\rangle=\partial E/\partial\varpi$ \citep{Feynman1939Forces}. If the body moves slowly and smoothly, and the state starts in a level that stays nondegenerate and separated from the others by a gap along the path, the adiabatic theorem keeps the state in the instantaneous eigenstate with the same quantum numbers, up to corrections that vanish as the speed goes to zero. At zeroth order in the speed, the force on the body is then $-\partial E/\partial\varpi=-\hbar\,\partial\omega/\partial\varpi$ at fixed quantum numbers, the form of Eqs.~\eqref{eq:force_energy} and~\eqref{eq:sgm} with the parameter in place of the position.

%===============================================================================
\section{Force in Planck units}
\label{sup:planck_force}
%===============================================================================

%-------------------------------------------------------------------------------
\subsection{The unit of force and force as a phase difference}
\label{sup:force_units}
%-------------------------------------------------------------------------------

The cancellation of $\hbar$ in Eq.~\eqref{eq:hbar_cancels_identity} can also be expressed in units of phase. Take a frequency gradient that is uniform across a separation $\Delta r$, measured along $-\nabla\omega$ (the direction of the force), and steady through a time $\Delta t$. Let the momentum field of the state be uniform, so that the convective term of Eq.~\eqref{eq:chain_convective} vanishes and the full gradient of the frequency field is the spectral gradient. Let the mixed second difference $\Delta_R\Delta_T\theta$ be the phase difference across $\Delta r$ at the end of $\Delta t$ minus the phase difference at the start. By Eqs.~\eqref{eq:phase_freq} and~\eqref{eq:mixed_orders}, $|\partial_t\nabla\theta|=|\nabla\omega|$, so $\Delta_R\Delta_T\theta=|\nabla\omega|\,\Delta r\,\Delta t$ radians. Taking magnitudes in Eq.~\eqref{eq:sgm} gives $F=\hbar\,|\nabla\omega|$ for the magnitude $F$ of the force, so the force that changes the phase difference by $\Delta_R\Delta_T\theta$ is
\begin{equation}
    F = \frac{\hbar \, \Delta_R\Delta_T\theta}{\Delta r \, \Delta t} \, ,
    \label{eq:force_phase}
\end{equation}
and one radian, $\Delta_R\Delta_T\theta=1$, corresponds to the force
\begin{equation}
    F = \frac{\hbar}{\Delta r \, \Delta t} \, .
    \label{eq:force_radian}
\end{equation}

By the fundamental theorem of calculus a difference quotient is the mean of the derivative over its span, so $\Delta_R\Delta_T\theta/(\Delta r\,\Delta t)$ is the mean, sign included, of $\partial_s\partial_t\theta$ over the separation and the interval $\Delta t$, where $s$ is the distance along the separation. Equation~\eqref{eq:force_phase} uses $\Delta r$ and $\Delta t$ only to form the difference quotient; the magnitude $F$ it gives is that of $-\hbar\nabla\omega$ at one point and one instant.

The constants $\hbar$, $c$, and $G$ form one length and one time, the Planck length and the Planck time,
\begin{align}
    \ell_P &= \sqrt{\hbar G/c^3} \, , \label{eq:planck_length} \\
    t_P &= \sqrt{\hbar G/c^5} \, , \label{eq:planck_time}
\end{align}
whose product is $\ell_P \, t_P=\hbar G/c^4$. The force that changes the phase difference across $\Delta r=\ell_P$ by one radian in $\Delta t=t_P$ is the \emph{Planck force} $F_P$:
\begin{align}
    F_P &= \frac{\hbar}{\ell_P \, t_P} \label{eq:force_unit_planck} \\
    &= \frac{\hbar \, c^4}{\hbar \, G} \notag \\
    &= \frac{c^4}{G} \qquad = 1.2103\times10^{44}~\text{N} \, . \label{eq:force_unit}
\end{align}
At $\Delta r=\ell_P$ and $\Delta t=t_P$, Eq.~\eqref{eq:force_phase} with Eq.~\eqref{eq:force_unit_planck} gives
\begin{equation}
    \begin{gathered}
    \boxed{\; \Delta_R\Delta_T\theta = \frac{F}{F_P} \;} \, . \\[-3pt]
    \text{\small\itshape (Force as the change of a phase difference)}
    \end{gathered}
    \label{eq:force_as_phase}
\end{equation}
The Planck force is also the Newtonian attraction of two bodies, each of the Planck mass $\sqrt{\hbar c/G}$, held one Planck length apart, since $G\,(\hbar c/G)/\ell_P^2=c^4/G$. For the gravitational force $F=mc^2\,|\nabla\alpha|$ (Eq.~\eqref{eq:hbar_cancels_identity}), Eqs.~\eqref{eq:force_unit} and~\eqref{eq:force_as_phase} give $\Delta_R\Delta_T\theta=mG\,|\nabla\alpha|/c^2$, which contains no $\hbar$.

The Planck force $F_P$ has not been shown to bound force in general: a quarter of it has been conjectured to bound the force between two bodies in general relativity \citep{Gibbons2002Maximum}. Counterexamples within general relativity have since been constructed \citep{JowseyVisser2021Counterexamples}. \citet{Hod2024Maximum} proved that in stable self-gravitating configurations with no horizon the radial pressure, multiplied by the area $4\pi r^2$ of the sphere of radius $r$, is at most $F_P$.

In the dimensions $\mathsf{M}$, $\mathsf{L}$, and $\mathsf{T}$ of mass, length, and time (Section~\ref{sec:hbar_cancels}), $\hbar$ has dimensions $\mathsf{M} \, \mathsf{L}^2 \, \mathsf{T}^{-1}$; $c$ has dimensions $\mathsf{L} \, \mathsf{T}^{-1}$, and $G$ has dimensions $\mathsf{M}^{-1} \, \mathsf{L}^3 \, \mathsf{T}^{-2}$. In matrix form, where $\mathsf{D}$ is the matrix of exponents whose rows are the three constants and whose columns are the three dimensions,
% M, L, T are set sans-serif per ISO 80000 so the dimension M is not read as a mass M.
\begin{align}
    \mathsf{D} &= \bordermatrix{ & \mathsf{M} & \mathsf{L} & \mathsf{T} \cr
        \hbar & 1 & 2 & -1 \cr
        c & 0 & 1 & -1 \cr
        G & -1 & 3 & -2 } . \notag
\end{align}
Since $\det\mathsf{D}=2$, the three constants are dimensionally independent, and the inverse exists and is
\begin{align}
    \mathsf{D}^{-1} &=
    \bordermatrix{ & \hbar & c & G \cr
        \mathsf{M} & 1/2 & 1/2 & -1/2 \cr
        \mathsf{L} & 1/2 & -3/2 & 1/2 \cr
        \mathsf{T} & 1/2 & -5/2 & 1/2 } \, .
    \label{eq:dimension_matrix}
\end{align}

A row of dimension exponents multiplied by $\mathsf{D}^{-1}$ gives the exponents of $\hbar$, $c$, and $G$ that form that dimension, so the rows of $\mathsf{D}^{-1}$ are the exponents of the unique mass, length, and time formed from the three constants: the Planck mass
\begin{equation}
    m_P = \sqrt{\hbar c/G} \, ,
    \label{eq:planck_mass}
\end{equation}
and the Planck length and Planck time of Eqs.~\eqref{eq:planck_length} and~\eqref{eq:planck_time}.

The cancellation of $\hbar$ from the Planck force of Eq.~\eqref{eq:force_unit} also follows from the dimension matrix: force has dimensions $\mathsf{M} \, \mathsf{L} \, \mathsf{T}^{-2}$, and the row $(1, 1, -2)$ times $\mathsf{D}^{-1}$ (Eq.~\eqref{eq:dimension_matrix}) gives $(0, 4, -1)$, an exponent of zero for $\hbar$. Since $\mathsf{D}$ is invertible, $F_P=c^4/G$ is the only product of powers of $\hbar$, $c$, and $G$ with the dimensions of force. The unit $F_P$ does not make force discrete, as the Planck mass does not make mass discrete: Eq.~\eqref{eq:sgm} holds for every value of the frequency gradient and so gives every value of the force.

%-------------------------------------------------------------------------------
\subsection{The constant of action and the gravitational sector}
\label{sup:sector_action}
%-------------------------------------------------------------------------------

An action has dimension $\mathsf{M}\,\mathsf{L}^2\,\mathsf{T}^{-1}$, and the row $(1, 2, -1)$ times $\mathsf{D}^{-1}$ gives $(1, 0, 0)$, with zero exponents for $c$ and $G$: the one action the three constants form is $\hbar$ itself, as the exponent count of Section~\ref{sec:hbar_cancels} found. A quantum phase is an action divided by $\hbar$, and the constants left in a gravitational force after $\hbar$ cancels, $c$ and $G$, form no action, so the gravitational force supplies no constant of action of its own (Proposition~\ref{prop:no_grav_action}).

An inverse-square law free of $\hbar$ does not by itself make a field classical: Coulomb's law contains no $\hbar$ either, and the electromagnetic field is quantized. The electrostatic coupling $e^2/(4\pi\varepsilon_0)$, built from the elementary charge $e$ and the vacuum permittivity $\varepsilon_0$, forms with $c$ the action $e^2/(4\pi\varepsilon_0c)=\alpha_\text{fs}\hbar$, where $\alpha_\text{fs}$ is the fine-structure constant, but no product of its powers and those of $c$ is a force. Conversely, $G$ and $c$ form the force of Eq.~\eqref{eq:force_unit} but no action. Adjoining a mass forms the action $G\,m^2/c$, but $m$ is not a constant of the sector. Every observed charge is an integer multiple of $e/3$, so $e$ belongs to the sector and to no one species; the particle masses are not integer multiples of any one mass unit. The one mass at which $G\,m^2/c$ equals $\hbar$ is the Planck mass of Eq.~\eqref{eq:planck_mass}, which already contains $\hbar$.

The standard quantization of the electromagnetic field also takes its constant of action from outside the sector, since $e^2/(4\pi\varepsilon_0c)$ is not $\hbar$. A quantization of the metric by a method the identity does not describe would not contradict the argument of Section~\ref{sec:hbar_cancels}. A curvature-squared term in the field law would contradict it, since its coefficient, combined with $c$, is a constant with the dimension of action that belongs to the gravitational sector, contrary to the second premise.

%-------------------------------------------------------------------------------
\subsection{One radian per light-crossing time}
\label{sec:resolution_floor}
%-------------------------------------------------------------------------------

The unit $\hbar/(\Delta r \, \Delta t)$ of Eq.~\eqref{eq:force_radian} depends on the length $\Delta r$ and the duration $\Delta t$ over which the phase difference is measured. A frequency difference is established across a separation $r$, and the phase at one end reaches the other end, no sooner than $r/c$ (Axiom~\ref{ax:causal_propagation}), so the shortest duration over which a gradient at scale $r$ can be measured at either end is the light-crossing time $r/c$. With $\Delta r=r$ and $\Delta t=r/c$, the force that changes the phase difference across $r$ by one radian in that time is, by Eq.~\eqref{eq:force_radian},
\begin{align}
    F_P(r) &= \frac{\hbar}{(r/c) \, r} \notag\\
    &= \frac{\hbar c}{r^2} ,
    \label{eq:force_floor}
\end{align}
the unit of Eq.~\eqref{eq:force_unit_planck} evaluated over the separation $r$ and its crossing time instead of one Planck length and one Planck time. The force $F_P(r)$ is a function of the separation, and $F_P$ without an argument is the constant of Eq.~\eqref{eq:force_unit_planck}. Over one crossing time a force $F$ changes the phase difference across $r$ by $F/F_P(r)$ radians (Eqs.~\eqref{eq:force_phase} and~\eqref{eq:force_floor}), so a force weaker than $F_P(r)$ needs more than one crossing time to change it by one radian. The Planck length divided by $c$ is the Planck time (Eqs.~\eqref{eq:planck_length} and~\eqref{eq:planck_time}), so $F_P(\ell_P)=F_P$. Equation~\eqref{eq:force_floor} itself contains no $G$; $G$ enters only through the length $\ell_P$ at which $F_P(r)$ equals $F_P$.

Newton's law falls with the same power of the separation (Eq.~\eqref{eq:newton}), so for two point masses $M$ and $m$ the radians of phase difference their mutual attraction accumulates in one crossing time do not depend on the separation,
\begin{equation}
    \frac{GMm/r^2}{F_P(r)} = \frac{GMm}{\hbar c} .
    \label{eq:radians_per_crossing}
\end{equation}
One radian per crossing requires $GMm=\hbar c$, met by two equal masses at the Planck mass (Eq.~\eqref{eq:planck_mass}). For two electrons the ratio is $1.75\times10^{-45}$ at every separation, so one radian of their mutual gravitational phase takes $5.7\times10^{44}$ crossing times to accumulate, which at a separation of one meter is $6\times10^{28}$~years. Test masses of $10^{-14}$~kg give $GMm/(\hbar c)\approx2\times10^{-13}$, a factor $10^{32}$ above the value for two electrons and a factor $5\times10^{12}$ below one radian per crossing. The proposal of \citet{Bose2017Spin} maintains the superposition for about $2.5$~s, some $4\times10^{12}$ light-crossing times of the $200~\mu\text{m}$ closest approach of the two masses, so a pair held at that separation for that time accumulates about $0.8$~radian of mutual gravitational phase.

By Remark~\ref{rem:lapse_test}, a mass $m$ held at separation $r$ from a mass $M$ oscillates, to first order in $GM/(c^2r)$, at
\begin{equation}
    \omega(r) = \omega_0 - \frac{GMm}{\hbar \, r} ,
    \label{eq:pair_frequency}
\end{equation}
where the next term is smaller than the offset by a further factor of order $GM/(c^2r)$. The offset is the Newtonian interaction energy $GMm/r$ divided by $\hbar$ and contains the two masses symmetrically, so interchanging them leaves it unchanged; in the joint two-body Hamiltonian $H_\text{grav,pair}$ of Eq.~\eqref{eq:h_grav_pair}, the pair factor $\tfrac{1}{2}$ prevents double counting so that the combined system's gravitational phase rate is $-GMm/(\hbar r)$ once, and an experiment measures a difference of offsets between branches held at different separations.

Over one light-crossing time $r/c$ the offset from $\omega_0$ accumulates a phase of $GMm/(\hbar c)$ radians in magnitude, the ratio of forces of Eq.~\eqref{eq:radians_per_crossing}. The two agree because $GMm/r$ is $r$ times the force $GMm/r^2$, and $r\,F_P(r)$ is $\hbar$ divided by the crossing time $r/c$ (Eq.~\eqref{eq:force_floor}). For a Yukawa potential of range $\lambda_Y$ the ratio of forces exceeds the accumulated phase by the factor $1+r/\lambda_Y$.

Both equations treat $M$ and $m$ as point masses, and Newton's shell theorem \citep{Newton1687Philosophiae} extends them to bodies of finite size. Two spherical masses of radii $R_M$ and $R_m$, held clear of each other, exert the same force on each other and have the same mutual potential energy as point masses at their centers, so Eqs.~\eqref{eq:radians_per_crossing} and~\eqref{eq:pair_frequency} remain exact at the center separation $r$. The extent enters through the crossing time alone: comparing the near and far faces instead of the centers gives crossing times in the interval $(r\pm(R_M+R_m))/c$, so the radian count has the fractional ambiguity $(R_M+R_m)/r$, and Eq.~\eqref{eq:radians_per_crossing} counts radians per crossing only when this ratio is small.

The ratio $GMm/(\hbar c)$ is the dimensionless gravitational coupling constant, the generalization to two arbitrary masses of the gravitational fine-structure constant that \citet{CarrRees1979Anthropic} define for two protons. For two equal masses it is the square of the mass in units of $m_P$. For any central force whose magnitude falls as $1/r^2$, the ratio of its magnitude to $F_P(r)$ is the same at every separation, since Eq.~\eqref{eq:force_floor} also falls as $1/r^2$. Because Eq.~\eqref{eq:force_floor} contains no coupling constant, this ratio is the dimensionless coupling constant of the force. With the potential energy set to zero at infinity, the ratio also equals the magnitude of the mutual phase that a pair at fixed separation accumulates in one light-crossing time of the separation.

Coulomb's law \citep{Coulomb1785Electricity} divided by $F_P(r)$ is
\begin{equation}
    \frac{e^2/(4\pi\varepsilon_0 \, r^2)}{F_P(r)} = \frac{e^2}{4\pi\varepsilon_0 \, \hbar c} \qquad = 1/137.036 ,
    \label{eq:alpha_crossing}
\end{equation}
the fine-structure constant \citep{Sommerfeld1916Spectral}: two unit charges accumulate one radian of mutual phase every $137$ crossing times, while the gravitational attraction of two electrons takes $5.7\times10^{44}$ (Eq.~\eqref{eq:radians_per_crossing}).

%===============================================================================
\section{The classical mass-shell condition, the ADM reduction, and source symmetry}
\label{sup:master_assembly}
%===============================================================================

At order $\hbar^0$ the frequency of Eq.~\eqref{eq:master} is $\hbar\omega=V+\alpha_n\sqrt{(mc^2)^2+c^2W^{-2}\boldsymbol{\pi}^2}-\mathbf{A}_g\cdot\boldsymbol{\pi}$ \citep{HutchisonWavefn2026}. It is the positive root of the mass-shell condition $g^{\mu\nu}P_\mu P_\nu=-m^2c^2$ of the gauge-covariant phase gradient $P_\mu=\hbar\,\partial_\mu\theta-qA_\mu$, with $q\Phi$ extended to the whole additive potential $V$, in the metric $g_{\mu\nu}$ of Eq.~\eqref{eq:lapse_shift_metric}, where $g^{\mu\nu}$ is its inverse. When $V$, $\mathbf{A}$, and $V_\text{sc}$ vanish, it is therefore the exact Hamiltonian of geodesic motion in that stationary metric \citep{MTW1973Gravitation}. The terms that the Dirac Hamiltonian adds contain $\hbar$ and derivatives of $\alpha_n$, $W$, and $\mathbf{A}_g$: at first order in $\hbar$ the spin couplings, which contain first derivatives, and at second order the gravitational Darwin terms, which contain second derivatives (Table~\ref{tab:forces}) \citep{ObukhovSilenkoTeryaev2009Spin,ObukhovSilenkoTeryaev2011Dirac}.

Let $h_{ij}$ be the spatial metric and $h^{ij}$ its inverse (not to be confused with the Planck constant $h$) \citep{MTW1973Gravitation}. In the ADM form of a stationary metric the Hamiltonian of a particle is
\begin{equation}
    H = q\Phi + \alpha_n\sqrt{m^2c^4+c^2 \, h^{ij} \, \pi_i \, \pi_j} - A_g^i \, \pi_i ,
    \label{eq:adm}
\end{equation}
where $A_g^i$ and $\pi_i$ are the components of $\mathbf{A}_g$ and $\boldsymbol{\pi}$. For $h^{ij}=W^{-2}\delta^{ij}$, and with $q\Phi$ extended to $V$ in the same way, $H$ is $\hbar\omega$. For an anisotropic $h_{ij}$, an orthonormal spatial triad replaces $W^{-1}\hat{\pi}_k$ \citep{HutchisonWavefn2026}. For a free particle in flat spacetime the root equals $\gamma mc^2$ at every speed: $\hbar\mathbf{k}=\gamma m\mathbf{v}$, where $\gamma=(1-v^2/c^2)^{-1/2}$ is the Lorentz factor, gives $mc^2\sqrt{1+\gamma^2v^2/c^2}=\gamma mc^2$.

The Painlev\'e--Gullstrand form of the Schwarzschild metric has $\alpha_n=W=1$ and the radial shift $\mathbf{A}_g=\sqrt{2GM/r}\,\hat{\mathbf{r}}$ \citep{MartelPoisson2001Regular}, whose curl vanishes, so the spin-rotation term is absent, and at order $\hbar^0$ the shift term $-\sqrt{2GM/r}\,\hat{\mathbf{r}}\cdot\boldsymbol{\pi}$ is the only term through which gravity acts. The change to these coordinates is a function of $r$ alone and preserves the conserved energy and angular momentum of the Schwarzschild coordinates \citep{MartelPoisson2001Regular}. Hamilton's equations of Eq.~\eqref{eq:master} in these coordinates therefore give, at leading order in $GM/(c^2r)$, the perihelion advance $6\pi GM/(c^2a_\text{orb}(1-e_\text{orb}^2))$ per orbit, for semimajor axis $a_\text{orb}$ and eccentricity $e_\text{orb}$, and, at $m=0$, the deflection $4GM/(c^2b_\text{imp})$ of light at impact parameter $b_\text{imp}$ \citep{MTW1973Gravitation}. Where the shift vanishes and the spatial slices are conformally flat, as in the isotropic coordinates of Eq.~\eqref{eq:isotropic} ($h_{ij}=W^2\delta_{ij}$), gravity acts through both $\alpha_n$ and $W^{-2}$ in Eq.~\eqref{eq:master}, giving the same deflection $4GM/(c^2b_\text{imp})$ of light (Section~\ref{sec:light}) and the same 1PN perihelion advance (Table~\ref{tab:forces}).

The fields in Eq.~\eqref{eq:master} are produced by other matter. A rigid source at $\mathbf{x}_\text{s}$ produces a field that depends on position through the displacement $\mathbf{x}-\mathbf{x}_\text{s}$. The offset $\hbar(\omega-\omega_0)$ that Eq.~\eqref{eq:master} assigns the particle is then a function of that displacement, so the gradient taken at the source is the negative of the gradient taken at the particle, $\nabla_{\mathbf{x}_\text{s}}\,\omega=-\nabla\omega$. Where the coupling depends on position alone, the offset is the interaction energy of the pair, one quantity that, by Axiom~\ref{ax:planck}, appears in the frequency of each particle, and Eq.~\eqref{eq:sgm} gives two mechanical forces of equal size and opposite sign, the exchange symmetry that Eq.~\eqref{eq:newton} has for gravity. For a momentum-dependent coupling, or for a configuration that changes in time so that Axiom~\ref{ax:causal_propagation} separates the instants at which the two gradients are evaluated, the two forces need not cancel, and the momentum of the field makes up the difference. The fields of a moving source follow by a Lorentz transformation, which mixes $\Phi$ with $\mathbf{A}$, and mixes $\alpha_n$, $\mathbf{A}_g$ and the spatial metric with one another, but never mixes a potential of one class of Table~\ref{tab:classes} with a potential of the other.
% Source note.  The displacement dependence is translation invariance of the
% source's own field, so grad_{x_s} = -grad_x is the chain rule on f(x - x_s);
% verification/sgm_check.py, the block for eq:adm, check (f).

%===============================================================================
\section{Experiments in detail}
\label{sup:experiment}
%===============================================================================

The neutron interferometer of Colella, Overhauser, and Werner (Section~\ref{sec:experiment}) splits the beam into two paths that enclose a parallelogram of area $\mathcal{A}_\text{par}$ and recombine. Rotating the crystal by an angle $\vartheta_\text{tilt}$ about the incident beam tilts the parallelogram, giving the paths a vertical separation proportional to $\sin\vartheta_\text{tilt}$. The fringe phase they predict is $\Delta\theta_\text{COW}\sin\vartheta_\text{tilt}$, where $\Delta\theta_\text{COW}$ is a phase in radians (they write it $q_\text{grav}$), $\Delta\theta_\text{COW}=4\pi \, \lambda_\text{neu} \, g \, h^{-2} \, m_\text{neu}^{2} \, d \, (d+t_\text{slab}\cos\vartheta_B)\tan\vartheta_B$. Here $\lambda_\text{neu}$ is the neutron wavelength, $m_\text{neu}$ the neutron mass, $h$ the Planck constant, $d=3.5$~cm the spacing of the crystal slabs, $t_\text{slab}=0.2$~cm their thickness, and $\vartheta_B=22.1^{\circ}$ the Bragg angle. The parallelogram has area $\mathcal{A}_\text{par}=2 \, d \, (d+t_\text{slab}\cos\vartheta_B)\tan\vartheta_B$, the neutron moves at the de~Broglie speed $v=h/(m_\text{neu} \, \lambda_\text{neu})$ (Eq.~\eqref{eq:debroglie}), and on the spacetime area $\mathcal{A}_{zt}=\mathcal{A}_\text{par} \, \sin\vartheta_\text{tilt}/v$ that the tilted parallelogram sweeps at that speed, the magnitude of Eq.~\eqref{eq:interferometer} is
\begin{align}
    \frac{m_\text{neu} \, g}{\hbar} \, \frac{\mathcal{A}_\text{par} \, \sin\vartheta_\text{tilt}}{v} &= \frac{2\pi \, \lambda_\text{neu} \, g \, m_\text{neu}^{2} \, \mathcal{A}_\text{par}}{h^{2}} \, \sin\vartheta_\text{tilt} \notag\\
    &= \Delta\theta_\text{COW} \, \sin\vartheta_\text{tilt} .
    \label{eq:cow_phase}
\end{align}
At $\lambda_\text{neu}=1.445$~\AA{} the speed is $v=2.74\times10^{3}$~m/s, the area is $\mathcal{A}_\text{par}=10.5$~cm$^{2}$, and $\Delta\theta_\text{COW}=59.6$~radians, about nineteen fringes as the crystal rotates through $180^{\circ}$. They measured the bending of the interferometer under its own weight separately, by rotating it and probing it with x rays. A residual discrepancy of a few percent between neutron-interferometric and macroscopic determinations of $g$ remained unresolved decades later \citep{Peters1999Measurement}.

In the cesium fountain of Peters, Chung, and Chu (Section~\ref{sec:experiment}), a $\pi/2$ pulse of two vertical counterpropagating laser beams drives a two-photon Raman transition, splitting each atom between two states whose momenta differ by the two-photon recoil $\hbar \, k_\text{eff}$, where the effective wavenumber $k_\text{eff}$ is the sum of the two wavenumbers. After the interval $T_\text{p}$ between pulses, the halves have separated by $\hbar \, k_\text{eff} \, T_\text{p}/m_\text{Cs}$ at the recoil speed $\hbar \, k_\text{eff}/m_\text{Cs}$. A $\pi$ pulse then exchanges the two states, and a final $\pi/2$ pulse recombines them after another interval $T_\text{p}$. The separation $\Delta z(t)$ opens and closes along a triangle of base $2 \, T_\text{p}$ and height $\hbar \, k_\text{eff} \, T_\text{p}/m_\text{Cs}$, of spacetime area $\mathcal{A}_{zt}=\hbar \, k_\text{eff} \, T_\text{p}^{2}/m_\text{Cs}$, so the magnitude of Eq.~\eqref{eq:interferometer} is
\begin{equation}
    \frac{m_\text{Cs} \, g}{\hbar} \, \frac{\hbar \, k_\text{eff}}{m_\text{Cs}} \, T_\text{p}^{2} = k_\text{eff} \, g \, T_\text{p}^{2} ,
    \label{eq:atom_phase}
\end{equation}
their working formula \citep{Kasevich1992Measurement}. With the Raman beams near the cesium D$_2$ wavelength of $852$~nm, $k_\text{eff}=1.475\times10^{7}$ radians per meter. At their longest interval $T_\text{p}=160$~ms the halves separate by just over a millimeter, and their phase difference at recombination is $k_\text{eff} \, g \, T_\text{p}^{2}=3.70\times10^{6}$~radians. The fringes in their figure are the $588{,}638$th and $588{,}639$th, a count proportional to the local value of $g$; each fringe spans $1.7\times10^{-6}$ of $g$, and the value of $g$ from their least-squares fit agreed with a falling corner-cube gravimeter to $(7\pm7)$ parts in $10^{9}$.

By Eq.~\eqref{eq:sgm}, each instrument measures a force as $\hbar$ times its phase divided by its spacetime area. The neutron parallelogram at $\sin\vartheta_\text{tilt}=1$ has $\mathcal{A}_{zt}=\mathcal{A}_\text{par}/v=3.83\times10^{-7}$~m\,s, the area used in Section~\ref{sec:experiment}, and on it the measured $54.3$~radians give $1.50\times10^{-26}$~N, short of the neutron's weight by the $9\%$ that Colella, Overhauser, and Werner ascribed to the bending of the interferometer. The cesium triangle has $\mathcal{A}_{zt}=1.80\times10^{-4}$~m\,s, and $\hbar$ times its phase divided by this area gives the weight of one atom found in Section~\ref{sec:experiment}.

\citet{Dobkowski2026Observation} interpret the phase shift between the falling and the held packet of Section~\ref{sec:intro} as a test of the equivalence principle, the local equivalence of a freely falling frame and an inertial one. The atoms are rubidium-87, of mass $m_\text{Rb}$. The falling packet is launched upward with speed $v_0=g \, T_\text{fall}$ and returns to the held packet after a free-fall time $2 \, T_\text{fall}$. The held packet stays at the height $z=0$, so its kinetic and gravitational terms vanish. The magnetic energy that holds it is left out, as in the authors' free-fall prediction. In the laboratory frame the Lagrangian of the falling packet relative to the held packet therefore integrates to
\begin{align}
    S &= \int_0^{2T_\text{fall}}\Big[\tfrac12 m_\text{Rb}(v_0-gt)^2 - m_\text{Rb}\,g\big(v_0t-\tfrac12 gt^2\big)\Big]dt \notag\\
    &= -\tfrac13\,m_\text{Rb}\, g^2 T_\text{fall}^3 ,
    \label{eq:galileo_action}
\end{align}
and Eq.~\eqref{eq:action_phase} gives the relative phase $\Delta\theta_g=-m_\text{Rb}\,g^2T_\text{fall}^3/(3\hbar)$, the cubic dependence on the free-fall time that the experiment observed. \citet{AsenbaumOverstreet2025Magnetic} identify the phase as a magnetic recoil phase that depends on the inertial mass alone, and \citet{Dobkowski2025Reply} reply that the levitation condition, which weighs the atoms, makes the phase depend on the gravitational acceleration.

%===============================================================================
\section{The additive class: phenomena and measurements}
\label{sup:additive}
%===============================================================================

%-------------------------------------------------------------------------------
\subsection{Aharonov--Bohm and related phases}
\label{sup:ab_readings}
%-------------------------------------------------------------------------------

Adding $\partial_k\lambda$ to $A_k$ for a scalar function $\lambda$, together with the right multiplication
\begin{equation}
    \psi \to \psi \, \exp\!\Big({-}\frac{q}{\hbar} \, \lambda \, e_3\Big) ,
    \label{eq:gauge_phase}
\end{equation}
leaves $D_k\psi$ multiplied by that same right factor: $e_3$ commutes with the exponential, and the derivative of the exponential gives $-\tfrac{q}{\hbar}(\partial_k\lambda) \, \psi \, e_3$, the shifted potential gives $+\tfrac{q}{\hbar}(\partial_k\lambda) \, \psi \, e_3$, and the two cancel. A $\lambda$ that depends on time requires the time component to change with it, $\Phi\to\Phi-\partial_t\lambda$: in Eq.~\eqref{eq:phase_connection_time} with $V=q\Phi$, the derivative of the exponential gives $-\tfrac{q}{\hbar}(\partial_t\lambda) \, \psi \, e_3$ and the shifted potential gives $+\tfrac{q}{\hbar}(\partial_t\lambda) \, \psi \, e_3$, which cancel, so $D_t\psi$ is multiplied by the same right factor.

The field strengths of Eqs.~\eqref{eq:phase_curvature} and~\eqref{eq:time_space_curvature} are unchanged as well: $\mathbf{B}=\nabla\times\mathbf{A}$ because the curl of a gradient vanishes, and $\mathbf{E}=-\nabla\Phi-\partial_t\mathbf{A}$ because the term $+\nabla\partial_t\lambda$ from $\Phi$ cancels the term $-\partial_t\nabla\lambda$ from $\mathbf{A}$. The two arms of an interferometer end at the same point, where their wave functions are multiplied by the same factor, so Eq.~\eqref{eq:ab} is unchanged.

The time component of the four-potential gives the electric form of the Aharonov--Bohm effect. The additive potential enters Eq.~\eqref{eq:phase_connection_time} once, as the term $-(V/\hbar) \, \psi \, e_3$, and in Eq.~\eqref{eq:master} it is added to $\hbar\omega$ outside the root, while the electric field appears only in the commutator of $D_t$ with $D_k$ (Eq.~\eqref{eq:time_space_curvature}). A particle inside a conducting tube whose potential is raised and lowered while the particle is far from either end crosses a region where $\mathbf{E}=\mathbf{0}$ throughout its transit, so no force acts on it, and its frequency is offset by $V/\hbar$ for as long as the potential is applied. Two such tubes, one on each arm of an interferometer and held at different potentials, shift the fringe by the difference of $\hbar^{-1}\int V \, dt$ between the two arms, which is the second experiment Aharonov and Bohm propose \citep{Aharonov1959Significance}.

In general the phase is $(q/\hbar)\oint A_\mu \, dx^\mu$ around a closed circuit in spacetime, where $A_\mu$ is the four-potential of Remark~\ref{rem:additivity}. Equation~\eqref{eq:ab} integrates its space components, the connection of Eq.~\eqref{eq:right_connection}, around a loop in space, and the pair of tubes integrates its time component, the connection of Eq.~\eqref{eq:phase_connection_time}, over the transit.

The commutator that gives the field in Eq.~\eqref{eq:phase_curvature} also quantizes the Landau levels \citep{Landau1930Diamagnetismus,HutchisonWavefn2026}, so both effects count flux quanta, on different loops: the phase around a loop enclosing one flux quantum $h/q$ is $2\pi$, and the Landau degeneracy is one state per flux quantum through the sample. The Aharonov--Casher phase of a magnetic moment about a line charge \citep{Cimmino1989Observation} and the He--McKellar--Wilkens phase of an electric dipole in a magnetic field \citep{Lepoutre2012HeMcKellar} arise in the same way from other additive couplings.

On the multiplicative side, the gravitational phase of a neutron interferometer \citep{Colella1975Observation} and the phase of an atom interferometer in the Compton-clock interpretation \citep{Mueller2010Precision,Lan2013Clock} are phases of the rest frequency $\omega_0=mc^2/\hbar$ (Eq.~\eqref{eq:rest_freq}); the phase is proportional to the mass, and the motion it generates does not depend on the mass \citep{Kajari2010Inertial}.

%-------------------------------------------------------------------------------
\subsection{The two classes at second order}
\label{sup:second_order}
%-------------------------------------------------------------------------------

A charged particle at rest in a static potential and a static metric has frequency $\omega=\omega_0\alpha+q\Phi/\hbar$. It remains at rest where $\nabla(\omega_0\alpha+q\Phi/\hbar)=0$, the gravitational force $-mc^2\nabla\alpha$ (Eq.~\eqref{eq:geometric_force}) balanced by the electric force $-q\nabla\Phi=+mc^2\nabla\alpha$. The potential in that expression is the coordinate-frame potential. A static local observer measures the potential $\Phi_\text{loc}$ through the energy $q\Phi_\text{loc}$ of the charge. An energy measured in the observer's proper time is $\alpha$ times smaller in coordinate time, as the rest energy is in Eq.~\eqref{eq:omega_alpha}, so $\Phi=\alpha\Phi_\text{loc}$. In the master equation (Eq.~\eqref{eq:master}) the coordinate-frame potential $V$ is added outside the root, and the lapse multiplies the root alone. Redshifting the local potential and then multiplying the whole frequency by $\alpha$ would count the factor twice, an error of $q\Phi(\alpha-1)$ in the energy, second order in the fields.

%===============================================================================
\section{Surface laws and the fall-off of the couplings}
\label{sup:surface_laws}
%===============================================================================

For a static charge density $\rho_q$, the potential $\Phi(\mathbf{x})=(4\pi\varepsilon_0)^{-1}\int\rho_q(\mathbf{y})\,|\mathbf{x}-\mathbf{y}|^{-1}\,d^3y$ satisfies $\nabla^2\Phi=-\rho_q/\varepsilon_0$, since $-1/(4\pi|\mathbf{x}-\mathbf{y}|)$ is the Green function of the Laplacian (Remark~\ref{rem:reach}). For a \emph{harmonic} function $w$, one with $\nabla^2w=0$ everywhere, the divergence of $w\,\nabla\Phi-\Phi\,\nabla w$ is $w\,\nabla^2\Phi$, and the divergence theorem gives
\begin{equation}
    \oint_{\partial\mathcal{V}}\big(w\,\nabla\Phi-\Phi\,\nabla w\big)\cdot d\boldsymbol{\mathcal{S}} =-\frac{1}{\varepsilon_0}\int_{\mathcal{V}}w\,\rho_q\,d^3y ,
    \label{eq:moment_law}
\end{equation}
where the right side is taken over the volume $\mathcal{V}$ enclosed by the surface $\partial\mathcal{V}$. The left side is evaluated on the surface alone, and the right side depends only on the charge inside, so the value is the same for every surface that encloses the same distribution of charge, whatever lies outside. Equation~\eqref{eq:moment_law} with $w=1$ is Gauss's law for charge. The weight $w=x_j$, where $x_j$ is the $j$th coordinate of $\mathbf{x}$, gives the $j$th component of the enclosed dipole moment $\int\mathbf{y}\,\rho_q\,d^3y$, the harmonic product $w=x_1x_2$ gives a component of the quadrupole moment, and each further harmonic polynomial gives one more moment.

The time-dilation factor obeys Eq.~\eqref{eq:moment_law} with $\Phi$ replaced by $\alpha$ and $\rho_q/\varepsilon_0$ by $-4\pi G\rho_m/c^2$ (Eq.~\eqref{eq:static_density_clock_field}), so $w=1$ gives Eq.~\eqref{eq:gauss_alpha} and $w=x_j$ the enclosed mass times the $j$th coordinate of its center of mass.

For a steady current density $\mathbf{J}_q$, the Amp\`ere--Maxwell law of Section~\ref{sec:maxwell} is $\nabla\times\mathbf{B}=\mu_0\,\mathbf{J}_q$. Together with $\mathbf{B}=\nabla\times\mathbf{A}$ (Eq.~\eqref{eq:lorentz}) and the gauge $\nabla\cdot\mathbf{A}=0$, it gives $\nabla^2\mathbf{A}=-\mu_0\,\mathbf{J}_q$, so each component of $\mathbf{A}$ obeys Eq.~\eqref{eq:moment_law} with $\rho_q/\varepsilon_0$ replaced by $\mu_0$ times that component of $\mathbf{J}_q$. Then $w=1$ gives the net enclosed current, which is zero for a steady current that stays inside the surface, and $w=x_j$ gives $\int y_j\,\mathbf{J}_q\,d^3y$, whose antisymmetric part is the enclosed magnetic moment $\tfrac{1}{2}\int\mathbf{y}\times\mathbf{J}_q\,d^3y$, which includes the moment of the spin through its magnetization current.

By Eq.~\eqref{eq:no_monopole} the flux of $\mathbf{B}$ through every closed surface is zero, for any kernel that gives $\mathbf{A}$ from the currents. The flux of $\mathbf{B}_g=\nabla\times\mathbf{A}_g$ (Section~\ref{sec:spinor}) is zero as well, since the divergence of a curl is zero, whatever the source law of the shift, which is not derived here.

Table~\ref{tab:falloff} separates the terms of Eq.~\eqref{eq:nr_expansion} by how they depend on the distance $r$ from a compact source, outside it. In the expansion of $1/|\mathbf{x}-\mathbf{y}|$ in powers of $|\mathbf{y}|/|\mathbf{x}|$, the charge and the mass of the source give parts of $\Phi$ and $\alpha$ that fall as $1/r$, its dipole moments parts that fall as $1/r^2$, and each moment of the next order a part that falls one power faster; the moments are those of Eq.~\eqref{eq:moment_law}. A field, the gradient or curl of a potential, falls one power faster than the potential. A part of the frequency that falls as $1/r$ advances the phase, in one light-crossing time $r/c$, by a number of radians that does not depend on $r$: $GMm/(\hbar c)$ for two masses (Eq.~\eqref{eq:radians_per_crossing}), and $q_1q_2/(4\pi\varepsilon_0\hbar c)$ for two charges $q_1$ and $q_2$, which for two electrons is the fine-structure constant $\alpha_\text{fs}$.

Where the shift vanishes, $\omega_\text{geom}=\omega_0\,\alpha$, so the flux of $\nabla\omega_\text{geom}$ through a closed surface is $\omega_0$ times that of $\nabla\alpha$ (Eq.~\eqref{eq:gauss_alpha}), and with $\omega_0=mc^2/\hbar$,
\begin{equation}
    \frac{1}{4\pi c}\oint_{\partial\mathcal{V}}\nabla\omega_\text{geom}\cdot d\boldsymbol{\mathcal{S}} =\frac{G \, M_\text{enc} \, m}{\hbar c} .
    \label{eq:phase_flux}
\end{equation}
The right side is the number of radians per light-crossing time of a particle of mass $m$ paired with the enclosed mass, the same for every distribution of that mass and every enclosing surface.

Let $f_W=f-2m_0c^2\ln W$ be the dressing potential $f$ of Section~\ref{sec:expansion_orders} with the spatial scale factor included. The Darwin term is proportional to $\nabla^2f_W$ at the particle's own position ($\nabla^2V=-q\,\rho_q/\varepsilon_0$ and $\nabla^2(\mathsf{G}-2m_0c^2\ln W)=(1+2\gamma_\text{PPN})4\pi G m_0\rho_m$), so it is a contact coupling to the charge and mass density of the other particles at that point and vanishes where that density vanishes. By Proposition~\ref{prop:reach} the density is generically nonzero at every distance, but at large distance it falls as the source's $|\psi|^2$ does, exponentially for a bound eigenstate below the threshold of the continuum.

Outside a long solenoid of flux $\Phi_B$, the vector potential is $\Phi_B/(2\pi r_\perp)$ along each circle about the axis, at distance $r_\perp$ from it, so the loop integral of Eq.~\eqref{eq:ab} is $\Phi_B$ at every $r_\perp$: the Aharonov--Bohm phase does not depend on the size of the loop, although the potential falls as $1/r_\perp$. Outside the solenoid the force is zero, and at a given kinematic momentum $\boldsymbol{\pi}$ the frequency has its free value (Section~\ref{sec:ab}).

An angular-momentum eigenstate confined to the circle of radius $r_\perp$ has the phase $n\varphi$, where $\varphi$ is the azimuth about the axis, and it is single-valued only if $n$ is an integer. Its kinematic momentum along the circle is therefore $(nh-q\Phi_B)/(2\pi r_\perp)$. A change of $\Phi_B$ by $h/q$, with $n$ relabeled as $n+1$, leaves the set of momenta unchanged, so the set of frequencies is periodic in $\Phi_B$ with period $h/q$.

\begin{table}[htbp]
\centering\fontsize{9.5}{11}\selectfont
\caption{The terms of the Taylor expansion of the master frequency equation (Eq.~\eqref{eq:nr_expansion} and Table~\ref{tab:forces}), sorted by the coefficient that depends on the other particles and by how that coefficient falls with the distance $r$ outside a compact source that is static or moves slowly. The surface law is Eq.~\eqref{eq:moment_law} with weight $w$, or the same law with $\Phi$ replaced by $\alpha$ or by a component of $\mathbf{A}$. At fixed values of a term's other factors, a coefficient proportional to $r^{-p}$, for $p=1$, $2$ or $3$, advances the phase in one light-crossing time $r/c$ by a number of radians proportional to $r^{1-p}$. For a central potential $f_W(r)$ the spin-orbit term is proportional to $r^{-1}\,df_W/dr$, and the Darwin term is proportional to $\nabla^2f_W$. The $1/r$ entry for $\mathbf{A}$ is for a charge or mass that moves as a whole: a steady current that stays inside the source has zero net current, so its $\mathbf{A}$ begins at $1/r^2$. The free kinetic and relativistic terms depend on the other particles only through $\mathbf{A}$ inside $\boldsymbol{\pi}$, and a product such as $q^2\mathbf{A}^2/(2m_0)$ falls as the product of the fall-offs of its factors. The entries for $\mathbf{A}_g$ and $\mathbf{B}_g$, apart from the zero flux of $\mathbf{B}_g$, assume that the mass current produces the shift through the $1/r$ kernel, as the charge current produces $\mathbf{A}$; that source law is not derived here.}
\label{tab:falloff}
\begin{tabular}{p{0.25\linewidth}p{0.11\linewidth}p{0.26\linewidth}p{0.24\linewidth}}
\toprule
Term & Coefficient & Fall-off outside the source & Surface law \\
\midrule
Geometric, $mc^2\alpha_n$; 1PN kinetic $(\alpha_nW^{-2}-1)\boldsymbol{\pi}^2/2m_0$ & $\alpha_n-1$, $W-1$ & $1/r$ (mass); $1/r^2$ (1PN $\Phi_g^2$); $1/r^3$ (mass quadrupole, with $r$ from the center of mass) & $w=1$: enclosed mass (Eq.~\eqref{eq:gauss_alpha}) \\
Extrinsic, $V$ & $\Phi$ & $1/r$ (charge); $1/r^2$ (electric dipole) & $w=1$: enclosed charge \\
Kinetic cross term $-q\,\mathbf{v}\cdot\mathbf{A}$; gravitomagnetic, $-\mathbf{A}_g\cdot\boldsymbol{\pi}$ & $\mathbf{A}$, $\mathbf{A}_g$ & $1/r$ (slowly moving charge or mass); $1/r^2$ (magnetic moment, spin) & $w=1$: net current; $w=x_j$: moment \\
Pauli; spin-rotation & $\mathbf{B}$, $\mathbf{B}_g$ & $1/r^2$ (slowly moving charge or mass); $1/r^3$ (magnetic moment, spin) & zero flux, for every kernel \\
Spin-orbit ($\omega_\text{so}$) & $r^{-1}\,df_W/dr$ & $1/r^3$ (charge or mass) & none \\
Darwin ($\omega_\text{so}$) & $\nabla^2f_W$ & zero outside the source's charge and mass density & none; $\nabla^2V=-q\,\rho_q/\varepsilon_0$ at the point \\
Free kinetic and relativistic at $\mathbf{A}=\mathbf{0}$, $\alpha_n=W=1$ & none & independent of the other particles & none \\
\bottomrule
\end{tabular}
\end{table}

%===============================================================================
\section{Worked derivations for configuration-space entanglement, 1PN cross-momentum coupling, and radiative mean-field exclusion}
\label{sup:dynamical_unquantized_gravity}
%===============================================================================

\emph{BMV reduced density matrix and purity (Proposition~\ref{prop:config_space_gravity}(i)).}
Let $|\Psi(\tau)\rangle = \tfrac{1}{2}\sum_{a,b\in\{0,1\}} e^{i\theta_{ab}}|a,b\rangle$ be the state produced by $\exp(-i H_\text{grav,pair}\tau/\hbar)$ from $|\Psi(0)\rangle = \tfrac{1}{2}(|0\rangle+|1\rangle)\otimes(|0\rangle+|1\rangle)$, where $\theta_{ab} = G m_1 m_2 \tau / (\hbar r_{ab})$ (Eq.~\eqref{eq:h_grav_pair}). Tracing over the second mass yields the $2\times 2$ reduced density matrix
\begin{equation}
    \hat{\rho}_1(\tau) = \frac{1}{2}\begin{pmatrix} 1 & \tfrac{1}{2}\big(e^{i(\theta_{00}-\theta_{10})} + e^{i(\theta_{01}-\theta_{11})}\big) \\ \tfrac{1}{2}\big(e^{-i(\theta_{00}-\theta_{10})} + e^{-i(\theta_{01}-\theta_{11})}\big) & 1 \end{pmatrix} .
    \label{eq:sup_rho1}
\end{equation}
With $\Delta\theta_\text{grav} = (\theta_{00}-\theta_{10}) - (\theta_{01}-\theta_{11})$, the off-diagonal modulus squared is $|(\hat{\rho}_1)_{01}|^2 = \tfrac{1}{16}\big|1 + e^{-i\Delta\theta_\text{grav}}\big|^2 = \tfrac{1}{4}\cos^2(\Delta\theta_\text{grav}/2)$, so $\operatorname{tr}(\hat{\rho}_1^2) = \tfrac{1}{4}+\tfrac{1}{4}+2|(\hat{\rho}_1)_{01}|^2 = \tfrac{1}{2}+\tfrac{1}{2}\cos^2(\Delta\theta_\text{grav}/2) = 1 - \tfrac{1}{2}\sin^2(\Delta\theta_\text{grav}/2)$, which is Eq.~\eqref{eq:gie_phase}.

\emph{Two-site Schr\"odinger--Newton overlap drift (Proposition~\ref{prop:config_space_gravity}(ii)).}
Under Eq.~\eqref{eq:sn_mean_field}, $\frac{d}{dt}|c_a|^2 = c_a^*\dot{c}_a + \dot{c}_a^* c_a = 0$, so $p_1 \equiv |c_1|^2 = p$ and $p_0 \equiv |c_0|^2 = 1-p$ are exact constants of motion. Integrating Eq.~\eqref{eq:sn_mean_field} gives $c_0(t) = c_0(0)\exp\big(i[U_\text{self}(1-p)+U_\text{cross} p]t/\hbar\big)$ and $c_1(t) = c_1(0)\exp\big(i[U_\text{self} p + U_\text{cross}(1-p)]t/\hbar\big)$. For two states $|\psi^{(A)}\rangle$ and $|\psi^{(B)}\rangle$ with $c_0^{(X)}(0)=\sqrt{1-p_X}$ and $c_1^{(X)}(0)=\sqrt{p_X}$ ($X\in\{A,B\}$), the inner product is
\begin{equation}
    \langle\psi^{(A)}(t)\mid\psi^{(B)}(t)\rangle = e^{i\phi_0(t)}\Big[\sqrt{(1-p_A)(1-p_B)} + \sqrt{p_A p_B}\,e^{i\Delta\omega_\text{SN}\,t}\Big] ,
    \label{eq:sup_sn_inner}
\end{equation}
where $\phi_0(t) = (U_\text{self}-U_\text{cross})(p_A-p_B)t/\hbar$ is an overall phase and $\Delta\omega_\text{SN} = \frac{2(U_\text{self}-U_\text{cross})}{\hbar}(p_B-p_A)$. Taking the modulus squared yields Eq.~\eqref{eq:sn_overlap_drift}. At $(p_A,p_B)=(\tfrac{1}{4},\tfrac{3}{4})$, $(1-p_A)(1-p_B)+p_A p_B = 2\sqrt{p_A(1-p_A)p_B(1-p_B)} = \tfrac{3}{8}$ and $\Delta\omega_\text{SN}=(U_\text{self}-U_\text{cross})/\hbar$, reducing Eq.~\eqref{eq:sn_overlap_drift} to $\tfrac{3}{4}\cos^2\big(\frac{(U_\text{self}-U_\text{cross})t}{2\hbar}\big)$, which vanishes at $t_*^\text{SN} = \pi\hbar/(U_\text{self}-U_\text{cross})$.

\emph{Legendre sign and Lense--Thirring velocity of $H_\text{EIH}^\text{cross}$ (Proposition~\ref{prop:eih_1pn}).}
At $1$PN order, the Einstein--Infeld--Hoffmann cross-velocity Lagrangian is $L_\text{EIH}^\text{cross} = -\sum_{i<j}\frac{G m_i m_j}{2c^2 r_{ij}}\big(7\mathbf{v}_i\cdot\mathbf{v}_j + (\mathbf{v}_i\cdot\hat{\mathbf{r}}_{ij})(\mathbf{v}_j\cdot\hat{\mathbf{r}}_{ij})\big)$ \citep{LandauLifshitz1975Classical}. Writing $L = \sum_i \tfrac{1}{2}m_i \mathbf{v}_i^2 + L_\text{EIH}^\text{cross}$ and Legendre-transforming to canonical momenta $\mathbf{p}_i = m_i\mathbf{v}_i + O(c^{-2})$ yields $H = \sum_i \frac{\mathbf{p}_i^2}{2m_i} - L_\text{EIH}^\text{cross}(\mathbf{v}\to\mathbf{p}/m) + O(c^{-4})$, flipping the sign to $+H_\text{EIH}^\text{cross}$ in Eq.~\eqref{eq:eih_1pn}. Because $\hat{p}_{i,a}$ and $\hat{p}_{j,b}$ act on distinct particle coordinates for $i\neq j$, the symmetrized sandwich $\tfrac{1}{2}(\hat{p}_{i,a}\mathcal{W}_{ij}^{ab}\hat{p}_{j,b} + \hat{p}_{j,b}\mathcal{W}_{ij}^{ab}\hat{p}_{i,a})$ is Hermitian on $\mathcal{H}_\text{mat}$ without needing any operator ordering approximation between $\mathbf{x}_i$ and $\hat{p}_i$ alone, and Hamilton's equation $\dot{x}_1^a = \partial H/\partial p_{1,a}$ at $\mathbf{p}_1=\mathbf{0}$ gives
\begin{equation}
    \dot{x}_1^a\big|_{\mathbf{p}_1=\mathbf{0}} = \sum_{b=1}^{3}\mathcal{W}_{12}^{ab}(\mathbf{r})\,p_{2,b} = \frac{G}{2c^2 r}\Big(7\,p_{2,a} + (\mathbf{p}_2\cdot\hat{\mathbf{r}})\,\hat{r}^a\Big) .
    \label{eq:sup_lt_drag}
\end{equation}

\emph{Quadrupole trace algebra and Bell GW-power bias (Proposition~\ref{prop:rad_meanfield_signaling}).}
For $E^{(0)}=\operatorname{diag}(1,-1,0)$ and $E^{(1)}=\bigl(\begin{smallmatrix}0&1&0\\1&0&0\\0&0&0\end{smallmatrix}\bigr)$, the trace-free oscillating quadrupole $\hat{I}_{ij}^\text{TF}(t)=\mathcal{Q}_0\cos(\Omega_\text{quad}t)\big(E_{ij}^{(0)}|0\rangle\langle 0|+E_{ij}^{(1)}|1\rangle\langle 1|\big)$ has third time derivative $\dddot{I}_{ij}^\text{TF}(t) = \mathcal{Q}_0\Omega_\text{quad}^3\sin(\Omega_\text{quad}t)\,\langle E_{ij}\rangle$. Since $\langle\sin^2(\Omega_\text{quad}t)\rangle_t = \tfrac{1}{2}$, Einstein's quadrupole power formula $P_\text{GW} = \frac{G}{5c^5}\langle\dddot{I}_{ij}^\text{TF}\dddot{I}_{ij}^\text{TF}\rangle_t$ evaluates to
\begin{equation}
    P_\text{GW}[\langle E\rangle] = \frac{G\,\mathcal{Q}_0^2\,\Omega_\text{quad}^6}{10\,c^5}\,\operatorname{tr}\big(\langle E\rangle^2\big) .
    \label{eq:sup_gw_power}
\end{equation}
When Alice measures $|\Phi^+\rangle = \tfrac{1}{\sqrt{2}}(|00\rangle+|11\rangle)$ in the $Z$ basis, Bob's state collapses to $|0\rangle$ or $|1\rangle$ with probability $\tfrac{1}{2}$ each, giving $\operatorname{tr}((E^{(0)})^2)=\operatorname{tr}((E^{(1)})^2)=2$ and $P_\text{GW}^{(Z)} = \frac{G\mathcal{Q}_0^2\Omega_\text{quad}^6}{5c^5} = P_0$. When Alice measures in the $X$ basis, Bob's state collapses to $|\pm\rangle = \tfrac{1}{\sqrt{2}}(|0\rangle\pm|1\rangle)$ with probability $\tfrac{1}{2}$ each; both branches yield the same expectation tensor $\langle\pm|E|\pm\rangle = \tfrac{1}{2}(E^{(0)}+E^{(1)})$ with $\operatorname{tr}\big((\tfrac{1}{2}(E^{(0)}+E^{(1)}))^2\big) = \tfrac{1}{4}(2+0+2) = 1$, giving $P_\text{GW}^{(X)} = \tfrac{1}{2}P_0$ and Eq.~\eqref{eq:gw_power_bias}.

\end{document}
